\chapter{Preliminaries from Riemannian geometry}
\label{chap:prelim}
\label{sectprelim}

In this chapter we will recall some basic facts in Riemannian
geometry. For more details we refer the reader to \cite{doCarmo} and
\cite{Petersen}. Throughout, we always adopt Einstein's summation
convention on repeated indices and `manifold' means a paracompact,
Hausdorff, smooth manifold.

%%% ----------------------------------------------------------------------

\section{Riemannian metrics and the Levi-Civita connection}

Let $M$ be a manifold and let $p$ be a point of $M$. Then $TM$
denotes the tangent bundle of $M$ and $T_pM$ is the tangent space at
$p$. Similarly, $T^*M$ denotes the cotangent bundle of $M$ and
$T^*_pM$ is the cotangent space at $p$. For any vector bundle
${\mathcal V}$ over $M$ we denote by $\Gamma({\mathcal V})$ the
vector space of smooth sections of ${\mathcal V}$.

\begin{definition}[Riemannian metric]
\label{def:riemannian-metric}
\lean{MorganTianLib.RiemannianMetric}
\leanok
Let $M$ be an $n$-dimensional manifold. A \it Riemannian metric \rm
$g$ on $M$ is a smooth section of $T^*M\otimes T^*M$ defining a
positive definite symmetric bilinear form on $T_pM$ for each $p\in
M$. In local coordinates $(x^1,\cdots, x^n)$, one has a natural
local basis $\{\partial_1,\cdots,\partial_n\}$ for $TM$, where
$\partial_i=\frac{\partial}{\partial x^i}$. The metric tensor
$g=g_{ij} dx^i\otimes dx^j$ is represented by a smooth matrix-valued
function
\[ g_{ij}=g(\partial_i,\partial_j). \]
 The
pair $(M,g)$ is a {\sl Riemannian manifold}. We denote by $(g^{ij})$
the inverse of the matrix $(g_{ij})$.
\end{definition}

Using a partition of unity one can easily see that any manifold
admits a Riemannian metric. A Riemannian metric on $M$ allows us to
measure lengths of smooth paths in $M$ and hence to define a
distance function by setting $d(p,q)$ equal to the infimum of the
lengths of smooth paths from $p$ to $q$. This makes $M$ a metric
space. For a point $p$ in a Riemannian manifold $(M,g)$ and for $r>
0$ we denote the metric ball of radius $r$ centered at $p$ in $M$ by
$B(p,r)$ or by $B_g(p,r)$ if the metric needs specifying or
emphasizing. It is defined by
$$B(p,r)=\{q\in M\left|\right. d(p,q)<r\}.$$

\begin{theorem}[Levi-Civita connection]
\label{thm:levi-civita-connection}
\uses{def:riemannian-metric}
\lean{MorganTianLib.existsUnique_leviCivita}
\leanok
Given a Riemannian metric $g$ on $M$, there uniquely exists a
torsion-free connection on $TM$ making $g$ parallel, i.e., there is
a unique $\mathbb{R}$-linear mapping $\nabla\colon
\Gamma(TM)\rightarrow \Gamma(T^*M\otimes TM)$ satisfying the Leibniz
formula
\[ \nabla(fX)=df\otimes X+f\nabla X, \]
and the following two additional conditions for all vector fields
$X$ and $Y$:
\begin{itemize}
\item[$\bullet$] ($g$ orthogonal) $d(g(X,Y))=g(\nabla X,Y)+g(X,\nabla
Y)$.
\item[$\bullet$] (Torsion-free) $\nabla_XY-\nabla_YX-[X,Y]=0$
(where, as is customary, we denote $\nabla Y(X)$ by $\nabla_XY$);
\end{itemize}
\end{theorem}

\begin{proof}
\uses{def:riemannian-metric}
\leanok
We first show that any connection $\nabla$ satisfying the two
conditions obeys the \emph{Koszul formula}: for all vector fields
$X,Y,Z$,
\begin{eqnarray}\label{Koszul}
2g(\nabla_XY,Z) & = & X(g(Y,Z))+Y(g(X,Z))-Z(g(X,Y)) \\
& & +\, g([X,Y],Z)-g([X,Z],Y)-g([Y,Z],X). \nonumber
\end{eqnarray}
Indeed, evaluating the condition $d(g(X,Y))=g(\nabla X,Y)+g(X,\nabla
Y)$ on the appropriate vector fields gives the three identities
\begin{eqnarray*}
X(g(Y,Z)) & = & g(\nabla_XY,Z)+g(Y,\nabla_XZ),\\
Y(g(X,Z)) & = & g(\nabla_YX,Z)+g(X,\nabla_YZ),\\
Z(g(X,Y)) & = & g(\nabla_ZX,Y)+g(X,\nabla_ZY).
\end{eqnarray*}
Adding the first two, subtracting the third, and using the
torsion-free condition three times in the form $\nabla_XZ-\nabla_ZX =
[X,Z]$, $\nabla_YZ-\nabla_ZY=[Y,Z]$, and $\nabla_XY+\nabla_YX =
2\nabla_XY-[X,Y]$, we obtain
\begin{eqnarray*}
X(g(Y,Z))+Y(g(X,Z))-Z(g(X,Y)) & = &
2g(\nabla_XY,Z)-g([X,Y],Z)\\
& & +\,g([X,Z],Y)+g([Y,Z],X),
\end{eqnarray*}
which is Equation~(\ref{Koszul}).

\emph{Uniqueness.} The right-hand side of
Equation~(\ref{Koszul}) does not involve $\nabla$. Since $g$ is
positive definite, and hence non-degenerate, on each tangent space,
and since $Z$ is arbitrary, Equation~(\ref{Koszul}) determines the
vector field $\nabla_XY$ for every pair $X,Y$. Thus any two
connections satisfying the two conditions agree.

\emph{Existence.} Fix a coordinate chart $(x^1,\ldots,x^n)$ and
define the smooth functions $\Gamma^k_{ij}$ by
Equation~(\ref{Gamma}) below, i.e.,
$\Gamma^{k}_{ij}=\frac{1}{2}g^{kl}(\partial_i g_{lj}+\partial_j
g_{il}-\partial_l g_{ij})$. For vector fields $X=X^i\partial_i$ and
$Y=Y^j\partial_j$ with components defined on the chart, set
$$\nabla_XY=X^i\left(\partial_iY^k+Y^j\Gamma^k_{ij}\right)\partial_k.$$
This expression is $\mathbb{R}$-linear in $Y$, is $C^\infty$-linear
in $X$ (no derivatives of $X^i$ appear), so it defines a map
$Y\mapsto \nabla Y\in \Gamma(T^*M\otimes TM)$ over the chart, and it
satisfies the Leibniz formula: replacing $Y$ by $fY$ produces the
extra term $X^i(\partial_if)Y^k\partial_k=df(X)\,Y$, so
$\nabla(fY)=df\otimes Y+f\nabla Y$. We check the two conditions. From
the symmetry $g_{ij}=g_{ji}$ it is immediate from the defining
formula that $\Gamma^k_{ij}=\Gamma^k_{ji}$, and hence for the
coordinate fields
$\nabla_{\partial_i}\partial_j-\nabla_{\partial_j}\partial_i
=(\Gamma^k_{ij}-\Gamma^k_{ji})\partial_k=0=[\partial_i,\partial_j]$;
since the torsion $T(X,Y)=\nabla_XY-\nabla_YX-[X,Y]$ is
$C^\infty$-bilinear in $X$ and $Y$ (the Leibniz terms cancel), its
vanishing on coordinate fields gives $T\equiv 0$. For metric
compatibility, lowering an index in the defining formula gives
$$g_{lj}\Gamma^l_{ki}=\frac{1}{2}\left(\partial_kg_{ji}+\partial_ig_{kj}-\partial_jg_{ki}\right),
\qquad
g_{il}\Gamma^l_{kj}=\frac{1}{2}\left(\partial_kg_{ij}+\partial_jg_{ki}-\partial_ig_{kj}\right),$$
and adding these two identities yields
$$\partial_k g_{ij}=g_{lj}\Gamma^l_{ki}+g_{il}\Gamma^l_{kj}.$$
Since both sides of the identity $d(g(X,Y))=g(\nabla X,Y)+g(X,\nabla
Y)$ are tensorial in the direction of differentiation and satisfy
the same Leibniz expansion in $X$ and $Y$, the displayed identity
for coordinate fields implies the identity for all vector fields.
Thus on each chart there exists a connection with the two required
properties. Finally, by the uniqueness statement already proved
(applied on the intersection of two charts, where both locally
defined connections satisfy the two conditions), the locally defined
connections agree on overlaps and hence patch together to a globally
defined connection $\nabla$ on $TM$.
\end{proof}

 We call the
above connection the {\sl Levi-Civita connection} of the metric and $\nabla X$ the {\sl covariant
derivative} of $X$. On a Riemannian manifold we always use the
Levi-Civita connection.

In local coordinates $(x^1,\ldots,x^n)$ the Levi-Civita connection
$\nabla$ is given by the $\nabla_{\partial_i}(\partial_j) =
\Gamma^{k}_{ij}\partial_k$, where the Christoffel symbols
$\Gamma^k_{ij}$ are the smooth functions
\begin{equation}\label{Gamma}
\Gamma^{k}_{ij}= \frac{1}{2}g^{kl}(\partial_i g_{l j} +\partial_j
g_{il}-\partial_l g_{ij}). \end{equation}
 Note that the above two
additional conditions for the Levi-Civita connection $\nabla$
correspond respectively to

\smallskip

$\bullet$ $\Gamma^k_{ij}=\Gamma^k_{ji}$,

$\bullet$ $\partial_k g_{ij}=g_{l j}\Gamma^l_{ki}+g_{il
}\Gamma^l_{kj}$.

The covariant derivative extends to all tensors. In the special case
of a function $f$ we have $\nabla( f)=df$. Note that there is a
possible confusion between this and  the notation in the literature
since one often sees $\nabla f$ written for the gradient of $f$,
which is the vector field dual to $df$. We always use $\nabla f$ to
mean $df$, and we will denote the gradient of $f$ by $(\nabla f)^*$,

The covariant derivative allows us to define the Hessian of a smooth
function at any point, not just a critical point. Let $f$ be a
smooth real-valued function on $M$.  We define the
Hessian  of $f$, denoted $\Hess(f)$, as
follows:
\begin{equation}\label{Hessian}
\Hess(f)(X,Y)=X(Y(f))-\nabla_XY(f).\end{equation}
\begin{lemma}[Hessian is symmetric]
\label{lem:hessian-symmetric}
\label{Hessformula}
\uses{thm:levi-civita-connection}
\lean{MorganTianLib.hessian,MorganTianLib.hessian_symm,MorganTianLib.hessian_smul,MorganTianLib.hessianAt_symm,MorganTianLib.hessian_eq_metricInner_cov_gradientField,MorganTianLib.hessianAt_chartBasisVecFiber}
\leanok
 The Hessian is a contravariant, symmetric two-tensor, i.e., for
 vector fields $X$ and $Y$ we have
$$\Hess(f)(X,Y)=\Hess(f)(Y,X)$$
and
$$\Hess(f)(\phi X,\psi Y)=\phi\psi \Hess(f)(X,Y)$$
for all smooth functions $\phi,\psi$. Other formulas for the Hessian
are
$$\Hess(f)(X,Y)=\langle \nabla_X(\nabla f),Y\rangle=\nabla_X(\nabla_Y(f))=\nabla^2f(X,Y).$$
Also, in local coordinates  we have
$$\Hess(f)_{ij}=\partial_i\partial_jf-(\partial_kf)\Gamma^k_{ij}.$$
\end{lemma}

\begin{proof}
\leanok
\uses{thm:levi-civita-connection}
The proof of symmetry is direct from the torsion-free assumption:
$$\Hess(f)(X,Y)-\Hess(f)(Y,X)=[X,Y](f)-(\nabla_XY-\nabla_YX)(f)=0.$$
 The fact that
$\Hess(f)$ is a tensor is also established by direct
computation. The equivalence of the various formulas is also
immediate:
\begin{eqnarray}\label{hessform}
\langle\nabla_X(\nabla f),Y\rangle & = & X(\langle \nabla
f,Y\rangle)-\langle\nabla f,\nabla_XY\rangle
\\ & = & X(Y(f))-\nabla_XY(f)=\Hess(f)(X,Y).\nonumber
\end{eqnarray}
Since $df=(\partial_rf) dx^r$ and
$\nabla(dx^k)=-\Gamma_{ij}^kdx^i\otimes dx^j$, it follows that
$$\nabla(df)=\left(\partial_i\partial_jf-(\partial_kf)\Gamma^k_{ij}\right)dx^i\otimes
dx^j.$$ It is direct from the definition that
$$\Hess(f)_{ij}=\Hess(f)(\partial_i,\partial_j)=\partial_i\partial_jf-(\partial_kf)\Gamma_{ij}^k.$$
\end{proof}

When the metric that we are using to define the Hessian is not clear
from the context, we introduce it into the notation and write $\Hess_g(f)$ to denote the Hessian of $f$ with respect to the metric
$g$.


The Laplacian $\triangle f$ is defined as the
trace of the Hessian: That is to say, in local coordinates near $p$
we have
$$\triangle f(p)=\sum_{ij}g^{ij}\Hess(f)(\partial_i,\partial_j).$$
Thus, if $\{X_i\}$ is an orthonormal basis for $T_pM$ then
\begin{equation}\label{laplacformula}
\triangle f(p)=\sum_i\Hess(f)(X_i,X_i).\end{equation} Notice that this is
the form of the Laplacian that is non-negative at a local minimum, and
consequently has a non-positive spectrum.

\section{Curvature of a Riemannian manifold}

For the rest of this chapter $(M,g)$ is a Riemannian manifold.

\begin{definition}[Riemann curvature tensor]
\label{def:riemann-curvature-tensor}
\uses{thm:levi-civita-connection}
\lean{MorganTianLib.riemannCurvature,MorganTianLib.curvatureForm}
\leanok
The \textit{Riemann curvature tensor} of $M$ is the  $(1,3)$-tensor on $M$
\[ {\mathcal R}(X,Y)Z = \nabla^{2}_{X,Y}Z
- \nabla^{2}_{Y,X}Z = \nabla_{X}\nabla_{Y}Z-\nabla_{Y}\nabla_{X}Z
-\nabla_{[X,Y]}Z, \] where
$\nabla^2_{X,Y}Z=\nabla_{X}\nabla_{Y}Z-\nabla_{\nabla_XY}Z$.
\end{definition}

In local coordinates the curvature tensor can be represented as \[
{\mathcal R}(\partial_i,\partial_j)\partial_k=
{{{R}_{ij}}^l}_k\partial_l, \] where
$${{R_{ij}}^l}_k=
\partial_i \Gamma^{l}_{jk}- \partial_j
\Gamma^{l}_{ik}+
\Gamma^{s}_{jk}\Gamma^{l}_{is}-\Gamma^{s}_{ik}\Gamma^{l}_{js}.$$
Using the metric tensor $g$, we can change ${\mathcal R}$ to a
$(0,4)$-tensor as follows:
\[ {\mathcal R}(X,Y,Z,W)=g({\mathcal R}(X,Y)W,Z). \]
(Notice the change of order in the last two variables.) Notice that
we use the same symbol and the same name for both the $(1,3)$ tensor
and the $(0,4)$ tensor; which one we are dealing with in a given
context is indicated by the index structure or the variables to
which the tensor is applied.  In local coordinates, the Riemann
curvature tensor can be represented as
\begin{align*} {\mathcal R}(\partial_i,\partial_j,\partial_k,\partial_l) & = R_{ijkl}\\
&=g_{ks}{{R_{ij}}^s}_l\\
&=g_{ks}(\partial_i \Gamma^{s}_{jl}- \partial_j \Gamma^{s}_{il}+
\Gamma^{t}_{jl}\Gamma^{s}_{it}-\Gamma^{t}_{il}\Gamma^{s}_{jt}).
\end{align*}


One can easily verify the following:

\begin{lemma}[Curvature symmetries and Bianchi identities]
\label{claim:curvature-symmetries-bianchi}
\label{Bianchi}
\uses{def:riemann-curvature-tensor}
\lean{MorganTianLib.curvatureForm_symmetries,MorganTianLib.curvatureForm_secondBianchi}
\leanok
The Riemann curvature tensor ${\mathcal R}$ satisfies the following
properties:

$\bullet$ (Symmetry) $R_{ijkl}=-R_{jikl}$, $R_{ijkl}=-R_{ijlk}$,
$R_{ijkl}=R_{klij}$,

$\bullet$ (1st Bianchi identity) The sum of
$R_{ijkl}$ over the cyclic permutation of

\quad any three indices vanishes,

$\bullet$ (2nd Bianchi identity)
$R_{ijkl,h}+R_{ijlh,k}+R_{ijhk,l}=0$, where
$$R_{ijkl,h}=(\nabla_{\partial_h}{\mathcal R})_{ijkl}.$$
\end{lemma}

\begin{proof}
\uses{def:riemann-curvature-tensor,thm:levi-civita-connection}
\leanok
Throughout the proof we write $\langle\cdot,\cdot\rangle$ for $g$
and $\nabla_i$ for $\nabla_{\partial_i}$, and we use repeatedly that
all the expressions involved are tensorial, so that identities need
only be checked on coordinate vector fields, for which
$[\partial_a,\partial_b]=0$ and hence ${\mathcal
R}(\partial_i,\partial_j)=\nabla_i\nabla_j-\nabla_j\nabla_i
=[\nabla_i,\nabla_j]$ as operators on vector fields, and for which
torsion-freeness reads $\nabla_a\partial_b=\nabla_b\partial_a$.

\emph{Step 1: $R_{ijkl}=-R_{jikl}$.} The defining formula for
${\mathcal R}(X,Y)Z$ is antisymmetric under interchange of $X$ and
$Y$, so ${\mathcal R}(\partial_i,\partial_j)=-{\mathcal
R}(\partial_j,\partial_i)$ and therefore
$R_{ijkl}=\langle{\mathcal
R}(\partial_i,\partial_j)\partial_l,\partial_k\rangle
=-\langle{\mathcal
R}(\partial_j,\partial_i)\partial_l,\partial_k\rangle=-R_{jikl}$.

\emph{Step 2: $R_{ijkl}=-R_{ijlk}$.} It suffices to show that
$\langle {\mathcal R}(X,Y)Z,Z\rangle=0$ for all $X,Y,Z$; the stated
antisymmetry then follows by replacing $Z$ by $Z+W$ and expanding.
Take $X=\partial_i$, $Y=\partial_j$. Metric compatibility gives
$$\langle\nabla_i\nabla_jZ,Z\rangle
=\partial_i\langle\nabla_jZ,Z\rangle-\langle\nabla_jZ,\nabla_iZ\rangle
=\tfrac{1}{2}\partial_i\partial_j\langle
Z,Z\rangle-\langle\nabla_jZ,\nabla_iZ\rangle,$$ and the same formula
with $i$ and $j$ exchanged. Subtracting, the two right-hand sides
agree (second partial derivatives of functions commute), so
$\langle[\nabla_i,\nabla_j]Z,Z\rangle=0$, which is the assertion.

\emph{Step 3: first Bianchi identity for the cyclic triple
$(i,j,l)$.} We claim that
\begin{equation}\label{firstBianchi}
R_{ijkl}+R_{jlki}+R_{likj}=0.
\end{equation}
By tensoriality it is enough to prove the operator identity
${\mathcal R}(\partial_i,\partial_j)\partial_l+{\mathcal
R}(\partial_j,\partial_l)\partial_i+{\mathcal
R}(\partial_l,\partial_i)\partial_j=0$; pairing with $\partial_k$
then gives Equation~(\ref{firstBianchi}). Expanding the three
curvature operators and regrouping the six terms in pairs,
$$\sum_{\mathrm{cyclic}}\left(\nabla_i\nabla_j\partial_l
-\nabla_j\nabla_i\partial_l\right)
=\nabla_i\left(\nabla_j\partial_l-\nabla_l\partial_j\right)
+\nabla_j\left(\nabla_l\partial_i-\nabla_i\partial_l\right)
+\nabla_l\left(\nabla_i\partial_j-\nabla_j\partial_i\right),$$
and each parenthesized difference vanishes by torsion-freeness.

\emph{Step 4: $R_{ijkl}=R_{klij}$.} Write down
Equation~(\ref{firstBianchi}) four times, cyclically permuting the
roles of the four indices, i.e., for the (cyclic triple; fixed
third-slot index) choices $(i,j,l;k)$, $(j,l,k;i)$, $(l,k,i;j)$ and
$(k,i,j;l)$:
\begin{eqnarray*}
R_{ijkl}+R_{jlki}+R_{likj} & = & 0\\
R_{jlik}+R_{lkij}+R_{kjil} & = & 0\\
R_{lkji}+R_{kijl}+R_{iljk} & = & 0\\
R_{kilj}+R_{ijlk}+R_{jkli} & = & 0.
\end{eqnarray*}
Add all four identities. By Step 2 the pairs
$\{R_{ijkl},R_{ijlk}\}$, $\{R_{jlki},R_{jlik}\}$,
$\{R_{lkij},R_{lkji}\}$ and $\{R_{kijl},R_{kilj}\}$ cancel, leaving
$$R_{likj}+R_{kjil}+R_{iljk}+R_{jkli}=0.$$
By Steps 1 and 2, $R_{iljk}=-R_{lijk}=R_{likj}$ and
$R_{jkli}=-R_{kjli}=R_{kjil}$, so this reads
$2R_{likj}+2R_{kjil}=0$, i.e., $R_{likj}=-R_{kjil}$. Applying Step 2
to the right-hand side gives $R_{likj}=R_{kjli}$, which is the pair
symmetry $R_{abcd}=R_{cdab}$ for $(a,b,c,d)=(l,i,k,j)$; since the
indices are arbitrary, $R_{ijkl}=R_{klij}$.

\emph{Step 5: the cyclic sum over any three indices vanishes.}
Equation~(\ref{firstBianchi}) is the assertion for the three index
positions $1,2,4$ (holding the third-position index fixed). Applying
the antisymmetry of Step 2 to each term of
Equation~(\ref{firstBianchi}) (with $k$ and $l$ interchanged)
turns it into the assertion for positions $1,2,3$; applying the pair
symmetry of Step 4 to each term of these two identities turns them
into the assertions for positions $3,4,2$ and $3,4,1$ respectively.
This exhausts all four three-element subsets of the index positions.

\emph{Step 6: second Bianchi identity.} Since $\nabla g=0$,
covariant differentiation commutes with raising and lowering
indices, so
$R_{ijkl,h}=\langle(\nabla_h{\mathcal
R})(\partial_i,\partial_j)\partial_l,\partial_k\rangle$, where
$\nabla_h {\mathcal R}$ is the covariant derivative of the
$(1,3)$-tensor ${\mathcal R}$; moreover the symmetries of Steps 1,
2 and 4 are preserved by covariant differentiation (differentiate
each identity and use $\nabla g=0$). By the Leibniz rule defining
$\nabla_h{\mathcal R}$, as operators on vector fields,
$$(\nabla_h{\mathcal R})(\partial_i,\partial_j)
=\nabla_h\circ{\mathcal R}(\partial_i,\partial_j)
-{\mathcal R}(\partial_i,\partial_j)\circ\nabla_h
-{\mathcal R}(\nabla_h\partial_i,\partial_j)
-{\mathcal R}(\partial_i,\nabla_h\partial_j),$$
that is, writing ${\mathcal R}_{ij}={\mathcal
R}(\partial_i,\partial_j)=[\nabla_i,\nabla_j]$ and
$\nabla_h\partial_i=\Gamma^s_{hi}\partial_s$,
$$(\nabla_h{\mathcal R})(\partial_i,\partial_j)
=[\nabla_h,{\mathcal R}_{ij}]
-\Gamma^s_{hi}{\mathcal R}_{sj}-\Gamma^s_{hj}{\mathcal R}_{is}.$$
Now sum this identity over the cyclic permutations of $(h,i,j)$. The
first terms contribute
$$[\nabla_h,[\nabla_i,\nabla_j]]+[\nabla_i,[\nabla_j,\nabla_h]]
+[\nabla_j,[\nabla_h,\nabla_i]]=0$$
by the Jacobi identity for commutators of linear operators. The
remaining terms contribute
$$-\Gamma^s_{hi}{\mathcal R}_{sj}-\Gamma^s_{hj}{\mathcal R}_{is}
-\Gamma^s_{ij}{\mathcal R}_{sh}-\Gamma^s_{ih}{\mathcal R}_{js}
-\Gamma^s_{jh}{\mathcal R}_{si}-\Gamma^s_{ji}{\mathcal R}_{hs},$$
which vanishes in pairs by the symmetry
$\Gamma^s_{ab}=\Gamma^s_{ba}$ and the antisymmetry ${\mathcal
R}_{ab}=-{\mathcal R}_{ba}$: the first and fourth, second and fifth,
and third and sixth terms cancel in pairs. Hence
$$(\nabla_h{\mathcal R})(\partial_i,\partial_j)
+(\nabla_i{\mathcal R})(\partial_j,\partial_h)
+(\nabla_j{\mathcal R})(\partial_h,\partial_i)=0,$$
and pairing with the metric,
$$R_{ijkl,h}+R_{jhkl,i}+R_{hikl,j}=0$$
for all indices, i.e., the cyclic sum over (position 1, position 2,
differentiation index) vanishes. Finally, applying the pair symmetry
$R_{abcd}=R_{cdab}$ (valid for $\nabla{\mathcal R}$ as noted above)
to each term of this identity with $(a,b,c)=(k,l,h)$ gives
$$R_{ijkl,h}+R_{ijlh,k}+R_{ijhk,l}
=R_{klij,h}+R_{lhij,k}+R_{hkij,l}=0,$$
which is the second Bianchi identity as stated.
\end{proof}


There are many important related curvatures.


\begin{definition}[Sectional curvature]
\label{def:sectional-curvature}
\uses{def:riemann-curvature-tensor}
\lean{MorganTianLib.sectionalCurvatureAt}
\leanok
The \textit{sectional curvature} of a
2-plane $P\subset T_pM$ is defined as
\[ K(P) = {\mathcal R}(X,Y,X,Y), \]
where $\{X,Y\}$ is an orthonormal basis of $P$. We say that $(M,g)$
has \textit{positive sectional curvature} (resp., \textit{negative
sectional curvature}) if $K(P)>0$ (resp., $K(P)<0$) for every
2-plane $P$.  There are analogous notions of non-negative and
non-positive sectional curvature.
\end{definition}


In local coordinates, suppose that $X = X^{i}\partial_i$ and $Y =
Y^{i}\partial_i$. Then we have
\[ K(P) = R_{ijkl}X^{i}Y^{j}X^{k}Y^{l}.\]
A Riemannian manifold is said to have {\sl constant sectional
curvature} if $K(P)$ is the same for all $p\in M$ and all $2$-planes
$P\subset T_pM$.

\begin{lemma}[Curvature tensor of constant curvature manifolds]
\label{lem:constant-curvature-tensor}
\uses{def:sectional-curvature}
\lean{MorganTianLib.hasConstantSectionalCurvature_iff_curvatureFormAt_eq}
\leanok
A manifold $(M,g)$ has constant
sectional curvature $\lambda$ if and only if \[ R_{ijkl} =
\lambda(g_{ik}g_{jl}-g_{il}g_{jk}), \]
i.e., if and only if ${\mathcal R}(X,Y,Z,W)
=\lambda\left(\langle X,Z\rangle\langle Y,W\rangle
-\langle X,W\rangle\langle Y,Z\rangle\right)$ for all vector fields
$X,Y,Z,W$.
\end{lemma}

\begin{proof}
\uses{def:sectional-curvature,claim:curvature-symmetries-bianchi}
\leanok
If the displayed identity holds then for any orthonormal pair $X,Y$
spanning a $2$-plane $P$ we get $K(P)={\mathcal
R}(X,Y,X,Y)=\lambda(1\cdot 1-0)=\lambda$, so the sectional curvature
is constant.

Conversely, suppose $K(P)=\lambda$ for every $2$-plane at every
point. Define the $(0,4)$-tensor
$$T(X,Y,Z,W)={\mathcal R}(X,Y,Z,W)
-\lambda\left(\langle X,Z\rangle\langle Y,W\rangle
-\langle X,W\rangle\langle Y,Z\rangle\right).$$
One checks directly that the subtracted tensor enjoys all the
symmetries of ${\mathcal R}$ listed in
Claim~\ref{claim:curvature-symmetries-bianchi} --- antisymmetry in
the first two and in the last two arguments, symmetry under
interchange of the two pairs, and the first Bianchi identity (for
the latter, the cyclic sum of
$g_{ik}g_{jl}-g_{il}g_{jk}$ over the permutations of $(i,j,l)$
telescopes to zero) --- and hence so does $T$.

We first claim $T(X,Y,X,Y)=0$ for \emph{all} $X,Y$. If $X,Y$ are
orthonormal this is the hypothesis $K=\lambda$. If $X,Y$ are
linearly dependent it holds by antisymmetry. In general, replacing
$Y$ by $Y+cX$ leaves $T(X,Y,X,Y)$ unchanged (expand and use
antisymmetry in the last two, respectively first two, arguments),
and $T(aX,bY,aX,bY)=a^2b^2\,T(X,Y,X,Y)$; using these operations any
linearly independent pair can be reduced to an orthonormal one.

Now polarize twice. Replacing $X$ by $X+Z$ in $T(X,Y,X,Y)=0$ and
expanding gives $T(X,Y,Z,Y)+T(Z,Y,X,Y)=0$, and by the pair symmetry
the two terms are equal; hence
$$T(X,Y,Z,Y)=0\qquad\text{for all }X,Y,Z.$$
Replacing $Y$ by $Y+W$ in this identity and expanding gives
$$T(X,Y,Z,W)=-T(X,W,Z,Y),$$
i.e., $T$ is antisymmetric under interchange of its second and
fourth arguments. Combining this with the first Bianchi identity in
the form $T(X,Y,Z,W)+T(Y,W,Z,X)+T(W,X,Z,Y)=0$ (cyclic in the first,
second and fourth arguments): by the new antisymmetry and the
antisymmetry in the first two arguments,
$$T(Y,W,Z,X)=-T(Y,X,Z,W)=T(X,Y,Z,W),$$
and likewise $T(W,X,Z,Y)=-T(W,Y,Z,X)=T(Y,W,Z,X)=T(X,Y,Z,W)$. The
Bianchi identity thus reads $3\,T(X,Y,Z,W)=0$, so $T\equiv 0$, which
is the displayed identity.
\end{proof}

Of course, the sphere of
radius $r$ in $\Ar^{n+1}$ has constant sectional curvature $1/r^2$,
$\Ar^n$ with the Euclidean metric has constant sectional curvature
$0$, and the hyperbolic space $\HH^n$, which, in the
Poincar\'e model, is given by the unit disk with the metric
$$\frac{4(dx^{2}_{1}+\cdots+
dx^{2}_{n})}{(1-\abs{x}^{2})^{2}},$$ or in the upper half-space
model with coordinates $(x^1,\ldots,x^n)$ is given by
$$\frac{ds^2}{(x^n)^2}$$ has constant sectional curvature $-1$. In all
three cases we denote the constant curvature metric by $g_\mathrm{st}$.
These properties of the three model spaces --- constant curvature,
completeness and simple connectivity --- are proved in
Lemma~\ref{lem:model-spaces} below.


\begin{definition}[Curvature operator]
\label{def:curvature-operator}
\uses{def:riemann-curvature-tensor,def:sectional-curvature}
\lean{MorganTianLib.curvatureOperator,MorganTianLib.curvatureOperator_ιMulti,MorganTianLib.curvatureOperator_symm,MorganTianLib.sectionalCurvature_eq_curvatureOperator,MorganTianLib.HasPositiveCurvatureOperator,MorganTianLib.HasNonnegativeCurvatureOperator,MorganTianLib.sectionalCurvature_nonneg_of_hasNonnegativeCurvatureOperator}
\leanok
Using the metric, one can replace the Riemann curvature tensor
${\mathcal R}$ by  a symmetric bilinear form $\Rm$ on
$\wedge^2TM$. In local coordinates let
$\varphi=\varphi^{ij}\partial_i\wedge\partial _j$ and
$\psi=\psi^{kl}\partial _k\wedge\partial_l$ be local sections of
$\wedge^2TM$. The formula for $\Rm$ is
\[ \Rm(\varphi, \psi)=R_{ijkl}\varphi^{ij}\psi^{kl}. \]
We call $\Rm$ the \textit{curvature
operator}. We say $(M,g)$ has
\textit{positive curvature operator} if $\Rm(\varphi,\varphi)>0$ for any nonzero 2-form $\varphi =
\varphi^{ij}\partial_{i}\wedge
\partial_{j}$ and has \textit{nonnegative curvature operator} if
$\Rm(\varphi,\varphi)\geq0$ for any  $\varphi\in \wedge^2TM$.
\end{definition}

\begin{definition}[Norm of the curvature operator]
\label{def:curvature-operator-norm}
\uses{def:curvature-operator,def:sectional-curvature}
\lean{MorganTianLib.wedgeInner,MorganTianLib.HasCurvatureOperatorNormLe,MorganTianLib.abs_curvatureForm_le_of_hasCurvatureOperatorNormLe,MorganTianLib.exists_orthonormal_changeBasis,MorganTianLib.abs_sectionalCurvature_le_of_hasCurvatureOperatorNormLe}
\leanok
Endow each fiber $\wedge^2T_xM$ with the inner product determined
by
$$\langle X\wedge Y,Z\wedge W\rangle
=g(X,Z)g(Y,W)-g(X,W)g(Y,Z),$$
so that $X\wedge Y$ has unit length whenever $X,Y$ are orthonormal;
if $\{e_1,\ldots,e_n\}$ is an orthonormal basis of $T_xM$ then
$\{e_i\wedge e_j\}_{i<j}$ is an orthonormal basis of
$\wedge^2T_xM$. With this normalization, for every $2$-plane
$P\subset T_xM$ with orthonormal basis $\{X,Y\}$ the antisymmetries
of $R_{ijkl}$ give
$$\Rm(X\wedge Y,X\wedge Y)={\mathcal R}(X,Y,X,Y)=K(P).$$
For $K\ge0$ we write $|\Rm(x)|\le K$ to mean that the eigenvalues
of the symmetric operator on $\wedge^2T_xM$ associated to the
bilinear form $\Rm$ by this inner product lie in $[-K,K]$;
equivalently, since that operator is self-adjoint, that
$|\Rm(\varphi,\varphi)|\le K\,\langle\varphi,\varphi\rangle$ for every
$\varphi\in\wedge^2T_xM$ (the Rayleigh-quotient form, which is the one
formalized). In
particular $|\Rm(x)|\le K$ implies $|K(P)|\le K$ for every
$2$-plane $P\subset T_xM$.
\end{definition}

Clearly, if the curvature operator is a positive (resp.,
non-negative) operator then the manifold is positively (resp.,
non-negatively) curved.


\begin{definition}[Ricci curvature and scalar curvature]
\label{def:ricci-curvature}
\uses{def:riemann-curvature-tensor}
\lean{MorganTianLib.ricciAt,MorganTianLib.scalarCurvatureAt}
\leanok
The \textit{Ricci curvature tensor},
denoted $\Ric$ or $\Ric_g$ when it is necessary to specify
the metric, is a symmetric contravariant two-tensor. In local
coordinates it is defined by
\[ \Ric(X,Y)=g^{kl}R(X,\partial_k,Y,\partial_l). \]
The value of this tensor at a point $p\in M$ is  given by
$\sum_{i=1}^nR(X(p),e_i,Y(p),e_i)$ where $\{e_{1},\cdots,e_{n}\}$ is
an orthonormal basis of $T_pM$. Clearly $\Ric$ is a symmetric
bilinear form on $TM$, given in local coordinates by
$$\Ric=\Ric_{ij}dx^i\otimes dx^j,$$ where
$\Ric_{ij}=\Ric(\partial_i,\partial _j)$. The
\textit{scalar curvature} is defined by:
$$R=R_g=\tr_g \Ric=g^{ij}\Ric_{ij}.$$ We will say that $\Ric \geq
k$ (or $\leq k$) if all the eigenvalues of $\Ric$ are $\geq k$
(or $\leq k$).
\end{definition}

\begin{lemma}[Ricci curvature as a trace of the curvature operator]
\label{lem:ricci-trace-curvature-operator}
\uses{def:ricci-curvature,def:riemann-curvature-tensor,claim:curvature-symmetries-bianchi}
\lean{MorganTianLib.metricInner_riemannCurvature_eq_curvatureFormAt,MorganTianLib.sum_metricInner_riemannCurvature_self_eq_ricciAt}
\leanok
Let $(M,g)$ be a Riemannian manifold with Levi-Civita connection
$\nabla$, let $X$ be a smooth vector field, let $p\in M$, and let
$\{e_1,\ldots,e_n\}$ be an orthonormal basis of $T_pM$, each $e_i$
extended to a smooth vector field $E_i$. Then
$$\sum_{i=1}^n\left\langle{\mathcal R}(E_i,X)X,E_i\right\rangle(p)
=\Ric\left(X(p),X(p)\right),$$
where ${\mathcal R}$ is the curvature operator of
Definition~\ref{def:riemann-curvature-tensor}.
\end{lemma}

\begin{proof}
\uses{def:ricci-curvature,claim:curvature-symmetries-bianchi}
\leanok
Since the $(0,4)$ curvature tensor is defined by
$R(X,Y,Z,W)=\langle{\mathcal R}(X,Y)W,Z\rangle$ and its value at $p$
depends only on the values of its arguments at $p$ (tensoriality,
Definition~\ref{def:riemann-curvature-tensor}), each summand is
$\langle{\mathcal R}(E_i,X)X,E_i\rangle(p)=R(e_i,X(p),e_i,X(p))$. By
the antisymmetry of $R$ in its first pair and in its second pair of
arguments (Claim~\ref{claim:curvature-symmetries-bianchi}),
$$R(e_i,X(p),e_i,X(p))=(-1)^2R(X(p),e_i,X(p),e_i)
=R(X(p),e_i,X(p),e_i),$$
and summing over $i$ gives
$\sum_iR(X(p),e_i,X(p),e_i)=\Ric(X(p),X(p))$ by the
orthonormal-basis formula of Definition~\ref{def:ricci-curvature}.
\end{proof}

Clearly, the curvatures are natural in the sense that if $F\colon N
\rightarrow M$ is a diffeomorphism and if $g$ is a Riemannian metric
on $M$, then $F^*g$ is a Riemannian metric on $N$ and we have $\Rm(F^*g)=F^*(\Rm(g))$, $\Ric(F^{*}g) = F^{*}(\Ric(g))$, and $R(F^*g)=F^*(R(g))$.

\subsection{Consequences of the Bianchi identities}

There is one consequence of the second Bianchi identity that will be
important later. For any contravariant two-tensor $\omega$ on $M$
(such as $\Ric$ or $\Hess(f)$) we define the contravariant
one-tensor $\operatorname{div}(\omega)$ as follows: For any vector field $X$
we set
$$\operatorname{div}(\omega)(X)=\nabla^*\omega(X)=g^{rs}\nabla_r(\omega)(X,\partial_s).$$

\begin{lemma}[Divergence of Ricci from second Bianchi identity]
\label{lem:div-ricci-scalar-curvature}
\label{divRic}
\uses{claim:curvature-symmetries-bianchi,def:ricci-curvature}
\lean{MorganTianLib.dir_scalarCurvature_eq_two_divRicci,MorganTianLib.divRicciAt,MorganTianLib.covRicciAt,MorganTianLib.divRicciAt_eq_sum_orthonormalBasis,MorganTianLib.covRicci_eq_frame_sum,MorganTianLib.dir_scalarCurvatureAt_eq_two_frame_covRicci,MorganTianLib.covRicci_congr_apply}
\leanok
$$dR=2\operatorname{div}(\Ric)=2\nabla^*\Ric.$$
\end{lemma}

\begin{proof}
\leanok
\uses{claim:curvature-symmetries-bianchi,def:ricci-curvature}
Since $\nabla g=0$, contraction with $g$ (in either index position)
commutes with covariant differentiation; we use this freely. In
local coordinates the definitions read
$\Ric_{ab}=g^{cd}R_{acbd}$, $R=g^{ab}\Ric_{ab}=g^{ab}g^{cd}R_{acbd}$
and $\operatorname{div}(\Ric)(\partial_h)=g^{rs}\Ric_{hs,r}$.
(In the formalization the contractions are computed against a smooth
local {\em orthonormal} frame $\{E_i\}$, so $g^{ik}=\delta^{ik}$ in the
computation below; the frame-derivative corrections that implement
"contraction commutes with $\nabla$" cancel by the antisymmetry
$\langle\nabla_U E_j,E_k\rangle=-\langle\nabla_U E_k,E_j\rangle$ of the
connection coefficients, which is exactly $\nabla g=0$.)

Start from the second Bianchi identity
(Claim~\ref{claim:curvature-symmetries-bianchi})
$$R_{ijkl,h}+R_{ijlh,k}+R_{ijhk,l}=0$$
and contract it with $g^{ik}g^{jl}$, treating the three terms in
turn. Each symmetry of ${\mathcal R}$ in
Claim~\ref{claim:curvature-symmetries-bianchi} asserts that a fixed
permutation of the arguments of the $(0,4)$-tensor ${\mathcal R}$
changes it by at most a sign; covariant differentiation commutes
with such permutations of the arguments, so differentiating each
identity shows that $\nabla{\mathcal R}$ satisfies the same
symmetries in its first four indices:
$R_{ijkl,h}=-R_{jikl,h}$, $R_{ijkl,h}=-R_{ijlk,h}$ and
$R_{ijkl,h}=R_{klij,h}$.

For the first term, relabelling the contraction indices,
$$g^{ik}g^{jl}R_{ijkl,h}=\left(g^{ik}g^{jl}R_{ijkl}\right)_{,h}
=\left(g^{ab}g^{cd}R_{acbd}\right)_{,h}=R_{,h}=\partial_hR,$$
where the second equality is the relabelling
$(i,j,k,l)=(a,c,b,d)$, using the pair symmetry
$R_{acbd}=R_{bdac}$ to identify the two full contractions of the
curvature tensor with the scalar curvature.

For the second term, the antisymmetry in the last two indices gives
$R_{ijlh}=-R_{ijhl}$, and $g^{jl}R_{ijhl}=\Ric_{ih}$ directly from
the definition of $\Ric$. Hence
$$g^{ik}g^{jl}R_{ijlh,k}=-g^{ik}\Ric_{ih,k}
=-g^{ik}\Ric_{hi,k}=-\operatorname{div}(\Ric)(\partial_h),$$
using the symmetry of $\Ric$.

For the third term, the pair symmetry followed by the antisymmetry
in the last two indices gives $R_{ijhk}=R_{hkij}=-R_{hkji}$, and
$g^{ik}R_{hkji}=\Ric_{hj}$ from the definition of $\Ric$ (with
contraction indices $k,i$). Hence
$$g^{ik}g^{jl}R_{ijhk,l}=-g^{jl}\Ric_{hj,l}
=-\operatorname{div}(\Ric)(\partial_h).$$

Adding the three terms, the contracted second Bianchi identity
becomes
$$\partial_hR-2\operatorname{div}(\Ric)(\partial_h)=0$$
for every $h$, which is exactly $dR=2\operatorname{div}(\Ric)$.
\end{proof}


We shall also need a formula relating the connection Laplacian on
contravariant one-tensors with the Ricci curvature. Recall that for
a smooth function $f$, we defined the symmetric two-tensor
$\nabla^2f$ by
$$\nabla^2f(X,Y)=\nabla_X\nabla_Y(f)-\nabla_{\nabla_X(Y)}(f)=\Hess(f)(X,Y),$$
and then defined the Laplacian
$$\triangle f=\tr\nabla^2f=g^{ij}(\nabla^2f)_{ij}.$$ These operators extend to tensors of any
rank. Suppose that $\omega$ is a contravariant tensor of rank $k$. Then we
define $\nabla^2\omega$ to be a contravariant tensor of rank $k+2$ given by
$$\nabla^2\omega(\cdot,X,Y)=(\nabla_{X}\nabla_{Y}\omega)(\cdot)
-\nabla_{\nabla_X(Y)}\omega(\cdot).$$ This expression is not
symmetric in the vector fields $X,Y$ but the commutator is given by
evaluating the curvature operator ${\mathcal R}(X,Y)$ on $\omega$.
We define the connection Laplacian on the tensor $\omega$ to be
$$\triangle \omega=g^{ij}\nabla^2(\omega)(\partial_i,\partial_j).$$
Direct computation gives the standard Bochner formula relating these
Laplacians with the Ricci curvature; see for example Proposition
4.36 on page 168 of \cite{Gallot}.

We first record the case of vector fields, from which the tensor
cases follow by duality.

\begin{definition}[Second covariant derivative of a vector field]
\label{def:second-covariant-derivative-vector-field}
\uses{thm:levi-civita-connection}
\lean{MorganTianLib.secondCov}
\leanok
Let $M$ be a smooth manifold equipped with an affine connection
$\nabla$, and let $X$, $Y$, $Z$ be smooth vector fields on $M$. The
{\it second covariant derivative} of $Z$ in the directions $X,Y$ is
the smooth vector field
$$\nabla^2Z(X,Y)=\nabla_X\left(\nabla_YZ\right)-\nabla_{\nabla_X(Y)}Z.$$
\end{definition}

\begin{lemma}[Tensoriality of the second covariant derivative]
\label{lem:second-covariant-derivative-tensorial}
\uses{def:second-covariant-derivative-vector-field}
\lean{MorganTianLib.secondCov_add_left,MorganTianLib.secondCov_smul_left,MorganTianLib.secondCov_add_middle,MorganTianLib.secondCov_smul_middle,MorganTianLib.secondCov_congr_left}
\leanok
Let $\nabla$ be an affine connection on $M$ and let $Z$ be a smooth
vector field. The second covariant derivative $\nabla^2Z(X,Y)$ is
additive in each of $X$ and $Y$ separately, and for every smooth
function $\varphi$ on $M$,
$$\nabla^2Z(\varphi X,Y)=\varphi\,\nabla^2Z(X,Y)
=\nabla^2Z(X,\varphi Y).$$
In particular the value $\nabla^2Z(X,Y)(p)$ at a point $p$ depends
on the direction field $X$ only through its value $X(p)$.
\end{lemma}

\begin{proof}
\uses{thm:levi-civita-connection}
\leanok
Additivity in $X$ is the additivity of $\nabla_{(\cdot)}$ in its
direction slot, applied to both terms of the definition; additivity
in $Y$ follows from the additivity of $\nabla$ in both slots. For
homogeneity in $X$, both terms are $\mathcal{D}(M)$-homogeneous in
$X$: $\nabla_{\varphi X}(\nabla_YZ)=\varphi\,\nabla_X(\nabla_YZ)$
directly, and $\nabla_{\varphi X}(Y)=\varphi\,\nabla_X(Y)$, whence
$\nabla_{\nabla_{\varphi X}(Y)}Z
=\nabla_{\varphi\,\nabla_X(Y)}Z=\varphi\,\nabla_{\nabla_X(Y)}Z$. For
homogeneity in $Y$, the Leibniz rule gives
$$\nabla_X\left(\nabla_{\varphi Y}Z\right)
=\nabla_X\left(\varphi\,\nabla_YZ\right)
=\varphi\,\nabla_X\left(\nabla_YZ\right)+X(\varphi)\,\nabla_YZ,$$
while $\nabla_X(\varphi Y)=\varphi\,\nabla_X(Y)+X(\varphi)\,Y$ and
the $\mathcal{D}(M)$-linearity of the direction slot give
$$\nabla_{\nabla_X(\varphi Y)}Z
=\varphi\,\nabla_{\nabla_X(Y)}Z+X(\varphi)\,\nabla_YZ;$$
subtracting, the two Leibniz cross terms $X(\varphi)\,\nabla_YZ$
cancel, leaving $\varphi\,\nabla^2Z(X,Y)$. Finally, both terms of
$\nabla^2Z(X,Y)$ are covariant derivatives in the direction $X$ — of
$\nabla_YZ$ and, since $\nabla_X(Y)(p)$ itself depends on $X$ only
through $X(p)$, of $Z$ — and the value of $\nabla_XW$ at $p$ depends
only on $X(p)$ by the standard bump-function argument (write a
direction field vanishing at $p$ as a finite sum
$\sum_i\varphi_iW_i$ with $\varphi_i(p)=0$ near $p$ and use the
$\mathcal{D}(M)$-homogeneity just established).
\end{proof}

\begin{lemma}[Ricci commutation identity for vector fields]
\label{lem:ricci-commutation-vector-field}
\uses{def:second-covariant-derivative-vector-field,def:riemann-curvature-tensor}
\lean{MorganTianLib.secondCov_sub_swap}
\leanok
Let $\nabla$ be a torsion-free affine connection on $M$. For all
smooth vector fields $X$, $Y$, $Z$,
$$\nabla^2Z(X,Y)-\nabla^2Z(Y,X)={\mathcal R}(X,Y)Z,$$
where ${\mathcal R}(X,Y)Z
=\nabla_X\nabla_YZ-\nabla_Y\nabla_XZ-\nabla_{[X,Y]}Z$ is the
curvature operator of Definition~\ref{def:riemann-curvature-tensor}.
\end{lemma}

\begin{proof}
\uses{thm:levi-civita-connection,def:riemann-curvature-tensor}
\leanok
Subtracting the defining expressions,
$$\nabla^2Z(X,Y)-\nabla^2Z(Y,X)
=\nabla_X\nabla_YZ-\nabla_Y\nabla_XZ
-\left(\nabla_{\nabla_X(Y)}Z-\nabla_{\nabla_Y(X)}Z\right).$$
By additivity of the direction slot of $\nabla$, the parenthesized
difference is $\nabla_{\nabla_X(Y)-\nabla_Y(X)}Z$, and
torsion-freeness gives $\nabla_X(Y)-\nabla_Y(X)=[X,Y]$; so the
right-hand side is
$\nabla_X\nabla_YZ-\nabla_Y\nabla_XZ-\nabla_{[X,Y]}Z
={\mathcal R}(X,Y)Z$.
\end{proof}

\begin{lemma}[Laplacian of differential form]
\label{lem:laplacian-one-form}
\label{lapformula}
\uses{lem:hessian-symmetric,def:ricci-curvature}
\lean{MorganTianLib.oneFormLaplacianAt,MorganTianLib.oneFormLaplacianAt_eq_sum_frame,MorganTianLib.metricInner_secondCov_gradientField_symm,MorganTianLib.sum_metricInner_secondCov_along_frame_eq_dir_laplacianAt,MorganTianLib.sum_metricInner_riemannCurvature_frame_eq_ricciAt,MorganTianLib.oneFormLaplacianAt_gradientField_eq_dirTangent_laplacianAt_add_ricciAt,MorganTianLib.oneFormLaplacianAt_gradientField_eq_metricInner_gradientAt_add_ricciAt}
\leanok
 Let $f$ be a smooth function on a Riemannian manifold.
Then we have the following formula for contravariant one-tensors:
$$ \triangle df=d\triangle f+\Ric((\nabla f)^*,\cdot).$$
\end{lemma}

\begin{proof}
\uses{lem:hessian-symmetric,def:ricci-curvature,claim:curvature-symmetries-bianchi,def:second-covariant-derivative-vector-field,lem:ricci-commutation-vector-field}
\leanok
Write $H=\nabla(df)=\Hess(f)$, a symmetric two-tensor
(Lemma~\ref{lem:hessian-symmetric}), and write
$\nabla^2_{X,Y}=\nabla_X\nabla_Y-\nabla_{\nabla_XY}$ for the second
covariant derivative in the directions $X,Y$, as in the discussion
preceding the lemma. A direct expansion shows that for any one-form
$\omega$ and vector field $Z$,
$$\left(\nabla^2_{X,Y}\omega\right)(Z)=(\nabla_X(\nabla\omega))(Y,Z),$$
where $\nabla\omega$ is regarded as a two-tensor with
$(\nabla\omega)(Y,Z)=(\nabla_Y\omega)(Z)$; in particular
$\left(\nabla^2_{X,Y}df\right)(Z)=(\nabla_XH)(Y,Z)$. In these terms
the two sides of the asserted formula are, at a point, with
$\{\partial_i\}$ local coordinate fields,
$$\triangle(df)(Z)=g^{jk}\left(\nabla_{\partial_j}H\right)(\partial_k,Z),
\qquad
d(\triangle f)(Z)=Z\left(g^{jk}H_{jk}\right)
=g^{jk}\left(\nabla_ZH\right)(\partial_j,\partial_k),$$
the latter because $\nabla g=0$, so differentiating the trace
commutes with the contraction.

\emph{Commutation formula for one-forms.} We claim that for any
one-form $\omega$,
\begin{equation}\label{oneformcomm}
\left(\nabla^2_{X,Y}\omega-\nabla^2_{Y,X}\omega\right)(Z)
=-\omega\left({\mathcal R}(X,Y)Z\right).
\end{equation}
Indeed, apply the second covariant derivative to the function
$\omega(Z)$: the Leibniz rule gives
$$\nabla^2_{X,Y}\left(\omega(Z)\right)
=\left(\nabla^2_{X,Y}\omega\right)(Z)
+(\nabla_X\omega)(\nabla_YZ)+(\nabla_Y\omega)(\nabla_XZ)
+\omega\left(\nabla^2_{X,Y}Z\right).$$
Antisymmetrize in $X$ and $Y$. The left-hand side is the Hessian of
the function $\omega(Z)$, which is symmetric in $X,Y$ by
Lemma~\ref{lem:hessian-symmetric}; the two middle terms on the right
are symmetric as a pair; and
$\nabla^2_{X,Y}Z-\nabla^2_{Y,X}Z={\mathcal R}(X,Y)Z$ by the
definition of the curvature tensor. This yields
Equation~(\ref{oneformcomm}).

Applied to $\omega=df$, Equation~(\ref{oneformcomm}) says
\begin{equation}\label{Hcomm}
(\nabla_XH)(Y,Z)-(\nabla_YH)(X,Z)=-df\left({\mathcal
R}(X,Y)Z\right).
\end{equation}

\emph{Comparing the two sides.} Fix a coordinate direction
$\partial_i$. Using Equation~(\ref{Hcomm}) twice and the symmetry
of $H$ (hence of $\nabla_WH$) in its two arguments,
\begin{eqnarray*}
(\nabla_{\partial_j}H)(\partial_k,\partial_i)
-(\nabla_{\partial_i}H)(\partial_k,\partial_j)
& = & \left[(\nabla_{\partial_j}H)(\partial_k,\partial_i)
-(\nabla_{\partial_k}H)(\partial_j,\partial_i)\right]
+\left[(\nabla_{\partial_k}H)(\partial_i,\partial_j)
-(\nabla_{\partial_i}H)(\partial_k,\partial_j)\right]\\
& = & -df\left({\mathcal R}(\partial_j,\partial_k)\partial_i\right)
-df\left({\mathcal R}(\partial_k,\partial_i)\partial_j\right).
\end{eqnarray*}
Contract with $g^{jk}$. The first term on the right contracts to
zero, since $g^{jk}$ is symmetric in $j,k$ while ${\mathcal
R}(\partial_j,\partial_k)$ is antisymmetric. For the second term,
writing $df=f_{,s}dx^s$ and ${\mathcal
R}(\partial_k,\partial_i)\partial_j={{R_{ki}}^s}_j\partial_s$ with
${{R_{ki}}^s}_j=g^{sm}R_{kimj}$, the symmetries of the curvature
tensor give $R_{kimj}=-R_{ikmj}$, so that
$$-g^{jk}\,df\left({\mathcal R}(\partial_k,\partial_i)\partial_j\right)
=-g^{jk}g^{sm}f_{,s}R_{kimj}
=g^{sm}f_{,s}\,g^{jk}R_{ikmj}
=g^{sm}f_{,s}\Ric_{im}
=\Ric\left((\nabla f)^*,\partial_i\right),$$
using the definition $\Ric_{im}=g^{jk}R_{ikmj}$ (contraction on the
second and fourth indices) and the fact that $(\nabla
f)^*=g^{sm}f_{,s}\partial_m$ is the gradient of $f$. Hence
$$g^{jk}(\nabla_{\partial_j}H)(\partial_k,\partial_i)
-g^{jk}(\nabla_{\partial_i}H)(\partial_j,\partial_k)
=\Ric\left((\nabla f)^*,\partial_i\right),$$
which, by the identifications in the first paragraph, is exactly
$\triangle df=d\triangle f+\Ric((\nabla f)^*,\cdot)$ evaluated on
$\partial_i$.
\end{proof}




\subsection{First examples}

The most homogeneous Riemannian manifolds are those of constant
sectional curvature. These are classified by the following theorem
(cf.\ Corollary 10 of Chapter 5 on page 147 of \cite{Petersen}). The
proof given below uses the theory of geodesics, Jacobi fields and
the exponential map developed in the section on geodesics below ---
in particular the subsection on local isometries and constant
curvature --- to which the reader may wish to turn first.

\begin{theorem}[Uniformization Theorem]
\label{thm:uniformization}
\uses{def:sectional-curvature}
{\bf (Uniformization Theorem)} If
$(M^n,g)$, $n\ge2$, is a complete, simply-connected Riemannian
manifold of constant sectional curvature $\lambda$, then:
\begin{enumerate}
\item If $\lambda=0$, then $M^n$ is isometric to Euclidean $n$-space.
\item If $\lambda>0$ there is a diffeomorphism $\phi \colon M \rightarrow S^{n}$
 such that g = $\lambda^{-1}\phi^{*}(g_\mathrm{st})$ where $g_{st}$ is the usual metric on the unit sphere in $\Ar^{n+1}$.
\item If $\lambda<0$ there is a diffeomorphism $\phi \colon M \rightarrow \HH^{n}$
 such that g = $\abs{\lambda}^{-1}\phi^{*}(g_\mathrm{st})$ where $g_{st}$ is the Poincar\'e metric of constant curvature
$-1$ on $\HH^n$.
\end{enumerate}
\end{theorem}

\begin{proof}
\uses{def:sectional-curvature,thm:hopf-rinow,def:exponential-map,lem:constant-curvature-jacobi,lem:local-isometry-covering,lem:local-isometry-rigidity,lem:local-isometry-exponential,def:local-isometry,thm:levi-civita-connection,lem:model-spaces,lem:model-polar-isometry,lem:submanifold-connection}
(The hypothesis $n\ge2$ is essential: for $n=1$ there are no
$2$-planes, so a complete, simply-connected Riemannian
$1$-manifold --- necessarily isometric to $\mathbb{R}$, by
arc-length parameterization --- has constant sectional curvature
$\lambda$ vacuously for \emph{every} $\lambda$, and the conclusion
fails for $\lambda>0$.) We freely use the notation
$s_\lambda$, $D_\lambda$ and the results of the subsection on local
isometries and constant curvature.

\emph{Step 1: behavior under scaling.} For a constant $c>0$ the
metrics $g$ and $cg$ have the same Christoffel symbols (the factor
$c$ cancels between $g^{kl}$ and $\partial g_{ij}$ in
Equation~(\ref{Gamma})), hence the same geodesics and the same
$(1,3)$ curvature tensor; therefore $R_{ijkl}(cg)=c\,R_{ijkl}(g)$,
and evaluating on a $cg$-orthonormal basis $\{X/\sqrt c,Y/\sqrt
c\}$ of a $2$-plane $P$ gives $K_{cg}(P)=K_g(P)/c$. Moreover
$d_{cg}=\sqrt{c}\,d_g$, so completeness is preserved, and simple
connectivity is a topological property. Consequently: if $(M,g)$ is
complete and simply connected of constant curvature $\lambda$, then
$(M,cg)$ is complete and simply connected of constant curvature
$\lambda/c$.

\emph{Step 2: the cases $\lambda\le 0$.} Suppose $(M,g)$ is
complete and simply connected of constant curvature $\lambda\le 0$,
and fix $p\in M$. By Lemma~\ref{lem:model-polar-isometry} (again
from the subsection on local isometries and constant curvature,
where it is proved independently of the present theorem),
$\exp_p$ is an isometry
$$(T_pM,G)\ \cong\ (M,g),$$
where $G=\exp_p^*g$ is given in polar coordinates on
$T_pM\setminus\{0\}$ by
\begin{equation}\label{modelpolar}
G=dr^2+s_\lambda(r)^2\,g_{S^{n-1}}.
\end{equation}

If $\lambda=0$ then $s_0(r)=r$ and Equation~(\ref{modelpolar})
says that $G$ agrees with the Euclidean metric $g_{p}$ of $T_pM$ on
$T_pM\setminus\{0\}$, hence everywhere by continuity. Thus $M$ is
isometric to Euclidean $n$-space, proving (1).

Suppose $\lambda<0$. By Lemma~\ref{lem:model-spaces}(3) (proved,
like the other results of the subsection on local isometries and
constant curvature, independently of the present theorem),
$(\HH^n,g_\mathrm{st})$ is
a complete, simply-connected manifold of constant sectional
curvature $-1$; by Step 1, $\hat g:=|\lambda|^{-1}g_\mathrm{st}$
is complete, simply connected, of constant curvature
$-1/|\lambda|^{-1}=\lambda$. Applying the previous analysis to both
$(M,g)$ at $p$ and $(\HH^n,\hat g)$ at some $q$, we obtain
isometries $M\cong (T_pM,G_M)$ and $\HH^n\cong(T_q\HH^n,G_{\HH})$,
where both pullback metrics are given by the same polar formula
(\ref{modelpolar}). Any linear isometry $A\colon (T_pM,g_p)\to
(T_q\HH^n,\hat g_q)$ preserves $r$ and restricts to an isometry of
unit spheres, so $A^*G_{\HH}=G_M$. Composing,
$$\psi:=\exp_q^{\HH}\circ A\circ(\exp_p^M)^{-1}\colon M\to\HH^n$$
is a diffeomorphism with $g=\psi^*\hat
g=|\lambda|^{-1}\psi^*(g_\mathrm{st})$, proving (3) with
$\phi=\psi$.

\emph{Step 3: the case $\lambda>0$.} Let
$S:=(S^n,\lambda^{-1}g_\mathrm{st})$ and
$\rho:=1/\sqrt\lambda$, so that $D:=D_\lambda=\pi\rho$. Rescaling
the standard embedding shows $S$ is isometric to the sphere
$S_\rho\subset \mathbb{R}^{n+1}$ of radius $\rho$ with its induced
metric; by Lemma~\ref{lem:model-spaces}(2), $S$ is complete,
simply connected, and of constant curvature
$1/\rho^2=\lambda$. We
work with the embedded picture $S_\rho$.

We first record the classical description of the exponential map of
$S_\rho$. Fix $\bar q\in S_\rho$ and a unit tangent vector $u\in
T_{\bar q}S_\rho$ (so $\langle \bar q,u\rangle=0$ in
$\mathbb{R}^{n+1}$). The curve
$p(t)=\cos(t/\rho)\,\bar q+\rho\sin(t/\rho)\,u$ lies on $S_\rho$,
has unit speed, and satisfies $p''(t)=-\rho^{-2}p(t)$, which is
everywhere normal to the sphere; hence $p$ is a geodesic by
Lemma~\ref{lem:submanifold-connection} (a curve in a submanifold of
Euclidean space is a geodesic if and only if its second derivative
is everywhere normal to the submanifold). Hence
$$\exp_{\bar q}(ru)=\cos(r/\rho)\,\bar q+\rho\sin(r/\rho)\,u.$$
For $0<r<\pi\rho$ the point $x=\exp_{\bar q}(ru)$ determines
$r$ (via $\langle x,\bar q\rangle=\rho^2\cos(r/\rho)$ and the
injectivity of $\cos$ on $(0,\pi)$) and then $u$ (via
$u=(x-\cos(r/\rho)\bar q)/(\rho\sin(r/\rho))$), and every $x\in
S_\rho\setminus\{\pm\bar q\}$ arises. Thus
$\exp_{\bar q}$ maps $B(0,D)\subset T_{\bar q}S_\rho$ bijectively to
$S_\rho\setminus\{-\bar q\}$, and it is non-singular there by
Lemma~\ref{lem:constant-curvature-jacobi}(2) applied to $S_\rho$;
hence
$$\exp_{\bar q}\colon B(0,D)\to S_\rho\setminus\{-\bar q\}$$
is a diffeomorphism. By
Lemma~\ref{lem:constant-curvature-jacobi}(3), the pullback
$G_{\bar q}:=\exp_{\bar q}^*(g_{S_\rho})$ on $B(0,D)$ has the polar
form (\ref{modelpolar}).

Now let $(M,g)$ be complete and simply connected of constant
curvature $\lambda$, and fix $p\in M$. By
Lemma~\ref{lem:constant-curvature-jacobi}, $\exp_p$ is defined on
$T_pM$ and non-singular on $B(0,D)\subset T_pM$, and
$G_p:=\exp_p^*g$ has the polar form (\ref{modelpolar}) on $B(0,D)$.
Fix a linear isometry $A\colon T_{\bar q}S_\rho\to T_pM$; as in Step
2, $A^*G_p=G_{\bar q}$ on $B(0,D)$. Define
$$F_1:=\exp_p\circ A\circ \left(\exp_{\bar
q}\right)^{-1}\colon\ S_\rho\setminus\{-\bar q\}\longrightarrow M.$$
As a composition of the isometry $(\exp_{\bar q})^{-1}\colon
S_\rho\setminus\{-\bar q\}\to (B(0,D),G_{\bar q})$, the isometry
$A$, and the local isometry $\exp_p\colon (B(0,D),G_p)\to M$, the
map $F_1$ is a local isometry.

Next we extend $F_1$ across $-\bar q$. Choose $x_0\in
S_\rho\setminus\{\bar q,-\bar q\}$, set $p_1=F_1(x_0)$ and
$B_1=d(F_1)_{x_0}\colon T_{x_0}S_\rho\to T_{p_1}M$, a linear
isometry, and define, by exactly the same construction based at
$x_0$ and $p_1$,
$$F_2:=\exp_{p_1}\circ B_1\circ\left(\exp_{x_0}\right)^{-1}
\colon\ S_\rho\setminus\{-x_0\}\longrightarrow M,$$
again a local isometry. Since $d(\exp_{x_0})_0$ and
$d(\exp_{p_1})_0$ are the identity, $F_2(x_0)=p_1=F_1(x_0)$ and
$d(F_2)_{x_0}=B_1=d(F_1)_{x_0}$. The overlap
$W=S_\rho\setminus\{-\bar q,-x_0\}$ is connected (a sphere of
dimension $n\ge 2$ minus two points), so by
Lemma~\ref{lem:local-isometry-rigidity} applied to $F_1|_W$ and
$F_2|_W$ we get $F_1=F_2$ on $W$. Since $-\bar q\not=-x_0$, the
domains of $F_1$ and $F_2$ cover $S_\rho$, and the two maps glue to
a local isometry
$$F\colon S_\rho\longrightarrow M.$$
As $S_\rho$ is complete and $M$ is connected,
Lemma~\ref{lem:local-isometry-covering} shows $F$ is a covering
map; as $M$ is simply connected and $S_\rho$ is connected, the
covering is one-sheeted by classical covering-space theory, so $F$
is an isometry from $S_\rho$, i.e., from $(S^n,
\lambda^{-1}g_\mathrm{st})$, onto $(M,g)$. Then
$\phi:=F^{-1}\colon M\to S^n$ is a diffeomorphism with
$g=(F^{-1})^*\left(\lambda^{-1}g_\mathrm{st}\right)
=\lambda^{-1}\phi^*(g_\mathrm{st})$, proving (2).
\end{proof}

Of course, if $(M^n,g)$ is a complete manifold of constant sectional
curvature then its universal covering satisfies the hypothesis of
the theorem and hence is one of $S^n, \Ar^n$, or $\HH^n$, up
to a constant scale factor. This implies that $(M,g)$ is isometric
to a quotient of one of these simply connected spaces of constant
curvature by the free action of a discrete group of isometries. Such
a Riemannian manifold is called a
\textit{space-form}.


\begin{definition}[Einstein manifold]
\label{def:einstein-manifold}
\uses{def:ricci-curvature}
\lean{MorganTianLib.IsEinstein}
\leanok
 The Riemannian manifold
$(M,g)$ is said to be an \textit{Einstein manifold with Einstein
constant} $\lambda$ if $\Ric(g)=\lambda g$.
\end{definition}

\begin{example}[Einstein manifolds in dimensions 2 and 3]
\label{ex:einstein-dimension-2-3}
\uses{def:einstein-manifold,def:sectional-curvature}
\lean{MorganTianLib.IsEinstein.sectionalCurvature_const}
\leanok
 Let $M$ be an $n$-dimensional manifold with $n$ being either $2$ or $3$.
If $(M,g)$ is Einstein with Einstein constant $\lambda$, then $M$ has
constant sectional curvature $\frac{\lambda}{n-1}$.
\end{example}

\begin{proof}
\uses{def:einstein-manifold,def:sectional-curvature,def:ricci-curvature,claim:curvature-symmetries-bianchi}
\leanok
Fix $p\in M$ and let $P\subset T_pM$ be any $2$-plane. Choose an
orthonormal basis $\{e_1,\ldots,e_n\}$ of $T_pM$ whose first two
vectors span $P$, and write $K_{ij}={\mathcal R}(e_i,e_j,e_i,e_j)$
for $i\not=j$, the sectional curvature of the plane spanned by $e_i$
and $e_j$. From the formula for the Ricci tensor in an orthonormal
basis, and since ${\mathcal R}(e_i,e_i,e_i,e_i)=0$ by the
antisymmetries of the curvature tensor
(Claim~\ref{claim:curvature-symmetries-bianchi}),
$$\Ric(e_i,e_i)=\sum_{j=1}^n{\mathcal R}(e_i,e_j,e_i,e_j)
=\sum_{j\not=i}K_{ij}.$$
Note also that for $i\not=j$, the antisymmetries of ${\mathcal R}$
in its first two and in its last two arguments give
$$K_{ji}={\mathcal R}(e_j,e_i,e_j,e_i)
=-{\mathcal R}(e_i,e_j,e_j,e_i)
={\mathcal R}(e_i,e_j,e_i,e_j)=K_{ij}.$$

Suppose $n=2$. Then the Einstein condition
$\Ric(e_1,e_1)=\lambda\,g(e_1,e_1)=\lambda$ reads
$K_{12}=\lambda=\frac{\lambda}{n-1}$, and $K_{12}=K(P)$.

Suppose $n=3$. The Einstein condition applied to $e_1,e_2,e_3$ gives
the three equations
$$K_{12}+K_{13}=\lambda,\qquad K_{12}+K_{23}=\lambda,\qquad
K_{13}+K_{23}=\lambda.$$
Subtracting the second from the first gives $K_{13}=K_{23}$, and
subtracting the third from the second gives $K_{12}=K_{13}$. Hence
all three are equal and $2K_{12}=\lambda$, so
$K(P)=K_{12}=\frac{\lambda}{2}=\frac{\lambda}{n-1}$.

In both cases the value of $K(P)$ is the constant
$\frac{\lambda}{n-1}$, independent of the point $p$ and of the plane
$P$; that is, $(M,g)$ has constant sectional curvature
$\frac{\lambda}{n-1}$.
\end{proof}

\begin{remark}[Complete Einstein surfaces and $3$-manifolds are space forms]
\label{rem:einstein-space-form}
\uses{ex:einstein-dimension-2-3,thm:uniformization}
Continuing Example~\ref{ex:einstein-dimension-2-3}: if moreover
$(M,g)$ is complete, then, having constant sectional curvature
$\frac{\lambda}{n-1}$, it is a space-form by the discussion following
Theorem~\ref{thm:uniformization}. This last step relies on the
uniformization theorem and is not part of the formalized statement
above.
\end{remark}

\subsection{Cones}

Another class of examples that will play an important role in our
study of the Ricci flow is that of cones.

\begin{definition}[Open cone]
\label{def:cone}
\label{conedefn}
Let $(N,g)$ be a Riemannian manifold. We define the {\em open cone}
over $(N,g)$ to be the manifold $N\times (0,\infty)$ with the metric
$\tilde g$ defined as follows: For any $(x,s)\in N\times (0,\infty)$
we have
$$\tilde g(x,s)=s^2g(x)+ds^2.$$
\end{definition}

Fix local coordinates $(x^1,\ldots,x^{n-1})$ on $N$ (of dimension
$n-1$, as below). Let
$\Gamma_{ij}^k;\ 1\le i,j,k\le n-1$, be the Christoffel symbols for
the Levi-Civita connection on $N$. Set $x^0=s$. In the local
coordinates $(x^0,x^1,\ldots,x^{n-1})$ for the cone we have the
Christoffel symbols $\widetilde \Gamma_{ij}^k;\ 0\le i,j,k\le n-1$,
for $\tilde g$. The relation between the metrics gives the following
relations between the two sets of Christoffel symbols:
\begin{eqnarray*}\widetilde
\Gamma_{ij}^k & = & \Gamma_{ij}^k; \ \ \ \ 1\le i,j,k\le n-1 \\
\widetilde \Gamma_{ij}^0 & = & -sg_{ij}; \ \ \ \ 1\le i,j\le n-1 \\
\widetilde\Gamma_{i0}^j & = & \widetilde\Gamma_{0i}^j =
s^{-1}\delta^j_i; \ \ \ \ 1\le i,j\le n-1 \\
\widetilde \Gamma_{i0}^0 & = & 0;\ \ \ \ 0\le i\le n-1 \\
\widetilde \Gamma_{00}^i & = & 0;\ \  \ \ 0\le i\le n-1.
\end{eqnarray*}

Denote by ${\mathcal R}_g$ the curvature tensor for $g$ and by
${\mathcal R}_{\tilde g}$ the curvature tensor for $\tilde g$. Then
the above formulas lead directly to:
\begin{eqnarray*}
{\mathcal R}_{\tilde g}(\partial_i,\partial_j)(\partial_0) & = & 0;\ \ \ \  0\le i,j\le n-1 \\
{\mathcal R}_{\tilde g}(\partial_i,\partial_j)(\partial_i) & = &
{\mathcal
R}_g(\partial_i,\partial_j)(\partial_i)+g_{ii}\partial_j-g_{ji}\partial_i;\
\ \ \  1\le i,j\le n-1 \\
\end{eqnarray*}

This allows us to compute the Riemann curvatures of the cone in
terms of those of $N$.

\begin{proposition}[Curvature operator of cone]
\label{prop:cone-curvature}
\label{conecurv}
\uses{def:cone,def:curvature-operator}
Let $N$ be a Riemannian manifold of dimension $n-1$. Fix $(p,s)\in
c(N)=N\times (0,\infty)$. With respect to the orthogonal splitting
$$\wedge^2T_{(p,s)}c(N)=\wedge^2T_pN\oplus
\left(T_pN\wedge\partial_0\right),$$
where $\partial_0$ is the coordinate field of $x^0=s$, the curvature
operator $\Rm_{\tilde g}(p,s)$ of the cone is block-diagonal,
$$\begin{pmatrix} s^2(\Rm_g(p)-\wedge^2g(p)) & 0 \\ 0 & 0
\end{pmatrix},$$
where $\wedge^2g(p)$ is the symmetric form on $\wedge^2T_pN$ induced
by $g$. (The corresponding display in the original source places
the non-zero block in an off-diagonal corner, where it is not
symmetric; the block belongs on the diagonal, as the proof shows.)
\end{proposition}

\begin{proof}
\uses{def:cone,def:curvature-operator,thm:levi-civita-connection,def:riemann-curvature-tensor,claim:curvature-symmetries-bianchi}
Throughout, indices $i,j,k,l$ refer to the local coordinates on $N$
and $0$ to the coordinate $x^0=s$; the metric components of $\tilde
g=s^2g+ds^2$ are $\tilde g_{00}=1$, $\tilde g_{0i}=0$ and $\tilde
g_{ij}=s^2g_{ij}$, so $\tilde g^{00}=1$, $\tilde g^{0i}=0$ and
$\tilde g^{ij}=s^{-2}g^{ij}$.

\emph{The Christoffel symbols.} The relations displayed before the
proposition follow from Equation~(\ref{Gamma}) applied to $\tilde
g$. Indeed, since the only $s$-dependence of the components of
$\tilde g$ is the factor $s^2$ in $\tilde g_{ij}$:
$$\widetilde\Gamma^k_{ij}
=\tfrac{1}{2}s^{-2}g^{kl}\left(\partial_i(s^2g_{lj})
+\partial_j(s^2g_{il})-\partial_l(s^2g_{ij})\right)=\Gamma^k_{ij},$$
$$\widetilde\Gamma^0_{ij}
=\tfrac{1}{2}\left(\partial_i\tilde g_{0j}+\partial_j\tilde
g_{i0}-\partial_0(s^2g_{ij})\right)=-sg_{ij},\qquad
\widetilde\Gamma^j_{i0}
=\tfrac{1}{2}s^{-2}g^{jl}\,\partial_0(s^2g_{il})=s^{-1}\delta^j_i,$$
and $\widetilde\Gamma^0_{i0}=\widetilde\Gamma^0_{00}
=\widetilde\Gamma^k_{00}=0$ because $\tilde g_{00}$ is constant and
$\tilde g_{i0}=0$.

\emph{The curvature operator ${\mathcal R}_{\tilde g}$ on coordinate
fields.} Write $\widetilde\nabla_a=\widetilde\nabla_{\partial_a}$.
From the Christoffel relations, $\widetilde\nabla_i\partial_0
=s^{-1}\partial_i$, $\widetilde\nabla_0\partial_0=0$ and
$\widetilde\nabla_i\partial_j
=\Gamma^k_{ij}\partial_k-sg_{ij}\partial_0$. First,
$$\widetilde\nabla_i\left(\widetilde\nabla_j\partial_0\right)
=\widetilde\nabla_i\left(s^{-1}\partial_j\right)
=s^{-1}\left(\Gamma^k_{ij}\partial_k-sg_{ij}\partial_0\right),$$
which is symmetric in $i$ and $j$, so ${\mathcal
R}_{\tilde g}(\partial_i,\partial_j)\partial_0=0$. Also
$$\widetilde\nabla_0\left(\widetilde\nabla_i\partial_0\right)
=\widetilde\nabla_0\left(s^{-1}\partial_i\right)
=-s^{-2}\partial_i+s^{-1}\widetilde\nabla_0\partial_i
=-s^{-2}\partial_i+s^{-2}\partial_i=0,$$
and $\widetilde\nabla_i(\widetilde\nabla_0\partial_0)=0$, so
${\mathcal R}_{\tilde g}(\partial_i,\partial_0)\partial_0=0$.
Next, for tangential fields, using the two displayed covariant
derivatives,
\begin{eqnarray*}
\widetilde\nabla_i\left(\widetilde\nabla_j\partial_k\right)
& = & \partial_i\left(\Gamma^l_{jk}\right)\partial_l
+\Gamma^l_{jk}\left(\Gamma^m_{il}\partial_m-sg_{il}\partial_0\right)
-s\,\partial_i(g_{jk})\,\partial_0-sg_{jk}\cdot s^{-1}\partial_i\\
& = & \left(\partial_i\Gamma^l_{jk}
+\Gamma^m_{jk}\Gamma^l_{im}\right)\partial_l
-g_{jk}\partial_i
-s\left(\Gamma^l_{jk}g_{il}+\partial_ig_{jk}\right)\partial_0.
\end{eqnarray*}
Antisymmetrizing in $i$ and $j$: the $\partial_l$-terms give exactly
the coordinate expression for ${\mathcal
R}_g(\partial_i,\partial_j)\partial_k$; the middle terms give
$g_{ik}\partial_j-g_{jk}\partial_i$; and the $\partial_0$-terms
cancel, since by metric compatibility on $N$,
$\partial_ig_{jk}=g_{lk}\Gamma^l_{ij}+g_{jl}\Gamma^l_{ik}$, so that
$$\Gamma^l_{jk}g_{il}+\partial_ig_{jk}
-\Gamma^l_{ik}g_{jl}-\partial_jg_{ik}
=g_{lk}\Gamma^l_{ij}-g_{lk}\Gamma^l_{ji}=0.$$
Hence
$$
{\mathcal R}_{\tilde g}(\partial_i,\partial_j)\partial_k
={\mathcal R}_g(\partial_i,\partial_j)\partial_k
+g_{ik}\partial_j-g_{jk}\partial_i,
$$
which for $k=i$ is the second relation displayed before the
proposition.

\emph{The $(0,4)$-tensor.} Since ${\mathcal R}_{\tilde
g}(\partial_a,\partial_b)\partial_0=0$ for all $a,b$ (the cases
$(i,j)$ and $(i,0)$ were computed above, and the case $(0,i)$
follows by antisymmetry), every component
$\widetilde R_{abc0}=\tilde g({\mathcal
R}_{\tilde g}(\partial_a,\partial_b)\partial_0,\partial_c)$
vanishes. Consequently, by the symmetries of the curvature tensor
(Claim~\ref{claim:curvature-symmetries-bianchi}), a component of
$\widetilde R$ vanishes whenever the index $0$ appears in any
position: $\widetilde R_{ab0d}=-\widetilde R_{ab d0}$ vanishes by
the antisymmetry in the last two indices, and then $\widetilde
R_{0bcd}=\widetilde R_{cd0b}$ and $\widetilde R_{a0cd}=\widetilde
R_{cda0}$ vanish by the pair symmetry. For all-tangential indices,
pairing the displayed operator identity with $s^2g$,
\begin{eqnarray*}
\widetilde R_{ijkl}
& = & \tilde g\left({\mathcal R}_{\tilde
g}(\partial_i,\partial_j)\partial_l,\partial_k\right)
=s^2g\left({\mathcal R}_g(\partial_i,\partial_j)\partial_l
+g_{il}\partial_j-g_{jl}\partial_i,\ \partial_k\right)\\
& = & s^2\left(R^g_{ijkl}
-\left(g_{ik}g_{jl}-g_{il}g_{jk}\right)\right).
\end{eqnarray*}

\emph{The curvature operator.} Decompose
$\wedge^2T_{(p,s)}c(N)=\wedge^2T_pN\oplus\left(T_pN\wedge\partial_0\right)$,
identifying the second summand with $T_pN$. If
$\varphi\in\wedge^2T_{(p,s)}c(N)$ is arbitrary and
$\psi=\psi^{k0}\partial_k\wedge\partial_0$ lies in the second
summand, then $\Rm_{\tilde
g}(\varphi,\psi)$ is a sum of components of $\widetilde R$ each
carrying the index $0$, and hence vanishes; thus $T_pN\wedge
\partial_0$ pairs to zero with everything. On the first summand, for
$\varphi=\varphi^{ij}\partial_i\wedge\partial_j$ and
$\psi=\psi^{kl}\partial_k\wedge\partial_l$ in $\wedge^2T_pN$,
$$\Rm_{\tilde g}(\varphi,\psi)=\widetilde
R_{ijkl}\varphi^{ij}\psi^{kl}
=s^2\left(R^g_{ijkl}\varphi^{ij}\psi^{kl}
-\left(g_{ik}g_{jl}-g_{il}g_{jk}\right)\varphi^{ij}\psi^{kl}\right)
=s^2\left(\Rm_g-\wedge^2g\right)(\varphi,\psi),$$
since $(g_{ik}g_{jl}-g_{il}g_{jk})\varphi^{ij}\psi^{kl}$ is
precisely the symmetric form $\wedge^2g$ induced by $g$ on
$\wedge^2T_pN$. This is exactly the stated block decomposition of
$\Rm_{\tilde g}(p,s)$ with respect to the splitting
$\wedge^2T_pN\oplus T_pN$: the restriction to $\wedge^2T_pN$ is
$s^2(\Rm_g(p)-\wedge^2g(p))$ and the summand $T_pN$ lies in the null
space.
\end{proof}

\begin{corollary}[Eigenvalues of cone curvature operator]
\label{cor:cone-eigenvalues}
\uses{prop:cone-curvature}
For any $p\in N$ let $\lambda_1,\ldots,\lambda_{(n-1)(n-2)/2}$ be
the eigenvalues of $\Rm_g(p)$. Then for any $s>0$  there are
$(n-1)$ zero eigenvalues of $\Rm_{\tilde g}(p,s)$. The other
$(n-1)(n-2)/2$ eigenvalues of $\Rm_{\tilde g}(p,s)$ are
$s^{-2}(\lambda_i-1)$.
\end{corollary}


\begin{proof}
\uses{prop:cone-curvature}
Clearly from Proposition~\ref{prop:cone-curvature}, we see that under the
orthogonal decomposition $\wedge^2T_{(p,s)}c(N)=\wedge^2T_pN\oplus
T_pN$ the second subspace is contained in the null space of $ \Rm_{\tilde g}(p,s)$, and hence contributes $(n-1)$ zero
eigenvalues. Likewise, from this proposition we see that the
eigenvalues of the restriction of $\Rm_{\tilde g}(p,s)$ to the
subspace $\wedge^2T_pN$ are  given by
$s^{-4}(s^2(\lambda_i-1))=s^{-2}(\lambda_i-1)$.
\end{proof}


\section{Geodesics and the exponential map}

Here we review standard material about geodesics, Jacobi fields, and
the exponential map.

\subsection{Geodesics and the energy functional}

\begin{definition}[Geodesic]
\label{def:geodesic}
\uses{thm:levi-civita-connection}
\lean{MorganTianLib.IsGeodesic,MorganTianLib.IsGeodesicOn}
\leanok
 Let $I$ be an open interval.
A smooth curve $\gamma\colon I \rightarrow M$ is called a
\textit{geodesic} if $\nabla_{\dot \gamma}\dot
\gamma =0$.
\end{definition}

In local coordinates, we write $\gamma(t)=(x^1(t),\ldots,x^n(t))$
and this equation becomes
$$0=\nabla_{\dot \gamma}\dot \gamma(t) =
\left(\sum\limits_k\left(\ddot x^k(t) +\dot x^i(t)\dot
x^j(t)\Gamma^{k}_{ij}(\gamma(t))\right)\partial_k\right).$$ This is
a system of $2^{nd}$ order ODE's. The local existence, uniqueness
and smoothness of a geodesic through any point $p\in M$ with initial
velocity vector $v\in T_pM$ follow from the classical ODE theory.
We call geodesics that minimize the length between their endpoints
{\em minimizing geodesics}. On a complete manifold, any two points
are joined by a minimizing geodesic, and all geodesics are defined
for all time; this is the Hopf--Rinow theorem
(Theorem~\ref{thm:hopf-rinow}), which we state and prove in the
subsection on the exponential mapping below, once the exponential
map is available.


Geodesics are critical points of the energy functional. Let $(M,g)$
be a complete Riemannian manifold. Consider the space of $C^1$-paths
in $M$ parameterized by the unit interval. On this space we have the
energy functional
$$E(\gamma)=\frac{1}{2}\int_{0}^{1}\langle \gamma'(t),\gamma'(t)\rangle dt.$$
 Suppose that we have a one-parameter family of  paths parameterized by
 $[0,1]$,
 all having the same initial point $p$ and the same final point $q$.
 By this we mean that we have a surface $\tilde\gamma(t,u)$ with the property that for each
 $u$ the path $\gamma_u=\tilde\gamma(\cdot,u)$ is a path from $p$ to $q$ parameterized by $[0,1]$.
 Let $\widetilde X=\partial\tilde \gamma/\partial t$ and
 $\widetilde Y=\partial\tilde \gamma/\partial u$ be the corresponding vector
 fields along the surface swept out by $\tilde \gamma$, and denote by $X$ and $Y$ the restriction
 of these vector fields along $\gamma_0$.
 We compute
 \begin{eqnarray*}
 \frac{dE(\gamma_u)}{du}\Bigl|_{u=0}\Bigr.& = & \left( \int_0^1\langle \nabla_{\widetilde Y}\widetilde X,\widetilde
 X\rangle dt\right)\left|_{u=0}\right. \\
 & = & \left(\int_0^1\langle\nabla_{\widetilde X}\widetilde Y,\widetilde
 X\rangle dt \right)\left|_{u=0}\right.\\
 & = & -\left( \int_0^1\langle\nabla_{\widetilde X}\widetilde X,\widetilde
 Y\rangle dt\right)\left|_{u=0}\right.=-\int_0^1\langle\nabla_XX,Y\rangle,
\end{eqnarray*}
where the first equality in the last line comes from integration by
parts and the fact that $\widetilde Y$ vanishes at the endpoints.
Given any vector field $Y$ along $\gamma_0$ there is a one-parameter
family $\tilde \gamma(t,u)$ of paths from $p$ to $q$ with $\tilde
\gamma(t,0)=\gamma_0$ and with $\widetilde Y(t,0)=Y$. Thus, from the
above expression we see that $\gamma_0$ is a critical point for the
energy functional on the space of paths from $p$ to $q$
parameterized by the interval $[0,1]$ if and only if $\gamma_0$ is a
geodesic.

Notice that it follows immediately from the geodesic equation that
the length of a tangent vector along a geodesic is constant. Thus,
if a geodesic is parameterized by $[0,1]$ we have
$$E(\gamma)=\frac{1}{2}L(\gamma)^2.$$
It is immediate from the Cauchy-Schwarz inequality that for any
curve $\mu$ parameterized by $[0,1]$ we have
$$E(\mu)\ge \frac{1}{2}L(\mu)^2$$
with equality if and only if $|\mu'|$ is constant. In particular, a
curve parameterized by $[0,1]$ minimizes distance between its
endpoints if it is a minimum for the energy functional on all paths
parameterized by $[0,1]$ with the given endpoints.






\subsection{Families of geodesics and Jacobi fields}


Consider a family of geodesics $\tilde\gamma(u,t)=\gamma_u(t)$
parameterized by the interval  $[0,1]$ with $\gamma_u(0)=p$ for all
$u$. Here, unlike the discussion above, we allow $\gamma_u(1)$ to
vary with $u$. As before define vector fields along the surface
swept out by $\tilde \gamma$:  $\widetilde X=\partial
\tilde\gamma/\partial t$ and let $\tilde
Y=\partial\tilde\gamma/\partial u$.  We denote by $X$ and $Y$ the
restriction of these vector fields to the geodesic
$\gamma_0=\gamma$.
 Since each
$\gamma_u$ is a geodesic, we have $\nabla_{\widetilde X}{\widetilde
X}=0$. Differentiating this equation in the $\widetilde Y$-direction
yields $\nabla_{\widetilde Y}\nabla_{\widetilde X}\widetilde X=0$.
Interchanging the order of differentiation, using
$\nabla_{\widetilde X}\widetilde Y=\nabla_{\widetilde Y}{\widetilde
X}$, and then restricting to $\gamma$, we get the {\em Jacobi
equation}:
$$\nabla_X\nabla_XY+{\mathcal R}(Y,X)X=0.$$
Notice that the left-hand side of the equation depends only on the
value of $Y$ along $\gamma$, not on the entire family. We denote the
left-hand side of this equation by $\Jac(Y)$, so that the Jacobi
equation now reads
$$\Jac(Y)=0.$$
 The fact that all the geodesics begin at the same point at time
$0$ means that $Y(0)=0$. A vector field $Y$ along a geodesic
$\gamma$ is said to be a {\em Jacobi field} if
it satisfies this equation and vanishes at the initial point $p$. A
Jacobi field is determined by its first derivative at $p$, i.e., by
$\nabla_XY(0)$. We have just seen that this is the equation
describing, to first order, variations of $\gamma$ by a family of
geodesics with the same starting point.

Jacobi fields are also determined by the energy functional. Consider
the space of paths parameterized by $[0,1]$ starting at a given
point $p$ but free to end anywhere in the manifold. Let $\gamma$ be
a geodesic (parameterized by $[0,1]$) from $p$ to $q$. Associated to
any one-parameter family $\tilde \gamma(t,u)$ of paths parameterized
by $[0,1]$ starting at $p$ we associate the second derivative of the
energy at $u=0$. Straightforward computation gives
$$\frac{d^2E(\gamma_u)}{du^2}\Bigl|_{u=0}\Bigr.=\langle \nabla_XY(1),Y(1)\rangle+\langle
X(1),\nabla_Y\widetilde Y(1,0)\rangle-\int_0^1\langle \Jac(Y),Y\rangle dt.$$

Notice that the first term is a boundary term from the integration
by parts, and it depends not just on the value of $Y$ (i.e., on
$\widetilde Y$ restricted to $\gamma$) but also on the first-order
variation of $\widetilde Y$ in the $Y$ direction. There is the
associated bilinear form that comes from two-parameter families
$\tilde \gamma(t,u_1,u_2)$ whose value at $u_1=u_1=0$ is $\gamma$.
It is
$$\frac{d^2E}{du_1du_2}\Bigl|_{u_1=u_2=0}\Bigr.=\langle \nabla_XY_1(1),Y_2(1)\rangle+\langle
X(1),\nabla_{Y_1}\widetilde Y_2(1,0)\rangle-\int_0^1\langle \Jac(Y_1),Y_2\rangle dt.$$ Notice that restricting to the space of
vector fields that vanish at both endpoints, the second derivatives
depend only on $Y_1$ and $Y_2$ and the formula is
$$\frac{d^2E}{du_1du_2}\Bigl|_{u_1=u_2=0}\Bigr.=-\int_0^1\langle \Jac(Y_1),Y_2\rangle dt,$$
so that this expression is symmetric in $Y_1$ and $Y_2$. The
associated quadratic form on the space of vector fields along
$\gamma$ vanishing at both endpoints
$$-\int_0^1\langle \Jac(Y),Y\rangle dt$$
is the second derivative of the energy function at $\gamma$ for any
one-parameter family whose value at $0$ is $\gamma$ and whose first
variation is given by $Y$.

The interchange of covariant derivatives used in the derivation of
the Jacobi equation above, and the Jacobi equation itself for a
family of geodesics, are recorded precisely — in local coordinates —
in the following lemma.

\begin{lemma}[Covariant commutation and the Jacobi equation for families]
\label{lem:covariant-commutation-jacobi}
\uses{thm:levi-civita-connection,def:riemann-curvature-tensor,def:geodesic}
\lean{MorganTianLib.covDerivAlong,MorganTianLib.christoffelCurvature,MorganTianLib.covDerivAlong_comm,MorganTianLib.covDerivAlong_fderiv_symm,MorganTianLib.covDerivAlong_geodesic_family_jacobi,MorganTianLib.chartChristoffelBilin,MorganTianLib.chartCurvature,MorganTianLib.chart_geodesic_family_jacobi}
\leanok
Let $U$ be a coordinate chart of $M$, with Christoffel symbols
$\Gamma^k_{ij}$ of the Levi-Civita connection
(Theorem~\ref{thm:levi-civita-connection}). For $x\in U$ write
$\Gamma_x(X,Y)$ for the symmetric bilinear map with components
$\Gamma^k_{ij}(x)X^iY^j$. Let $u=u(s,t)$ be a $C^2$ two-parameter
family of points of $U$, written in coordinates, and for a $C^2$
vector field $V(s,t)$ along $u$ (in coordinates) set
$$\nabla_sV=\partial_sV+\Gamma_u(\partial_su,V),\qquad
\nabla_tV=\partial_tV+\Gamma_u(\partial_tu,V).$$
Then:
\begin{enumerate}
\item ({\em Curvature commutation})
$\nabla_s\nabla_tV-\nabla_t\nabla_sV
={\mathcal R}(\partial_su,\partial_tu)V$, where ${\mathcal R}$ is
the curvature tensor of
Definition~\ref{def:riemann-curvature-tensor} in its coordinate
expression.
\item ({\em Torsion-freeness})
$\nabla_s\partial_tu=\nabla_t\partial_su$.
\item ({\em Jacobi equation}) If $u$ is $C^3$ and every line
$t\mapsto u(s,t)$ is a geodesic, then $Y=\partial_su$ satisfies
$$\nabla_t\nabla_tY+{\mathcal R}(Y,\partial_tu)\partial_tu=0.$$
\end{enumerate}
\end{lemma}

\begin{proof}
\leanok
\uses{thm:levi-civita-connection,def:riemann-curvature-tensor,def:geodesic}
The symmetry $\Gamma_x(X,Y)=\Gamma_x(Y,X)$ holds because
$\Gamma^k_{ij}=\Gamma^k_{ji}$, the Levi-Civita connection being
torsion-free.

(1) Abbreviate $u_s=\partial_su$ and $u_t=\partial_tu$, all
functions being evaluated at a fixed parameter point. Denote by
$(\partial_X\Gamma)_x(Y,Z)$ the derivative of the map
$x\mapsto\Gamma_x(Y,Z)$ (with $Y,Z$ held fixed) at $x$ in the
direction $X$; it is again bilinear in $(Y,Z)$, with components
$(\partial_m\Gamma^k_{ij})(x)X^mY^iZ^j$. Expanding the outer
covariant derivative and applying the chain and product rules to the
inner one,
\begin{align*}
\nabla_s\nabla_tV
&=\partial_s\bigl(\partial_tV+\Gamma_u(u_t,V)\bigr)
  +\Gamma_u\bigl(u_s,\partial_tV+\Gamma_u(u_t,V)\bigr)\\
&=\partial_s\partial_tV
  +(\partial_{u_s}\Gamma)_u(u_t,V)
  +\Gamma_u(\partial_s\partial_tu,V)
  +\Gamma_u(u_t,\partial_sV)\\
&\qquad+\Gamma_u(u_s,\partial_tV)
  +\Gamma_u\bigl(u_s,\Gamma_u(u_t,V)\bigr).
\end{align*}
The same formula with $s$ and $t$ interchanged gives
$\nabla_t\nabla_sV$. In the difference the terms
$\partial_s\partial_tV=\partial_t\partial_sV$ and
$\Gamma_u(\partial_s\partial_tu,V)=\Gamma_u(\partial_t\partial_su,V)$
cancel by the symmetry of second partial derivatives (Schwarz — for
which $C^2$ suffices), and the terms
$\Gamma_u(u_t,\partial_sV)+\Gamma_u(u_s,\partial_tV)$ cancel
pairwise. There remains
$$\nabla_s\nabla_tV-\nabla_t\nabla_sV
=(\partial_{u_s}\Gamma)_u(u_t,V)-(\partial_{u_t}\Gamma)_u(u_s,V)
+\Gamma_u\bigl(u_s,\Gamma_u(u_t,V)\bigr)
-\Gamma_u\bigl(u_t,\Gamma_u(u_s,V)\bigr),$$
whose $l$-th component is
$$\bigl(\partial_i\Gamma^l_{jk}-\partial_j\Gamma^l_{ik}
+\Gamma^m_{jk}\Gamma^l_{im}-\Gamma^m_{ik}\Gamma^l_{jm}\bigr)
u_s^iu_t^jV^k={{R_{ij}}^l}_ku_s^iu_t^jV^k,$$
by the coordinate expression of the curvature tensor following
Definition~\ref{def:riemann-curvature-tensor}. This is
${\mathcal R}(u_s,u_t)V$.

(2) By definition
$\nabla_s\partial_tu=\partial_s\partial_tu+\Gamma_u(u_s,u_t)$ and
$\nabla_t\partial_su=\partial_t\partial_su+\Gamma_u(u_t,u_s)$; these
agree by Schwarz symmetry and the symmetry of $\Gamma_u$.

(3) Since each $t$-line is a geodesic, $\nabla_tu_t=0$ identically
(Definition~\ref{def:geodesic} in coordinates), hence also
$\nabla_s\nabla_tu_t=0$. Since $u$ is $C^3$, the field $V=u_t$ is
$C^2$ and part~(1) applies to it:
$$0=\nabla_s\nabla_tu_t
=\nabla_t\nabla_su_t+{\mathcal R}(u_s,u_t)u_t.$$
By part~(2), $\nabla_su_t=\nabla_tu_s=\nabla_tY$, so
$$\nabla_t\nabla_tY+{\mathcal R}(Y,u_t)u_t=0,$$
as claimed.
\end{proof}

Along a fixed coordinate curve the Jacobi equation, viewed as an
equation for the pair $(J,\nabla J)$, is a first-order linear
system, and the basic solvability theory follows.

\begin{lemma}[Jacobi fields along a coordinate curve]
\label{lem:jacobi-field-coordinates}
\uses{lem:covariant-commutation-jacobi,def:riemann-curvature-tensor,thm:levi-civita-connection,lem:second-order-linear-ode}
\lean{MorganTianLib.chartCurvatureEndo,MorganTianLib.IsJacobiFieldOn,MorganTianLib.jacobiPairCoeffCoord,MorganTianLib.exists_isJacobiFieldOn_Icc,MorganTianLib.exists_isJacobiFieldOn_Icc_of_curve,MorganTianLib.continuousOn_jacobiPairCoeffCoord,MorganTianLib.IsJacobiFieldOn.eqOn_of_left,MorganTianLib.IsJacobiFieldOn.eqOn_of_right,MorganTianLib.IsJacobiFieldOn.add,MorganTianLib.IsJacobiFieldOn.const_smul,MorganTianLib.IsJacobiFieldOn.eqOn_zero}
\leanok
Let $u\colon[a,b]\to U$ be a $C^1$ curve in a coordinate chart $U$,
with $\Gamma$ and ${\mathcal R}$ as in
Lemma~\ref{lem:covariant-commutation-jacobi}, the coefficient
functions being continuous along $u$. Say that a pair of curves
$J,DJ\colon[a,b]\to\mathbb{R}^n$ is a {\em Jacobi field along $u$}
(with covariant derivative $DJ$) if
$$J'=DJ-\Gamma_u(\dot u,J),\qquad
DJ'=-{\mathcal R}(J,\dot u)\dot u-\Gamma_u(\dot u,DJ)$$
on $[a,b]$ (one-sided derivatives at the endpoints) — equivalently,
$\nabla J=DJ$ and $\nabla DJ+{\mathcal R}(J,\dot u)\dot u=0$, the
Jacobi equation. Then:
\begin{enumerate}
\item (Existence) for every pair $(J_0,DJ_0)$ there is a Jacobi
field along $u$ with $J(a)=J_0$, $DJ(a)=DJ_0$.
\item (Uniqueness) two Jacobi fields along $u$ agreeing (in both
components) at $a$, or at $b$, agree on all of $[a,b]$.
\item (Superposition) sums and scalar multiples of Jacobi fields
along $u$ are Jacobi fields; a Jacobi field with $J(a)=0$,
$DJ(a)=0$ vanishes identically, so one with nontrivial initial
data does not vanish identically.
\end{enumerate}
\end{lemma}

\begin{proof}
\leanok
\uses{lem:covariant-commutation-jacobi,lem:second-order-linear-ode}
On $\mathbb{R}^n\times\mathbb{R}^n$ define
$$A(t)(x,w)=\bigl(w-\Gamma_{u(t)}(\dot u(t),x),\;
-{\mathcal R}(x,\dot u(t))\dot u(t)-\Gamma_{u(t)}(\dot u(t),w)\bigr).$$
Each $A(t)$ is linear in $(x,w)$ — $\Gamma_x(\cdot,\cdot)$ is
bilinear and $x\mapsto{\mathcal R}(x,v)v$ is linear, being composed
of the bilinear maps $\Gamma$ and the directional derivative
$(\partial_X\Gamma)$ from the proof of
Lemma~\ref{lem:covariant-commutation-jacobi} — hence bounded, and
$t\mapsto A(t)$ is continuous, since $\Gamma$ and its first
derivatives are continuous on the chart and $u,\dot u$ are
continuous. By construction $(J,DJ)$ is a Jacobi field along $u$ if
and only if $W=(J,DJ)$ solves the first-order linear system
$W'=A(t)W$ on $[a,b]$.

Existence with prescribed $W(a)$, and forward and backward
uniqueness, now follow exactly as in the proof of
Lemma~\ref{lem:second-order-linear-ode}: the Picard--Lindel\"of
contraction on subintervals of length $1/(2K)$ (where
$\|A\|\le K$ on $[a,b]$, by continuity and compactness) glued
finitely many times gives existence, and Gr\"onwall's inequality
applied to the difference of two solutions gives uniqueness from
either endpoint — that argument used only the boundedness and
continuity of the coefficient of the first-order system, not its
special shape. Superposition is immediate from the linearity of
$A(t)$; the zero pair is a solution, so a Jacobi field with
vanishing initial data coincides with it by uniqueness.
\end{proof}

The coefficient operator ${\mathcal R}$ of the chart Jacobi system is
not merely an ad hoc combination of Christoffel symbols: it {\em is}
the Riemann curvature tensor of
Definition~\ref{def:riemann-curvature-tensor} read in the chart, and
in particular it inherits pointwise bounds from the sectional
curvature. This is the bridge through which the geometric hypothesis
``$K(P)\le K$'' enters the Jacobi equation.

\begin{lemma}[The curvature tensor in chart coordinates]
\label{lem:chart-curvature-coordinates}
\uses{def:riemann-curvature-tensor,thm:levi-civita-connection,def:sectional-curvature,lem:covariant-commutation-jacobi,claim:curvature-symmetries-bianchi}
\lean{MorganTianLib.curvatureFormAt_chartFrame,MorganTianLib.chartCurvature_basis,MorganTianLib.exists_chartFrame_leviCivita_christoffel_nhds,MorganTianLib.exists_contMDiff_eventuallyEq,MorganTianLib.cov_congr_apply_right,MorganTianLib.alg_curvature_le_of_sectionalCurvature_le,MorganTianLib.chartCurvature_pairing_le_of_sectionalCurvatureAt_le,MorganTianLib.chartCurvature_pairing_le_of_sectionalCurvatureAt_le'}
\leanok
Let $(U,x)$ be a coordinate chart of $(M,g)$, $p\in U$, and let
$G_{ml}=g(\partial_m,\partial_l)$ be the Gram matrix of the
coordinate frame. Let ${\mathcal R}$ be the coefficient curvature of
Lemma~\ref{lem:covariant-commutation-jacobi} built from the
Christoffel symbols $\Gamma^k_{ij}$ of $g$ in this chart, with
components
$${\mathcal R}^m_{ijk}
=\partial_i\Gamma^m_{jk}-\partial_j\Gamma^m_{ik}
+\sum_r\bigl(\Gamma^r_{jk}\Gamma^m_{ir}
-\Gamma^r_{ik}\Gamma^m_{jr}\bigr).$$
Then:
\begin{enumerate}
\item ({\em The bridge.}) For all coordinate vectors
$v,w,z,t\in\mathbb{R}^n$, the Riemann curvature $(0,4)$-tensor
${\mathcal R}iem$ of Definition~\ref{def:riemann-curvature-tensor}
satisfies, at $p$,
$${\mathcal R}iem\Bigl(\sum_a v^a\partial_a,\ \sum_b w^b\partial_b,\
\sum_c z^c\partial_c,\ \sum_d t^d\partial_d\Bigr)
\;=\;-\bigl\langle{\mathcal R}(v,w)z,\;t\bigr\rangle_G
\;=\;-\sum_{m,l}{\mathcal R}^m_{\,\cdot}v^iw^jz^k\,G_{ml}\,t^l,$$
where $\langle a,b\rangle_G=\sum_{m,l}G_{ml}a^mb^l$ is the Gram
pairing and the middle expression contracts
${\mathcal R}(v,w)z=\sum_m{\mathcal R}^m_{ijk}v^iw^jz^k\,e_m$.
\item ({\em The sectional bound.}) If $K(P)\le K$ for every
$2$-plane $P\subset T_pM$, then for all $J,u\in\mathbb{R}^n$
$$\bigl\langle{\mathcal R}(J,u)u,\;J\bigr\rangle_G\ \le\
K\bigl(\langle J,J\rangle_G\langle u,u\rangle_G
-\langle J,u\rangle_G^2\bigr),$$
and if moreover $K\ge0$ then
$\langle{\mathcal R}(J,u)u,J\rangle_G\le
K\,\langle J,J\rangle_G\,\langle u,u\rangle_G$.
\end{enumerate}
\end{lemma}

\begin{proof}
\leanok
\uses{thm:levi-civita-connection,claim:curvature-symmetries-bianchi,def:sectional-curvature,lem:covariant-commutation-jacobi,def:riemann-curvature-tensor}
(1) Both sides are multilinear in $(v,w,z,t)$ — the left by the
tensoriality of the curvature tensor
(Definition~\ref{def:riemann-curvature-tensor}: the value at $p$
depends only on the values of the arguments at $p$, and is linear in
each), the right because ${\mathcal R}(v,w)z$ is built from the
bilinear maps $\Gamma$, the directional derivative
$(\partial_v\Gamma)$, and compositions thereof, all linear in each
slot, and $\langle\cdot,\cdot\rangle_G$ is bilinear. It therefore
suffices to prove the identity on the coordinate frame
$v=e_i,\,w=e_j,\,z=e_k,\,t=e_l$.

Extend the coordinate fields to global smooth vector fields
$X_1,\ldots,X_n$ agreeing with $\partial_1,\ldots,\partial_n$ on a
neighbourhood $V\subseteq U$ of $p$ (multiply by a bump function
equal to $1$ near $p$ and supported in $U$). Since the covariant
derivative and the Lie bracket at a point depend only on the fields
near that point, we may compute with the $X_a$ as if they were the
coordinate fields: on $V$,
$$[X_i,X_j]=0,\qquad
\nabla_{X_i}X_j=\sum_m\Gamma^m_{ij}\,X_m,$$
the first because coordinate fields commute, the second by the
characterization of the Levi-Civita connection through the
Christoffel symbols (Theorem~\ref{thm:levi-civita-connection},
Equation~\eqref{Gamma}). Differentiating the second identity along
$X_j$ at $p$ and using the Leibniz rule,
$$\nabla_{X_j}\nabla_{X_i}X_k
=\sum_m\bigl(\partial_j\Gamma^m_{ik}\bigr)X_m
+\sum_m\Gamma^m_{ik}\,\nabla_{X_j}X_m
=\sum_m\Bigl(\partial_j\Gamma^m_{ik}
+\sum_r\Gamma^r_{ik}\Gamma^m_{jr}\Bigr)X_m
\quad\text{at }p,$$
where the germ-locality of $\nabla$ in its differentiated slot
justifies substituting the local formula before differentiating.
Hence, with the convention of
Definition~\ref{def:riemann-curvature-tensor} and vanishing bracket,
$${\mathcal R}iem(X_i,X_j)X_k
=\nabla_{X_i}\nabla_{X_j}X_k-\nabla_{X_j}\nabla_{X_i}X_k
=\sum_m{\mathcal R}^m_{ijk}\,X_m\quad\text{at }p,$$
by the very definition of ${\mathcal R}^m_{ijk}$. Pairing with
$X_l$ and using
${\mathcal R}iem(X,Y,Z,W)=g({\mathcal R}iem(X,Y)W,Z)
=-g({\mathcal R}iem(X,Y)Z,W)$ (the antisymmetry of
Claim~\ref{claim:curvature-symmetries-bianchi} in the last two
slots) gives
${\mathcal R}iem(e_i,e_j,e_k,e_l)
=-\sum_m{\mathcal R}^m_{ijk}G_{ml}
=-\langle{\mathcal R}(e_i,e_j)e_k,e_l\rangle_G$, as claimed.

(2) Write $B(v,w,z,t)$ for the left side of (1), an algebraic
curvature form on $(T_pM,g_p)$ by
Claim~\ref{claim:curvature-symmetries-bianchi}. By (1) and the
antisymmetry in the last two slots,
$$\langle{\mathcal R}(J,u)u,J\rangle_G=-B(J,u,u,J)=B(J,u,J,u).$$
If $J,u$ span a $2$-plane $P$, then by
Definition~\ref{def:sectional-curvature}
$$B(J,u,J,u)=K(P)\cdot
\bigl(\langle J,J\rangle\langle u,u\rangle-\langle J,u\rangle^2\bigr)
\le K\bigl(\langle J,J\rangle\langle u,u\rangle
-\langle J,u\rangle^2\bigr),$$
since the wedge factor is positive by Cauchy--Schwarz (with equality
iff dependent). If $J,u$ are linearly dependent, both sides vanish:
the wedge factor vanishes by the equality case of Cauchy--Schwarz,
and $B(J,u,J,u)=0$ because either $u=0$, or $J=au$ and then
$B(au,u,au,u)=a\,B(u,u,au,u)=0$ since $B(u,u,\cdot,\cdot)=0$ by
antisymmetry in the first two slots. Here the inner products are
those of $g_p$, i.e.\ exactly the Gram pairings
$\langle\cdot,\cdot\rangle_G$ in coordinates. Finally, for $K\ge0$
discard the nonnegative subtracted term:
$K(\langle J,J\rangle_G\langle u,u\rangle_G-\langle J,u\rangle_G^2)
\le K\langle J,J\rangle_G\langle u,u\rangle_G$.
\end{proof}

The Jacobi equation is best analyzed in a parallel frame along the
geodesic, whose existence and rigidity are recorded in the next
lemma.

\begin{lemma}[Parallel transport and parallel frames]
\label{lem:parallel-frame}
\uses{def:riemannian-metric,thm:levi-civita-connection}
\lean{MorganTianLib.IsParallelSolOn,MorganTianLib.exists_parallelFrame_Icc,MorganTianLib.IsParallelSolOn.eqOn_of_left,MorganTianLib.IsParallelSolOn.chartMetricInner_eq,MorganTianLib.chartMetricInner_parallelFrame_eq}
\leanok
Let $(M,g)$ be a Riemannian manifold with Levi-Civita connection
$\nabla$, and let $\gamma\colon[a,b]\to M$ be a smooth curve whose
image lies in the domain of a single coordinate chart. A vector field
$V$ along $\gamma$ is {\em parallel} if $\nabla_{\gamma'}V=0$ on
$[a,b]$ (one-sided derivatives at the endpoints). Then:
\begin{enumerate}
\item for every $v\in T_{\gamma(a)}M$ there is a unique parallel
field $V$ along $\gamma$ with $V(a)=v$; more generally, for any
family $(v_i)_i$ of vectors in $T_{\gamma(a)}M$ there is a family
$(e_i)_i$ of parallel fields with $e_i(a)=v_i$;
\item the inner product of two parallel fields is constant:
$\langle V(t),W(t)\rangle_{\gamma(t)}
=\langle V(a),W(a)\rangle_{\gamma(a)}$ for all $t\in[a,b]$;
\item consequently a parallel family whose initial values form an
orthonormal basis of $T_{\gamma(a)}M$ remains an orthonormal basis
of $T_{\gamma(t)}M$ for every $t$.
\end{enumerate}
\end{lemma}

\begin{proof}
\uses{thm:levi-civita-connection}
\leanok
In the coordinates of the chart, with $u$ the coordinate image of
$\gamma$, a field $V$ is parallel if and only if it solves the
first-order linear system $V'=-\Gamma(u',V)(u)$, where
$\Gamma$ is the Christoffel contraction of the Levi-Civita
connection (Theorem~\ref{thm:levi-civita-connection}); the
coefficient $A(t)=-\Gamma(u'(t),\cdot)(u(t))$ is linear in $V$,
continuous in $t$ and, the interval being compact, bounded, say
$\|A\|\le K$. Existence with prescribed value $V(a)=v$ follows from
the Picard--Lindel\"of theorem applied on successive subintervals of
length at most $1/(2K)$, glued finitely many times; a family of
initial vectors is handled componentwise. Uniqueness follows from
Gr\"onwall's inequality: the difference $Z$ of two solutions
satisfies $\|Z'\|=\|A(t)Z\|\le K\|Z\|$ and $Z(a)=0$, hence
$\|Z(t)\|\le\|Z(a)\|e^{K(t-a)}=0$. (Along a compact curve crossing
several charts, cover $[a,b]$ by finitely many subintervals each
mapped into one chart; solutions in overlapping charts agree on
overlaps by this uniqueness, so they patch to a parallel field
along the whole curve, with the same uniqueness and inner-product
constancy. This patched form is what is used when the curve is a
whole geodesic, as in
Lemma~\ref{lem:exponential-differential-jacobi}.)

For (2), metric compatibility of the Levi-Civita connection gives,
at each interior time,
$$\frac{d}{dt}\langle V,W\rangle
=\langle\nabla_{\gamma'}V,W\rangle
+\langle V,\nabla_{\gamma'}W\rangle=0,$$
and $t\mapsto\langle V(t),W(t)\rangle$ is continuous on $[a,b]$, so
it is constant on $[a,b]$. For (3), applying (2) to each pair
$(e_i,e_j)$ shows that the Gram matrix
$\langle e_i(t),e_j(t)\rangle$ is constant in $t$; initial
orthonormality is therefore preserved, and an orthonormal family of
$n=\dim M$ vectors is automatically a basis.
\end{proof}

The analytic content of the theory of Jacobi fields is the following
standard lemma about second-order linear equations, which we state
over an arbitrary real Banach space; it will be applied both to
vector-valued solutions (Jacobi fields, in a parallel frame along the
geodesic) and to operator-valued ones (matrices of Jacobi fields, as
in Lemma~\ref{lem:geodesic-polar-form}).

\begin{lemma}[Second-order linear ODE: existence, uniqueness,
superposition]
\label{lem:second-order-linear-ode}
\lean{MorganTianLib.IsJacobiSolOn,MorganTianLib.exists_isJacobiSolOn_Icc,MorganTianLib.IsJacobiSolOn.eqOn_of_left,MorganTianLib.IsJacobiSolOn.eqOn_of_right,MorganTianLib.IsJacobiSolOn.add,MorganTianLib.IsJacobiSolOn.const_smul,MorganTianLib.IsJacobiSolOn.sub,MorganTianLib.IsJacobiSolOn.eqOn_zero}
\leanok
Let $F$ be a real Banach space, let $a\le b$, and let
$R\colon[a,b]\to L(F)$ be a continuous family of bounded operators.
Say that a pair of curves $y,v\colon[a,b]\to F$ {\em solves the
second-order linear equation} $y''+R(t)\,y=0$ on $[a,b]$ if $y'=v$
and $v'=-R(t)\,y$ on $[a,b]$ (one-sided derivatives at the
endpoints). Then:
\begin{enumerate}
\item (Existence) for every $(y_0,v_0)\in F\times F$ there is a
solution with $y(a)=y_0$ and $v(a)=v_0$.
\item (Uniqueness) two solutions that agree at $a$ (both in $y$ and
in $v$) agree on all of $[a,b]$; likewise two solutions that agree
at $b$.
\item (Superposition) sums, differences and scalar multiples of
solutions are solutions; in particular the solution with $y(a)=0$,
$v(a)=0$ is identically zero, so a solution with nontrivial initial
data does not vanish identically.
\end{enumerate}
\end{lemma}

\begin{proof}
\leanok
On the product space $F\times F$ with the norm
$\|(y,v)\|=\max(\|y\|,\|v\|)$, define $A(t)(y,v)=(v,-R(t)\,y)$. Then
$A(t)$ is a bounded operator with $\|A(t)\|\le\max(1,\|R(t)\|)$:
indeed
$\|A(t)(y,v)\|=\max(\|v\|,\|R(t)\,y\|)
\le\max(1,\|R(t)\|)\,\|(y,v)\|$. The map $t\mapsto A(t)$ is
continuous since $R$ is, and $(y,v)$ solves the second-order
equation if and only if the curve $W=(y,v)$ solves the first-order
linear system $W'=A(t)\,W$ on $[a,b]$.

By continuity of $R$ on the compact interval $[a,b]$ there is $C$
with $\|R\|\le C$ there, so $\|A\|\le K:=\max(1,C)$ on $[a,b]$.
Existence for $W'=A(t)\,W$ with prescribed $W(a)$ follows from the
Picard--Lindel\"of theorem: on any subinterval of length at most
$1/(2K)$ the standard contraction argument produces a solution with
prescribed left-endpoint value, and finitely many such solutions
glue to a solution on all of $[a,b]$. Forward and backward
uniqueness follow from Gr\"onwall's inequality: if $W_1,W_2$ are
two solutions and $Z=W_1-W_2$, then $\|Z'\|=\|A(t)Z\|\le K\|Z\|$,
so $\|Z(t)\|\le\|Z(a)\|\,e^{K(t-a)}$ going forward and
$\|Z(t)\|\le\|Z(b)\|\,e^{K(b-t)}$ going backward (for the latter
apply the forward estimate to the time-reversed pair, which solves
the system with coefficient $-A(b+a-t)$); either vanishes when the
two solutions agree at the corresponding endpoint. Superposition is
immediate from the linearity of the equation and of the derivative;
the zero pair is a solution, so a solution with vanishing initial
data coincides with it by uniqueness.
\end{proof}

\subsection{Minimal geodesics}

\begin{definition}[Conjugate point]
\label{def:conjugate-point}
\uses{def:geodesic,lem:jacobi-field-coordinates}
\lean{MorganTianLib.chartVectorRep,MorganTianLib.IsJacobiFieldAlongOn,MorganTianLib.IsConjugatePointAt}
\leanok
Let $\gamma$ be a geodesic beginning at $p\in M$. A {\sl Jacobi field
along $\gamma$} is a field of tangent vectors $J(t)\in T_{\gamma(t)}M$
along $\gamma$, together with its covariant derivative
$\nabla J(t)\in T_{\gamma(t)}M$, such that near every time, read in the
coordinates of a chart containing the nearby piece of $\gamma$, the pair
$(J,\nabla J)$ satisfies the chart Jacobi system of
Lemma~\ref{lem:jacobi-field-coordinates} — equivalently, $J$ satisfies
the Jacobi equation
$\nabla_{\dot\gamma}\nabla_{\dot\gamma}J
+{\mathcal R}(J,\dot\gamma)\dot\gamma=0$. The notion is chart-local, so
it is meaningful for geodesics that leave any single chart. For any
$t>0$ we say that $q=\gamma(t)$ is a {\sl conjugate point along
$\gamma$} if there is a non-zero Jacobi field
along $\gamma$ vanishing at $p$ and at $\gamma(t)$. (Morgan--Tian leave
the vanishing at $p$ implicit in the phrase ``geodesic beginning at
$p$''; their subsequent arguments use it — e.g.\ the proofs of
Proposition~\ref{prop:minimal-geodesic-no-conjugate} and
Lemma~\ref{lem:conjugate-sturm} — and we make it explicit.)
\end{definition}


\begin{proposition}[Minimal geodesics have no conjugate points]
\label{prop:minimal-geodesic-no-conjugate}
\label{jacmin}
\uses{def:conjugate-point}
Suppose that $\gamma\colon [0,1]\to M$ is a minimal geodesic. Then
for any $t<1$ the restriction of $\gamma$ to $[0,t]$ is the unique
minimal geodesic between its endpoints and there are no conjugate
points on $\gamma([0,1))$, i.e., there is no non-zero Jacobi field
along $\gamma$ vanishing at any $t\in [0,1)$.
\end{proposition}

\begin{proof}
\uses{def:geodesic,def:conjugate-point,claim:second-variation-minimal-geodesic,def:exponential-map,claim:curvature-symmetries-bianchi}
Write $p=\gamma(0)$, $q=\gamma(1)$ and $L=L(\gamma)=d(p,q)$; since
$\gamma$ is a geodesic its speed is constant, $|\gamma'(t)|\equiv
L$. We use throughout that for a piecewise smooth path
$\sigma$ parameterized by $[0,1]$ one has $E(\sigma)\ge
\frac{1}{2}L(\sigma)^2$, with equality if and only if $|\sigma'|$ is
constant, and that consequently $E(\gamma)=\frac12 L^2\le
\frac12 L(\sigma)^2\le E(\sigma)$ for every piecewise smooth path
$\sigma$ from $p$ to $q$: that is, a minimal geodesic minimizes
energy among all piecewise smooth paths with the same endpoints,
parameterized by $[0,1]$. (The first variation and energy
computations earlier in this section extend to piecewise smooth
paths and variations by applying them on each smooth piece and
summing.)

\emph{Part 1: for $t_0<1$ the restriction $\gamma|_{[0,t_0]}$ is the
unique minimal geodesic between its endpoints.} Fix $0<t_0<1$.
First, $\gamma|_{[0,t_0]}$ is minimal: if some rectifiable curve
$\sigma$ from $\gamma(0)$ to $\gamma(t_0)$ had
$L(\sigma)<L(\gamma|_{[0,t_0]})$, then the concatenation of $\sigma$
with $\gamma|_{[t_0,1]}$ would be a curve from $p$ to $q$ of length
$L(\sigma)+(1-t_0)L<L$, contradicting $L=d(p,q)$.

Now suppose that $\mu\colon [0,t_0]\to M$ is a minimal geodesic from
$\gamma(0)$ to $\gamma(t_0)$ distinct from $\gamma|_{[0,t_0]}$. Being
minimal geodesics between the same endpoints, both have length
$d(\gamma(0),\gamma(t_0))=t_0L$, and both have constant speed, equal
to $L$. If we had $\mu'(t_0)=\gamma'(t_0)$, then $\mu$ and
$\gamma|_{[0,t_0]}$ would be two solutions of the (second-order)
geodesic equation agreeing in position and velocity at $t=t_0$, and
by the uniqueness statement of classical ODE theory they would
coincide, contrary to the assumption that they are distinct. Hence
$$\Delta:=\gamma'(t_0)-\mu'(t_0)\not=0.$$
Let $c\colon [0,1]\to M$ be the concatenation, equal to $\mu$ on
$[0,t_0]$ and to $\gamma$ on $[t_0,1]$. Then $c$ is a piecewise
smooth path from $p$ to $q$ of constant speed $L$, so
$$E(c)=\frac{1}{2}L^2=\frac{1}{2}d(p,q)^2,$$
i.e., $c$ minimizes energy among piecewise smooth paths from $p$ to
$q$. Choose a piecewise smooth vector field $Y$ along $c$ with
$Y(0)=0$, $Y(1)=0$ and $Y(t_0)=\Delta$ (for instance $\phi(t)V(t)$
with $V$ the parallel field along each piece with
$V(t_0)=\Delta$ and $\phi$ a smooth bump function vanishing at $0,1$
with $\phi(t_0)=1$), and consider the variation
$\tilde c(t,u)=\exp_{c(t)}(uY(t))$, defined for $|u|$ small. On each
of the two smooth pieces the first-variation computation from the
beginning of this section applies: writing $\widetilde X=\partial
\tilde c/\partial t$ and $\widetilde Y=\partial\tilde c/\partial u$,
and using $\nabla_{\widetilde Y}\widetilde X=\nabla_{\widetilde
X}\widetilde Y$ and integration by parts on each piece,
$$\frac{dE(c_u)}{du}\Bigl|_{u=0}\Bigr.
=\left(\int_0^{t_0}+\int_{t_0}^1\right)
\langle\nabla_{X}Y,X\rangle\,dt
=\Bigl[\langle Y,X\rangle\Bigr]_0^{t_0^-}
+\Bigl[\langle Y,X\rangle\Bigr]_{t_0^+}^{1}
-\int_0^1\langle Y,\nabla_XX\rangle\,dt.$$
The integral vanishes since both pieces of $c$ are geodesics, and
the boundary terms give
$$\frac{dE(c_u)}{du}\Bigl|_{u=0}\Bigr.
=\langle Y(t_0),\mu'(t_0)\rangle-\langle
Y(t_0),\gamma'(t_0)\rangle=-\langle \Delta,\Delta\rangle
=-|\Delta|^2<0.$$
Hence $E(c_u)<E(c)=\frac{1}{2}d(p,q)^2$ for all sufficiently small
$u>0$, and therefore $L(c_u)\le\sqrt{2E(c_u)}<d(p,q)$ for such $u$
--- a contradiction, since $c_u$ is a path from $p$ to $q$. Thus no
such $\mu$ exists, proving Part 1.

For the second assertion of the proposition we shall use the
following claim, together with a refinement of the second-variation
formula valid for piecewise smooth variations.

\begin{lemma}[Second variation is zero iff Jacobi field]
\label{claim:second-variation-minimal-geodesic}
\uses{def:geodesic}
Suppose that $\gamma$ is a minimal geodesic and $Y$ is a field
vanishing at both endpoints. Let $\tilde\gamma(t,u)$ be any
one-parameter family of curves parameterized by $[0,1]$, with
$\gamma_0=\gamma$ and with $\gamma_u(0)=\gamma_0(0)$ and
$\gamma_u(1)=\gamma_0(1)$ for all $u$.
Suppose that the first-order variation of $\tilde \gamma$ at $u=0$
is given by $Y$. Then
$$\frac{d^2E(\gamma_u)}{du^2}\Bigl|_{u=0}\Bigr.=0$$
if and only if $Y$ is a Jacobi field.
\end{lemma}

\begin{proof}
\uses{def:geodesic}
Suppose that $\tilde\gamma(u,t)$ is a one-parameter family of curves
from $\gamma(0)$ to $\gamma(1)$ with $\gamma_0=\gamma$ and $Y$ is
the first-order variation of this family along $\gamma$.  Since
$\gamma$ is a minimal geodesic we have
$$-\int_0^1\langle \Jac(Y),Y\rangle dt=\frac{d^2E(\gamma_u)}{du^2}\Bigl|_{u=0}\Bigr.\ge 0.$$
 The associated symmetric bilinear form is
$$B_\gamma(Y_1,Y_2)=-\int_\gamma\langle \Jac(Y_1),Y_2\rangle dt$$
is symmetric when $Y_1$ and $Y_2$ are constrained to vanish at both
endpoints. Since the associated quadratic form is non-negative, we
see by the usual argument for symmetric bilinear forms that
$B_\gamma(Y,Y)=0$ if and only if $B_\gamma(Y,\cdot)=0$ as a linear
functional on the space of vector fields along $\gamma$ vanishing at
point endpoints. This of course occurs if and only if $\Jac(Y)=0$.
\end{proof}


\emph{Part 2: there are no conjugate points on $\gamma([0,1))$.}
For a piecewise smooth vector field $Y$ along $\gamma$ vanishing at
both endpoints, define the \emph{index form}
$$I(Y_1,Y_2)=\int_0^1\left(\langle\nabla_XY_1,\nabla_XY_2\rangle
-\langle{\mathcal R}(Y_1,X)X,Y_2\rangle\right)dt,$$
the integral being taken piecewise; $I$ is a symmetric bilinear form
on the vector space of such fields (symmetry of the curvature term
is the pair symmetry of the curvature tensor). For \emph{smooth}
$Y_1,Y_2$ vanishing at the endpoints, integration by parts gives
$I(Y_1,Y_2)=-\int_0^1\langle \Jac(Y_1),Y_2\rangle
dt=B_\gamma(Y_1,Y_2)$, the second-variation form of
Claim~\ref{claim:second-variation-minimal-geodesic}.

We first extend the second-variation formula: for every piecewise
smooth $Y$ along $\gamma$ vanishing at both endpoints, the variation
$\tilde\gamma(t,u)=\exp_{\gamma(t)}(uY(t))$ satisfies
\begin{equation}\label{secvarindex}
\frac{d^2E(\gamma_u)}{du^2}\Bigl|_{u=0}\Bigr.=I(Y,Y).
\end{equation}
Indeed, with $\widetilde X=\partial\tilde\gamma/\partial t$ and
$\widetilde Y=\partial\tilde\gamma/\partial u$ we have, on each
smooth piece,
$$\frac{d^2E(\gamma_u)}{du^2}
=\frac{d}{du}\int_0^1\langle\nabla_{\widetilde X}\widetilde
Y,\widetilde X\rangle\,dt
=\int_0^1\left(\langle\nabla_{\widetilde Y}\nabla_{\widetilde
X}\widetilde Y,\widetilde X\rangle
+|\nabla_{\widetilde X}\widetilde Y|^2\right)dt,$$
using $\nabla_{\widetilde Y}\widetilde X=\nabla_{\widetilde
X}\widetilde Y$ twice. Commuting the covariant derivatives,
$\nabla_{\widetilde Y}\nabla_{\widetilde X}\widetilde
Y=\nabla_{\widetilde X}\nabla_{\widetilde Y}\widetilde Y+{\mathcal
R}(\widetilde Y,\widetilde X)\widetilde Y$. For each fixed $t$ the
curve $u\mapsto\exp_{\gamma(t)}(uY(t))$ is a geodesic with velocity
field $\widetilde Y$, so $\nabla_{\widetilde Y}\widetilde Y=0$
identically. Evaluating at $u=0$ (where $\widetilde X=X=\gamma'$ and
$\widetilde Y=Y$) and using
$\langle{\mathcal R}(Y,X)Y,X\rangle=-\langle{\mathcal
R}(Y,X)X,Y\rangle$, which follows from the symmetries of the
curvature tensor, we obtain
Equation~(\ref{secvarindex}). Moreover, since $\gamma$ minimizes
energy among piecewise smooth paths from $p$ to $q$ and
$\gamma_u$ is such a path for each small $u$, the function
$u\mapsto E(\gamma_u)$ has a minimum at $u=0$, so
\begin{equation}\label{indexnonneg}
I(Y,Y)\ge 0
\end{equation}
for every piecewise smooth $Y$ along $\gamma$ vanishing at both
endpoints. By Equations~(\ref{secvarindex}) and
(\ref{indexnonneg}) and the Cauchy--Schwarz inequality for
positive semi-definite symmetric bilinear forms, if $I(Y,Y)=0$ then
$I(Y,Z)=0$ for every piecewise smooth $Z$ vanishing at both
endpoints. (For smooth fields, Equation~(\ref{secvarindex})
together with
Claim~\ref{claim:second-variation-minimal-geodesic} recovers the
statement that the null space of $I$ on smooth fields is the space
of Jacobi fields.)

Now suppose that for some $0<t_0<1$ the point $\gamma(t_0)$ is a
conjugate point along $\gamma$, so that there is a non-zero Jacobi
field $Y$ along $\gamma|_{[0,t_0]}$ with $Y(0)=Y(t_0)=0$. Note that
$\nabla_XY(t_0)\not=0$: the Jacobi equation is a second-order linear
ODE for $Y$ along $\gamma$, so if both $Y(t_0)=0$ and
$\nabla_XY(t_0)=0$ then $Y\equiv 0$ by ODE uniqueness. Extend $Y$ to
a field $\hat Y$ along all of $\gamma$ by setting $\hat Y=0$ on
$[t_0,1]$; then $\hat Y$ is piecewise smooth (smooth except possibly
at $t_0$, where it is continuous) and vanishes at both endpoints.
On $[0,t_0]$, integration by parts and $\Jac(Y)=0$ give
$$\int_0^{t_0}\left(|\nabla_XY|^2-\langle{\mathcal
R}(Y,X)X,Y\rangle\right)dt
=\Bigl[\langle\nabla_XY,Y\rangle\Bigr]_0^{t_0}
-\int_0^{t_0}\langle \Jac(Y),Y\rangle\,dt=0,$$
the boundary terms vanishing because $Y(0)=Y(t_0)=0$. Since $\hat
Y\equiv 0$ on $[t_0,1]$, we conclude $I(\hat Y,\hat Y)=0$, and hence
by the previous paragraph $I(\hat Y,Z)=0$ for every piecewise smooth
$Z$ vanishing at the endpoints. On the other hand, computing
$I(\hat Y,Z)$ by integration by parts on the two pieces, and using
$\Jac(\hat Y)=0$ on each piece,
$$I(\hat Y,Z)
=\Bigl[\langle\nabla_X\hat Y,Z\rangle\Bigr]_0^{t_0^-}
+\Bigl[\langle\nabla_X\hat Y,Z\rangle\Bigr]_{t_0^+}^{1}
-\int_0^1\langle \Jac(\hat Y),Z\rangle\,dt
=\langle\nabla_XY(t_0),Z(t_0)\rangle,$$
since $Z(0)=Z(1)=0$ and $\nabla_X\hat Y=0$ on $(t_0,1]$. Choosing
$Z$ piecewise smooth, vanishing at the endpoints, with
$Z(t_0)=\nabla_XY(t_0)$ (a bump function times a parallel field, as
in Part 1), we get $I(\hat Y,Z)=|\nabla_XY(t_0)|^2>0$, a
contradiction. Hence there is no conjugate point $\gamma(t_0)$ with
$t_0<1$. (There is no conjugate point at $t_0=0$ by definition,
conjugate points occurring only at parameters $t>0$.)
\end{proof}





\subsection{The exponential mapping}


\begin{definition}[Exponential map]
\label{def:exponential-map}
\uses{def:geodesic}
\lean{MorganTianLib.expMap,MorganTianLib.expDomain}
\leanok
For any $p\in M$, we can define the \textit{exponential
map} at $p$, $\exp_{p}$. It is
defined on an open neighborhood $O_{p}$ of the origin in $T_{p}M$
and is defined by $\exp_{p}(v)=\gamma_v(1)$, the endpoint of
the unique geodesic $\gamma_v\colon [0,1]\rightarrow M$ starting
from $p$ with initial velocity vector $v$. We always take
$O_{p}\subset T_{p}M$ to be the maximal domain on which $\exp_p$ is defined, so that $O_{p}$ is a star-shaped open
neighborhood of $0\in T_pM$.
\end{definition}

By the Hopf--Rinow Theorem (Theorem~\ref{thm:hopf-rinow} below),
if $M$ is complete then the exponential map is defined on all of
$T_pM$.

Near the origin of $T_pM$ the exponential map is a diffeomorphism
onto a metric ball, uniformly over compact sets of centers. (The
forward reference to the Gauss Lemma,
Lemma~\ref{lem:geodesic-polar-form} in the section on curvature
comparison below, involves no circularity: that lemma depends only
on the material of the present section and of
Section~2.)

\begin{lemma}[Normal balls, uniformly on compact sets]
\label{lem:normal-ball}
\uses{def:exponential-map,def:geodesic}
Let $(M,g)$ be a Riemannian manifold, not assumed complete, and let
$K\subset M$ be compact. There is $\delta=\delta(K)>0$ such that
for every $x\in K$ the exponential map $\exp_x$ is defined on
$B(0,\delta)\subset T_xM$ and:
\begin{enumerate}
\item $\exp_x$ is a diffeomorphism from $B(0,\delta)$ onto the
metric ball $B(x,\delta)$, and $\exp_x(B(0,s))=B(x,s)$ for every
$0<s\le\delta$;
\item $d(x,\exp_x(v))=|v|$ for all $|v|<\delta$;
\item every piecewise smooth curve from $x$ to $\exp_x(v)$,
$|v|<\delta$, has length $\ge|v|$, with equality if and only if it
is a monotone reparameterization of the radial segment
$t\mapsto\exp_x(tv)$, $t\in[0,1]$; and every piecewise smooth
curve starting at $x$ and not staying in $B(x,s)$, $0<s\le\delta$,
has length $\ge s$.
\end{enumerate}
\end{lemma}

\begin{proof}
\uses{lem:geodesic-polar-form,def:exponential-map,def:geodesic}
Throughout, $|v|=\sqrt{g_x(v,v)}$ for $v\in T_xM$, and a monotone
reparameterization of the radial segment means a curve of the form
$t\mapsto\exp_x(u(t)v)$ with $u\colon[0,1]\to[0,1]$ continuous and
monotone. Since the distance $d$ is the infimum of the lengths of
smooth paths, and smooth paths are in particular piecewise smooth,
any lower bound for the lengths of all piecewise smooth paths
between two points is also a lower bound for their distance.

\emph{Step 1: a uniform radius $\delta_1$.} In local coordinates
the geodesic equation is a second-order ODE system with smooth
coefficients, so by the classical theorems on local existence,
uniqueness and smooth dependence on initial conditions, every
point of $M$ has an open neighborhood $U_0$ for which there are
$\epsilon_0,\eta_0>0$ such that for every $y\in U_0$ and $v\in
T_yM$ with $|v|<\epsilon_0$ the geodesic $\gamma_v$ with
$\gamma_v(0)=y$ and $\gamma_v'(0)=v$ is defined for $|t|<\eta_0$,
and $(y,v,t)\mapsto\gamma_v(t)$ is smooth. For $s>0$ the curve
$t\mapsto\gamma_v(st)$ is again a geodesic, with initial velocity
$sv$, so by uniqueness of solutions
$$\gamma_{sv}(t)=\gamma_v(st)$$
whenever the right-hand side is defined. Replacing $\epsilon_0$ by
$\epsilon_0\eta_0/2$ and using this homogeneity with $s=\eta_0/2$,
we may assume the geodesics above are defined for $|t|<2$. Hence
$\exp_y(v)=\gamma_v(1)$ is defined for $y\in U_0$ and
$|v|<\epsilon_0$, and the map
$$E(y,v)=(y,\exp_y(v))$$
is defined and smooth on an open neighborhood of the zero section
of $TM$, with values in $M\times M$.

Fix $x\in M$ and compute $dE$ at $(x,0)$ in a chart trivializing
$TM$ near $x$. Since $\exp_y(0)=y$ for every $y$, the derivative
of $y\mapsto E(y,0)=(y,y)$ is $\xi\mapsto(\xi,\xi)$; and since
$\exp_x(t\zeta)=\gamma_{t\zeta}(1)=\gamma_\zeta(t)$ by homogeneity,
we get $d(\exp_x)_0=\mathrm{Id}$, so the derivative of $E$ in the
fiber direction at $(x,0)$ is $\zeta\mapsto(0,\zeta)$. Thus
$dE_{(x,0)}(\xi,\zeta)=(\xi,\xi+\zeta)$, which is invertible. By
the inverse function theorem there is an open neighborhood $W_x$
of $(x,0)$ in $TM$ on which $E$ is a diffeomorphism onto an open
subset of $M\times M$. Shrinking, we may take $W_x$ to be a tube
$\{(y,v)\,|\,y\in U_x,\ v\in T_yM,\ |v|<\epsilon_x\}$ with open
$U_x\ni x$: in a local trivialization of $TM$ over a chart, the
sets $\{y\in U',\ |v|_e<\epsilon'\}$, with $|\cdot|_e$ a Euclidean
fiber norm, form a neighborhood basis of $(x,0)$, and on a compact
neighborhood of $x$ one has $|v|_e\le C\,|v|$ by continuity and
compactness, so the corresponding $g$-tubes form a neighborhood
basis as well.

Cover the compact set $K$ by finitely many sets
$U_{x_1},\ldots,U_{x_N}$ as above and put
$$\delta_1=\min_{1\le i\le N}\epsilon_{x_i}>0.$$
Given $x\in K$, choose $i$ with $x\in U_{x_i}$; then
$\{x\}\times B(0,\delta_1)\subset W_{x_i}$, and we conclude, for
every $x\in K$:
\begin{itemize}
\item $\exp_x$ is defined and smooth on $B(0,\delta_1)\subset
T_xM$;
\item $\exp_x$ is injective on $B(0,\delta_1)$: if
$\exp_x(v)=\exp_x(w)$ with $|v|,|w|<\delta_1$ then
$E(x,v)=E(x,w)$ with both arguments in $W_{x_i}$, so $v=w$;
\item $\exp_x$ is non-singular at every $v\in B(0,\delta_1)$:
restricting $E$ to the fiber through $(x,v)$ shows that
$dE_{(x,v)}$ maps the vertical vector $(0,\xi)$ to
$(0,d(\exp_x)_v\xi)$, so $d(\exp_x)_v\xi=0$ forces $\xi=0$ by
invertibility of $dE_{(x,v)}$.
\end{itemize}
Being an injective local diffeomorphism, $\Phi:=
\exp_x|_{B(0,\delta_1)}$ is an open map, hence a diffeomorphism
from $B(0,\delta_1)$ onto the open set
$\Phi(B(0,\delta_1))\subseteq M$. For $0<\rho\le\delta_1$ write
$U_\rho=\exp_x(B(0,\rho))$; for $\rho<\delta_1$ this is an open
subset of $M$.

\emph{Step 2: length estimates in polar form.} Fix $x\in K$ and
let $\tilde g=\Phi^*g$ on $B(0,\delta_1)$; since $\Phi$ is
non-singular there, $\tilde g$ is a Riemannian metric, and by the
definition of the pullback,
$L(\Phi\circ\tilde c)=L_{\tilde g}(\tilde c)$ for every piecewise
smooth curve $\tilde c$ in $B(0,\delta_1)$, where $L_{\tilde g}$
denotes length with respect to $\tilde g$. Apply
Lemma~\ref{lem:geodesic-polar-form}(1) with $p=x$, with $V$ the
whole unit sphere of $T_xM$ and with $r_0=\delta_1$; the
hypotheses hold by Step 1. Thus around each point of
$B(0,\delta_1)\setminus\{0\}$ there are polar coordinates
$(r,\theta)$, with $r$ the norm function, in which
$$\tilde g=dr^2+g_{ij}(r,\theta)\,d\theta^i\otimes d\theta^j,$$
and the two-tensor $g_{ij}\,d\theta^i\otimes d\theta^j$ is
positive definite on the tangent spaces of the spheres
$\{r=\mathrm{const}\}$, being the restriction of the Riemannian
metric $\tilde g$ to the span of the linearly independent
coordinate fields $\partial_{\theta^i}$.

\emph{Claim A.} Let $\tilde c\colon[a,b]\to
B(0,\delta_1)\setminus\{0\}$ be piecewise smooth and set
$r(t)=|\tilde c(t)|$. Then
$L_{\tilde g}(\tilde c)\ge|r(b)-r(a)|$, with equality if and only
if $\tilde c(t)=r(t)\,w$ for a fixed unit vector $w\in T_xM$ and
$r$ is monotone.

Indeed, $r$ is piecewise smooth (the norm is smooth away from the
origin), and at every point of smoothness, computing in a polar
chart around $\tilde c(t)$,
$$|\tilde c'(t)|_{\tilde g}^2=\dot r(t)^2
+g_{ij}\,\dot\theta^i\dot\theta^j\ \ge\ \dot r(t)^2,$$
with equality at $t$ if and only if the angular velocity vanishes
at $t$. Hence
$$L_{\tilde g}(\tilde c)=\int_a^b|\tilde c'|_{\tilde g}\,dt
\ \ge\ \int_a^b|\dot r|\,dt\ \ge\
\Big|\int_a^b\dot r\,dt\Big|=|r(b)-r(a)|.$$
Suppose equality holds throughout. Then first
$g_{ij}\,\dot\theta^i\dot\theta^j=0$ at every point of smoothness
(a continuous non-negative function with vanishing integral), so
the angular path $w(t)=\tilde c(t)/r(t)$ is locally constant on
each smooth piece, hence constant on $[a,b]$ by continuity; call
it $w$. Second, writing $P=\int_a^b\max(\dot r,0)\,dt$ and
$N=\int_a^b\max(-\dot r,0)\,dt$, equality gives $P+N=|P-N|$, so
$P=0$ or $N=0$; by continuity of $\dot r$ on each smooth piece,
$r$ is monotone. Conversely, if $\tilde c=r(\cdot)\,w$ with $r$
monotone, both inequalities above are equalities. This proves
Claim A.

\emph{Claim B.} Let $\tilde c\colon[a,b]\to B(0,\delta_1)$ be
piecewise smooth with $\tilde c(a)=0$, and set
$\rho=|\tilde c(b)|$. Then $L_{\tilde g}(\tilde c)\ge\rho$.
Moreover, if $\rho>0$ and $L_{\tilde g}(\tilde c)=\rho$, then
$$\tilde c(t)=u(t)\,\tilde c(b),\qquad t\in[a,b],$$
for a continuous non-decreasing $u\colon[a,b]\to[0,1]$ with
$u(a)=0$ and $u(b)=1$.

For the inequality we may assume $\rho>0$. Given
$0<\epsilon<\rho$, the set $\{t\,|\,|\tilde c(t)|=\epsilon\}$ is
non-empty (by the intermediate value theorem) and closed; let
$t_\epsilon$ be its maximum. On $(t_\epsilon,b]$ the continuous
function $|\tilde c|$ never takes the value $\epsilon$ and equals
$\rho>\epsilon$ at $b$, so $|\tilde c|>\epsilon$ there; in
particular $\tilde c\ne0$ on $[t_\epsilon,b]$. Claim A gives
$$L_{\tilde g}(\tilde c)\ \ge\
L_{\tilde g}\big(\tilde c|_{[t_\epsilon,b]}\big)\ \ge\
\rho-\epsilon,$$
and letting $\epsilon\to0$ yields the bound.

Now suppose $L_{\tilde g}(\tilde c)=\rho>0$. The set
$\{t\,|\,\tilde c(t)=0\}$ is closed, non-empty and does not
contain $b$; let $t_*<b$ be its maximum, so $\tilde c\ne0$ on
$(t_*,b]$. Fix $t\in(t_*,b]$ and write $r(t)=|\tilde c(t)|$. The
bound just proved, applied on $[a,t]$, and Claim A, applied on
$[t,b]$, give
$$\rho=L_{\tilde g}\big(\tilde c|_{[a,t]}\big)
+L_{\tilde g}\big(\tilde c|_{[t,b]}\big)\ \ge\
r(t)+|\rho-r(t)|\ \ge\ \rho.$$
Hence $r(t)\le\rho$, both inequalities are equalities, and the
equality case of Claim A holds on $[t,b]$: there $\tilde
c=r(\cdot)\,w$ with $r$ monotone and, evaluating at $b$,
$w=\tilde c(b)/\rho$; since $r(t)\le\rho=r(b)$, $r$ is
non-decreasing on $[t,b]$. As $t\in(t_*,b]$ was arbitrary, on
$(t_*,b]$ we have $\tilde c(t)=r(t)\,\tilde c(b)/\rho$ with $r$
non-decreasing. Finally, letting $t\downarrow t_*$ in
$L_{\tilde g}(\tilde c|_{[a,t]})=r(t)$ and using
$r(t)\to|\tilde c(t_*)|=0$ gives
$L_{\tilde g}(\tilde c|_{[a,t_*]})=0$, so $\tilde c'\equiv0$ on
$[a,t_*]$ and $\tilde c\equiv\tilde c(t_*)=0$ there. Setting
$u=r/\rho$ on $(t_*,b]$ and $u=0$ on $[a,t_*]$ proves Claim B.

\emph{Step 3: exit estimate.} Let $x\in K$, let $0<\rho<\delta_1$,
and let $c\colon[0,1]\to M$ be piecewise smooth with $c(0)=x$ and
$c(t_1)\notin U_\rho$ for some $t_1$. Then $L(c)\ge\rho$.

The set $\{t\,|\,c(t)\notin U_\rho\}$ is closed ($U_\rho$ is open
since $\rho<\delta_1$) and non-empty; let $\tau$ be its minimum.
Since $c(0)=\exp_x(0)\in U_\rho$ we have $\tau>0$, and $c(t)\in
U_\rho$ for $0\le t<\tau$. Because $\rho<\delta_1$, the closed
ball $\overline{B(0,\rho)}\subset T_xM$ is a compact subset of
$B(0,\delta_1)$, so $\exp_x\big(\overline{B(0,\rho)}\big)$ is a
compact, hence closed, subset of $M$ containing $U_\rho$.
Therefore
$$c(\tau)=\lim_{t\to\tau^-}c(t)\in\overline{U_\rho}\subseteq
\exp_x\big(\overline{B(0,\rho)}\big),$$
say $c(\tau)=\exp_x(u)$ with $|u|\le\rho$; and $|u|=\rho$, since
$|u|<\rho$ would give $c(\tau)\in U_\rho$. In particular
$c([0,\tau])\subseteq\Phi(B(0,\delta_1))$, so
$\tilde c:=\Phi^{-1}\circ c|_{[0,\tau]}$ is a piecewise smooth
curve in $B(0,\delta_1)$ with $\tilde c(0)=0$ and
$|\tilde c(\tau)|=\rho$. By Claim B,
$$L(c)\ \ge\ L\big(c|_{[0,\tau]}\big)
=L_{\tilde g}(\tilde c)\ \ge\ \rho.$$

\emph{Step 4: conclusion, with $\delta=\delta_1/2$.} Set
$\delta=\delta_1/2$ and fix $x\in K$. By Step 1, $\exp_x$ is
defined on $B(0,\delta)$ and is a diffeomorphism from $B(0,\delta)$
(indeed from $B(0,\delta_1)$) onto its open image.

\emph{Lower bound, and (2).} Let $|v|<\delta$ and let
$c\colon[0,1]\to M$ be piecewise smooth from $x$ to $\exp_x(v)$.
Fix $\rho$ with $|v|<\rho<\delta_1$. If $c$ does not stay in
$U_\rho$, Step 3 gives $L(c)\ge\rho>|v|$. Otherwise
$c([0,1])\subseteq U_\rho\subseteq\Phi(B(0,\delta_1))$, and
$\tilde c=\Phi^{-1}\circ c$ is piecewise smooth with
$\tilde c(0)=0$ and $\Phi(\tilde c(1))=\exp_x(v)=\Phi(v)$, whence
$\tilde c(1)=v$ by injectivity; Claim B gives
$L(c)=L_{\tilde g}(\tilde c)\ge|v|$. This proves the first
assertion of (3), and it shows $d(x,\exp_x(v))\ge|v|$. For the
reverse inequality consider the radial segment
$\sigma(t)=\exp_x(tv)$, $t\in[0,1]$, a smooth path from $x$ to
$\exp_x(v)$: its lift $t\mapsto tv$ has, for $t>0$, constant
angular coordinates and radial coordinate $r(t)=t|v|$, so
$|\sigma'(t)|=|\dot r(t)|=|v|$ by the polar form of $\tilde g$;
since $|\sigma'|$ is continuous on $[0,1]$, $L(\sigma)=|v|$.
Hence $d(x,\exp_x(v))=|v|$, proving (2).

\emph{(1).} Let $0<s\le\delta$. If $|v|<s$ then
$d(x,\exp_x(v))=|v|<s$ by (2), so
$\exp_x(B(0,s))\subseteq B(x,s)$. Conversely, let $q\in B(x,s)$.
By the definition of the distance there is a smooth path
$c\colon[0,1]\to M$ from $x$ to $q$ with $L(c)<s$; choose $s'$
with $L(c)<s'<s$. Since $0<s'<\delta_1$ and $L(c)<s'$, Step 3
shows that $c$ stays in $U_{s'}$; in particular
$q\in U_{s'}\subseteq\exp_x(B(0,s))$. Hence
$\exp_x(B(0,s))=B(x,s)$ for all $0<s\le\delta$, and taking
$s=\delta$ shows that $\exp_x$ is a diffeomorphism from
$B(0,\delta)$ onto the metric ball $B(x,\delta)$.

\emph{(3), exit part.} Let $0<s\le\delta$ and let $c$ be a
piecewise smooth curve with $c(0)=x$ not staying in
$B(x,s)=U_s$ (the equality by (1)). Since $s\le\delta<\delta_1$,
Step 3 with $\rho=s$ gives $L(c)\ge s$.

\emph{(3), equality case.} If $v=0$ the radial segment is the
constant path at $x$, and a piecewise smooth curve from $x$ to $x$
has length $0$ if and only if its velocity vanishes identically,
i.e., if and only if it is constant; so the assertion holds. Let
$0<|v|<\delta$ and let $c$ be a piecewise smooth curve from $x$ to
$\exp_x(v)$ with $L(c)=|v|$. For $\rho=(|v|+\delta_1)/2$ the
curve stays in $U_\rho$, since otherwise $L(c)\ge\rho>|v|$ by
Step 3. Lifting as above, $\tilde c=\Phi^{-1}\circ c$ satisfies
$\tilde c(0)=0$, $\tilde c(1)=v$ and
$L_{\tilde g}(\tilde c)=|v|$. By the equality case of Claim B,
$\tilde c(t)=u(t)\,v$ with $u\colon[0,1]\to[0,1]$ continuous,
non-decreasing, $u(0)=0$ and $u(1)=1$; that is,
$c(t)=\exp_x(u(t)v)$ is a monotone reparameterization of the
radial segment.

Conversely, suppose the piecewise smooth curve $c$ from $x$ to
$\exp_x(v)$ has the form $c(t)=\exp_x(u(t)v)$ with
$u\colon[0,1]\to[0,1]$ continuous and monotone. Then
$\tilde c(t)=u(t)v=\Phi^{-1}(c(t))$ is piecewise smooth, and hence
so is $u=\langle\tilde c,v\rangle/|v|^2$. From $c(0)=x=\Phi(0)$
and $c(1)=\exp_x(v)=\Phi(v)$, injectivity of $\Phi$ gives $u(0)=0$
and $u(1)=1$; being monotone with $u(0)<u(1)$, $u$ is
non-decreasing, so $\dot u\ge0$ on each smooth piece. Let
$t_0=\max\{t\,|\,u(t)=0\}$; then $u\equiv0$ on $[0,t_0]$ and $u>0$
on $(t_0,1]$. On $[0,t_0]$ the curve is constant and contributes
no length; on each smooth piece contained in $(t_0,1]$ the lift
has constant angular coordinates and radial coordinate $u(t)|v|$,
so $|c'(t)|=\dot u(t)\,|v|$ by the polar form. Summing over the
pieces,
$$L(c)=\int_{t_0}^1\dot u(t)\,|v|\,dt=|v|\,(u(1)-u(t_0))=|v|.$$
This completes the equality case, and the proof.
\end{proof}

In particular, for every $p\in M$ there exists
$r_0=r_{0}(p,M)>0$ such that the restriction of $\exp_p$ to the
ball $B_{g|T_pM}(0,r_0)$ in $T_pM$ is a diffeomorphism onto the
metric ball $B_g(p,r_0)$. Fix $g$-orthonormal linear coordinates
on $T_pM$. Transferring these coordinates via $\exp_p$ to
coordinates on $B(p,r_0)$ gives us {\sl Gaussian normal
coordinates} on $B(p,r_0)\subset M$.

The differential of the exponential mapping is given by Jacobi
fields, as follows.

\begin{lemma}[The differential of the exponential map via Jacobi fields]
\label{lem:exponential-differential-jacobi}
\uses{def:exponential-map,def:geodesic,def:conjugate-point}
% Partial Lean progress (existence/uniqueness of Y_Z, first paragraph of the
% proof, via chart-change covariance of the Jacobi pair system; the chart
% partition of the compact geodesic and the abstract composition engine for
% chaining the one-chart flow-step / chart-junction derivatives; the flow-step
% link exists_geodesic_flow_step_jacobiTransport identifying the one-chart flow
% derivative with the Jacobi variational-pair transport, the even links of the
% chain; and now the MANIFOLD-FORM flow-step link
% exists_geodesic_flow_step_jacobiTransport_manifold, which identifies the
% flow's internal base geodesic with the actual manifold geodesic gamma via
% intrinsic geodesic uniqueness on a preconnected interval and so transports
% the chart Jacobi variational pair directly along gamma -- the (f,L,x,p)
% one-step datum consumed by the composition engine); and now the ARBITRARY-BASE
% flow-step link exists_geodesic_flow_step_jacobiTransport_manifold_at, which
% bases a flow step at any interior time a (not just t=0) and lands it at any
% prescribed nearby time b, transporting the chart Jacobi variational pair from
% a to b -- the shifting is powered by translation invariance of the geodesic
% equation (IsGeodesicOn.comp_const_add), so the flow's 0-clock base geodesic
% is identified with the time-shifted curve s |-> gamma (a+s).  This is exactly
% the per-boundary datum a partition of [0,1] feeds to the composition engine.
% RUNG 2b NOW COMPLETE: the BALL-UNIFORM flow-step link
% exists_geodesic_flow_step_jacobiTransport_manifold_ball anchors ONE shared
% geodesic flow at a center time s0 and serves EVERY base a <= b in a whole time
% window (s0-rho, s0+rho); this makes the cover elements genuine open intervals
% (any subinterval admits a step), breaking the finiteness circularity, and
% exists_geodesic_flowstep_partition then refines the cover of [0,1] to a finite
% chart/flow-step partition with per-piece flow-step data and both-sided chart
% membership at every boundary.  The chart junction is packaged in flow-step
% shape by exists_geodesic_junction_step (odd links), and a global Jacobi field
% restricts to each nondegenerate piece by IsJacobiFieldAlongOn.mono.  RUNG 3 NOW
% COMPLETE: exists_geodesic_jacobiTransport_chain combines each piece into a
% single link junction_i . flowStep_i (identity flow part on degenerate boundary
% pieces, where tau_i = tau_{i+1} once gamma has reached time 1), indexed
% directly by the piece -- no parity index -- and chains them through the
% marked-chain-factored engine hasStrictFDerivAt_comp_chain', producing ONE
% composed map F_0 strictly differentiable at the initial chart state whose
% single derivative D_0 transports the Jacobi variational pair of EVERY manifold
% Jacobi field along gamma from time 0 to time 1.
%
% RUNG 4 IN PROGRESS (run 0110 session 0006).  The obstruction, now identified and
% removed: EVERY layer of the flow chain pinned the endpoint map down only AT the
% single base state of gamma (flowEnd (state at tau_i) = state at tau_{i+1}),
% whereas d(exp_p)_v is a derivative IN v, so HasFDerivAt of exp_p is a
% NEIGHBOURHOOD statement in v and cannot be extracted from base-point data alone.
% The nearby-state data was in fact already present in the construction --
% exists_geodesic_flow_step builds a flow Z z for EVERY z in a ball around the base
% state, with sprayBase beta (Z z) a genuine geodesic -- but it was discarded at the
% first manifold wrapper (exists_geodesic_flow_step_jacobiTransport_manifold).
% It is now re-exported and threaded through the whole chain:
%   MorganTianLib.sprayFlow_eq_chartState — the semantic core: a spray flow state IS
%       the chart state of its own base curve, Z z sigma = (phi_beta(c), (phi_beta . c)').
%   exists_geodesic_flow_step_jacobiTransport_manifold_ball — now ALSO yields a
%       neighbourhood W of the base chart state at a, such that every z in W is
%       realized by a geodesic c_z (chart state z at time a, confined to chart beta),
%       and flowEnd z is exactly the chart state of c_z at time b.
%   exists_geodesic_flowstep_partition — threads that datum per piece, (Wnb i, mwin i).
%   MorganTianLib.hasStrictFDerivAt_comp_chain_nbhd — the composition engine now
%       intersects the finitely many preimages (f_{i-1} . ... . f_0)^{-1}(W i) into ONE
%       neighbourhood W_0 of the initial state (each partial composite is continuous
%       at x_0 and sends x_0 to x_i), and exposes, for every z in W_0, the ORBIT
%       orb i = (f_{i-1} . ... . f_0)(z), with orb i in W i and F_0 z = orb n.
% The per-piece geodesics attached to an orbit are then recognized as pieces of ONE
% geodesic NOT by a sSup gluing walk but by UNIQUENESS, via the new bridge
% MorganTianLib/Ch01/GeodesicOfState.lean (assumes M complete, as Morgan-Tian do
% throughout -- a documented strengthening of the ambient hypotheses of this lemma):
%   MorganTianLib.exists_geodesic_chartState — on a complete manifold, every chart-beta
%       state (foot x, velocity w) at every time a is realized by a geodesic on all
%       of R.  Hopf-Rinow prescribes the velocity in the chart AT x, so w is first
%       converted by the tangent coordinate change and the result translated in time;
%       the cocycle C_{x->beta} . C_{beta->x} = id returns the chart-beta reading.
%   MorganTianLib.eq_of_chartState_eq — two global geodesics sharing ONE chart-beta
%       state coincide (OpenGALib's IsGeodesicOn.eqOn_of_deriv_chartReading_eq on the
%       preconnected set univ; it compares velocities in an ARBITRARY chart, which is
%       exactly the chart-state vocabulary the flow chain speaks).
% RUNG 4 NOW COMPLETE (run 0110 session 0010) -- the lemma is fully formalized.
%   (i) exists_geodesic_jacobiTransport_chain_nbhd (MorganTianLib/Ch01/FlowChainNbhd.lean)
%       SHRINKS each per-piece window to W_i = Wnb_i cap flowPart_i^{-1}(foot^{-1}(source
%       beta_{i+1})), so a NEARBY geodesic's foot at tau_{i+1} lands in both chart sources
%       and the junction stateTransition_chartState (which quantifies over an ARBITRARY
%       geodesic) serves the odd links too.  Still a neighbourhood of the base state: the
%       flow part is continuous there (strictly differentiable), the chart inverse is
%       continuous at the image point, and gamma(tau_{i+1}) does lie in the open set.
%       Then the ORBIT INDUCTION runs: orb i = chart-beta_i state of the single global
%       geodesic c at tau_i, each piece's LOCAL geodesic being identified with c on the
%       piece's open time window by intrinsic uniqueness (they share a chart state there)
%       -- no sSup gluing walk.  Conclusion: F_0 computes the time-1 chart-zeta state of
%       the geodesic emanating from EVERY nearby initial state, not just that of gamma.
%       Note this needs NO completeness hypothesis: c is universally quantified, so only
%       uniqueness of geodesics is used, never existence.
%   (ii)-(iv) MorganTianLib/Ch01/ExpDifferential.lean.  Compose F_0 with the initial-state
%       map A(w) = (phi_alpha(p), C w), C = tangentCoordChange p->alpha : E ->L E, which
%       sends a tangent vector w at p to the chart-alpha state at time 0 of gamma_w.  A is
%       AFFINE (a constant plus a continuous linear map), so d A = (0, C .).  Every
%       gamma_w starts at p and is global, so the neighbourhood clause applies to each,
%       giving F_0(A w) = (phi_zeta(exp_p w), ...) for all w near v; hence phi_zeta . exp_p
%       agrees with w |-> (F_0 (A w))_1 near v and the chain rule gives
%       d(phi_zeta . exp_p)_v = fst . D_0 . (0, C .)  [hasFDerivAt_chartReading_expMapGlobal].
%       Evaluating on Z: the initial variational pair of Y_Z (Y_Z(0)=0, nabla_X Y_Z(0)=Z)
%       is exactly (0, C Z) by jacobiVarPair_of_left_eq_zero, so D_0 carries it to the
%       chart-zeta variational pair of Y_Z at time 1, whose FIRST component is the chart
%       reading of Y_Z(1).  That is d(exp_p)_v(Z) = Y_Z(1)  [expDifferential_eq_jacobiField].
%
% TWO DOCUMENTED DEVIATIONS from the statement as written above.
%   (a) M is assumed COMPLETE (as Morgan-Tian do throughout), so O_p = T_pM and the
%       "maximal domain" hypothesis v in O_p is vacuous.  Completeness is what supplies the
%       global geodesic gamma_w for every w (Hopf-Rinow).
%   (b) The Lean uses MorganTianLib.expMapGlobal (MorganTianLib/Ch01/GlobalExp.lean), NOT
%       OpenGALib's expMap.  Reason: OpenGALib defines expMap g p v = maximalGeodesic g p v 1,
%       and maximalGeodesic is an integral curve of geodesicVectorFieldChart g p -- the
%       geodesic spray written in the chart AT p, a formula that is meaningless outside
%       (chartAt H p).source.  So maximalGeodesicInterval g p v stops when the geodesic
%       leaves that single chart (OpenGALib's own HopfRinow.lean notes it already fails on
%       the round circle), and expMap is correct only for v small enough that gamma_v([0,1])
%       stays in one chart -- it is the junk value p otherwise.  Since this lemma concerns
%       exp_p at an ARBITRARY v, whose geodesic may cross many charts, it cannot be stated
%       against expMap.  expMapGlobal is therefore built the way Morgan-Tian define exp_p:
%       the time-1 value of the GLOBAL geodesic with initial data (p,v) (Hopf-Rinow), pinned
%       chart-independently by geodesic uniqueness (globalGeodesic_eq), so the Classical
%       choice made in its definition is immaterial.  Repairing maximalGeodesic to be
%       chart-independent belongs to OpenGALib and is filed there, not here.
%
% The two derivative statements are read in the terminal chart zeta (the Frechet derivative
% of w |-> phi_zeta(exp_p w) at v), which is the coordinate form of d(exp_p)_v; the
% identification with the invariant mfderiv is a routine chart-reading step, not the
% mathematical content, and is not separately formalized.
\lean{MorganTianLib.IsJacobiFieldOn.transfer,MorganTianLib.chartCurvature_coordChange,MorganTianLib.exists_isJacobiFieldAlongOn,MorganTianLib.IsJacobiFieldAlongOn.eqOn_of_initial,MorganTianLib.IsJacobiFieldAlongOn.isJacobiFieldOn_of_mem_source,MorganTianLib.exists_geodesic_chart_partition,MorganTianLib.hasStrictFDerivAt_comp_chain,MorganTianLib.hasStrictFDerivAt_comp_chain',MorganTianLib.jacobiVarPair_of_left_eq_zero,MorganTianLib.exists_geodesic_flow_step_jacobiTransport,MorganTianLib.exists_geodesic_flow_step_jacobiTransport_manifold,MorganTianLib.IsGeodesicOn.comp_const_add,MorganTianLib.exists_geodesic_flow_step_jacobiTransport_manifold_at,MorganTianLib.exists_geodesic_flow_step_jacobiTransport_manifold_ball,MorganTianLib.exists_geodesic_flowstep_partition,MorganTianLib.exists_geodesic_junction_step,MorganTianLib.IsJacobiFieldAlongOn.mono,MorganTianLib.exists_geodesic_jacobiTransport_chain,MorganTianLib.sprayFlow_eq_chartState,MorganTianLib.hasStrictFDerivAt_comp_chain_nbhd,MorganTianLib.exists_geodesic_chartState,MorganTianLib.eq_of_chartState_eq,MorganTianLib.stateTransition_chartState,MorganTianLib.exists_geodesic_jacobiTransport_chain_nbhd,MorganTianLib.globalGeodesic,MorganTianLib.globalGeodesic_eq,MorganTianLib.expMapGlobal,MorganTianLib.expMapGlobal_eq_of_isGeodesic,MorganTianLib.hasFDerivAt_chartReading_expMapGlobal,MorganTianLib.expDifferential_eq_jacobiField}
\leanok
Let $p\in M$ and let $v\in O_p$, the maximal domain of $\exp_p$.
Let $\gamma=\gamma_v\colon[0,1]\to M$ and $X=\gamma'$. For
$Z\in T_pM$ let $Y_Z$ be the unique Jacobi field along $\gamma$
with $Y_Z(0)=0$ and $\nabla_XY_Z(0)=Z$. Then, under the natural
identification of $T_v(T_pM)$ with $T_pM$,
$$d(\exp_p)_v(Z)=Y_Z(1)\qquad\text{for all }Z\in T_pM.$$
In particular, $d(\exp_p)_v$ is singular if and only if
$\exp_p(v)$ is a conjugate point along $\gamma_v$.
\end{lemma}

\begin{proof}
\uses{def:conjugate-point,def:geodesic,thm:levi-civita-connection,def:exponential-map,lem:second-order-linear-ode,lem:parallel-frame,lem:covariant-commutation-jacobi,lem:jacobi-field-coordinates}
\leanok
We first record the existence and uniqueness of $Y_Z$. Let
$e_1,\ldots,e_n$ be a parallel frame along $\gamma$
(Lemma~\ref{lem:parallel-frame}), obtained by
solving the first-order linear ODE $\nabla_Xe_i=0$ with initial
conditions a basis of $T_pM$; since the metric is parallel, the
inner products $\langle e_i,e_j\rangle$ are constant in $t$, so the
$e_i(t)$ remain a basis of $T_{\gamma(t)}M$ for all $t$. Writing a
vector field along $\gamma$ as $Y=\sum_iy^ie_i$, we have
$\nabla_XY=\sum_i\dot y^ie_i$ and
$\nabla_X\nabla_XY=\sum_i\ddot y^ie_i$, so the Jacobi equation
$$\nabla_X\nabla_XY+{\mathcal R}(Y,X)X=0$$
from the discussion of families of geodesics in this section
becomes the second-order linear system
$$\ddot y^i(t)+\sum_jR^i_j(t)\,y^j(t)=0,$$
where the smooth functions $R^i_j(t)$ are the matrix entries of
$Y\mapsto{\mathcal R}(Y,X)X$ in the frame $(e_i)$. By existence and
uniqueness for second-order linear equations
(Lemma~\ref{lem:second-order-linear-ode}), for each $Z\in T_pM$
there is exactly one
solution $Y$ on $[0,1]$ with $Y(0)=0$ and $\nabla_XY(0)=Z$; since
it vanishes at the initial point, it is a Jacobi field, and it is
the field $Y_Z$ of the statement. In particular, if $Z\neq0$ then
$Y_Z$ is not identically zero, since the zero field has vanishing
initial covariant derivative.

Fix $Z\in T_pM$. Since $O_p$ is an open star-shaped neighborhood
of $0\in T_pM$, there is $\varepsilon>0$ such that $v+sZ\in O_p$
for $|s|<\varepsilon$, and then $t(v+sZ)\in O_p$ for all
$t\in[0,1]$. We claim that for any $w\in O_p$ and $t\in[0,1]$ we
have $\exp_p(tw)=\gamma_w(t)$: the curve $c(u)=\gamma_w(tu)$,
$u\in[0,1]$, satisfies $c(0)=p$, $c'(0)=tw$ and
$\nabla_{c'}c'=t^2\nabla_{\gamma_w'}\gamma_w'=0$, so by the ODE
uniqueness of geodesics $c=\gamma_{tw}$, and hence
$\exp_p(tw)=\gamma_{tw}(1)=c(1)=\gamma_w(t)$. Therefore the map
$$\tilde\gamma(t,s)=\exp_p\bigl(t(v+sZ)\bigr),\qquad
(t,s)\in[0,1]\times(-\varepsilon,\varepsilon),$$
satisfies $\tilde\gamma(\cdot,s)=\gamma_{v+sZ}$; it is a family of
geodesics parameterized by $[0,1]$, all starting at $p$, with
$\tilde\gamma(\cdot,0)=\gamma$. It is smooth, since solutions of
the geodesic equation depend smoothly on their initial conditions.

Set $J(t)=\partial\tilde\gamma/\partial s(t,0)$. By
Lemma~\ref{lem:covariant-commutation-jacobi}(3), $J$ satisfies the Jacobi
equation along $\gamma$, and $J(0)=0$ because $\tilde\gamma(0,s)=p$
for all $s$. We compute $\nabla_XJ(0)$. For the smooth map
$\tilde\gamma$ the covariant derivatives of the two coordinate
fields may be interchanged,
$$\nabla_{\partial_t}\partial_s\tilde\gamma
=\nabla_{\partial_s}\partial_t\tilde\gamma,$$
since in local coordinates both sides have components
$\partial_s\partial_t\tilde\gamma^k
+\Gamma^k_{ij}\,\partial_t\tilde\gamma^i\,\partial_s\tilde\gamma^j$;
here mixed partial derivatives commute, and
$\Gamma^k_{ij}=\Gamma^k_{ji}$ because the Levi-Civita connection of
Theorem~\ref{thm:levi-civita-connection} is torsion-free and the
coordinate fields have vanishing brackets. (This is the same
interchange used in the derivation of the Jacobi equation.) Hence
$$\nabla_XJ(0)=\nabla_{\partial_t}\partial_s\tilde\gamma(0,0)
=\nabla_{\partial_s}\partial_t\tilde\gamma(0,0).$$
Now $\partial_t\tilde\gamma(0,s)=\gamma_{v+sZ}'(0)=v+sZ$ is a
vector field along the constant curve $s\mapsto p$. For a vector
field $W$ along a constant curve $c(s)\equiv p$, the covariant
derivative in local coordinates is
$\nabla_{\partial_s}W=(\dot W^k+\Gamma^k_{ij}\dot c^iW^j)
\partial_k=\dot W^k\partial_k$, i.e., the ordinary derivative of
$W$ as a curve in the fixed vector space $T_pM$. Therefore
$$\nabla_XJ(0)=\frac{d}{ds}\Bigl|_{s=0}\Bigr.(v+sZ)=Z.$$
Thus $J$ is a Jacobi field along $\gamma$ with $J(0)=0$ and
$\nabla_XJ(0)=Z$, so $J=Y_Z$ by the uniqueness established in the
first paragraph.

On the other hand, $\tilde\gamma(1,s)=\exp_p(v+sZ)$, and under the
natural identification of $T_v(T_pM)$ with $T_pM$ the vector $Z$
corresponds to the velocity at $s=0$ of the curve $s\mapsto v+sZ$
in $T_pM$. Hence
$$d(\exp_p)_v(Z)=\frac{d}{ds}\Bigl|_{s=0}\Bigr.\exp_p(v+sZ)
=J(1)=Y_Z(1),$$
which is the asserted formula.

Finally, note that $\exp_p(v)=\gamma_v(1)$ with $1>0$. If
$d(\exp_p)_v$ is singular, there is $Z\neq0$ with
$Y_Z(1)=d(\exp_p)_v(Z)=0$. As observed above, $Y_Z$ is then a
non-zero Jacobi field along $\gamma_v$ vanishing at $\gamma_v(1)$,
so by Definition~\ref{def:conjugate-point} the point $\exp_p(v)$
is a conjugate point along $\gamma_v$. Conversely, if $\exp_p(v)$
is a conjugate point along $\gamma_v$, there is a non-zero Jacobi
field $Y$ along $\gamma_v$ vanishing at $\gamma_v(1)$; being a
Jacobi field, $Y$ also vanishes at the initial point, $Y(0)=0$.
Set $Z=\nabla_XY(0)$. Then $Z\neq0$, since otherwise $Y\equiv0$ by
uniqueness, and $Y=Y_Z$. Hence
$d(\exp_p)_v(Z)=Y_Z(1)=Y(1)=0$ with $Z\neq0$, so $d(\exp_p)_v$ is
singular.
\end{proof}

\begin{corollary}[Exponential map is local diffeomorphism]
\label{cor:exponential-local-diffeomorphism}
\label{star}
\uses{prop:minimal-geodesic-no-conjugate}
Suppose that $\gamma$ is a minimal geodesic parameterized by $[0,1]$
starting at $p$. Let $X(0)=\gamma'(0)\in T_pM$. Then for each
$t_0<1$ the restriction $\gamma|_{[0,t_0]}$ is a minimal geodesic
and $\exp_p\colon T_pM\to M$ is a local diffeomorphism near
$t_0X(0)$.
\end{corollary}

\begin{proof}
\uses{prop:minimal-geodesic-no-conjugate,lem:exponential-differential-jacobi}
Of course, $\exp_p(t_0X(0))=\gamma(t_0)$. By
Lemma~\ref{lem:exponential-differential-jacobi} (applied to the
geodesic $\gamma|_{[0,t_0]}=\gamma_{t_0X(0)}$, reparameterized by
$[0,1]$), the differential of the
exponential mapping at $t_0X(0)$ is singular if and only if there
is a non-zero
Jacobi field along $\gamma|_{[0,t_0]}$ vanishing at $0$ and at
$\gamma(t_0)$. According
to  Proposition~\ref{prop:minimal-geodesic-no-conjugate} the only such Jacobi field is the
trivial one. Hence, the differential of $\exp_p$ at $t_0X(0)$
is an isomorphism, completing the proof.
\end{proof}

Two fundamental global consequences of this local theory are the
existence of minimizing geodesics and the Hopf--Rinow theorem.

\begin{lemma}[Existence of minimizing geodesics]
\label{lem:minimizing-geodesic-exists}
\uses{def:geodesic,def:exponential-map}
Let $(M,g)$ be a Riemannian manifold, not assumed complete or
connected (with the convention that $d(p,q)=\infty$ when there is
no smooth path from $p$ to $q$), let $p\in M$ and $D>0$. Suppose
that every unit-speed
geodesic starting at $p$ is defined on $[0,D]$. Then every $q\in
M$ with $d(p,q)\le D$ is joined to $p$ by a geodesic of length
$d(p,q)$; equivalently, $q=\exp_p(w)$ for some $w\in T_pM$ with
$|w|=d(p,q)$.
\end{lemma}

\begin{proof}
\uses{lem:normal-ball,def:geodesic,def:exponential-map}
Throughout we use that $d(a,b)\le L(\sigma)$ for every piecewise
smooth curve $\sigma$ from $a$ to $b$, so that $d(a,b)$ is also the
infimum of the lengths of piecewise smooth curves from $a$ to $b$:
precomposing $\sigma$ with a smooth nondecreasing surjection of
$[0,1]$ onto itself all of whose derivatives vanish at the finitely
many corner times produces a curve with the same endpoints and, by
the change-of-variables formula, the same length, and this curve is
smooth (in local coordinates near a corner every one-sided
derivative of the reparameterized curve tends to $0$ at the corner
time).

For $x\in M$ write $\delta_x>0$ for the constant furnished by
Lemma~\ref{lem:normal-ball} applied to the compact set $\{x\}$:
$\exp_x$ is a diffeomorphism of $B(0,\delta_x)\subset T_xM$ onto
the metric ball $B(x,\delta_x)$, with $d(x,\exp_x(v))=|v|$ for
$|v|<\delta_x$, and every piecewise smooth curve starting at $x$
that does not stay in $B(x,s)$, $s\le\delta_x$, has length at least
$s$. In particular, for $0<\rho<\delta_x$ we have
$\exp_x(B(0,\rho))=B(x,\rho)$: the inclusion $\subset$ holds
because $d(x,\exp_x(v))=|v|<\rho$, and conversely every $y\in
B(x,\rho)\subset B(x,\delta_x)$ is $\exp_x(v)$ for a (unique) $v$
with $|v|=d(x,y)<\rho$. We also record that $d(x,y)>0$ whenever
$x\not=y$: either $y\in B(x,\delta_x)$, in which case
$y=\exp_x(v)$ with $v\not=0$ (as $\exp_x$ is injective on
$B(0,\delta_x)$ and $\exp_x(0)=x\not=y$), so $d(x,y)=|v|>0$; or
$y\notin B(x,\delta_x)$, so $d(x,y)\ge\delta_x$.

\emph{Step 0: small spheres realize the distance.} For $x\in M$
and $0<\rho<\delta_x$ set
$$S_\rho(x)=\exp_x\left(\{v\in T_xM\,|\,|v|=\rho\}\right),$$
a compact subset of $M$, each of whose points $y$ satisfies
$d(x,y)=\rho$. We claim that for every $q\in M$ with
$\rho\le d(x,q)<\infty$,
$$d(x,q)=\rho+\min_{y\in S_\rho(x)}d(y,q),$$
the minimum being attained because $d(\cdot,q)$ is continuous
(indeed $1$-Lipschitz, by the triangle inequality) and $S_\rho(x)$
is compact and nonempty: since $0<\rho\le d(x,q)<\infty$, some
piecewise smooth curve of finite length joins $x$ to $q\not=x$;
such a nonconstant curve cannot exist in a $0$-dimensional
manifold, so $\dim M\ge1$ and the sphere $\{|v|=\rho\}$ in $T_xM$
is nonempty.

The inequality $\le$ is immediate: for every $y\in S_\rho(x)$,
$$d(x,q)\le d(x,y)+d(y,q)=\rho+d(y,q).$$
For the inequality $\ge$, let $c\colon[0,1]\to M$ be any piecewise
smooth curve from $x$ to $q$; it suffices to bound $L(c)$ from
below by the right-hand side. First note that the frontier of
$B(x,\rho)$ in $M$ is contained in $S_\rho(x)$: since
$\rho<\delta_x$, the set
$C:=\exp_x\left(\overline{B(0,\rho)}\right)
=B(x,\rho)\cup S_\rho(x)$ is compact, hence closed in $M$, so
$\overline{B(x,\rho)}\subset C$; and $S_\rho(x)\cap
B(x,\rho)=\emptyset$ because the points of $S_\rho(x)$ are at
distance exactly $\rho$ from $x$; hence
$\overline{B(x,\rho)}\setminus B(x,\rho)\subset S_\rho(x)$. Now
$c(0)=x\in B(x,\rho)$ while $c(1)=q\notin B(x,\rho)$; let
$$t_1=\inf\{t\in[0,1]\,|\,c(t)\notin B(x,\rho)\}.$$
Then $t_1>0$, $c(t)\in B(x,\rho)$ for $0\le t<t_1$, and
$c(t_1)\notin B(x,\rho)$: if $c(t_1)$ lay in the open set
$B(x,\rho)$, then $c(t)\in B(x,\rho)$ for all $t$ in a
neighborhood of $t_1$, contradicting the definition of $t_1$ as an
infimum. Being the limit of the points $c(t)\in B(x,\rho)$ as
$t\uparrow t_1$, the point $y_1:=c(t_1)$ lies in
$\overline{B(x,\rho)}\setminus B(x,\rho)\subset S_\rho(x)$. The
curve $c|_{[0,t_1]}$ starts at $x$ and does not stay in
$B(x,\rho)$, so by the exit estimate of
Lemma~\ref{lem:normal-ball} its length is at least $\rho$; and
$c|_{[t_1,1]}$ joins $y_1\in S_\rho(x)$ to $q$, so its length is
at least $d(y_1,q)\ge\min_{y\in S_\rho(x)}d(y,q)$. Adding,
$$L(c)=L\left(c|_{[0,t_1]}\right)+L\left(c|_{[t_1,1]}\right)
\ge\rho+\min_{y\in S_\rho(x)}d(y,q),$$
and taking the infimum over all such $c$ proves the claim.

\emph{Step 1: setup.} Let $d:=d(p,q)\le D$. If $d=0$ then $q=p$ by
the positivity established above, and the constant geodesic at $p$
(with $w=0$) satisfies the conclusion. So assume $d>0$ and fix
$0<\rho_0<\min(\delta_p,d)$. By Step 0 applied at $x=p$ (note
$d(p,q)=d>\rho_0$) there is $x_0\in S_{\rho_0}(p)$ with
$$d(p,q)=\rho_0+d(x_0,q).$$
Write $x_0=\exp_p(\rho_0u)$ with $u\in T_pM$, $|u|=1$, and let
$\gamma(t)=\exp_p(tu)$ be the unit-speed geodesic with
$\gamma(0)=p$ and $\gamma'(0)=u$; by hypothesis $\gamma$ is
defined on $[0,D]\supseteq[0,d]$. Set
$$A=\{t\in[\rho_0,d]\,|\,d(\gamma(t),q)=d-t\}.$$
Since $\gamma(\rho_0)=x_0$ and $d(x_0,q)=d-\rho_0$, we have
$\rho_0\in A$. Moreover $A$ is the zero set of the continuous
function $t\mapsto d(\gamma(t),q)-(d-t)$ on the compact interval
$[\rho_0,d]$, so $A$ is a nonempty compact set and $t^*:=\sup A$
belongs to $A$.

\emph{Step 2: a minimizing broken path.} Suppose $t^*<d$, and set
$p'=\gamma(t^*)$, so that $d(p',q)=d-t^*>0$. Fix
$0<\rho'<\min(\delta_{p'},d-t^*)$. By Step 0 applied at $x=p'$
(note $d(p',q)>\rho'$) there is $x'\in S_{\rho'}(p')$ with
$$d(p',q)=\rho'+d(x',q),\qquad\text{i.e.,}\qquad
d(x',q)=d-t^*-\rho'.$$
By the triangle inequality $d(p,q)\le d(p,x')+d(x',q)$, so
$$d(p,x')\ge d-(d-t^*-\rho')=t^*+\rho'.$$
Write $x'=\exp_{p'}(\rho'v')$ with $v'\in T_{p'}M$, $|v'|=1$, and
let $\alpha(s)=\exp_{p'}(sv')$, $0\le s\le\rho'$, be the
unit-speed radial geodesic from $p'$ to $x'$. Set $T=t^*+\rho'$
and define the piecewise smooth path $c\colon[0,1]\to M$ of
constant speed $T$ by
$$c(t)=\begin{cases}\gamma(Tt),&0\le t\le t_0:=t^*/T,\\
\alpha(Tt-t^*),&t_0\le t\le1.\end{cases}$$
(Affine reparameterizations of geodesics are geodesics, as is
immediate from the geodesic equation, so both pieces of $c$ are
geodesics.) The path $c$ runs from $p$ to $x'$ and has length
$t^*+\rho'=T$; combined with $d(p,x')\ge T$ this gives
$$d(p,x')=T=L(c),$$
so $c$ realizes the distance from $p$ to $x'$.

\emph{Step 3: a minimizing broken geodesic has no corner.} We
claim that the two one-sided velocities of $c$ at $t_0$ coincide,
i.e., that
$$\Delta:=c'(t_0^+)-c'(t_0^-)=T\left(v'-\gamma'(t^*)\right)=0.$$
Suppose $\Delta\not=0$. Since $|c'|\equiv T$,
$$E(c)=\frac{1}{2}\int_0^1|c'(t)|^2\,dt=\frac{1}{2}T^2
=\frac{1}{2}d(p,x')^2,$$
while by the Cauchy--Schwarz inequality every piecewise smooth
path $\sigma$ from $p$ to $x'$ parameterized by $[0,1]$ satisfies
$E(\sigma)\ge\frac{1}{2}L(\sigma)^2\ge\frac{1}{2}d(p,x')^2$; thus
$c$ minimizes energy among all such paths. Now choose a piecewise
smooth vector field $Y$ along $c$ as follows: let $V$ be the
vector field along $c$ obtained by parallel transport of
$\Delta\in T_{c(t_0)}M$ along each of the two smooth pieces of $c$
(so $V$ is smooth on each piece and continuous, with
$V(t_0)=\Delta$), and set $Y=\phi V$ with $\phi$ a smooth function
on $[0,1]$ satisfying $\phi(0)=\phi(1)=0$ and $\phi(t_0)=1$.
Consider the variation
$$\tilde c(t,u)=\exp_{c(t)}(uY(t)),$$
defined and smooth on each of $[0,t_0]\times(-u_0,u_0)$ and
$[t_0,1]\times(-u_0,u_0)$ for some $u_0>0$, because the domain of
the exponential map is an open neighborhood of the zero section of
$TM$ and $c([0,1])$ is compact. Since $Y(0)=Y(1)=0$, each path
$c_u=\tilde c(\cdot,u)$ runs from $p$ to $x'$. Write $\widetilde
X=\partial\tilde c/\partial t$ and $\widetilde Y=\partial\tilde
c/\partial u$; then $\widetilde X(t,0)=c'(t)$, and $\widetilde
Y(t,0)=Y(t)$ because $u\mapsto\exp_{c(t)}(uY(t))$ is the geodesic
with initial velocity $Y(t)$. On each smooth piece, compatibility
of the Levi-Civita connection with the metric, the symmetry
$\nabla_{\widetilde Y}\widetilde X=\nabla_{\widetilde X}\widetilde
Y$ (torsion-freeness applied to the coordinate fields of the
parameterized surface $\tilde c$), and integration by parts give
$$\frac{d}{du}\,\frac{1}{2}\int\langle\widetilde X,\widetilde
X\rangle\,dt=\int\langle\nabla_{\widetilde X}\widetilde
Y,\widetilde X\rangle\,dt=\Bigl[\langle\widetilde Y,\widetilde
X\rangle\Bigr]-\int\langle\widetilde Y,\nabla_{\widetilde
X}\widetilde X\rangle\,dt.$$
At $u=0$ both pieces of $c$ are geodesics, so the integrals on the
right vanish and, summing the two pieces and using $Y(0)=Y(1)=0$,
$$\frac{dE(c_u)}{du}\Bigl|_{u=0}\Bigr.
=\Bigl[\langle Y,c'\rangle\Bigr]_0^{t_0^-}
+\Bigl[\langle Y,c'\rangle\Bigr]_{t_0^+}^{1}
=\langle Y(t_0),c'(t_0^-)\rangle-\langle Y(t_0),c'(t_0^+)\rangle
=-|\Delta|^2<0.$$
Hence $E(c_u)<E(c)=\frac{1}{2}d(p,x')^2$ for all sufficiently
small $u>0$, and for such $u$
$$L(c_u)\le\sqrt{2E(c_u)}<d(p,x'),$$
a contradiction, since $c_u$ is a piecewise smooth path from $p$
to $x'$. Therefore $\Delta=0$, i.e.,
$\alpha'(0)=v'=\gamma'(t^*)$. Now $\alpha$ and
$s\mapsto\gamma(t^*+s)$ are geodesics with the same initial
position $p'$ and the same initial velocity, so by the uniqueness
statement of classical ODE theory (the geodesic equation is a
second-order ODE) they coincide on $[0,\rho']$; note that
$t^*+\rho'<d\le D$, so $\gamma$ is defined there. In particular
$$x'=\alpha(\rho')=\gamma(t^*+\rho').$$

\emph{Step 4: conclusion.} By Steps 2 and 3,
$$d\left(\gamma(t^*+\rho'),q\right)=d(x',q)=d-(t^*+\rho'),$$
and $\rho_0\le t^*<t^*+\rho'<d$, so $t^*+\rho'\in A$,
contradicting $t^*=\sup A$. Hence $t^*=d$, that is,
$d(\gamma(d),q)=0$, and by the positivity of $d$ established above
$\gamma(d)=q$. Thus $\gamma|_{[0,d]}$ is a unit-speed geodesic
from $p$ to $q$ of length $d=d(p,q)$, and it is minimizing since
every curve from $p$ to $q$ has length at least $d(p,q)$. Finally,
set $w=d\,u\in T_pM$, so that $|w|=d$; the curve
$t\mapsto\gamma(td)$, $t\in[0,1]$, is a geodesic with initial
velocity $w$, whence $\exp_p(w)=\gamma(d)=q$.
\end{proof}

We can now state and prove the classical theorem of Hopf and Rinow
(compare Theorem 16 of Chapter 5 on p.~137 of \cite{Petersen}).

\begin{theorem}[Hopf-Rinow Theorem]
\label{thm:hopf-rinow}
\uses{def:geodesic}
(Hopf-Rinow) Let $(M,g)$ be a Riemannian manifold which is
complete as a metric space; for parts (2) and (3) assume in
addition that $M$ is connected. Then:
\begin{enumerate}
\item every geodesic extends to a geodesic defined for all time;
\item any two points $p,q\in M$ are joined by a minimizing
geodesic, of length $d(p,q)$;
\item every closed and bounded subset of $M$ is compact.
\end{enumerate}
\end{theorem}

\begin{proof}
\uses{def:geodesic,def:exponential-map,lem:minimizing-geodesic-exists,thm:levi-civita-connection}
Throughout we use that in local coordinates the geodesic equation
is a second-order ODE system with smooth coefficients, so that
classical ODE theory provides local existence, uniqueness, and
smooth dependence on initial conditions for geodesics. We record
two elementary observations.

(a) \emph{Geodesics have constant speed:} if $\gamma$ is a
geodesic then, since the Levi-Civita connection makes $g$
parallel,
$$\frac{d}{dt}\langle \gamma',\gamma'\rangle
=2\langle \nabla_{\gamma'}\gamma',\gamma'\rangle=0.$$

(b) \emph{Affine reparameterizations of geodesics are geodesics:}
if $\gamma$ is a geodesic and $\alpha,\beta\in\Ar$, then
$t\mapsto\gamma(\alpha t+\beta)$ is a geodesic, since the
coordinate form of the geodesic equation for the reparameterized
curve is $\alpha^2$ times the corresponding expression for
$\gamma$.

(1) Let $\gamma$ be a geodesic. By (a) its speed is a constant
$\lambda\ge0$. If $\lambda=0$ then $\gamma$ is constant, and the
same constant curve defined on all of $\Ar$ extends it. If
$\lambda>0$, then $c(t)=\gamma(t/\lambda)$ is a unit-speed
geodesic by (b), and if $c$ extends to a geodesic defined for all
time then so does $\gamma$, since $t\mapsto c(\lambda t)$ is a
geodesic extending $\gamma$. So we may assume $\gamma$ has unit
speed.

Any two geodesics that agree, together with their velocities, at
a single point agree on the intersection of their domains: the
set of parameters where they agree to first order is nonempty,
open by ODE uniqueness, and closed by continuity, and the
intersection of two open intervals is connected. Consequently the
union of all geodesics extending $\gamma$ is a geodesic, again
denoted $\gamma$, defined on an open interval $(a,b)$ and
admitting no extension to a larger interval. We must show
$(a,b)=\Ar$. Suppose $b<\infty$; the argument for the left
endpoint is the same.

For $s\le t$ in $(a,b)$ the restriction $\gamma|_{[s,t]}$ is a
smooth path of length $t-s$, so $d(\gamma(s),\gamma(t))\le t-s$.
Hence for any sequence $t_i\uparrow b$ the points $\gamma(t_i)$
form a Cauchy sequence, which by completeness converges to some
$\bar q\in M$; letting $i\rightarrow\infty$ in
$d(\gamma(t),\bar q)\le(t_i-t)+d(\gamma(t_i),\bar q)$ gives
$$d(\gamma(t),\bar q)\le b-t\qquad\text{for all } t<b.$$

Fix a coordinate chart $\varphi\colon U\rightarrow \Ar^n$ with
$\varphi(\bar q)=0$ and $\rho>0$ so small that
$K=\varphi^{-1}(\{|x|\le\rho\})$ is a compact subset of $U$. Since
the function $(x,v)\mapsto\langle v,v\rangle$ is continuous and
positive on the compact set of pairs with $x\in K$ and
$|d\varphi(v)|=1$, there is $c>0$ with
$$\langle v,v\rangle\ge c^2\,|d\varphi(v)|^2
\qquad\text{for all } v\in T_xM,\ x\in K,$$
where $|\cdot|$ denotes the Euclidean norm on $\Ar^n$. We claim
that every $q$ with $d(\bar q,q)<c\rho$ lies in $K$. Indeed, let
$\sigma\colon[0,1]\rightarrow M$ be a smooth path with
$\sigma(0)=\bar q$ and $\sigma(1)=q\notin K$, and let $[0,t^*]$ be
the largest initial interval with $\sigma([0,t^*])\subset K$. If
$|\varphi(\sigma(t^*))|<\rho$ then $\sigma$ would remain in $K$
slightly beyond $t^*$, so $|\varphi(\sigma(t^*))|=\rho$ and
$$L(\sigma)\ge\int_0^{t^*}|\sigma'|\,dt
\ge c\int_0^{t^*}\left|\frac{d}{dt}\varphi(\sigma(t))\right|dt
\ge c\,|\varphi(\sigma(t^*))-\varphi(\sigma(0))|=c\rho.$$
Taking the infimum over such $\sigma$ gives $d(\bar q,q)\ge
c\rho$, proving the claim.

Next we find $\epsilon>0$ such that every unit-speed geodesic
starting at a point of $K$ is defined at least on
$(-\epsilon,\epsilon)$. Each point of $K$ has a compact
neighborhood contained in a chart domain; finitely many of these,
say $N_1,\ldots,N_m$ with charts $\varphi_j\colon
U_j\rightarrow\Ar^n$, cover $K$. Over $U_j$ the geodesic equation
is the first-order system $x'=y$,
$(y^k)'=-\Gamma^{k}_{ij}(x)y^iy^j$ on the open set
$\varphi_j(U_j)\times\Ar^n$, and the set of initial conditions
$$S_j=\left\{(x,y)\ :\ x\in\varphi_j(K\cap N_j),\
g_{ij}(x)y^iy^j=1\right\}$$
is compact: it is closed, its first factor ranges in a compact
set, and over that set positive definiteness of $g$ bounds $|y|$
exactly as above. Solutions of a smooth ODE with initial
conditions in a compact subset of its domain are defined on a
common time interval (classical ODE theory: on a compact
neighborhood the right-hand side is bounded and Lipschitz, giving
a uniform Picard existence time), so there is $\epsilon_j>0$ such
that the solution with any initial condition in $S_j$ is defined
on $(-\epsilon_j,\epsilon_j)$; take
$\epsilon=\min_j\epsilon_j$.

Now choose $t_0\in(a,b)$ with $b-t_0<\min(\epsilon/2,c\rho)$.
Then $d(\gamma(t_0),\bar q)\le b-t_0<c\rho$, so
$\gamma(t_0)\in K$, and the geodesic $\sigma$ with
$\sigma(0)=\gamma(t_0)$ and $\sigma'(0)=\gamma'(t_0)$ is defined
on $(-\epsilon,\epsilon)$. As noted above,
$\sigma(s)=\gamma(t_0+s)$ wherever both sides are defined, so
$\gamma$ and $t\mapsto\sigma(t-t_0)$ patch to a geodesic on
$(a,t_0+\epsilon)$, and $t_0+\epsilon>b+\epsilon/2>b$. This
contradicts the maximality of $(a,b)$; hence $b=\infty$, and
likewise $a=-\infty$.

(2) Since $M$ is a connected manifold, any two points $p,q$ are
joined by a smooth path ($M$ is locally path-connected, hence
path-connected, and within charts a continuous path may be
replaced by a smooth one), so $d(p,q)<\infty$. Set
$D=\max(d(p,q),1)>0$. By (1) every
unit-speed geodesic from $p$ extends to a geodesic defined for
all time; in particular it is defined on $[0,D]$. Since
$d(p,q)\le D$, Lemma~\ref{lem:minimizing-geodesic-exists} gives a
geodesic from $p$ to $q$ of length $d(p,q)$. Since by definition of $d$ no
path from $p$ to $q$ has length less than $d(p,q)$, this geodesic
minimizes the length between its endpoints, i.e., it is a
minimizing geodesic.

(3) Let $C\subset M$ be closed and bounded; the empty set is
compact, so assume $C\not=\emptyset$ and fix $p\in C$ and
$R<\infty$ with $d(p,x)\le R$ for all $x\in C$. By (1) the
geodesic $\gamma_v$ is defined at time $1$ for every $v\in T_pM$,
so the exponential map $\exp_p$ of
Definition~\ref{def:exponential-map} is defined on all of
$T_pM$; it is continuous (indeed smooth) by the smooth dependence
of solutions of the geodesic equation on their initial
conditions. We claim that
$$\{x\in M\ :\ d(p,x)\le R\}\subset
\exp_p\left(\overline{B(0,R)}\right),$$
where $\overline{B(0,R)}=\{w\in T_pM\,:\,|w|\le R\}$. Given such
an $x$, part (2) provides a minimizing geodesic from $p$ to $x$
of length $d(p,x)$; parameterize it at unit speed as
$\sigma\colon[0,d(p,x)]\rightarrow M$ and set
$w=d(p,x)\,\sigma'(0)$ (if $x=p$, take $w=0$). Then
$|w|=d(p,x)\le R$, and $t\mapsto\sigma(t\,d(p,x))$ is a geodesic
by (b) with initial point $p$ and initial velocity $w$, so by
uniqueness it equals $\gamma_w$ and
$$\exp_p(w)=\gamma_w(1)=\sigma(d(p,x))=x.$$
The ball $\overline{B(0,R)}$ is a closed and bounded subset of
the finite-dimensional inner product space
$(T_pM,\langle\cdot,\cdot\rangle)$, hence compact, so its image
$\exp_p\left(\overline{B(0,R)}\right)$ under the continuous map
$\exp_p$ is compact. Finally, $C$ is a closed subset of $M$
contained in this compact set, hence compact.
\end{proof}

Suppose now that $M$ is complete, and fix a point $p\in M$. By
Theorem~\ref{thm:hopf-rinow}(2), every $q\in M$ is the endpoint
of a minimizing geodesic from $p$ parameterized by $[0,1]$;
consequently, $\exp_p\colon T_pM\to M$ is onto.

\begin{definition}[Cut locus]
\label{def:cut-locus}
\uses{def:exponential-map}
Let $M$ be complete and connected, and let $p\in M$. Let $U_p\subset T_pM$ denote
the set of all $v\in T_pM$ for which: (i) $\gamma_v$ is the unique
minimal geodesic from $p$ to $\gamma_v(1)$, and (ii) $\exp_p$ is
a local diffeomorphism at $v$. We set ${\mathcal C}_p\subset M$
equal to $M\setminus \exp_p(U_p)$. Then ${\mathcal C}_p$ is
called the {\sl cut locus from $p$}.
\end{definition}

The basic properties of $U_p$ and ${\mathcal C}_p$ --- that $U_p$
is an open, star-shaped neighborhood of the origin and that the
cut locus is closed and of measure zero --- are established in the
following lemma. Its proof rests on the localized cut-locus
apparatus of Lemma~\ref{lem:localized-cut-locus} in the section on
curvature comparison below; there is no circularity: although that
lemma is typeset later, neither it nor anything in its chain of
dependencies makes reference to the cut locus or to the present
lemma.

\begin{lemma}[Properties of the cut locus]
\label{lem:cut-locus-properties}
\uses{def:cut-locus}
Let $M$ be complete and connected, and let $p\in M$. Then:
\begin{enumerate}
\item $U_p$ is an open, star-shaped neighborhood of $0$ in
$T_pM$, and $d(p,\exp_p(w))=|w|$ for every $w\in U_p$;
\item $\exp_p(U_p)$ is open and the cut locus
${\mathcal C}_p=M\setminus\exp_p(U_p)$ is closed;
\item ${\mathcal C}_p$ has measure zero.
\end{enumerate}
\end{lemma}

\begin{proof}
\uses{def:cut-locus,lem:localized-cut-locus,thm:hopf-rinow,lem:normal-ball}
For $R>0$ let $U_p^{<R}$ be the set of
Lemma~\ref{lem:localized-cut-locus}(3). Comparing the defining
conditions, $U_p^{<R}=U_p\cap B(0,R)$: for $0<|v|<R$ conditions
(i) and (ii) there are literally conditions (i) and (ii) of
Definition~\ref{def:cut-locus}, and $0$ belongs to both sets
($\gamma_0$ is the unique minimal --- indeed the unique ---
geodesic from $p$ to $p$ of length $0$ parameterized by $[0,1]$,
and $\exp_p$ is a local diffeomorphism at $0$ by
Lemma~\ref{lem:normal-ball}). By the Hopf--Rinow theorem
(Theorem~\ref{thm:hopf-rinow}(3)) every ball $B(p,R)$ has compact
closure in $M$, so Lemma~\ref{lem:localized-cut-locus} applies for
every $R>0$.

(1) By Lemma~\ref{lem:localized-cut-locus}(3), each
$U_p\cap B(0,R)=U_p^{<R}$ is an open, star-shaped neighborhood of
$0$ with $d(p,\exp_p(w))=|w|$ for $w\in U_p^{<R}$; taking the
union over $R>0$ gives (1).

(2) By condition (ii), $\exp_p$ is a local diffeomorphism at
every point of the open set $U_p$, so
$\exp_p(U_p)$ is open, and its complement ${\mathcal C}_p$ is
closed.

(3) Fix $r>0$ and let $R>r$. By (1) and
$d(p,\exp_p(w))=|w|$,
$$\exp_p\left(U_p^{<R}\cap B(0,r)\right)
=\exp_p\left(U_p\cap B(0,r)\right)\subseteq\exp_p(U_p),$$
so
$${\mathcal C}_p\cap B(p,r)\ \subseteq\
B(p,r)\setminus\exp_p\left(U_p^{<R}\cap B(0,r)\right),$$
which has measure zero by
Lemma~\ref{lem:localized-cut-locus}(4). Writing ${\mathcal C}_p$
as the countable union of the sets ${\mathcal C}_p\cap B(p,r)$,
$r\in\mathbb{N}$, shows that ${\mathcal C}_p$ has measure zero.
\end{proof}

\begin{proposition}[Exponential map to cut locus complement]
\label{prop:exponential-diffeomorphism-cut-locus}
\uses{def:cut-locus}
The map
$$\exp_p\colon U_p\to M\setminus {\mathcal C}_p$$
is a diffeomorphism. \end{proposition}

\begin{proof}
\uses{def:cut-locus,def:exponential-map,def:geodesic,lem:cut-locus-properties}
By the definition of the cut locus, ${\mathcal C}_p=M\setminus
\exp_p(U_p)$, so $\exp_p(U_p)=M\setminus{\mathcal C}_p$ and
the map in question is onto.

It is injective: suppose $v,w\in U_p$ satisfy
$\exp_p(v)=\exp_p(w)=q$. By condition (i) in the
definition of $U_p$, the geodesic $\gamma_v\colon[0,1]\to M$ is the
unique minimal geodesic from $p$ to $q$, and likewise for
$\gamma_w$; hence, both being minimal geodesics from $p$ to $q$
parameterized by $[0,1]$, uniqueness forces $\gamma_v=\gamma_w$.
Differentiating at $t=0$ gives $v=\gamma_v'(0)=\gamma_w'(0)=w$.

Finally, by condition (ii) in the definition of $U_p$, the map
$\exp_p$ is a local diffeomorphism at every point of $U_p$: every
$v\in U_p$ has an open neighborhood $N_v$ in $T_pM$ carried
diffeomorphically by $\exp_p$ onto an open neighborhood of
$\exp_p(v)$ in $M$; since $U_p$ is open
(Lemma~\ref{lem:cut-locus-properties}), we may shrink $N_v$ so
that $N_v\subseteq U_p$. Then $W_v:=\exp_p(N_v)$ is an open
neighborhood of $\exp_p(v)$ in $M\setminus{\mathcal C}_p$, and on
$W_v$ the smooth local inverse $(\exp_p|_{N_v})^{-1}$ takes
values in $U_p$; by the injectivity of $\exp_p$ on $U_p$ proved
above, it therefore agrees on $W_v$ with the global
(set-theoretic) inverse of $\exp_p|_{U_p}$. Thus the global
inverse is smooth near each point of $M\setminus{\mathcal C}_p$,
and $\exp_p\colon U_p\to
M\setminus{\mathcal C}_p$ is a diffeomorphism.
\end{proof}

The domain $U_p$ has the following radial structure, which will be
used repeatedly in the volume arguments below.

\begin{lemma}[Radial structure of the domain $U_p$]
\label{lem:cut-time-star-shaped}
\uses{def:cut-locus,lem:cut-locus-properties}
Let $M$ be complete and connected, let $p\in M$, let
$U_p\subset T_pM$ be the
set of Definition~\ref{def:cut-locus} (open, by
Lemma~\ref{lem:cut-locus-properties}), and let $S$ denote the
unit sphere of $(T_pM,g_p)$. For $v\in S$ set
$$c(v)=\sup\{t>0\ |\ tv\in U_p\}\in(0,\infty].$$
Then:
\begin{enumerate}
\item $\{t>0\ |\ tv\in U_p\}=(0,c(v))$; in particular $U_p$ is
star-shaped about the origin;
\item for every $w\in U_p$ the geodesic $\gamma_w$ is the unique
minimal geodesic from $p$ to $\exp_p(w)$, and
$d(p,\exp_p(w))=|w|$;
\item for $v\in S$ the unit-speed radial geodesic
$\gamma_v(t)=\exp_p(tv)$ is minimizing on $[0,t]$ for every
$t<c(v)$, and there is no conjugate point along $\gamma_v$ on
$[0,c(v))$.
\end{enumerate}
\end{lemma}

\begin{proof}
\uses{def:cut-locus,lem:cut-locus-properties,prop:minimal-geodesic-no-conjugate,def:conjugate-point,lem:exponential-differential-jacobi}
(2) By condition (i) of Definition~\ref{def:cut-locus}, $\gamma_w$
is the unique minimal geodesic from $p$ to $\exp_p(w)$; being
minimal, its length $|w|$ equals $d(p,\exp_p(w))$.

(1) Since $U_p$ is an open neighborhood of $0$
(Lemma~\ref{lem:cut-locus-properties}(1)), the set
$\{t>0\,|\,tv\in U_p\}$ is non-empty and open in $(0,\infty)$, so
$c(v)>0$; it therefore suffices to
show that if $tv\in U_p$ and $0<s<t$ then $sv\in U_p$ (downward
closedness together with openness gives exactly the interval
$(0,c(v))$). By
condition (i), $\gamma_{tv}\colon[0,1]\to M$ is the unique minimal
geodesic from $p$ to its endpoint. By
Proposition~\ref{prop:minimal-geodesic-no-conjugate}, its
restriction to $[0,s/t]$ --- which is $\gamma_{sv}$ after
reparameterization --- is the unique minimal geodesic between its
endpoints, so condition (i) holds at $sv$; and there are no
conjugate points along $\gamma_{tv}|_{[0,1)}$, so by
Lemma~\ref{lem:exponential-differential-jacobi},
$\exp_p$ is a local diffeomorphism at $sv$, which is
condition (ii). Hence $sv\in U_p$.

(3) Given $t<c(v)$, pick $t'\in(t,c(v))$; then $t'v\in U_p$, and
the argument for (1) shows that $\gamma_v|_{[0,t]}$ is minimizing
and that there is no conjugate point along $\gamma_v$ at any
parameter in $[0,t')\supseteq[0,t]$.
\end{proof}




\begin{definition}[Injectivity radius]
\label{def:injectivity-radius}
\uses{def:exponential-map,def:cut-locus}
The {\sl injectivity radius $\inj_M(p)$ of $M$ at $p$} is the
supremum of the $r> 0$ for which the restriction of $\exp_p\colon T_pM\to M$ to the ball $B(0,r)$ of radius $r$ in $T_pM$
is a diffeomorphism into $M$. Clearly, $\inj_M(p)$ is the
distance in $T_pM$ from $0$ to the frontier of $U_p$. It is also the
distance in $M$ from $p$ to the cut locus ${\mathcal C}_p$.
\end{definition}

Suppose that  $\inj_M(p)=r$. There are two possibilities:
Either there is a broken, closed geodesic through $p$, broken only
at $p$, of length $2r$, or there is a geodesic $\gamma$ of length
$r$ emanating from $p$ whose endpoint is a conjugate point along
$\gamma$. The first case happens when the exponential mapping is not
one-to-one of the closed ball of radius $r$ in $T_pM$, and the
second happens when there is a tangent vector in $T_pM$ of length
$r$ at which $\exp_p$ is not a local diffeomorphism.
This dichotomy (due to Klingenberg) is made precise by the
following lemma.

\begin{lemma}[Klingenberg: short geodesic loop at the injectivity radius]
\label{lem:klingenberg-loop}
\uses{def:injectivity-radius,def:exponential-map,def:conjugate-point,def:geodesic,thm:hopf-rinow}
Let $M$ be complete and connected, $p\in M$, and suppose
$\rho:=\inj_M(p)<\infty$. If no geodesic from $p$ has a conjugate point at distance
$\le\rho$ along it, then there is a geodesic loop $\sigma\colon
[0,2\rho]\to M$ with $\sigma(0)=\sigma(2\rho)=p$, parameterized at
unit speed, smooth on $(0,2\rho)$ (i.e., broken at most at $p$), of
length exactly $2\rho$.
\end{lemma}

\begin{proof}
\uses{def:injectivity-radius,def:conjugate-point,thm:hopf-rinow,lem:exponential-differential-jacobi}
Throughout, we use that $\exp_p$ is injective on the open ball
$B(0,\rho)\subset T_pM$ (any two points of $B(0,\rho)$ lie in some
$B(0,r)$ with $r<\rho$, on which $\exp_p$ is a diffeomorphism by
the definition of the injectivity radius), and that
$d(\exp_p)$ is non-singular at every $v$ with $|v|\le\rho$ (by
Lemma~\ref{lem:exponential-differential-jacobi}, its
kernel at $v$ is non-trivial exactly when there is a non-zero
Jacobi field along $t\mapsto
\exp_p(tv)$ vanishing at both ends, i.e., a conjugate point at
distance $|v|\le\rho$, which is excluded by hypothesis).

\emph{Step 1: two distinct vectors with the same image.} For every
$\epsilon>0$ the map $\exp_p$ fails to be a diffeomorphism from
$B(0,\rho+\epsilon)$ onto its image. For small $\epsilon$ it is a
local diffeomorphism at each point of
$\overline{B(0,\rho+\epsilon)}$: non-singularity holds on the
compact set $\{|v|\le\rho\}$ by hypothesis, hence on a slightly
larger ball, since the set where $d\exp_p$ is singular is
closed. Therefore, for all small $\epsilon>0$, $\exp_p$ is not
injective on $B(0,\rho+\epsilon)$: there are $a_\epsilon\not=
b_\epsilon$ in $B(0,\rho+\epsilon)$ with
$\exp_p(a_\epsilon)=\exp_p(b_\epsilon)$. Let
$\epsilon\to0$ and pass to convergent subsequences $a_\epsilon\to
v$, $b_\epsilon\to v'$, so $|v|,|v'|\le\rho$ and
$\exp_p(v)=\exp_p(v')=:x$. If $v=v'$ then, $\exp_p$
being a local diffeomorphism at $v$, it is injective on a
neighborhood of $v$ containing both $a_\epsilon$ and $b_\epsilon$
for small $\epsilon$ --- a contradiction. So $v\not=v'$. Both are also non-zero: if, say, $v'=0$, then
$t\mapsto\exp_p(tv)$ is a geodesic from $p$ to $p$ of length
$|v|\le\rho<2\rho$, smooth on the open parameter interval, and the
midpoint argument of Step 4 below (which uses nothing from Steps
2--3) already contradicts the injectivity of $\exp_p$ on
$B(0,\rho)$.

Let $\gamma_1(t)=\exp_p(tv)$ and $\gamma_2(t)=\exp_p(tv')$,
$t\in[0,1]$, two distinct geodesics from $p$ to $x$, of lengths
$\rho_1=|v|$ and $\rho_2=|v'|$, and let $T_1=\gamma_1'(1)/\rho_1$
and $T_2=\gamma_2'(1)/\rho_2$ be their unit terminal velocities at
$x$.

\emph{Step 2: the terminal velocities satisfy $T_1=\pm T_2$.}
Suppose not; then $T_1\not=T_2$ and $T_1\not=-T_2$, so
$w:=-(T_1+T_2)\not=0$ satisfies $\langle w,T_1\rangle
=-1-\langle T_1,T_2\rangle<0$ and likewise $\langle
w,T_2\rangle<0$. Let $c(s)=\exp_x(sw)$. Since
$d(\exp_p)$ is non-singular at $v$ and at $v'$, the inverse
function theorem provides smooth curves $v(s),v'(s)$ in $T_pM$, for
$s$ small, with $v(0)=v$, $v'(0)=v'$ and
$\exp_p(v(s))=\exp_p(v'(s))=c(s)$. Consider the variation
$\tilde\gamma(t,s)=\exp_p(t\,v(s))$ through geodesics from $p$;
its energy is $E(s)=\frac12|v(s)|^2$. By the first-variation
computation of Section~3 (integration by parts, with the
$t$-integral vanishing since each $\tilde\gamma(\cdot,s)$ is a
geodesic, and no boundary term at $t=0$ since the initial point is
fixed),
$$\frac{dE}{ds}\Bigl|_{s=0}\Bigr.
=\Bigl[\langle \widetilde Y,\widetilde
X\rangle\Bigr]_{t=0}^{t=1}
=\langle w,\gamma_1'(1)\rangle=\rho_1\langle w,T_1\rangle<0,$$
and likewise for $v'(s)$. Hence for small $s>0$ both $|v(s)|<
|v|\le\rho$ and $|v'(s)|<|v'|\le\rho$, while
$v(s)\not=v'(s)$ (they are close to the distinct vectors $v,v'$)
and $\exp_p(v(s))=\exp_p(v'(s))$. This contradicts the
injectivity of $\exp_p$ on $B(0,\rho)$. So $T_1=\pm T_2$.

\emph{Step 3: the case $T_1=T_2$ is impossible.} Let $\beta$ be the
unit-speed geodesic with $\beta(0)=x$, $\beta'(0)=-T_1$; it is
defined for all time by completeness. Reversing $\gamma_1$ and
$\gamma_2$ shows $\beta(\rho_1)=p$ and $\beta(\rho_2)=p$. If
$\rho_1=\rho_2$ then $\gamma_1$ and $\gamma_2$ are geodesics with
the same endpoint and same terminal velocity, hence equal (ODE
uniqueness backwards from $t=1$), contradicting $v\not=v'$. So,
say, $\rho_2<\rho_1$, and $\sigma_0:=\beta|_{[\rho_2,\rho_1]}$ is a
geodesic (in particular smooth on the open interval) from $p$ to
$p$ of length $L_0=\rho_1-\rho_2\in(0,\rho]$. The midpoint
argument of Step 4 below, applied to $\sigma_0$, contradicts
injectivity on $B(0,\rho)$ (as $L_0\le \rho<2\rho$). Hence
$T_1=-T_2$.

\emph{Step 4: conclusion.} Since $T_1=-T_2$, the concatenation
$\sigma$ of $\gamma_1$ with the reversal of $\gamma_2$,
parameterized at unit speed on $[0,\rho_1+\rho_2]$, is smooth at
$x$ (the velocities match there) and is therefore a geodesic loop
based at $p$, smooth on the open interval, of length
$L=\rho_1+\rho_2\le2\rho$. We claim $L=2\rho$, which finishes the
proof. Suppose $L<2\rho$ (this includes the situation of the loop
$\sigma_0$ from Step 3, with $L=L_0\le\rho$). The two vectors
$$u_1=\tfrac{L}{2}\,\sigma'(0)\in T_pM,\qquad
u_2=-\tfrac{L}{2}\,\sigma'(L)\in T_pM$$
satisfy $|u_1|=|u_2|=L/2<\rho$ and
$\exp_p(u_1)=\sigma(L/2)=\exp_p(u_2)$ (traverse the loop
from its two ends to its midpoint; $\sigma|_{[L/2,L]}$ reversed is
the geodesic $t\mapsto\exp_p(-t\sigma'(L))$). If $u_1=u_2$,
then the geodesics $t\mapsto\sigma(t)$ and $t\mapsto\sigma(L-t)$
have the same initial data, hence
$\sigma(t)=\sigma(L-t)$ for all $t$; differentiating at $t=L/2$
gives $\sigma'(L/2)=0$, which is absurd for a unit-speed curve
(note that $\sigma$ is smooth at the parameter $L/2$, the loop
being smooth away from its endpoints). So $u_1\not=u_2$ are
distinct points of $B(0,\rho)$ with the same image under
$\exp_p$, contradicting injectivity. Hence $L=2\rho$.
\end{proof}

\subsection{Local isometries and constant curvature}

In this subsection we collect the facts about local isometries and
about manifolds of constant curvature that are used in the proof of
the Uniformization Theorem (Theorem~\ref{thm:uniformization}).

\begin{definition}[Local isometry]
\label{def:local-isometry}
\uses{def:riemannian-metric}
A smooth map $F\colon (N,h)\to (M,g)$ between Riemannian manifolds
of the same dimension is a {\sl local isometry} if $F^*g=h$, i.e.,
if $dF_x\colon T_xN\to T_{F(x)}M$ is a linear isometry for every
$x\in N$. By the inverse function theorem, a local isometry
restricts to an isometry from a sufficiently small neighborhood of
each point onto its image. An {\sl isometry} is a bijective local
isometry; its inverse is then also an isometry.
\end{definition}

\begin{lemma}[Local isometries and the exponential map]
\label{lem:local-isometry-exponential}
\uses{def:local-isometry,def:geodesic,def:exponential-map}
A local isometry $F\colon N\to M$ carries geodesics to geodesics.
If in addition $N$ is complete (as a metric space), then for every
$x\in N$, every $v\in T_xN$ and every $t\in \mathbb{R}$ the point
$\exp_{F(x)}\left(t\,dF_x(v)\right)$ is defined and
$$F\left(\exp_x(tv)\right)=\exp_{F(x)}\left(t\,dF_x(v)\right).$$
\end{lemma}

\begin{proof}
\uses{thm:levi-civita-connection,thm:hopf-rinow}
Being a geodesic is a local property, and on a small neighborhood
$U$ of any point a local isometry is an isometry onto its image. An
isometry carries the Levi-Civita connection to the Levi-Civita
connection: the pushforward under an isometry of the Levi-Civita
connection of the source is a torsion-free connection making the
target metric parallel, hence is the Levi-Civita connection of the
target by the uniqueness in
Theorem~\ref{thm:levi-civita-connection}. Consequently
$\nabla_{\dot\gamma}\dot\gamma=0$ implies $\nabla_{(F\circ
\gamma)^{\textstyle .}}(F\circ\gamma)^{\textstyle .}=0$ on $U$, and
geodesics map to geodesics.

Now suppose $N$ is complete. By the Hopf--Rinow Theorem
(Theorem~\ref{thm:hopf-rinow}), the geodesic
$t\mapsto\exp_x(tv)$ is defined for all $t\in\mathbb{R}$. Its image
$t\mapsto F(\exp_x(tv))$ is then a geodesic of $M$, defined for all
$t$, with initial position $F(x)$ and initial velocity $dF_x(v)$. By
the uniqueness of geodesics with given initial conditions, it
coincides with $t\mapsto \exp_{F(x)}(t\,dF_x(v))$ wherever the
latter is defined; in particular the maximal geodesic with this
initial data is defined for all $t$, and the displayed identity
holds for all $t$.
\end{proof}

Combining this naturality with the uniform normal balls of
Lemma~\ref{lem:normal-ball} shows that a local isometry is, at a
uniform scale over any compact set, an isometry of metric balls.

\begin{lemma}[Local isometries on normal balls]
\label{lem:local-isometry-normal-balls}
\uses{def:local-isometry,def:exponential-map}
Let $F\colon(N,h)\to(M,g)$ be a local isometry and let $K\subset
N$ be compact. Then there is $\delta>0$ such that:
\begin{enumerate}
\item for every $x\in K$, $F$ maps the metric ball $B(x,\delta)$
diffeomorphically onto the metric ball $B(F(x),\delta)$, and
$$F\left(\exp_x(v)\right)=\exp_{F(x)}\left(dF_x(v)\right)
\qquad(|v|<\delta);$$
\item $d_M\left(F(a),F(b)\right)=d_N(a,b)$ for all $a,b\in
B(x,\delta/4)$, $x\in K$;
\item if $c$ is a continuous curve of finite length in $M$ and
$\tilde c$ is a continuous curve in $N$ with $F\circ\tilde c=c$
and $\tilde c([0,1])\subseteq K$, then the lengths agree:
$L_h(\tilde c)=L_g(c)$.
\end{enumerate}
\end{lemma}

\begin{proof}
\uses{lem:normal-ball,lem:local-isometry-exponential,def:local-isometry,def:geodesic}
Apply Lemma~\ref{lem:normal-ball} to the compact set $K\subset N$
and to the compact set $F(K)\subset M$ (the continuous image of a
compact set), and let $\delta>0$ be the smaller of the two
resulting radii. Then for every $x\in K$ and every $0<s\le\delta$
the map $\exp_x$ is a diffeomorphism of $B(0,s)\subset T_xN$ onto
the metric ball $B(x,s)$: the restriction to $B(0,s)$ of the
diffeomorphism provided by Lemma~\ref{lem:normal-ball} remains a
diffeomorphism onto its image, this image consists of points
$\exp_x(v)$ with $d_N(x,\exp_x(v))=|v|<s$, and conversely every
$q\in B(x,s)$ equals $\exp_x(v)$ for some $v$ with
$|v|=d_N(x,q)<s$. The same statements hold at every point of
$F(K)$, with $d_M$ in place of $d_N$. Moreover the balls $B(x,s)$,
$x\in K$, $0<s\le\delta$, are open subsets of $N$: $\exp_x$ is
continuous and injective on $B(0,\delta)$, so by invariance of
domain (applied in charts of $N$) the image of the open set
$B(0,s)$ is open.

Since $dF_x\colon T_xN\to T_{F(x)}M$ is a linear isometry for
every $x\in N$ (Definition~\ref{def:local-isometry}), every
piecewise smooth curve $\sigma$ in $N$ satisfies
$|(F\circ\sigma)'|=|dF_\sigma(\sigma')|=|\sigma'|$ pointwise, so
$L_g(F\circ\sigma)=L_h(\sigma)$. Taking the infimum over piecewise
smooth curves from $a$ to $b$ yields
$$d_M\left(F(a),F(b)\right)\le d_N(a,b)
\qquad\text{for all }a,b\in N,$$
an inequality used repeatedly below.

(1) Fix $x\in K$ and $v\in T_xN$ with $|v|<\delta$. Since $v$ lies
in the domain of $\exp_x$, the geodesic $\gamma_v$ with initial
position $x$ and initial velocity $v$ is defined on $[0,1]$, and
$\exp_x(tv)=\gamma_v(t)$ for $t\in[0,1]$: indeed
$s\mapsto\gamma_v(ts)$ is a geodesic with initial velocity $tv$,
hence coincides with $\gamma_{tv}$ by the uniqueness of solutions
of the geodesic equation (Definition~\ref{def:geodesic}) with
given initial conditions, and one evaluates at $s=1$. By the first
assertion of Lemma~\ref{lem:local-isometry-exponential}, which
involves no completeness hypothesis, $\alpha:=F\circ\gamma_v$ is a
geodesic of $M$ on $[0,1]$, with $\alpha(0)=F(x)$ and, by the
chain rule, $\alpha'(0)=dF_x(v)$. On the other hand
$|dF_x(v)|=|v|<\delta$, so by the choice of $\delta$ the map
$t\mapsto\exp_{F(x)}\left(t\,dF_x(v)\right)$ is defined on $[0,1]$
and, by the same rescaling as above, is the geodesic of $M$ with
initial position $F(x)$ and initial velocity $dF_x(v)$. Two
geodesics with the same initial position and velocity coincide on
their common interval of definition, again by ODE uniqueness;
evaluating at $t=1$ gives
$$F\left(\exp_x(v)\right)=\exp_{F(x)}\left(dF_x(v)\right),
\qquad|v|<\delta.$$
Consequently, on the ball
$B(x,\delta)=\exp_x\left(B(0,\delta)\right)$,
$$F|_{B(x,\delta)}=\exp_{F(x)}\circ\,dF_x\circ
\left(\exp_x|_{B(0,\delta)}\right)^{-1}.$$
Here $\left(\exp_x|_{B(0,\delta)}\right)^{-1}$ is a diffeomorphism
of $B(x,\delta)$ onto $B(0,\delta)\subset T_xN$; the linear
isometry $dF_x$ maps $B(0,\delta)\subset T_xN$ diffeomorphically
onto $B(0,\delta)\subset T_{F(x)}M$; and, since $F(x)\in F(K)$,
$\exp_{F(x)}$ maps the latter ball diffeomorphically onto the
metric ball $B(F(x),\delta)$. Hence $F$ maps $B(x,\delta)$
diffeomorphically onto $B(F(x),\delta)$.

(2) Fix $x\in K$ and $a,b\in B(x,\delta/4)$; then $d_N(a,b)\le
d_N(a,x)+d_N(x,b)<\delta/2$. In view of the displayed inequality
above, only $d_N(a,b)\le d_M(F(a),F(b))$ requires proof. Let
$0<\epsilon<\delta/4$ and choose a piecewise smooth curve
$\sigma\colon[0,1]\to M$ from $F(a)$ to $F(b)$ with
$$L_g(\sigma)<d_M\left(F(a),F(b)\right)+\epsilon\le
d_N(a,b)+\epsilon<\frac{3\delta}{4}.$$
For every $t\in[0,1]$,
$$d_M\left(F(x),\sigma(t)\right)\le
d_M\left(F(x),F(a)\right)+d_M\left(F(a),\sigma(t)\right)\le
d_N(x,a)+L_g(\sigma)<\frac{\delta}{4}+\frac{3\delta}{4}=\delta,$$
since $d_M(F(a),\sigma(t))\le L_g(\sigma|_{[0,t]})\le
L_g(\sigma)$; hence $\sigma$ lies in $B(F(x),\delta)$. By part
(1), $\tilde\sigma:=\left(F|_{B(x,\delta)}\right)^{-1}\circ\sigma$
is a piecewise smooth curve in $B(x,\delta)$ with
$F\circ\tilde\sigma=\sigma$, so $L_h(\tilde\sigma)=L_g(\sigma)$.
Since $F$ is injective on $B(x,\delta)$ and $a,b\in B(x,\delta)$,
the endpoints of $\tilde\sigma$, namely the preimages in
$B(x,\delta)$ of $F(a)$ and of $F(b)$, are $a$ and $b$. Therefore
$$d_N(a,b)\le L_h(\tilde\sigma)=L_g(\sigma)
<d_M\left(F(a),F(b)\right)+\epsilon,$$
and letting $\epsilon\to0$ completes the proof of (2).

(3) For a merely continuous curve the length is, by definition,
the supremum over all partitions $0=t_0\le t_1\le\dots\le t_k=1$
of the sums $\sum_i d\left(c(t_i),c(t_{i+1})\right)$ of the
distances between consecutive points; by the triangle inequality
these sums do not decrease when the partition is refined.

Since the balls $B(y,\delta/8)$, $y\in K$, are open in $N$, the
sets $U_t:=\tilde c^{-1}\left(B(\tilde c(t),\delta/8)\right)$,
$t\in[0,1]$, form an open cover of the compact space $[0,1]$
(here $\tilde c(t)\in K$). Let $\eta>0$ be a Lebesgue number of
this cover. If $|s-s'|\le\eta$, then $s$ and $s'$ lie in a common
$U_t$, whence
$$d_N\left(\tilde c(s),\tilde c(s')\right)\le
d_N\left(\tilde c(s),\tilde c(t)\right)
+d_N\left(\tilde c(t),\tilde c(s')\right)<\frac{\delta}{4}.$$
Call a partition of $[0,1]$ \emph{fine} if its consecutive points
differ by at most $\eta$. For a fine partition $\{t_i\}$ we have
$\tilde c(t_i)\in K$ and
$\tilde c(t_i),\tilde c(t_{i+1})\in B(\tilde c(t_i),\delta/4)$, so
part (2), applied with $x=\tilde c(t_i)$, together with
$F\circ\tilde c=c$, gives
$$d_N\left(\tilde c(t_i),\tilde c(t_{i+1})\right)
=d_M\left(c(t_i),c(t_{i+1})\right)$$
for every $i$: for a fine partition the two partition sums agree
term by term. Finally, every partition admits a fine refinement,
and refining never decreases the partition sums, so the suprema
defining $L_h(\tilde c)$ and $L_g(c)$ are unchanged when
restricted to fine partitions, over which the sums agree term by
term. Hence $L_h(\tilde c)=L_g(c)$; in particular $\tilde c$ has
finite length.
\end{proof}

\begin{lemma}[Rigidity of local isometries]
\label{lem:local-isometry-rigidity}
\uses{def:local-isometry}
Let $N$ be connected and let $F_1,F_2\colon N\to M$ be local
isometries with $F_1(x_0)=F_2(x_0)$ and $d(F_1)_{x_0}=d(F_2)_{x_0}$
for some $x_0\in N$. Then $F_1=F_2$.
\end{lemma}

\begin{proof}
\uses{lem:local-isometry-exponential,def:exponential-map,lem:normal-ball}
Let $A=\{x\in N\,|\, F_1(x)=F_2(x)\text{ and }
d(F_1)_x=d(F_2)_x\}$. It is non-empty by hypothesis, and closed,
since $F_1,F_2$ and their differentials are continuous. It is also
open: given $x\in A$, choose $r>0$ so small that $\exp_x$ is a
diffeomorphism of $B(0,r)\subset T_xN$ onto $B(x,r)$ and both
geodesics $t\mapsto\exp_x(tv)$, $|v|<r$, $t\in[0,1]$, and their
images under $F_1,F_2$ are defined (such $r$ exists by
Lemma~\ref{lem:normal-ball} applied to $\{x\}\subset N$ and to
$\{F_1(x)\}\subset M$). Exactly as in the proof of
Lemma~\ref{lem:local-isometry-exponential}, for $|v|<r$,
$$F_i\left(\exp_x(v)\right)=\exp_{F_i(x)}\left(d(F_i)_x
(v)\right),\qquad i=1,2,$$
and since $F_1(x)=F_2(x)$ and $d(F_1)_x=d(F_2)_x$ the two right-hand
sides agree. Hence $F_1=F_2$ on the open set $B(x,r)$, and then also
$d(F_1)=d(F_2)$ there, so $B(x,r)\subset A$. Since $N$ is connected,
$A=N$.
\end{proof}

\begin{lemma}[Local isometries from complete manifolds are coverings]
\label{lem:local-isometry-covering}
\uses{def:local-isometry}
Let $F\colon N\to M$ be a local isometry, where $N$ is complete and
$M$ is connected. Then $F$ is a surjective covering map.
\end{lemma}

\begin{proof}
\uses{lem:local-isometry-exponential,thm:hopf-rinow,def:exponential-map}
Throughout we use the identity of
Lemma~\ref{lem:local-isometry-exponential} (``naturality''), valid
for all $t$ since $N$ is complete.

\emph{Surjectivity.} The image $F(N)$ is open, $F$ being a local
diffeomorphism. It is also closed: let $y\in\overline{F(N)}$ and
choose $\epsilon>0$ such that $\exp_y$ is a diffeomorphism of
$B(0,\epsilon)\subset T_yM$ onto the neighborhood $B(y,\epsilon)$
of $y$. Choose $z=F(x')\in B(y,\epsilon)$. There is a geodesic
$\sigma\colon [0,1]\to M$ from $z$ to $y$ (the reversal of the
radial geodesic from $y$ to $z$ in $B(y,\epsilon)$). Let
$v=d F_{x'}^{-1}\left(\sigma'(0)\right)$ and set
$c(t)=\exp_{x'}(tv)$; by naturality $F\circ c$ is the geodesic with
initial data $(z,\sigma'(0))$, i.e., $F(c(t))=\sigma(t)$. Then
$F(c(1))=y$, so $y\in F(N)$. Since $M$ is connected, $F(N)=M$.

\emph{Even covering.} Fix $y\in M$ and $\epsilon>0$ as above. For
each $x\in F^{-1}(y)$ define
$$\varphi_x=\exp_x\circ\, dF_x^{-1}\circ
\left(\exp_y|_{B(0,\epsilon)}\right)^{-1}
\colon B(y,\epsilon)\longrightarrow N,$$
which is defined on all of $B(y,\epsilon)$ because $N$ is complete.
By naturality $F\circ\varphi_x=\mathrm{id}_{B(y,\epsilon)}$.
Consequently $d(\varphi_x)$ is everywhere injective, hence (equal
dimensions) $\varphi_x$ is a local diffeomorphism, and it is
injective, having a left inverse. Therefore
$V_x:=\varphi_x(B(y,\epsilon))$ is open and $\varphi_x\colon
B(y,\epsilon)\to V_x$ is a diffeomorphism with inverse $F|_{V_x}$.
We show that $F^{-1}(B(y,\epsilon))$ is the disjoint union of the
$V_x$.

First, given $z'\in F^{-1}(B(y,\epsilon))$, write
$F(z')=\exp_y(u)$ with $|u|<\epsilon$, let $\sigma(t)=\exp_y((1-t)u)$
be the geodesic from $F(z')$ to $y$, and let $c$ be its lift from
$z'$, i.e., the geodesic $c(t)=\exp_{z'}\left(t\,
dF_{z'}^{-1}(\sigma'(0))\right)$, which satisfies $F\circ c=\sigma$
by naturality. Then $x^*:=c(1)\in F^{-1}(y)$. The reversed geodesic
$\bar c(t)=c(1-t)$ runs from $x^*$ to $z'$ and satisfies $F(\bar
c(t))=\sigma(1-t)=\exp_y(tu)$; differentiating at $t=0$ gives
$dF_{x^*}(\bar c'(0))=u$, i.e., $\bar c'(0)=dF_{x^*}^{-1}(u)$.
Hence
$$z'=\bar c(1)=\exp_{x^*}\left(dF_{x^*}^{-1}(u)\right)
=\varphi_{x^*}\left(\exp_y(u)\right)=\varphi_{x^*}(F(z'))\in
V_{x^*}.$$

Second, the $V_x$ are disjoint: suppose $z'\in V_{x_1}\cap V_{x_2}$
with $x_1,x_2\in F^{-1}(y)$, and let $u=\exp_y^{-1}(F(z'))$. For
$i=1,2$, from $z'=\varphi_{x_i}(F(z'))
=\exp_{x_i}(dF_{x_i}^{-1}(u))$ we see, reversing the geodesic
$t\mapsto \exp_{x_i}\left(t\,dF_{x_i}^{-1}(u)\right)$ exactly as
above, that $x_i$ is the endpoint at $t=1$ of the geodesic starting
at $z'$ with initial velocity $dF_{z'}^{-1}(\tau'(0))$, where
$\tau(t)=\exp_y((1-t)u)$ is a curve depending only on $y$ and
$F(z')$. This description does not involve $x_i$, so $x_1=x_2$.

Thus every point of $M$ has an evenly covered neighborhood, and $F$
is a covering map.
\end{proof}

We shall also need the classical Gauss formula identifying the
Levi-Civita connection of a submanifold of Euclidean space.

\begin{lemma}[Levi-Civita connection of a Euclidean submanifold]
\label{lem:submanifold-connection}
\uses{def:riemannian-metric,thm:levi-civita-connection,def:geodesic}
Let $M\subset\Ar^N$ be an embedded smooth submanifold with the
induced Riemannian metric, and for $x\in M$ let
$P_x\colon\Ar^N\to T_xM$ denote the orthogonal projection. Then
the Levi-Civita connection of $M$ is given by
$$\nabla_XY(x)=P_x\left(D_XY(x)\right),$$
where $D_XY$ is the directional derivative in $\Ar^N$ of (any
smooth extension of) $Y$ along $X$. Consequently, a smooth curve
$\gamma$ in $M$ is a geodesic if and only if its Euclidean second
derivative $\gamma''(t)$ is normal to $T_{\gamma(t)}M$ for all
$t$.
\end{lemma}

\begin{proof}
\uses{thm:levi-civita-connection,def:geodesic,def:riemannian-metric}
Set $n=\dim M$. For $x\in M$ we identify $T_xM$, via the
differential of the inclusion $M\hookrightarrow\Ar^N$, with an
$n$-dimensional linear subspace of $\Ar^N$, and we regard a vector
field on $M$ as a smooth map $Y\colon M\to\Ar^N$ with $Y(x)\in
T_xM$ for all $x$. If $F\colon U\to\Ar^N$, with $U\subset\Ar^n$
open, is a local parametrization of $M$ (the inverse of a chart,
viewed as a map into $\Ar^N$), the coordinate fields are
$\partial_i=F_i:=\partial F/\partial x^i$; writing
$Y=Y^i\partial_i=Y^iF_i$ shows that our notion of vector field
agrees with the usual notion of smooth section of $TM$, since the
frame $F_1,\ldots,F_n$ is smooth in both senses. The induced
metric is $g(u,v)=\langle u,v\rangle$ for $u,v\in T_xM$, the
restriction of the Euclidean inner product; as $dF$ is injective,
$g_{ij}=\langle F_i,F_j\rangle$ is smooth and positive definite,
so the induced metric is a Riemannian metric in the sense of
Definition~\ref{def:riemannian-metric}.

\emph{Step 1: the tangential derivative is well defined.} Since
$M$ is embedded, every point of $M$ has a slice chart: an open set
$\Omega\subset\Ar^N$ and a diffeomorphism $\Phi$ from $\Omega$
onto a product $V\times W$ of open balls in
$\Ar^n\times\Ar^{N-n}$ with $\Phi(M\cap\Omega)=V\times\{0\}$.
Then $r:=\Phi^{-1}\circ\pi\circ\Phi\colon\Omega\to M\cap\Omega$,
where $\pi(v,w)=(v,0)$, is a smooth retraction onto $M\cap\Omega$,
and for any vector field $Y$ on $M$ the map $\widetilde Y:=Y\circ
r$ is a smooth extension of $Y|_{M\cap\Omega}$ to $\Omega$. For
$x\in M$, $v\in T_xM$ and any smooth extension $\widetilde Y$ of
$Y$ near $x$, set $D_vY(x):=D\widetilde Y(x)[v]$. This does not
depend on the choice of extension: choosing a smooth curve
$c\colon(-\epsilon,\epsilon)\to M$ with $c(0)=x$ and $c'(0)=v$
(for instance the image of a short line segment under a
parametrization), the chain rule gives
$$D\widetilde Y(x)[v]=\frac{d}{dt}\Big|_{t=0}\widetilde Y(c(t))
=\frac{d}{dt}\Big|_{t=0}Y(c(t)),$$
and the right-hand side involves only $Y$ along $c$. For vector
fields $X,Y$ on $M$ we write $D_XY(x):=D_{X(x)}Y(x)$; by the
above this is well defined and depends on $X$ only through the
value $X(x)$, linearly. Moreover $x\mapsto P_x$ is smooth:
applying the Gram--Schmidt process to a local frame
$F_1,\ldots,F_n$ produces a smooth orthonormal frame
$e_1,\ldots,e_n$ of $TM$, and
$P_x(w)=\sum_{i=1}^n\langle w,e_i(x)\rangle\,e_i(x)$. Hence
$$\nabla^T_XY(x):=P_x\bigl(D_XY(x)\bigr)$$
defines a smooth vector field $\nabla^T_XY$ on $M$, and since
$D_vY(x)$ is linear in $v\in T_xM$, the assignment
$Y\mapsto\nabla^TY$ is an $\mathbb{R}$-linear map
$\Gamma(TM)\to\Gamma(T^*M\otimes TM)$.

\emph{Step 2: $\nabla^T$ satisfies the Leibniz formula and the
two conditions of Theorem~\ref{thm:levi-civita-connection}.}
First, the Leibniz formula: for a smooth function $f$ on $M$ and
a curve $c$ as in Step 1 with $c'(0)=X(x)$,
$$D_X(fY)(x)=\frac{d}{dt}\Big|_{t=0}f(c(t))\,Y(c(t))
=X(f)(x)\,Y(x)+f(x)\,D_XY(x);$$
applying $P_x$ and using $P_x(Y(x))=Y(x)$ (as $Y(x)\in T_xM$)
gives $\nabla^T(fY)=df\otimes Y+f\,\nabla^TY$.

Next, the $g$ orthogonal condition. For vector fields $X,Y,Z$ on
$M$ and $x\in M$, differentiating
$t\mapsto\langle Y(c(t)),Z(c(t))\rangle$ at $t=0$ gives
$$X(g(Y,Z))(x)=\langle D_XY(x),Z(x)\rangle
+\langle Y(x),D_XZ(x)\rangle.$$
Since $Z(x)\in T_xM$ while $D_XY(x)-P_x(D_XY(x))$ is orthogonal
to $T_xM$, we have $\langle D_XY(x),Z(x)\rangle
=\langle P_x(D_XY(x)),Z(x)\rangle=g(\nabla^T_XY,Z)(x)$, and
likewise for the second term. As $X$ and $x$ are arbitrary,
$d(g(Y,Z))=g(\nabla^TY,Z)+g(Y,\nabla^TZ)$.

Finally, torsion-freeness. Fix $x_0\in M$, a slice-chart
neighborhood $\Omega$ as in Step 1, and the extensions
$\widetilde X:=X\circ r$, $\widetilde Y:=Y\circ r$. In the linear
coordinates $z^1,\ldots,z^N$ of $\Ar^N$, the Lie bracket of
vector fields on the open set $\Omega$ has components
$[\widetilde X,\widetilde Y]^k
=\widetilde X^j\partial_j\widetilde Y^k
-\widetilde Y^j\partial_j\widetilde X^k$ (the second-order terms
cancel by the symmetry of second partial derivatives), that is,
$$[\widetilde X,\widetilde Y]
=D_{\widetilde X}\widetilde Y-D_{\widetilde Y}\widetilde X
\quad\text{on }\Omega.$$
We claim that $[\widetilde X,\widetilde Y](x)=[X,Y](x)$ for $x\in
M\cap\Omega$, where $[X,Y]$ denotes the Lie bracket of $X$ and
$Y$ as vector fields on the manifold $M$; in particular the
left-hand side is then tangent to $M$. Indeed, for any smooth
function $f$ on $\Omega$ and $x\in M\cap\Omega$, since
$\widetilde Y(x)=Y(x)\in T_xM$ we get
$$(\widetilde Yf)(x)=df(x)[Y(x)]=d(f|_M)(x)[Y(x)]
=\bigl(Y(f|_M)\bigr)(x);$$
applying the same identity to $\widetilde X$ and the smooth
function $\widetilde Yf$, and then with the roles of $X$ and $Y$
exchanged,
$$([\widetilde X,\widetilde Y]f)(x)
=\bigl(X(Y(f|_M))-Y(X(f|_M))\bigr)(x)
=\bigl([X,Y](f|_M)\bigr)(x)=df(x)\bigl[[X,Y](x)\bigr].$$
Taking for $f$ the coordinate functions $z^1,\ldots,z^N$ shows
$[\widetilde X,\widetilde Y](x)=[X,Y](x)$, as claimed. Since
$D_{\widetilde X}\widetilde Y(x)=D\widetilde Y(x)[X(x)]
=D_XY(x)$ and likewise with $X$ and $Y$ exchanged, we conclude
that $D_XY(x)-D_YX(x)=[X,Y](x)$ lies in $T_xM$, and applying
$P_x$ yields
$$\nabla^T_XY-\nabla^T_YX=[X,Y].$$

\emph{Step 3: identification of the connections.} By Steps 1 and
2, $\nabla^T$ is a torsion-free connection on $TM$ making $g$
parallel. By the uniqueness assertion of
Theorem~\ref{thm:levi-civita-connection}, $\nabla^T=\nabla$, the
Levi-Civita connection of the induced metric. This proves the
displayed formula.

\emph{Step 4: the geodesic criterion.} Let $\gamma\colon I\to M$
be a smooth curve and $t_0\in I$. Choose a local parametrization
$F\colon U\to\Ar^N$ around $\gamma(t_0)$ and write
$\gamma(t)=F(x(t))$ for $t$ near $t_0$, with
$x=(x^1,\ldots,x^n)$ smooth. For the coordinate fields
$\partial_j=F_j$, the curve $s\mapsto F(u+se_i)$ passes through
$F(u)$ with velocity $F_i(u)$, so by Step 1
$$D_{\partial_i}\partial_j\,(F(u))
=\frac{d}{ds}\Big|_{s=0}F_j(u+se_i)=F_{ij}(u),\qquad
F_{ij}:=\frac{\partial^2F}{\partial x^i\partial x^j},$$
and therefore, by Step 3 and the coordinate description
$\nabla_{\partial_i}\partial_j=\Gamma^k_{ij}\partial_k$ of the
Levi-Civita connection in terms of the Christoffel symbols
$\Gamma^k_{ij}$ of the induced metric in this chart,
$$P\bigl(F_{ij}\bigr)=\nabla_{\partial_i}\partial_j
=\Gamma^k_{ij}F_k$$
at each point of $U$ (evaluation points suppressed). By the
chain rule in $\Ar^N$,
$$\gamma'(t)=\dot x^i(t)\,F_i(x(t)),\qquad
\gamma''(t)=\ddot x^k(t)\,F_k(x(t))
+\dot x^i(t)\dot x^j(t)\,F_{ij}(x(t)),$$
so applying $P_{\gamma(t)}$, which fixes each $F_k(x(t))$,
$$P_{\gamma(t)}\bigl(\gamma''(t)\bigr)
=\Bigl(\ddot x^k(t)+\Gamma^k_{ij}(x(t))\,
\dot x^i(t)\dot x^j(t)\Bigr)F_k(x(t)).$$
Since $F_1,\ldots,F_n$ form a basis of $T_{\gamma(t)}M$, the
projection $P_{\gamma(t)}(\gamma''(t))$ vanishes if and only if
$\ddot x^k+\Gamma^k_{ij}\dot x^i\dot x^j=0$ for all $k$, which is
exactly the local-coordinate form of the geodesic equation
$\nabla_{\dot\gamma}\dot\gamma=0$ of
Definition~\ref{def:geodesic}. On the other hand,
$P_{\gamma(t)}(\gamma''(t))=0$ if and only if $\gamma''(t)$ is
orthogonal to $T_{\gamma(t)}M$. Since every $t_0\in I$ lies in
the domain of such a chart and the geodesic condition is
chart-independent, $\gamma$ is a geodesic if and only if
$\gamma''(t)$ is normal to $T_{\gamma(t)}M$ for all $t\in I$.
\end{proof}

We now solve the Jacobi equation explicitly in constant curvature.
For $\lambda\in\mathbb{R}$ let $s_\lambda\colon
\mathbb{R}\to\mathbb{R}$ be the solution of
$$s_\lambda''+\lambda\, s_\lambda=0,\qquad s_\lambda(0)=0,\quad
s_\lambda'(0)=1,$$
that is, $s_\lambda(r)=\frac{1}{\sqrt\lambda}\sin(\sqrt\lambda\,r)$
for $\lambda>0$, $s_0(r)=r$, and
$s_\lambda(r)=\frac{1}{\sqrt{-\lambda}}\sinh(\sqrt{-\lambda}\,r)$
for $\lambda<0$. Set
$$D_\lambda=\begin{cases}\pi/\sqrt{\lambda}&\lambda>0\\
+\infty&\lambda\le 0,\end{cases}$$
so that $s_\lambda>0$ on $(0,D_\lambda)$.

\begin{lemma}[Jacobi fields and the exponential map in constant curvature]
\label{lem:constant-curvature-jacobi}
\uses{lem:constant-curvature-tensor,def:geodesic,def:exponential-map,thm:hopf-rinow}
Let $(M,g)$ be a complete Riemannian $n$-manifold of constant
sectional curvature $\lambda$ and let $p\in M$.
\begin{enumerate}
\item Let $\gamma$ be a unit-speed geodesic with $\gamma(0)=p$ and
$X=\gamma'$. The Jacobi fields along $\gamma$ vanishing at $0$ are
exactly the fields
$$Y(t)=a\,t\,X(t)+s_\lambda(t)E(t),$$
where $a\in \mathbb{R}$ and $E$ is a parallel field along $\gamma$
with $E(0)$ perpendicular to $X(0)$; here
$\nabla_XY(0)=aX(0)+E(0)$.
\item $\exp_p$ is defined on all of $T_pM$ and is non-singular at
every $v\in T_pM$ with $|v|<D_\lambda$.
\item On $B(0,D_\lambda)\setminus\{0\}\subset T_pM$, in polar
coordinates $(r,\theta)$ (where $\theta$ denotes coordinates on the
unit sphere of $(T_pM,g_p)$),
$$\exp_p^*g=dr^2+s_\lambda(r)^2\,g_{S^{n-1}},$$
where $g_{S^{n-1}}$ is the round metric on the unit sphere of
$(T_pM,g_p)$; moreover $\exp_p^*g$ is a smooth metric on all of
$B(0,D_\lambda)$.
\end{enumerate}
\end{lemma}

\begin{proof}
\uses{lem:constant-curvature-tensor,claim:curvature-symmetries-bianchi,thm:hopf-rinow}
By completeness and the Hopf--Rinow Theorem
(Theorem~\ref{thm:hopf-rinow}), $\exp_p$ is defined on all of
$T_pM$.

(1) By Lemma~\ref{lem:constant-curvature-tensor}, for any vector
field $Y$ along the unit-speed geodesic $\gamma$,
$$\langle{\mathcal R}(Y,X)X,W\rangle={\mathcal R}(Y,X,W,X)
=\lambda\left(\langle Y,W\rangle\langle X,X\rangle
-\langle Y,X\rangle\langle X,W\rangle\right)$$
for every $W$, so ${\mathcal R}(Y,X)X=\lambda\left(Y-\langle
Y,X\rangle X\right)=\lambda Y^\perp$, where $Y^\perp$ denotes the
component of $Y$ perpendicular to $X$. The Jacobi equation
$\nabla_X\nabla_XY+{\mathcal R}(Y,X)X=0$ therefore reads
$\nabla_X\nabla_XY=-\lambda Y^\perp$. Write $u(t)=\langle
Y,X\rangle$; since $X$ is parallel along $\gamma$,
$u''=\langle\nabla_X\nabla_XY,X\rangle=-\lambda\langle
Y^\perp,X\rangle=0$, so $u(t)=u(0)+u'(0)t$. Choosing a parallel
orthonormal frame $E_1=X,E_2,\ldots,E_n$ along $\gamma$ and writing
$Y=uX+\sum_{i\ge 2}y_iE_i$, the perpendicular components satisfy
$y_i''+\lambda y_i=0$. For a Jacobi field with $Y(0)=0$ we get
$u(t)=at$ with $a=\langle \nabla_XY(0),X(0)\rangle$ and
$y_i(t)=y_i'(0)s_\lambda(t)$, which is the asserted form with
$E=\sum_{i\ge2}y_i'(0)E_i$. Conversely each such field solves the
Jacobi equation and vanishes at $0$.

(2) As recalled in the discussion of the exponential mapping above,
for $v,w\in T_pM$ the differential of $\exp_p$ at $v$
applied to $w$ equals $Y_w(1)$, where $Y_w$ is the Jacobi field
along the geodesic $t\mapsto\exp_p(tv)$ with $Y_w(0)=0$ and
$\nabla Y_w(0)=w$. Fix $v$ with $0<r:=|v|<D_\lambda$ and let
$\gamma(s)=\exp_p(s\,v/r)$ be the unit-speed reparameterization. If
$\widetilde Y$ is a Jacobi field along $\gamma$ then
$t\mapsto\widetilde Y(rt)$ is a Jacobi field along
$t\mapsto\exp_p(tv)$, and every Jacobi field arises this way; under
this correspondence $\nabla Y(0)=r\,\nabla\widetilde Y(0)$.
Writing $w=cv/r+w^\perp$ with $w^\perp\perp v$, part (1) gives
\begin{equation}\label{dexpformula}
d(\exp_p)_v(w)=Y_w(1)=\widetilde Y(r)
=c\,\gamma'(r)+\frac{s_\lambda(r)}{r}\,P_\gamma(w^\perp),
\end{equation}
where $P_\gamma$ denotes parallel transport along $\gamma$ from $p$
to $\gamma(r)$ (extended linearly). Since $P_\gamma(w^\perp)$ is
perpendicular to $\gamma'(r)$ (parallel transport preserves inner
products and $w^\perp\perp\gamma'(0)$), the right side vanishes only
if $c=0$ and $s_\lambda(r)w^\perp=0$; as $s_\lambda(r)>0$ for
$0<r<D_\lambda$, the differential is injective. At $v=0$,
$d(\exp_p)_0=\mathrm{id}$.

(3) Fix a unit vector $v$ and $0<r<D_\lambda$. In polar coordinates,
the coordinate field $\partial_r$ at the point $rv$ corresponds to
the vector $v\in T_pM$, and the tangent space to the sphere of
radius $r$ at $rv$ consists of the vectors $w\perp v$. By
Equation~(\ref{dexpformula}) (with $c=1,w^\perp=0$, respectively
$c=0$),
$$d(\exp_p)_{rv}(v)=\gamma'(r),\qquad
d(\exp_p)_{rv}(w)=\frac{s_\lambda(r)}{r}P_\gamma(w)\quad (w\perp
v).$$
Hence, using that parallel transport is a linear isometry and
$\gamma'(r)\perp P_\gamma(w)$:
$$(\exp_p^*g)(v,v)=|\gamma'(r)|^2=1,\qquad
(\exp_p^*g)(v,w)=0,\qquad
(\exp_p^*g)(w_1,w_2)
=\frac{s_\lambda(r)^2}{r^2}\langle w_1,w_2\rangle$$
for $w,w_1,w_2\perp v$. A tangent vector to the unit sphere
$\theta\mapsto$ direction, scaled to the sphere of radius $r$, has
Euclidean length $r$ times its unit-sphere length; the displayed
formulas say exactly that
$\exp_p^*g=dr^2+s_\lambda(r)^2g_{S^{n-1}}$. Finally $\exp_p^*g$ is
smooth on $B(0,D_\lambda)$, being the pullback of a smooth metric
under a smooth map, and it is a metric (positive definite) there
because $d\exp_p$ is injective at each point of $B(0,D_\lambda)$ by
part (2).
\end{proof}

The following consequence packages the case $\lambda\le0$ of the
preceding lemma for repeated use.

\begin{lemma}[Polar model of nonpositive constant curvature]
\label{lem:model-polar-isometry}
\uses{def:sectional-curvature,def:exponential-map,def:local-isometry,lem:constant-curvature-jacobi}
Let $(M,g)$ be a complete, simply-connected Riemannian $n$-manifold
of constant sectional curvature $\lambda\le0$ and let $p\in M$.
Then $G:=\exp_p^*g$ is a smooth metric on $T_pM$, given in
polar coordinates on $T_pM\setminus\{0\}$ by
$$G=dr^2+s_\lambda(r)^2\,g_{S^{n-1}},$$
and $\exp_p\colon(T_pM,G)\to(M,g)$ is an isometry.
\end{lemma}

\begin{proof}
\uses{lem:constant-curvature-jacobi,lem:local-isometry-covering,def:local-isometry}
Since $D_\lambda=\infty$,
Lemma~\ref{lem:constant-curvature-jacobi} says that $\exp_p$ is
defined on all of $T_pM$, is everywhere non-singular, and the
pullback $G=\exp_p^*g$ is a smooth metric on $T_pM$ with the
stated polar form. By construction $\exp_p\colon (T_pM,G)\to
(M,g)$ is a local isometry.

We claim $(T_pM,G)$ is complete. First, for any piecewise smooth
curve $c$ in $T_pM$, the polar form gives
$|c'|_G\ge |\frac{d}{dt}(r\circ c)|$ wherever $c\not=0$, so
$L_G(c)\ge |r(c(1))-r(c(0))|$; hence any $G$-Cauchy sequence lies in
some set $K_R=\{r\le R\}$, and moreover any curve of $G$-length
$<1$ starting in $K_R$ stays in $K_{R+1}$. Second, the function
$s_\lambda(r)/r$ extends to a positive continuous function on
$[0,R+1]$ (with value $1$ at $r=0$), so there are constants $0<m\le
M'$ with $m^2\,g_{\mathrm{euc}}\le G\le M'^2\,g_{\mathrm{euc}}$ on
$K_{R+1}$, comparing the polar form with
$g_{\mathrm{euc}}=dr^2+r^2g_{S^{n-1}}$. Hence a $G$-Cauchy sequence
in $K_R$ is Cauchy for the Euclidean metric, so it converges in the
closed set $K_R$ to some $x_\infty$, and comparing lengths of
straight segments (which lie in $K_{R+1}$) shows
$d_G(x_k,x_\infty)\le M'\,|x_k-x_\infty|\to 0$. So $(T_pM,G)$ is
complete.

By Lemma~\ref{lem:local-isometry-covering}, $\exp_p\colon
(T_pM,G)\to M$ is a covering map. Since $M$ is simply connected and
$T_pM$ is connected, classical covering-space theory says the
covering is one-sheeted, i.e., $\exp_p$ is a bijective local
isometry, hence an isometry.
\end{proof}

We can now verify the properties of the three model spaces asserted
after Definition~\ref{def:sectional-curvature}.

\begin{lemma}[The model spaces]
\label{lem:model-spaces}
\uses{def:sectional-curvature,def:geodesic,lem:constant-curvature-tensor}
Let $n\ge2$.
\begin{enumerate}
\item $(\Ar^n,g_{\mathrm{euc}})$ is complete, simply connected, of
constant sectional curvature $0$.
\item For $\rho>0$ the sphere $S^n_\rho\subset\Ar^{n+1}$ of radius
$\rho$, with the induced metric $\rho^2g_{\mathrm{st}}$, is
complete, simply connected, of constant sectional curvature
$1/\rho^2$.
\item The upper half-space $\{x\in\Ar^n\,|\,x^n>0\}$ with the
metric $g_{ij}=(x^n)^{-2}\delta_{ij}$ is complete, simply
connected, of constant sectional curvature $-1$. (This is the
upper half-space presentation of $(\HH^n,g_{\mathrm{st}})$
recalled after Definition~\ref{def:sectional-curvature}; the
Poincar\'e disk presentation is carried to it by the classical
inversion map, as verified at the end of the proof.)
\end{enumerate}
\end{lemma}

\begin{proof}
\uses{lem:constant-curvature-tensor,prop:cone-curvature,def:cone,def:curvature-operator,thm:levi-civita-connection,claim:curvature-symmetries-bianchi,def:riemann-curvature-tensor}
We use throughout that multiplying a metric by a constant $c>0$
leaves the Christoffel symbols (Equation~(\ref{Gamma})), hence the
$(1,3)$ curvature tensor, unchanged, multiplies the $(0,4)$-tensor
by $c$, and multiplies each sectional curvature by $1/c$
(evaluate on a rescaled orthonormal basis); it also multiplies
distances by $\sqrt c$, so preserves completeness.

(1) In Cartesian coordinates $g_{ij}=\delta_{ij}$, so
$\Gamma^k_{ij}=0$ and $R_{ijkl}=0$; hence every sectional
curvature vanishes. Geodesics are straight lines, defined for all
time; $(\Ar^n,d_{\mathrm{euc}})$ is a complete metric space, and
$\Ar^n$ is convex, hence simply connected.

(2) By the scaling rule it suffices to treat $\rho=1$. In polar
coordinates the Euclidean metric on $\Ar^{n+1}\setminus\{0\}$ is
$dr^2+r^2g_{S^n}$, which is exactly the cone metric
(Definition~\ref{def:cone}) over $(S^n,g_{\mathrm{st}})$. By (1)
its curvature operator vanishes, so
Proposition~\ref{prop:cone-curvature} (applied with $N=S^n$) gives
$$0=s^2\left(\Rm_{g_{\mathrm{st}}}-\wedge^2g_{\mathrm{st}}\right)
\qquad\text{on }\wedge^2T_qS^n,$$
i.e., $\Rm_{g_{\mathrm{st}}}=\wedge^2g_{\mathrm{st}}$ as bilinear
forms on $\wedge^2TS^n$. A $(0,4)$-tensor antisymmetric in its
first two and in its last two indices (as the curvature tensor is,
by Claim~\ref{claim:curvature-symmetries-bianchi}) is determined by
the bilinear form it induces on $\wedge^2$, so
$R_{ijkl}=g_{ik}g_{jl}-g_{il}g_{jk}$ on $S^n$, and
Lemma~\ref{lem:constant-curvature-tensor} gives constant sectional
curvature $1$. The sphere is compact, hence complete as a metric
space, and it is simply connected for $n\ge2$ (classical
topology).

(3) Write $g=e^{2\phi}\delta$ with $\phi=-\log x^n$, so
$\partial_i\phi=-\delta_{in}/x^n$. From
Equation~(\ref{Gamma}),
$$\Gamma^k_{ij}=\delta_{kj}\,\partial_i\phi
+\delta_{ik}\,\partial_j\phi-\delta_{ij}\,\partial_k\phi
=\frac{-\delta_{kj}\delta_{in}-\delta_{ik}\delta_{jn}
+\delta_{ij}\delta_{kn}}{x^n},$$
so that
$$\nabla_{\partial_i}\partial_j=\frac{-\delta_{jn}\partial_i
-\delta_{in}\partial_j+\delta_{ij}\partial_n}{x^n},\qquad
\nabla_{\partial_i}\partial_n=-\frac{\partial_i}{x^n}.$$
A direct computation of ${\mathcal
R}(\partial_i,\partial_j)\partial_k
=\nabla_{\partial_i}\nabla_{\partial_j}\partial_k
-\nabla_{\partial_j}\nabla_{\partial_i}\partial_k$ from these
formulas (expanding
$\nabla_{\partial_i}\left(\nabla_{\partial_j}\partial_k\right)$ by
the Leibniz rule, with $\partial_i(1/x^n)=-\delta_{in}/(x^n)^2$;
all terms carrying two indices equal to $n$ cancel in the
antisymmetrization) leaves
$${\mathcal R}(\partial_i,\partial_j)\partial_k
=\frac{\delta_{ik}\,\partial_j-\delta_{jk}\,\partial_i}{(x^n)^2}.$$
Hence, using $g_{ab}=\delta_{ab}/(x^n)^2$,
$$R_{ijkl}=g\left({\mathcal
R}(\partial_i,\partial_j)\partial_l,\partial_k\right)
=\frac{\delta_{il}\delta_{jk}-\delta_{jl}\delta_{ik}}{(x^n)^4}
=-\left(g_{ik}g_{jl}-g_{il}g_{jk}\right),$$
and Lemma~\ref{lem:constant-curvature-tensor} gives constant
sectional curvature $-1$. For completeness: for any piecewise
smooth curve $c$,
$$L_g(c)=\int\frac{|c'|_{\mathrm{euc}}}{x^n(c)}\,dt
\ \ge\ \int\left|\frac{d}{dt}\log x^n(c(t))\right|dt
\ \ge\ \left|\log\frac{x^n(c(1))}{x^n(c(0))}\right|,$$
so points of a $d_g$-bounded set have $x^n$ bounded above and away
from $0$. Let $(x_k)$ be a $d_g$-Cauchy sequence; it lies in a
slab $S_{a,b}=\{a\le x^n\le b\}$ with $0<a\le b<\infty$. Applying
the displayed estimate to initial segments shows that every
piecewise smooth curve of $g$-length $<1$ starting in $S_{a,b}$
stays in the enlarged slab $S_{a/e,\,be}$, on which
$(be)^{-2}\delta\le g\le(a/e)^{-2}\delta$. Hence, for $k,l$ large,
curves from $x_k$ to $x_l$ of $g$-length
$<\min\left(2\,d_g(x_k,x_l),1\right)$ stay in $S_{a/e,\,be}$ and
give $|x_k-x_l|_{\mathrm{euc}}\le 2be\,d_g(x_k,x_l)$: the sequence
is Euclidean-Cauchy, so it converges to some $x_\infty$ in the
closed slab $S_{a,b}$. The straight segments from $x_k$ to
$x_\infty$ lie in the convex slab $S_{a,b}$, where $g\le
a^{-2}\delta$, so $d_g(x_k,x_\infty)\le
a^{-1}|x_k-x_\infty|_{\mathrm{euc}}\to0$. Thus $(M,g)$ is
complete. The half-space is convex, hence simply
connected.

Finally we verify that the Poincar\'e disk presentation is carried
to the upper half-space presentation by the classical inversion.
Let
$$\iota(x)=\frac{2\,(x+e_n)}{|x+e_n|^2}-e_n,$$
the inversion in the sphere of radius $\sqrt2$ centered at $-e_n$;
it is an involution of $\Ar^n\setminus\{-e_n\}$. From the
elementary identity
$\left|\frac{a}{|a|^2}-\frac{b}{|b|^2}\right|=\frac{|a-b|}{|a|\,|b|}$
(expand both sides), applied with $a=x+e_n$, $b=x'+e_n$,
\begin{equation}\label{invconf}
|\iota(x)-\iota(x')|
=\frac{2\,|x-x'|}{|x+e_n|\,|x'+e_n|}.
\end{equation}
Direct computation of the $n$-th coordinate gives, using
$|x+e_n|^2=|x|^2+2x^n+1$,
$$\iota(x)^n=\frac{2(x^n+1)}{|x+e_n|^2}-1
=\frac{1-|x|^2}{|x+e_n|^2},$$
so $\iota$ maps the open unit disk into the upper half-space; and
since $\iota(e_n)=0$, Equation~(\ref{invconf}) with $x'=e_n$ gives
$|\iota(x)|=|x-e_n|/|x+e_n|$, which is $<1$ exactly when $x^n>0$,
so $\iota$ maps the upper half-space into the disk. Being an
involution, $\iota$ therefore restricts to mutually inverse
diffeomorphisms between the disk and the half-space. By
Equation~(\ref{invconf}), $d\iota_x$ multiplies all lengths by the
factor $2/|x+e_n|^2$, so $\iota^*\delta=
\left(2/|x+e_n|^2\right)^2\delta$ as quadratic forms; hence
$$\iota^*\left(\frac{\delta}{(y^n)^2}\right)
=\frac{\left(\frac{2}{|x+e_n|^2}\right)^2\,\delta}
{\left(\frac{1-|x|^2}{|x+e_n|^2}\right)^2}
=\frac{4\,\delta}{\left(1-|x|^2\right)^2},$$
which is precisely the Poincar\'e metric of the disk presentation.
\end{proof}

\section{Computations in Gaussian normal coordinates}

In this section we compute the metric and the Laplacian (on
functions) in local Gaussian coordinates. A direct computation shows
that in Gaussian normal coordinates on a metric ball about $p\in M$
the metric takes the form
\begin{eqnarray}\label{metricexp}
g_{ij}(x)& = & \delta_{ij} + \frac{1}{3}R_{iklj}x^kx^l +
\frac{1}{6}R_{iklj,s}x^kx^lx^s \\
& & +(\frac{1}{20}R_{iklj,st} +\frac{2}{45}\sum_m
R_{iklm}R_{jstm})x^kx^lx^sx^t + O(r^{5}), \nonumber
\end{eqnarray}
where $r$ is the distance from $p$. (See, for example Proposition
3.1 on page 41 of~\cite{Sakai}, with the understanding that, with
the conventions there, the quantity $R_{ijkl}$ there differs by sign
from ours.) The fourth-order coefficients in
Equation~(\ref{metricexp}) are a classical computation which we do
not reproduce here, citing~\cite{Sakai} for them; the second- and
third-order terms, which are the ones used in the sequel, we now
derive from the Jacobi equation.

\begin{lemma}[Leading terms of the metric in normal coordinates]
\label{lem:normal-coordinates-expansion}
\uses{def:exponential-map,def:riemann-curvature-tensor}
In Gaussian normal coordinates $(x^1,\ldots,x^n)$ centered at $p$,
$$g_{ij}(x)=\delta_{ij}+\frac13R_{iklj}x^kx^l
+\frac16R_{iklj,s}x^kx^lx^s+O(|x|^4),$$
where the curvature tensor and its covariant derivative are
evaluated at $p$ in the orthonormal frame defining the coordinates.
\end{lemma}

\begin{proof}
\uses{lem:constant-curvature-jacobi,claim:curvature-symmetries-bianchi,def:geodesic}
Fix a unit vector $v\in T_pM$, let $\gamma(t)=\exp_p(tv)$, and
let $\{e_i\}$ be the orthonormal basis of $T_pM$ defining the
coordinates. As recalled in the discussion of the exponential map,
the coordinate field $\partial_i$ at the point $\exp_p(tv)$
equals $J^{(i)}(t)/t$, where $J^{(i)}$ is the Jacobi field along
$\gamma$ with $J^{(i)}(0)=0$ and $\nabla J^{(i)}(0)=e_i$ (this is
the description of $d(\exp_p)_{tv}$ by Jacobi fields, applied to
the radial rescaling as in the proof of
Lemma~\ref{lem:constant-curvature-jacobi}(2)). Hence
$$g_{ij}(tv)=\frac{1}{t^2}\left\langle
J^{(i)}(t),J^{(j)}(t)\right\rangle .$$

Write $R_t$ for the symmetric operator $X\mapsto{\mathcal
R}(X,\gamma'(t))\gamma'(t)$ on $\gamma'(t)^\perp\oplus\mathbb{R}
\gamma'(t)$, so that the Jacobi equation reads $J''=-R_t(J)$, where
primes denote covariant derivatives along $\gamma$. Differentiating
the Jacobi equation covariantly and evaluating at $t=0$, using
$J(0)=0$, $J'(0)=w$:
$$J''(0)=-R_0(J(0))=0,\qquad J'''(0)=-R_0'(J(0))-R_0(J'(0))
=-R_0(w),$$
$$J''''(0)=-R_0''(J(0))-2R_0'(J'(0))-R_0(J''(0))=-2R_0'(w),$$
where $R_0'=\left(\nabla_{\gamma'}{\mathcal
R}\right)(\cdot,\gamma')\gamma'$ evaluated at $0$, since $\gamma'$
is parallel. The covariant Taylor expansion in a parallel frame
therefore gives
$$J(t)=t\,w-\frac{t^3}{6}R_0(w)-\frac{t^4}{12}R_0'(w)+O(t^5),$$
with the remainder locally uniform in $v$. Taking inner products
(parallel transport preserves $\langle\cdot,\cdot\rangle$) and
using the symmetry of $R_0$ and $R_0'$, which follows from the pair
symmetry of the curvature tensor and of its covariant derivative
(Claim~\ref{claim:curvature-symmetries-bianchi}),
$$\left\langle J^{(i)}(t),J^{(j)}(t)\right\rangle
=t^2\delta_{ij}-\frac{t^4}{3}\left\langle
R_0(e_i),e_j\right\rangle-\frac{t^5}{6}\left\langle
R_0'(e_i),e_j\right\rangle+O(t^6),$$
so that
$$g_{ij}(tv)=\delta_{ij}-\frac{t^2}{3}\,{\mathcal
R}(e_i,v,e_j,v)-\frac{t^3}{6}\,\left(\nabla_v{\mathcal
R}\right)(e_i,v,e_j,v)+O(t^4),$$
using $\langle R_0(e_i),e_j\rangle=\langle{\mathcal
R}(e_i,v)v,e_j\rangle={\mathcal R}(e_i,v,e_j,v)$. Finally, setting
$x=tv$ and using the antisymmetry of the curvature tensor in its
last two arguments,
$$-\frac{t^2}{3}{\mathcal R}(e_i,v,e_j,v)
=+\frac13{\mathcal R}(e_i,x,x,e_j)=\frac13R_{iklj}x^kx^l,$$
and in the same way
$-\frac{t^3}{6}(\nabla_v{\mathcal R})(e_i,v,e_j,v)
=\frac16R_{iklj,s}x^kx^lx^s$. This gives the asserted expansion,
in particular confirming, with our conventions, the signs of the
second- and third-order terms of Equation~(\ref{metricexp}).
\end{proof}

Let $\gamma$ be a geodesic in $M$ emanating from $p$ in the direction $v$.
Choose local coordinates $\theta^1,\ldots,\theta^{n-1}$ on the unit sphere in
$T_pM$ in a neighborhood of $v/|v|$. Then $(r, \theta^{1}, ... ,\theta^{n-1})$
are local coordinates at any point of the ray emanating from the origin in the
$v$ direction (except at $p$). Transferring these via $\exp_p$ produces
local coordinates $(r,\theta^{1}, ... ,\theta^{n-1})$ along $\gamma$. Using
Gauss's lemma (Lemma 12 of Chapter 5 on p. 133 of \cite{Petersen}), we can
write the metric locally as
\[ g=dr^{2}+
r^{2}h_{ij}(r,\theta)d\theta^{i}\otimes d\theta^{j}. \] Then the
volume form
\begin{align*}
dV &= \sqrt{\det(g_{ij})}dr\wedge d\theta^{1}
\wedge\cdots\wedge
d\theta^{n-1}\\
& = r^{n-1}\sqrt{\det(h_{ij})}dr\wedge d\theta^{1}
\wedge\cdots\wedge d\theta^{n-1}.
\end{align*}



\begin{lemma}[Laplacian in local coordinates, Christoffel form]
\label{lem:laplacian-christoffel-formula}
\uses{lem:hessian-symmetric,thm:levi-civita-connection}
\lean{MorganTianLib.laplacianAt_eq_chart_formula,MorganTianLib.laplacianAt_eq_invGram_sum}
\leanok
In any coordinate chart, the Laplacian of a smooth function $f$
satisfies
\[ \triangle f
   = g^{ij}\left(\partial_i\partial_jf-\Gamma^k_{ij}\partial_kf\right)
   = g^{ij}\partial_i\partial_jf+b^k\partial_kf,
   \qquad b^k:=-g^{ij}\Gamma^k_{ij}, \]
where $(g^{ij})$ is the inverse of the Gram matrix $(g_{ij})$ of the
metric in the chart and $\Gamma^k_{ij}$ are the Christoffel symbols of
the Levi-Civita connection. In particular, in a chart the Laplacian is
a second-order operator whose leading coefficient matrix $(g^{ij})$ is
symmetric positive definite and whose coefficients $g^{ij}$, $b^k$ are
smooth.
\end{lemma}

\begin{proof}
\uses{lem:hessian-symmetric,thm:levi-civita-connection}
\leanok
By the definition of the Laplacian as the metric trace of the Hessian
(Equation~(\ref{laplacformula})), $\triangle f(p)=\sum_i
\Hess(f)(X_i,X_i)$ for any $g$-orthonormal basis $\{X_i\}$ of $T_pM$.
This trace equals the $g^{ij}$-weighted sum over the coordinate frame:
writing $C$ for the matrix expanding the orthonormal basis in the
coordinate frame, $X_i=\sum_aC_{ai}\partial_a$, orthonormality reads
$C^TGC=\mathrm{Id}$ with $G=(g_{ab})$ the Gram matrix, whence
$G^{-1}=CC^T$ and, for any bilinear form $B$,
$$\sum_iB(X_i,X_i)=\sum_{a,b}\Big(\sum_iC_{ai}C_{bi}\Big)B(\partial_a,\partial_b)
=\sum_{a,b}(CC^T)_{ab}B(\partial_a,\partial_b)
=g^{ab}B(\partial_a,\partial_b).$$
Hence $\triangle f=g^{ij}\Hess(f)(\partial_i,\partial_j)$. By the
local-coordinate formula for the Hessian
(Lemma~\ref{lem:hessian-symmetric}) and
$\nabla_{\partial_i}\partial_j=\Gamma^k_{ij}\partial_k$ for the
Levi-Civita connection (Theorem~\ref{thm:levi-civita-connection}),
$$\triangle
f=g^{ij}\Hess(f)(\partial_i,\partial_j)=g^{ij}\left(\partial_i\partial_j(f)-
\nabla_{\partial_i}\partial_jf\right)
=g^{ij}\partial_i\partial_jf-g^{ij}\Gamma_{ij}^k\partial_kf.$$
Positive definiteness of $(g^{ij})$: it is the inverse of the positive
definite symmetric matrix $G$, and for $\xi\neq0$, setting
$\eta:=G^{-1}\xi\neq0$ gives
$\xi^{T}G^{-1}\xi=(G\eta)^{T}G^{-1}(G\eta)=\eta^{T}G\eta>0$.
Smoothness of the coefficients follows from smoothness of the metric:
$G$ is smooth with positive determinant, so $g^{ij}$ (by Cramer's
rule) and $\Gamma^k_{ij}$ (differences of derivatives of $G$
contracted with $g^{ij}$) are smooth.
\end{proof}

\begin{lemma}[Laplacian in local coordinates]
\label{lem:laplacian-local-formula}
\uses{lem:laplacian-christoffel-formula}
The \textit{Laplacian operator} acting on scalar
functions on $M$ is given in local coordinates by
\[ \triangle  = \frac{1}{\sqrt{\det(g)}}\partial_i
\left(g^{ij}\sqrt{\det(g)}\partial_j\right). \]
\end{lemma}

\begin{proof}
\uses{lem:hessian-symmetric,lem:laplacian-christoffel-formula}
Let us compute the derivative at a point $p$. We have
$$\frac{1}{\sqrt{\det(g)}}\partial_i
\left(g^{ij}\sqrt{\det(g)}\partial_j\right)f=g^{ij}\partial_i\partial_jf+\partial_ig^{ij}\partial_jf+
\frac{1}{2}g^{ij}\partial_iTr(\tilde g)\partial_jf,$$ where $\tilde
g=g(p)^{-1}g$. On the other hand, by the Christoffel form of the
coordinate Laplacian (Lemma~\ref{lem:laplacian-christoffel-formula}),
$$\triangle
f=g^{ij}\partial_i\partial_jf-g^{ij}\Gamma_{ij}^k\partial_kf.$$ Thus,
to prove the claim it suffices to show that
$$g^{ij}\Gamma_{ij}^k=-(\partial_ig^{ik}+\frac{1}{2}g^{ik}\Tr(\partial_i\tilde g)).$$
From the definition of the Christoffel symbols we have
$$g^{ij}\Gamma_{ij}^k=\frac{1}{2}g^{ij}g^{kl}(\partial_ig_{jl}+\partial_jg_{il}-\partial_lg_{ij}).$$
Of course, $g^{ij}\partial_ig_{jl}=-\partial_ig^{ij}g_{jl}$, so that
 $g^{ij}g^{kl}\partial_ig_{jl}=-\partial_ig^{ik}$. It follows by symmetry that $g^{ij}g^{jl}\partial_jg_{il}
 =-\partial_ig^{ik}$. The last term is clearly $-\frac{1}{2}g^{ik}\Tr(\partial_i\tilde g).$
\end{proof}

Using Gaussian local coordinates near $p$, we have
\begin{align*}
\triangle r &=
\frac{1}{r^{n-1}\sqrt{\det(h)}}\partial_r\left(r^{n-1}\sqrt{\det(h)}\right)\\
&= \frac{n-1}{r}
 + \partial_r\log\left(\sqrt{\det(h)}\right).
\end{align*}
From this one computes directly that
 \[ \triangle r
 = \frac{n-1}{r} - \frac{r}{3}\Ric(v,v)+ O(r^{2}), \]
 where $v = \dot r(0)$, cf,  p.265-268 of \cite{Petersen}. So
 $$\triangle r \leq
 \frac{n-1}{r}\ \ \text{when } r \ll 1\ \ \text{and } \Ric > 0.$$
This local computation has a global analogue, in which the
inequality holds across the cut locus in the following weak sense.

\begin{remark}[Laplacian in weak sense]
\label{rem:laplacian-weak-sense}
\label{calabi}
\uses{lem:laplacian-local-formula}
Let $f$ be a locally Lipschitz function on $M$ and let $h$ be a
locally integrable function on $M$. The statement that $\triangle
f\leq h$ in {\sl the sense of distributions} (or equivalently {\sl
in the weak sense}) means that for any non-negative test function
$\phi$, that is to say for any non-negative compactly supported
$C^\infty$-function $\phi$, we have
$$\int_Mf\triangle \phi\, d\vol\le \int_M h\,\phi\, d\vol.$$
By Rademacher's theorem a locally Lipschitz function is
differentiable almost everywhere, its almost-everywhere defined
gradient $\nabla f$ is locally bounded, and it coincides with the
distributional gradient; in particular the restriction of $\nabla
f$ to any compact subset of $M$ is an $L^2$ one-form. Moreover, for
every test function $\phi$,
$$\int_Mf\triangle
\phi\, d\vol=-\int_M\langle \nabla f,\nabla\phi\rangle\, d\vol.$$
Indeed, covering the compact support of $\phi$ by finitely many
charts and using a subordinate partition of unity, it suffices to
prove this in a single chart; there, by the divergence form of the
Laplacian (Lemma~\ref{lem:laplacian-local-formula}),
$$\sqrt{\det g}\ \triangle\phi
=\partial_i\left(\sqrt{\det g}\,g^{ij}\,\partial_j\phi\right),$$
and the assertion becomes the integration-by-parts identity for the
locally Lipschitz function $f$ against the compactly supported
smooth vector field with components $\sqrt{\det
g}\,g^{ij}\partial_j\phi$, which is the statement that the
almost-everywhere gradient of $f$ represents its distributional
gradient (standard real analysis).
\end{remark}

\begin{example}[Distance function Laplacian bound]
\label{ex:distance-function-laplacian}
\uses{lem:laplacian-local-formula,def:ricci-curvature,rem:laplacian-weak-sense}
Let $(M,g)$ be a complete connected Riemannian manifold of
dimension $n\ge2$, let $p\in M$, and let $f(x)=d(p,x)$ be the
distance function from $p$. If $(M, g)$ has
$\Ric \geq 0$, then
\[ \triangle f \leq \frac{n-1}{f} \]
 in the sense of distributions (the meaning of which is explained
in Remark~\ref{rem:laplacian-weak-sense} above).
\end{example}

\begin{proof}
\uses{def:cut-locus,lem:cut-locus-properties,prop:exponential-diffeomorphism-cut-locus,lem:cut-time-star-shaped,lem:volume-element-comparison,lem:geodesic-polar-form,def:ricci-curvature,rem:laplacian-weak-sense}
(Completeness makes the apparatus of the cut locus available;
$n\ge2$ ensures that the boundary term at $t=0$ in the integration
by parts below vanishes.) The triangle inequality gives
$|f(x)-f(y)|\le d(x,y)$, so $f$ is Lipschitz and
Remark~\ref{rem:laplacian-weak-sense} applies: the assertion means
that for every non-negative test function $\phi$,
$$-\int_M\langle\nabla f,\nabla\phi\rangle\,d\vol
\le\int_M\frac{n-1}{f}\,\phi\,d\vol,$$
where we have already rewritten the left-hand side by the
integration-by-parts identity of that remark. Fix such a $\phi$,
and fix $R>0$ so large that $f\le R$ on the (compact) support of
$\phi$.

Let $U_p\subset T_pM$ be the open set of
Definition~\ref{def:cut-locus}, whose image misses only the
measure-zero set ${\mathcal C}_p$
(Lemma~\ref{lem:cut-locus-properties}), let
$c(v)\in(0,\infty]$ be the cut time of
Lemma~\ref{lem:cut-time-star-shaped}, so that $\{t>0\,|\,tv\in
U_p\}=(0,c(v))$ for unit vectors $v$, and let $\lambda(t,v)$ be the
volume element in polar coordinates. Here
Lemma~\ref{lem:geodesic-polar-form} applies around every ray of
$U_p$: since $U_p$ is open and star-shaped, for each unit vector
$v$ and $t<c(v)$ there are a neighborhood $V$ of $v$ in the unit
sphere and $r_0>t$ with $\{sw\,|\,w\in V,\ 0\le s<r_0\}\subset
U_p\cup\{0\}$ (the cut time is lower semicontinuous, $U_p$ being
open), and $\exp_p$ is non-singular on this cone. On
$\exp_p(U_p)\setminus\{p\}$ the function $f$ is smooth: indeed
$f(\exp_p(w))=|w|$ for $w\in U_p$, so $f$ corresponds under
$\exp_p$ to the radial coordinate $r$, and its gradient is the
radial field $\partial_r$
(Lemma~\ref{lem:geodesic-polar-form}(2)). Hence, writing $\tilde
\phi=\phi\circ\exp_p$, we have $\langle\nabla
f,\nabla\phi\rangle=\partial_t\tilde\phi$ on $U_p$, and by the
change of variables formula and Fubini's theorem,
$$\int_M\langle\nabla f,\nabla\phi\rangle\,d\vol
=\int_{S}\int_0^{c(v)}\partial_t\tilde\phi(tv)\,
\lambda(t,v)\,dt\,d\mathrm{vol}_S(v).$$

Before integrating by parts we record the bounds that make all the
integrals below finite. First, if $\tilde\phi(tv)\not=0$ then
$\exp_p(tv)$ lies in the support of $\phi$, so
$t=d(p,\exp_p(tv))\le R$ (Lemma~\ref{lem:cut-time-star-shaped}(2));
thus all integrands below vanish for $t>R$. Second,
$|\partial_t\tilde\phi(tv)|=|\langle\nabla\phi,\partial_r\rangle|
\le\sup_M|\nabla\phi|=:C$, since $|\partial_r|=1$. Third, the
radial geodesic $\gamma_v$ is minimizing and free of conjugate
points on $[0,c(v))$ (Lemma~\ref{lem:cut-time-star-shaped}(3)), and
$\Ric\ge0$, so Lemma~\ref{lem:volume-element-comparison} with $k=0$
applies: $t\mapsto\lambda(t,v)/t^{n-1}$ is non-increasing with
limit $1$; in particular $0\le\lambda(t,v)\le t^{n-1}$,
$\lambda(t,v)\to0$ as $t\to0$, $\lambda(\cdot,v)$ has a finite
limit at $c(v)$ when $c(v)<\infty$, and
$$\partial_t\lambda(t,v)=\Tr(A)\,\lambda(t,v)
\le\frac{n-1}{t}\,\lambda(t,v).$$
Consequently $|\partial_t\tilde\phi|\,\lambda\le C\,t^{n-1}$ and
$0\le\tilde\phi\,\frac{n-1}{t}\,\lambda\le(n-1)
\left(\sup\phi\right)t^{n-2}$, both supported in $t\le R$, so the
integrals of these two integrands over $(0,c(v))$, and their
iterated integrals over $S$, converge absolutely ($n\ge2$).

Fix a unit vector $v$; we claim
$$-\int_0^{c(v)}\partial_t\tilde\phi\,\lambda\,dt
\le\int_0^{c(v)}\tilde\phi\,\frac{n-1}{t}\,\lambda\,dt.$$
Integrate by parts on a compact subinterval
$[\epsilon,T]\subset(0,c(v))$ (all functions are smooth there) and
use the pointwise bound on $\partial_t\lambda$:
$$-\int_\epsilon^{T}\partial_t\tilde\phi\,\lambda\,dt
=-\Bigl[\tilde\phi\lambda\Bigr]_{\epsilon}^{T}
+\int_\epsilon^{T}\tilde\phi\,\partial_t\lambda\,dt
\le\tilde\phi(\epsilon v)\lambda(\epsilon,v)
-\tilde\phi(Tv)\lambda(T,v)
+\int_\epsilon^{T}\tilde\phi\,\frac{n-1}{t}\,\lambda\,dt.$$
As $\epsilon\to0$ the first boundary term tends to $0$ (as
$\lambda\to0$ and $\tilde\phi$ is bounded), and the term
$-\tilde\phi(Tv)\lambda(T,v)$ is $\le0$; letting $T\to c(v)$ (or
$T\to\infty$ when $c(v)=\infty$, the integrands vanishing for
$t>R$), the two integrals converge to their absolutely convergent
limits, proving the claim. Integrating
over $S$ and changing variables back,
$$-\int_M\langle\nabla f,\nabla\phi\rangle\,d\vol
\le\int_S\int_0^{c(v)}\frac{n-1}{t}\,
\tilde\phi(tv)\,\lambda(t,v)\,dt\,d\mathrm{vol}_S(v)
=\int_M\frac{n-1}{f}\,\phi\,d\vol,$$
since $t=f$ under the change of variables and ${\mathcal C}_p$ has
measure zero (Lemma~\ref{lem:cut-locus-properties}), while both
iterated integrals converge absolutely by the bounds recorded
above. This is the required inequality.
\end{proof}
\noindent [Compare \cite{Petersen}, p. 284 Lemma 42].

\medskip

Since $\abs{\nabla f} = 1$ and $\triangle f$ is the mean curvature
of the geodesic sphere $\partial B(x,r)$, $\Ric(v,v)$ measures
the difference of the mean curvature between the standard Euclidean
sphere and the geodesic sphere in the direction $v$. Another
important geometric object  is the shape operator associated to $f$,
denoted $S$. By definition it is the Hessian of $f$; i.e.,
$S=\nabla^2f=\Hess(f)$.

\section{Basic curvature comparison results}
In this section we will recall some of the basic curvature
comparison results in Riemannian geometry. The reader can refer to
\cite{Petersen}, Section 1 of Chapter 9 for details.

\smallskip

We fix a point $p\in M$. For any real number $k\ge 0$ let $H^n_k$ denote the
simply connected, complete Riemannian $n$-manifold of constant sectional
curvature $-k$. Fix a point $q_k\in H^n_k$, and consider the exponential map
$\exp_{q_k}\colon T_{q_k}(H^n_k)\to H^n_k$. This map is a global
diffeomorphism. Let us consider the pullback, $\tilde h_k$, of the Riemannian
metric on $H^n_k$ to $T_{q_k}H^n_k$. A formula for this tensor is easily given
in polar coordinates on $T_{q_k}(H^n_k)$ in terms of the following function.

\begin{definition}[Sinh comparison function]
\label{def:sinh-comparison-function}
\lean{MorganTianLib.snK}
\leanok
We define a function $\sn_k$ as follows:$$\sn_k(r)=\begin{cases}
r \ \  & \text{if }k=0 \\
\frac{1}{\sqrt{k}}\sinh(\sqrt{k}r) \ \ & \text{if } k>0.
\end{cases}$$
\end{definition}

The function $\sn_{k}(r)$ is the solution to the equation
\begin{align*}
 \varphi'' - k\varphi &= 0,\\
\varphi(0) &= 0,\\
\varphi'(0) &= 1.
\end{align*}
We define
 $\ct_{k}(r) = \frac{\sn_{k}^{'}(r)}{\sqrt{k}\sn_{k}(r)}$.

Before stating the comparison theorems we assemble the technical
apparatus on which all of them rest: the structure of the metric in
geodesic polar coordinates, the Riccati equation satisfied by the
shape operator of the geodesic spheres, a scalar Riccati comparison
lemma, and the elementary Sturm comparison bounding the location of
conjugate points from below.

The Gauss Lemma — the assertion that the polar coordinate spheres meet
the radial geodesics orthogonally — rests on a single fact about Jacobi
fields, which we isolate first because it is proved entirely along one
geodesic, with no reference to the exponential map or to polar
coordinates.

\begin{lemma}[Radial Jacobi fields stay orthogonal to the geodesic]
\label{lem:radial-jacobi-orthogonal}
\uses{def:geodesic,def:conjugate-point,def:riemann-curvature-tensor,claim:curvature-symmetries-bianchi}
\lean{MorganTianLib.chartCurvature_inner_velocity_eq_zero,MorganTianLib.innerVelocity,MorganTianLib.innerVelocity_eq_chartMetricInner_of_geodesicAt,MorganTianLib.IsJacobiFieldOn.hasDerivAt_chartInner_fst_velocity,MorganTianLib.IsJacobiFieldOn.hasDerivAt_chartInner_snd_velocity,MorganTianLib.IsJacobiFieldAlongOn.innerVelocity_snd_eq,MorganTianLib.IsJacobiFieldAlongOn.innerVelocity_fst_eq,MorganTianLib.IsJacobiFieldAlongOn.innerVelocity_fst_eq_zero,MorganTianLib.IsJacobiFieldAlongOn.innerVelocity_snd_eq_zero}
\leanok
Let $\gamma$ be a geodesic on $[a,b]$ in a Riemannian manifold $(M,g)$,
and let $J$ be a Jacobi field along $\gamma$, with covariant derivative
$\nabla J$. Then:
\begin{enumerate}
\item $\langle\nabla J,\dot\gamma\rangle$ is constant on $[a,b]$;
\item $\langle J,\dot\gamma\rangle$ is an affine function of $t$:
$$\langle J(t),\dot\gamma(t)\rangle
=\langle J(a),\dot\gamma(a)\rangle
+(t-a)\,\langle\nabla J(a),\dot\gamma(a)\rangle;$$
\item consequently, if $J$ is {\sl radial}, i.e.\ $J(a)=0$ and
$\langle\nabla J(a),\dot\gamma(a)\rangle=0$, then both
$\langle J(t),\dot\gamma(t)\rangle=0$ and
$\langle\nabla J(t),\dot\gamma(t)\rangle=0$ for every $t\in[a,b]$.
\end{enumerate}
\end{lemma}

\begin{proof}
\uses{claim:curvature-symmetries-bianchi,thm:levi-civita-connection,def:geodesic}
\leanok
Write $X=\dot\gamma$; since $\gamma$ is a geodesic, $\nabla_XX=0$.

(1) The Levi-Civita connection is metric-compatible, so
$$\frac{d}{dt}\langle\nabla J,X\rangle
=\langle\nabla_X\nabla J,X\rangle+\langle\nabla J,\nabla_XX\rangle
=\langle\nabla_X\nabla_XJ,X\rangle
=-\langle{\mathcal R}(J,X)X,X\rangle,$$
the last step by the Jacobi equation
$\nabla_X\nabla_XJ+{\mathcal R}(J,X)X=0$. Now
$\langle{\mathcal R}(J,X)X,X\rangle={\mathcal R}(J,X,X,X)$, and by the
antisymmetry of the curvature tensor in its last two arguments
(Claim~\ref{claim:curvature-symmetries-bianchi}, $R_{ijkl}=-R_{ijlk}$)
this quantity equals its own negative, hence vanishes. So
$\langle\nabla J,X\rangle$ has vanishing derivative at every interior
time; being continuous on $[a,b]$, it is constant there.

(2) Again by metric compatibility and $\nabla_XX=0$,
$$\frac{d}{dt}\langle J,X\rangle
=\langle\nabla J,X\rangle+\langle J,\nabla_XX\rangle
=\langle\nabla J,X\rangle,$$
which by (1) is the constant $\langle\nabla J(a),X(a)\rangle$. A
continuous function on $[a,b]$ whose derivative is the constant $c$ at
every interior time is affine with slope $c$; applying this to
$t\mapsto\langle J(t),X(t)\rangle$ gives the displayed formula.

(3) Under the stated hypotheses both coefficients in (2) vanish, so
$\langle J(t),X(t)\rangle=0$; and by (1),
$\langle\nabla J(t),X(t)\rangle=\langle\nabla J(a),X(a)\rangle=0$.
\end{proof}

A second consequence of the curvature symmetries, also needed below, is
that the operator appearing in the Jacobi equation is self-adjoint.

\begin{lemma}[The Jacobi operator is self-adjoint]
\label{lem:jacobi-operator-symmetric}
\uses{def:riemann-curvature-tensor,claim:curvature-symmetries-bianchi}
\lean{MorganTianLib.curvatureFormAt_jacobi_symm,MorganTianLib.chartCurvature_inner_symm,MorganTianLib.chartCurvatureEndo_inner_symm}
\leanok
Let $u\in T_qM$. The {\sl Jacobi operator} $R_u\colon T_qM\to T_qM$,
$R_u(X)={\mathcal R}(X,u)u$, is self-adjoint:
$$\langle R_u(X),Y\rangle=\langle X,R_u(Y)\rangle
\qquad\text{for all }X,Y\in T_qM.$$
\end{lemma}

\begin{proof}
\uses{claim:curvature-symmetries-bianchi}
\leanok
Unwinding the definitions, $\langle R_u(X),Y\rangle={\mathcal
R}(X,u,u,Y)$ and $\langle X,R_u(Y)\rangle={\mathcal R}(Y,u,u,X)$, so
the claim is that ${\mathcal R}(X,u,u,Y)={\mathcal R}(Y,u,u,X)$. By
antisymmetry in the last two arguments
(Claim~\ref{claim:curvature-symmetries-bianchi},
$R_{ijkl}=-R_{ijlk}$),
$${\mathcal R}(X,u,u,Y)=-{\mathcal R}(X,u,Y,u),\qquad
{\mathcal R}(Y,u,u,X)=-{\mathcal R}(Y,u,X,u),$$
and by the pair symmetry $R_{ijkl}=R_{klij}$ of the same claim,
${\mathcal R}(X,u,Y,u)={\mathcal R}(Y,u,X,u)$. Combining the three
identities gives ${\mathcal R}(X,u,u,Y)={\mathcal R}(Y,u,u,X)$.
\end{proof}

% Lean status of lem:geodesic-polar-form: part (1) — the Gauss-Lemma
% orthogonality, i.e. that the polar fields (Jacobi fields vanishing at p with
% initial covariant derivative tangent to the unit sphere) stay orthogonal to
% the radial direction — is formalized and checked as
% lem:radial-jacobi-orthogonal.
%
% Part (3): the FRAME FORM of the matrix Jacobi equation is now formalized and
% checked at the MANIFOLD level, along a geodesic that may cross arbitrarily many
% charts, in MorganTianLib/Ch01/FrameJacobiSystem.lean:
%   MorganTianLib.exists_orthonormalParallelFrameAlong  — a parallel g-orthonormal
%       frame exists along any geodesic (Gram matrix of a parallel family is
%       constant, so orthonormality propagates from the initial point);
%   MorganTianLib.IsJacobiFieldAlongOn.exists_frameJacobiSystem — in that frame the
%       Jacobi equation becomes the CLOSED linear scalar system c' = d,
%       d' = sum_j R_ij c_j with R_ij = R(E_j, gamma', gamma', E_i) symmetric,
%       i.e. c'' = R c, with no manifold vectors and no charts left in it.
% This is exactly the "J'' = -R(r) J in a parallel frame" step of the proof of (3)
% and the shape IsRadialJacobi consumes.  Supporting new lemmas, both intrinsic:
%   MorganTianLib.curvatureFormAt_jacobi_symm      (self-adjointness of R_u on T_pM)
%   MorganTianLib.metricInner_orthonormal_expansion (v = sum_i <v,e_i> e_i on T_qM)
%
% Part (3) is now COMPLETE at the manifold level: the three gaps flagged here
% previously (operator-valued assembly, curv_cont / curv_bound, and the
% Ioo-vs-Icc endpoint case) are all closed, and the resulting datum is packaged
% as the new node lem:jacobi-metric-comparison below.  Specifically:
%   MorganTianLib.continuousOn_frameCurv (MorganTianLib/Ch01/FrameCurvContinuity.lean)
%       — r |-> R_ij(r) is CONTINUOUS along a geodesic crossing arbitrarily many
%       charts.  This was the missing analytic input.  It cannot be proved by
%       unfolding: curvatureFormAt goes through a Classical.choose extension, and
%       a parallel field's E-representation jumps at every chart change (it is
%       carried at its own moving foot).  The proof routes through the chart
%       bridge curvatureFormAt_chartVectorRep, in which the two jumps cancel, and
%       glues the resulting single-chart continuity by a local-to-global argument.
%       Boundedness (curv_bound) is then compactness of [0,b].
%   MorganTianLib.exists_isRadialJacobi (MorganTianLib/Ch01/RadialJacobiExists.lean)
%       — the OPERATOR-valued J(r) exists: the matrix Jacobi equation is a linear
%       ODE in the Banach algebra E ->L E, so exists_isJacobiSolOn_Icc applies with
%       initial data (0, 1).  Its columns are vector Jacobi fields
%       (IsJacobiSolOn.apply), which is what identifies J(r)w with the geometry.
%   MorganTianLib.exists_isRadialJacobi_of_geodesic (Ch01/FrameRadialBridge.lean)
%       — the producer: a geodesic yields an honest IsRadialJacobi datum, AND
%       every Jacobi field along it with J(0)=0 is a column of it,
%       frameVec J t = J(t) . frameVec (nabla J) 0, by ODE uniqueness.
%       The endpoint case is handled by taking the geodesic on a strictly larger
%       interval [a,b], a < 0 < B < b, so every time of [0,B] is interior.
%
% Parts (2) (S = Hess(r) identified with the operator A) and (4) (the polar volume
% density lambda = det J) remain: they need the polar-coordinate description of
% exp_p, i.e. lem:exponential-differential-jacobi.
\begin{lemma}[Radial geometry in geodesic polar coordinates]
\label{lem:geodesic-polar-form}
\uses{def:exponential-map,def:geodesic,lem:hessian-symmetric,def:riemann-curvature-tensor,def:sectional-curvature,lem:radial-jacobi-orthogonal,lem:jacobi-operator-symmetric}
Let $p\in M$, let $V$ be an open set of unit vectors in $T_pM$, let
$0<r_0\le\infty$, and suppose that $\exp_p$ is defined on the cone
$C=\{tw\,|\,w\in V,\ 0\le t<r_0\}$ and is non-singular at every
point of $C$. Let $\tilde g=\exp_p^*g$ on $C$, let
$(r,\theta^1,\ldots,\theta^{n-1})$ be polar coordinates on $C$
(where $\theta$ are local coordinates on $V$), and let $\hat
g=\hat g_{ij}(\theta)d\theta^i\otimes d\theta^j$ denote the round
metric of the unit sphere of $T_pM$ in these coordinates. Then:
\begin{enumerate}
\item (Gauss Lemma) $\tilde g=dr^2+g_{ij}(r,\theta)\,
d\theta^i\otimes d\theta^j$; the coordinate fields
$\partial_{\theta^i}$ restrict along each radial geodesic
$\gamma_w(t)=\exp_p(tw)$ to Jacobi fields vanishing at $t=0$, and
the radial field $\partial_r$ restricts to $\gamma_w'$.
\item Let $f=r$ be the radial coordinate function and
$S=\Hess(f)$. Then the gradient of $f$ is the radial field
$\partial_r$, $S(\partial_r,\cdot)=0$,
$$S_{ij}:=S(\partial_{\theta^i},\partial_{\theta^j})
=\tfrac{1}{2}\partial_rg_{ij},$$
and the associated $g$-symmetric operator $A$ on the tangent space
of the coordinate spheres, defined by $\langle
A(X),Y\rangle=S(X,Y)$, equals $X\mapsto\nabla_X\partial_r$ and
satisfies along each radial geodesic the Riccati equation
$$\nabla_{\partial_r}A+A^2+R_{\partial_r}=0,\qquad
R_{\partial_r}(X):={\mathcal R}(X,\partial_r)\partial_r,$$
where for unit $X\perp\partial_r$ one has $\langle
R_{\partial_r}(X),X\rangle=K(X\wedge\partial_r)$, the sectional
curvature of the plane spanned by $X$ and $\partial_r$.
\item As $r\to0$, locally uniformly in $\theta$,
$$g_{ij}(r,\theta)=r^2\left(\hat g_{ij}(\theta)+O(r^2)\right),
\qquad A(r)=\frac{1}{r}\,\mathrm{Id}+O(r).$$
\item Setting $\lambda(r,\theta)
=\sqrt{\det\left(g_{ij}(r,\theta)\right)/\det\left(\hat
g_{ij}(\theta)\right)}$, the volume form of $\tilde g$ is
$\lambda\, dr\wedge d\mathrm{vol}_{S^{n-1}}$ and
$$\partial_r\log\lambda=\Tr(A),$$
where $\Tr$ is the trace of the operator $A$. Moreover
$\lambda(r,\theta)/r^{n-1}\to 1$ as $r\to 0$, locally uniformly in
$\theta$.
\end{enumerate}
\end{lemma}

\begin{proof}
\uses{def:conjugate-point,claim:curvature-symmetries-bianchi,lem:hessian-symmetric,thm:levi-civita-connection,lem:exponential-differential-jacobi,lem:second-order-linear-ode,lem:jacobi-small-time,lem:parallel-frame,lem:jacobi-matrix-inverse,lem:radial-jacobi-orthogonal}
(1) Fix $w\in V$ and consider the variation
$\tilde\gamma(t,\theta)=\exp_p(t\,v(\theta))$ through radial
geodesics, where $v(\theta)\in V$. As in the discussion of Jacobi
fields, each field
$J_i=\partial\tilde\gamma/\partial\theta^i$ is a Jacobi field along
the radial geodesic, vanishing at $t=0$, with $\nabla
J_i(0)=\partial v/\partial\theta^i$ perpendicular to $v$ (the
$\theta$-derivatives of a map into the unit sphere are tangent to
it). By the chain rule $\partial_r$ corresponds to
$\gamma_w'$, so $\tilde g(\partial_r,\partial_r)=1$. Each $J_i$ is
therefore a radial Jacobi field in the sense of
Lemma~\ref{lem:radial-jacobi-orthogonal}(3), whence
$\langle J_i,\gamma_w'\rangle\equiv0$, i.e.,
$\tilde g(\partial_r,\partial_{\theta^i})=0$. This proves the
displayed form of $\tilde g$.

(2) From (1), $dr(\partial_r)=1$ and $dr(\partial_{\theta^i})=0$
while $\tilde g(\partial_r,\cdot)=dr$, so the gradient of $f=r$ is
$\partial_r$. By the formula for the Hessian
(Lemma~\ref{lem:hessian-symmetric}),
$S(X,Y)=\langle\nabla_X\partial_r,Y\rangle$, so the operator $A$ is
$X\mapsto\nabla_X\partial_r$, and it is $g$-symmetric because the
Hessian is a symmetric two-tensor. Since the radial curves are
geodesics, $S(\partial_r,\cdot)
=\langle\nabla_{\partial_r}\partial_r,\cdot\rangle=0$. As
$[\partial_r,\partial_{\theta^i}]=0$,
$$\partial_rg_{ij}
=\langle\nabla_{\partial_r}\partial_{\theta^i},\partial_{\theta^j}\rangle
+\langle\partial_{\theta^i},\nabla_{\partial_r}\partial_{\theta^j}\rangle
=\langle\nabla_{\partial_{\theta^i}}\partial_r,\partial_{\theta^j}\rangle
+\langle\partial_{\theta^i},\nabla_{\partial_{\theta^j}}\partial_r\rangle
=2S_{ij}.$$
For the Riccati equation, write $U=\partial_r$ and let $X$ be a
coordinate field $\partial_{\theta^i}$, so $[U,X]=0$ and
$\nabla_UU=0$. Then, by the definition of the curvature tensor,
$$\nabla_U(A(X))=\nabla_U\nabla_XU
=\nabla_X\nabla_UU+{\mathcal R}(U,X)U=-{\mathcal R}(X,U)U
=-R_{\partial_r}(X),$$
while by the Leibniz rule
$\nabla_U(A(X))=(\nabla_UA)(X)+A(\nabla_UX)
=(\nabla_UA)(X)+A(\nabla_XU)=(\nabla_UA)(X)+A^2(X)$. Comparing the
two expressions and using that the coordinate fields span the
tangent spaces of the spheres gives the Riccati equation. Finally,
for unit $X\perp U$, $\langle R_{\partial_r}(X),X\rangle={\mathcal
R}(X,U,X,U)=K(X\wedge U)$ directly from the definition of sectional
curvature.

(3) Fix $w\in V$ and a parallel orthonormal frame
$E_1,\ldots,E_{n-1}$ of $w^\perp$ along $\gamma_w$
(Lemma~\ref{lem:parallel-frame}), and let
${\mathcal J}(r)$ be the matrix whose columns are the components in
this frame of the Jacobi fields with ${\mathcal J}(0)=0$, ${\mathcal
J}'(0)=\mathrm{Id}$ (derivatives are covariant derivatives along
$\gamma_w$, which in the parallel frame become ordinary
derivatives). The Jacobi equation reads ${\mathcal
J}''=-{\mathcal R}(r){\mathcal J}$ with ${\mathcal R}(r)$ the matrix
of $R_{\partial_r}$ in the frame (perpendicular Jacobi fields stay
perpendicular by the linearity argument in (1)), and $\|{\mathcal
R}(r)\|\le C$ on compact $r$-intervals and compact sets of
directions. By the small-time asymptotics of
Lemma~\ref{lem:jacobi-small-time} applied to the operator-valued
solution ${\mathcal J}$ (with ${\mathcal J}(0)=0$,
${\mathcal J}'(0)=\mathrm{Id}$), on any interval $[0,r_1]$ on which
$\|{\mathcal R}\|\le C$ we have, with $K=\max(1,C)$ and
$M=e^{Kr_1}$,
$$\|{\mathcal J}(r)-r\,\mathrm{Id}\|\le\tfrac16\,CM\,r^3,\qquad
\|{\mathcal J}'(r)-\mathrm{Id}\|\le\tfrac12\,CM\,r^2.$$
Since $\exp_p$ is non-singular on
$C$, ${\mathcal J}(r)$ is invertible for $0<r<r_0$ (a kernel vector
of ${\mathcal J}(r)$ is a Jacobi field vanishing at $0$ and $r$,
i.e., a kernel vector of $d(\exp_p)_{rw}$ by
Lemma~\ref{lem:exponential-differential-jacobi}). Because
$A(J)=\nabla_{\partial_r}J$ for every polar Jacobi field
($\nabla_{\partial_r}\partial_{\theta^i}
=\nabla_{\partial_{\theta^i}}\partial_r=A(\partial_{\theta^i})$),
the matrix of $A$ in the frame is ${\mathcal J}'{\mathcal J}^{-1}$,
and Lemma~\ref{lem:jacobi-matrix-inverse} gives
$$A(r)={\mathcal J}'(r){\mathcal J}(r)^{-1}
=\frac{1}{r}\mathrm{Id}+O(r).$$
Similarly, since the fields $J_i$ have $\nabla J_i(0)=\partial
v/\partial\theta^i$ and $\langle \partial v/\partial\theta^i,
\partial v/\partial\theta^j\rangle=\hat g_{ij}$, the vector-valued
case of Lemma~\ref{lem:jacobi-small-time} gives
$J_i(r)=r\,\partial v/\partial\theta^i+O(r^3)$, whence
$$g_{ij}(r,\theta)=\langle J_i(r),J_j(r)\rangle
=r^2\left(\hat g_{ij}(\theta)+O(r^2)\right).$$
All estimates are locally uniform in the direction, since the
constants in Lemma~\ref{lem:jacobi-small-time} depend only on the
bound $C$ for ${\mathcal R}$, on $r_1$, and on the norms of the
initial data — and both ${\mathcal R}$ and the initial derivatives
$\partial v/\partial\theta^i$ are bounded on compact sets of
directions, the latter by smoothness of $\theta\mapsto v(\theta)$.

(4) The volume form of $\tilde g=dr^2+g_{ij}d\theta^id\theta^j$ is
$\sqrt{\det g_{ij}}\,dr\wedge d\theta^1\wedge\cdots\wedge
d\theta^{n-1}=\lambda\,dr\wedge d\mathrm{vol}_{S^{n-1}}$, since
$d\mathrm{vol}_{S^{n-1}}=\sqrt{\det\hat g}\,d\theta^1
\wedge\cdots\wedge d\theta^{n-1}$. By (2) and the standard formula
for the derivative of a determinant,
$$\partial_r\log\lambda=\tfrac12\partial_r\log\det(g_{ij})
=\tfrac12\Tr\left(g^{-1}\partial_rg\right)
=\Tr\left(g^{-1}S\right)=\Tr(A).$$
The limit $\lambda/r^{n-1}\to1$ follows from the expansion of
$g_{ij}$ in (3).
\end{proof}

\begin{lemma}[Small-time asymptotics of Jacobi-type solutions]
\label{lem:jacobi-small-time}
\uses{lem:second-order-linear-ode}
\lean{MorganTianLib.IsJacobiSolOn.max_norm_le,MorganTianLib.IsJacobiSolOn.norm_snd_le,MorganTianLib.IsJacobiSolOn.norm_fst_le,MorganTianLib.IsJacobiSolOn.norm_snd_sub_le,MorganTianLib.IsJacobiSolOn.norm_fst_sub_le}
\leanok
Let $F$ be a real Banach space and let $R\colon[0,b]\to L(F)$ be
continuous with $\|R(t)\|\le C$ for all $t\in[0,b]$; set
$K=\max(1,C)$. Let $(y,v)$ solve the second-order linear equation
$y''+R(t)\,y=0$ on $[0,b]$ in the sense of
Lemma~\ref{lem:second-order-linear-ode}, with $y(0)=0$ and
$v(0)=v_0$, and put $M=\|v_0\|\,e^{Kb}$. Then for all $t\in[0,b]$:
$$\max\bigl(\|y(t)\|,\|v(t)\|\bigr)\le\|v_0\|\,e^{Kt},\qquad
\|y(t)\|\le M\,t,$$
$$\|v(t)-v_0\|\le\tfrac12\,CM\,t^2,\qquad
\|y(t)-t\,v_0\|\le\tfrac16\,CM\,t^3.$$
More generally, without assuming $y(0)=0$, any solution on $[a,b]$
satisfies the a-priori bound
$\max(\|y(t)\|,\|v(t)\|)\le
\max(\|y(a)\|,\|v(a)\|)\,e^{K(t-a)}$.
\end{lemma}

\begin{proof}
\uses{lem:second-order-linear-ode}
\leanok
The a-priori bound is Gr\"onwall's inequality for the pair
$W=(y,v)$: with $A(t)(y,v)=(v,-R(t)\,y)$ as in the proof of
Lemma~\ref{lem:second-order-linear-ode} we have
$\|W'\|=\|A(t)\,W\|\le K\,\|W\|$ on $[a,b]$, whence
$\|W(t)\|\le\|W(a)\|\,e^{K(t-a)}$ in the sup norm
$\|(y,v)\|=\max(\|y\|,\|v\|)$. With $a=0$ and $y(0)=0$ this reads
$\max(\|y(t)\|,\|v(t)\|)\le\|v_0\|\,e^{Kt}$; in particular
$\|v\|\le M$ on $[0,b]$.

By the fundamental theorem of calculus,
$y(t)=\int_0^tv(s)\,ds$, so $\|y(t)\|\le\int_0^t\|v\|\le M\,t$.
Next, $v(t)-v_0=-\int_0^tR(s)\,y(s)\,ds$, so
$$\|v(t)-v_0\|\le\int_0^t\|R(s)\|\,\|y(s)\|\,ds
\le\int_0^tCMs\,ds=\tfrac12\,CM\,t^2.$$
Finally $y(t)-t\,v_0=\int_0^t\bigl(v(s)-v_0\bigr)\,ds$, so
$$\|y(t)-t\,v_0\|\le\int_0^t\tfrac12\,CMs^2\,ds
=\tfrac16\,CM\,t^3.$$
\end{proof}

\begin{lemma}[Inverse asymptotics of the matrix Jacobi solution]
\label{lem:jacobi-matrix-inverse}
\uses{lem:second-order-linear-ode,lem:jacobi-small-time}
\lean{MorganTianLib.IsJacobiSolOn.norm_one_sub_inv_smul_fst_le,MorganTianLib.IsJacobiSolOn.isUnit_fst,MorganTianLib.IsJacobiSolOn.norm_snd_mul_inverse_fst_sub_le}
\leanok
Let $F$ be a unital real Banach algebra and let
$R\colon[0,b]\to L(F)$ be continuous with $\|R\|\le C$ on $[0,b]$;
set $K=\max(1,C)$ and $M=e^{Kb}$. Let $(y,v)$ solve $y''+R(t)\,y=0$
on $[0,b]$ in the sense of Lemma~\ref{lem:second-order-linear-ode},
with $y(0)=0$ and $v(0)=1$ (the unit of $F$). Then for
$t\in(0,b]$:
\begin{enumerate}
\item $\|1-t^{-1}y(t)\|\le\tfrac16\,CM\,t^2$; in particular, if
$CMt^2<6$ then $y(t)$ is invertible in $F$;
\item if $CMt^2\le3$, then
$$\bigl\|v(t)\,y(t)^{-1}-t^{-1}\cdot1\bigr\|\le 2\,CM\,t.$$
\end{enumerate}
Applied to the matrix Jacobi field ${\mathcal J}$ of
Lemma~\ref{lem:geodesic-polar-form}(3) (with $F$ the algebra of
operators on the orthogonal complement of the geodesic's direction),
this is the assertion that ${\mathcal J}(r)$ is invertible for small
$r>0$ with $A(r)={\mathcal J}'(r){\mathcal J}(r)^{-1}
=\tfrac1r\,\mathrm{Id}+O(r)$.
\end{lemma}

\begin{proof}
\uses{lem:jacobi-small-time}
\leanok
(1) By Lemma~\ref{lem:jacobi-small-time},
$\|y(t)-t\cdot1\|\le\tfrac16\,CM\,t^3$; dividing by $t>0$ gives the
displayed bound for $z:=t^{-1}y(t)$. If $CMt^2<6$ then
$\|1-z\|<1$, so the geometric series
$\sum_{n\ge0}(1-z)^n$ converges and inverts
$z=1-(1-z)$; hence $z$, and with it $y(t)=t\,z$, is invertible.

(2) Write $w=1-z$, so $\|w\|\le\tfrac16CMt^2\le\tfrac12$ by (1) and
the hypothesis. The geometric series gives
$\|z^{-1}\|\le(1-\|w\|)^{-1}\le2$. By
Lemma~\ref{lem:jacobi-small-time} again,
$\|v(t)-1\|\le\tfrac12CMt^2$, so
$$\|v(t)-z\|\le\|v(t)-1\|+\|1-z\|
\le\tfrac12CMt^2+\tfrac16CMt^2=\tfrac23\,CMt^2.$$
Since $y(t)^{-1}=t^{-1}z^{-1}$ and $1=z\,z^{-1}$,
$$v(t)\,y(t)^{-1}-t^{-1}\cdot1
=t^{-1}\bigl(v(t)-z\bigr)z^{-1},$$
whence
$\|v(t)y(t)^{-1}-t^{-1}\cdot1\|
\le t^{-1}\cdot\tfrac23CMt^2\cdot2=\tfrac43\,CMt\le2\,CMt$.
\end{proof}

\begin{lemma}[The radial shape operator and its Riccati equation]
\label{lem:radial-shape-riccati}
\uses{lem:second-order-linear-ode,lem:jacobi-matrix-inverse,lem:jacobi-small-time}
\lean{MorganTianLib.IsRadialJacobi,MorganTianLib.jacobiOpCoeff,MorganTianLib.shapeOp,MorganTianLib.wronskian_eq,MorganTianLib.shapeOp_symm,MorganTianLib.hasDerivAt_shapeOp,MorganTianLib.tendsto_shapeOp_sub_inv_smul_id,MorganTianLib.IsRadialJacobi.isUnit_fst}
\leanok
Let $E$ be a real Hilbert space, let $b>0$, and let
$\mathcal R\colon[0,b]\to L(E)$ be continuous with
$\|\mathcal R\|\le C$ and each $\mathcal R(t)$ symmetric. Let
$({\mathcal J},{\mathcal J}')$ solve the matrix Jacobi equation
${\mathcal J}''+{\mathcal R}(t){\mathcal J}=0$ on $[0,b]$ in the sense
of Lemma~\ref{lem:second-order-linear-ode}, with ${\mathcal J}(0)=0$
and ${\mathcal J}'(0)=\mathrm{Id}$. Let $0<r_0\le b$ and suppose
${\mathcal J}(r)$ is invertible for every $r\in(0,r_0)$. Define the
\emph{radial shape operator}
$$A(r)={\mathcal J}'(r)\,{\mathcal J}(r)^{-1},\qquad r\in(0,r_0).$$
Then:
\begin{enumerate}
\item (Wronskian identity) $\langle{\mathcal J}'(t)a,{\mathcal
J}(t)c\rangle=\langle{\mathcal J}(t)a,{\mathcal J}'(t)c\rangle$ for all
$t\in[0,b]$ and $a,c\in E$; consequently $A(r)$ is symmetric for every
$r\in(0,r_0)$.
\item (Riccati equation) $A$ is differentiable on $(0,r_0)$ and
$$A'+A^2+{\mathcal R}=0 .$$
\item (Asymptotics) $A(r)-\tfrac1r\,\mathrm{Id}\to0$ in operator norm
as $r\to0^+$.
\end{enumerate}
\end{lemma}

\begin{proof}
\uses{lem:jacobi-matrix-inverse,lem:jacobi-small-time}
\leanok
(1) Fix $a,c\in E$ and set
$w(t)=\langle{\mathcal J}'(t)a,{\mathcal J}(t)c\rangle
-\langle{\mathcal J}(t)a,{\mathcal J}'(t)c\rangle$. Differentiating and
using ${\mathcal J}''=-{\mathcal R}{\mathcal J}$,
$$w'=\langle-{\mathcal R}({\mathcal J}a),{\mathcal J}c\rangle
+\langle{\mathcal J}'a,{\mathcal J}'c\rangle
-\langle{\mathcal J}'a,{\mathcal J}'c\rangle
-\langle{\mathcal J}a,-{\mathcal R}({\mathcal J}c)\rangle
=-\langle{\mathcal R}({\mathcal J}a),{\mathcal J}c\rangle
+\langle{\mathcal J}a,{\mathcal R}({\mathcal J}c)\rangle=0,$$
since ${\mathcal R}(t)$ is symmetric. As ${\mathcal J}(0)=0$ we have
$w(0)=0$, so $w\equiv0$ on $[0,b]$. For the symmetry of $A(r)$, let
$X,Y\in E$ and put $a={\mathcal J}(r)^{-1}X$, $c={\mathcal
J}(r)^{-1}Y$; then
$$\langle A(r)X,Y\rangle=\langle{\mathcal J}'(r)a,{\mathcal J}(r)c\rangle
=\langle{\mathcal J}(r)a,{\mathcal J}'(r)c\rangle=\langle X,A(r)Y\rangle .$$

(2) Inversion is differentiable on the units of the Banach algebra
$L(E)$, with $({\mathcal J}^{-1})'=-{\mathcal J}^{-1}{\mathcal
J}'{\mathcal J}^{-1}$. Hence, by the product rule and the Jacobi
equation,
$$A'={\mathcal J}''{\mathcal J}^{-1}
+{\mathcal J}'\bigl(-{\mathcal J}^{-1}{\mathcal J}'{\mathcal
J}^{-1}\bigr)
=-{\mathcal R}\bigl({\mathcal J}{\mathcal J}^{-1}\bigr)
-\bigl({\mathcal J}'{\mathcal J}^{-1}\bigr)^2
=-{\mathcal R}-A^2 .$$

(3) This is Lemma~\ref{lem:jacobi-matrix-inverse}(2), which gives
$\|A(r)-\tfrac1r\,\mathrm{Id}\|\le2\,CM\,r$ for small $r>0$.
\end{proof}

\begin{lemma}[Radial comparison in a parallel frame]
\label{lem:radial-comparison}
\uses{lem:radial-shape-riccati,lem:operator-riccati-upper,lem:trace-riccati-comparison,lem:ratio-monotone,lem:jacobi-small-time,def:sinh-comparison-function}
\lean{MorganTianLib.shapeOp_inner_le,MorganTianLib.norm_jacobi_sq_le,MorganTianLib.trace_shapeOp_le,MorganTianLib.tendsto_snK_div_self,MorganTianLib.tendsto_norm_jacobi_div_self}
\leanok
In the setting of Lemma~\ref{lem:radial-shape-riccati}, fix $k\ge0$.
\begin{enumerate}
\item If $\langle{\mathcal R}(r)X,X\rangle\ge-k\,|X|^2$ for all
$r\in(0,r_0)$ and $X\in E$, then
$$\langle A(r)X,X\rangle\ \le\ \frac{\sn_k'(r)}{\sn_k(r)}\,|X|^2 .$$
\item Under the same hypothesis, for every $w\in E$,
$$|{\mathcal J}(r)w|^2\ \le\ \sn_k^2(r)\,|w|^2\qquad(0<r<r_0).$$
\item If instead $\Tr{\mathcal R}(r)\ge-m\,k$ for all $r\in(0,r_0)$,
where $m=\dim E<\infty$, then
$$\Tr A(r)\ \le\ m\,\frac{\sn_k'(r)}{\sn_k(r)}\qquad(0<r<r_0).$$
\end{enumerate}
\end{lemma}

\begin{proof}
\uses{lem:radial-shape-riccati,lem:operator-riccati-upper,lem:trace-riccati-comparison,lem:ratio-monotone,lem:jacobi-small-time}
\leanok
(1) By the Riccati equation of
Lemma~\ref{lem:radial-shape-riccati}(2),
$$\langle A'(r)X,X\rangle=-\langle{\mathcal R}(r)X,X\rangle
-\langle A(r)^2X,X\rangle\ \le\ k\,|X|^2-\langle A(r)^2X,X\rangle,$$
which is exactly the hypothesis of
Lemma~\ref{lem:operator-riccati-upper}; its remaining hypotheses —
symmetry of $A(r)$ and $A(r)-\tfrac1r\mathrm{Id}\to0$ — are
Lemma~\ref{lem:radial-shape-riccati}(1) and (3). The conclusion is (1).

(2) Fix $w$ and set $h(r)=|{\mathcal J}(r)w|^2$. Since
${\mathcal J}'(r)w=A(r)\bigl({\mathcal J}(r)w\bigr)$, writing
$Y={\mathcal J}(r)w$ we get
$$h'(r)=2\langle{\mathcal J}'(r)w,{\mathcal J}(r)w\rangle
=2\langle A(r)Y,Y\rangle\ \le\ 2\,\frac{\sn_k'(r)}{\sn_k(r)}\,|Y|^2
=2\,\frac{\sn_k'(r)}{\sn_k(r)}\,h(r)$$
by part (1). Hence $h/\sn_k^2$ is non-increasing on $(0,r_0)$ by
Lemma~\ref{lem:ratio-monotone}. By Lemma~\ref{lem:jacobi-small-time},
${\mathcal J}(r)=r\,\mathrm{Id}+O(r^3)$, so $|{\mathcal J}(r)w|/r\to|w|$
as $r\to0^+$; and $\sn_k(r)/r\to1$, since $\sn_k(0)=0$ and
$\sn_k'(0)=1$. Therefore $h(r)/\sn_k^2(r)\to|w|^2$ as $r\to0^+$, and
the monotonicity gives $h(r)/\sn_k^2(r)\le|w|^2$ for $0<r<r_0$, which
is (2).

(3) Taking traces in the Riccati equation and using linearity of the
trace,
$$\Tr A'(r)+\Tr\bigl(A(r)^2\bigr)=-\Tr{\mathcal R}(r)\ \le\ m\,k,$$
which is the hypothesis of
Lemma~\ref{lem:trace-riccati-comparison}. Moreover the trace is a
continuous linear functional on $L(E)$ (as $\dim E<\infty$) and
$\Tr(\mathrm{Id})=m$, so Lemma~\ref{lem:radial-shape-riccati}(3) gives
$\Tr A(r)-m/r\to0$ as $r\to0^+$. Together with the symmetry of $A(r)$,
Lemma~\ref{lem:trace-riccati-comparison} applies and yields (3).
\end{proof}

\begin{lemma}[Derivative of the radial volume element]
\label{lem:radial-volume-element}
\uses{lem:radial-shape-riccati}
\lean{MorganTianLib.volumeElement,MorganTianLib.hasDerivAt_volumeElement}
\leanok
In the setting of Lemma~\ref{lem:radial-shape-riccati}, suppose in addition
that $\dim E=m<\infty$, and set
$$\lambda(r)=\det{\mathcal J}(r).$$
Then $\lambda$ is differentiable on $(0,r_0)$ and
$$\lambda'(r)=\lambda(r)\,\Tr A(r),$$
equivalently $\partial_r\log\lambda=\Tr A$ wherever $\lambda\ne0$.
\end{lemma}

\begin{proof}
\uses{lem:radial-shape-riccati}
\leanok
Fix a basis of $E$ and identify endomorphisms with their matrices; the
determinant and the trace of an endomorphism are those of its matrix, and
composition becomes matrix multiplication. The determinant is a polynomial
in the matrix entries, hence differentiable, and its Fréchet derivative is
Jacobi's formula
$$d(\det)_M(B)=\det M\cdot\Tr\bigl(M^{-1}B\bigr)\qquad(M\ \text{invertible}).$$
Since ${\mathcal J}(r)$ is invertible for $r\in(0,r_0)$, the chain rule
applied to $r\mapsto\det{\mathcal J}(r)$ gives
$$\lambda'(r)=\det{\mathcal J}(r)\cdot
\Tr\bigl({\mathcal J}(r)^{-1}{\mathcal J}'(r)\bigr)
=\lambda(r)\,\Tr\bigl({\mathcal J}'(r){\mathcal J}(r)^{-1}\bigr)
=\lambda(r)\,\Tr A(r),$$
where the middle equality is the cyclicity of the trace,
$\Tr(XY)=\Tr(YX)$.
\end{proof}

\begin{lemma}[Scalar Riccati comparison]
\label{lem:scalar-riccati-comparison}
\uses{def:sinh-comparison-function}
\lean{MorganTianLib.scalar_riccati_comparison,MorganTianLib.scalar_riccati_comparison_Ioi}
\leanok
Fix $k\ge 0$ and $0<r_0\le\infty$. Suppose $\phi\colon
(0,r_0)\to\mathbb{R}$ is differentiable with
$$\phi'+\phi^2\le k\quad\text{on }(0,r_0),\qquad\text{and}\qquad
\phi(r)-\frac1r\to0\ \text{ as } r\to0^+.$$
Then
$$\phi(r)\le\frac{\sn_k'(r)}{\sn_k(r)}\qquad\text{for all }r\in
(0,r_0).$$
\end{lemma}

\begin{proof}
\uses{def:sinh-comparison-function}
\leanok
Let $a(r)=\sn_k'(r)/\sn_k(r)$. From $\sn_k''=k\,\sn_k$ we get
$a'+a^2=k$, and from the Taylor expansion of $\sn_k$ at $0$,
$a(r)-\frac1r\to0$ as $r\to0^+$. Set
$\psi(r)=\sn_k^2(r)\left(\phi(r)-a(r)\right)$. Then
$\psi(r)\to0$ as $r\to0^+$ (as $\sn_k^2(r)\to 0$ and $\phi-a\to0$),
and
$$\psi'=2\,\sn_k\sn_k'\,(\phi-a)+\sn_k^2(\phi'-a')
\le 2\,\sn_k^2\,a\,(\phi-a)-\sn_k^2(\phi-a)(\phi+a)
=-\sn_k^2\,(\phi-a)^2\le0,$$
using $\phi'-a'\le(k-\phi^2)-(k-a^2)=-(\phi-a)(\phi+a)$. Hence
$\psi\le0$ on $(0,r_0)$, and since $\sn_k^2>0$ there,
$\phi\le a$.
\end{proof}

We also record an elementary companion, used to convert the shape
operator bounds of the comparison theorems into bounds on the metric
coefficients.

\begin{lemma}[Ratio monotonicity under a logarithmic-derivative bound]
\label{lem:ratio-monotone}
\lean{MorganTianLib.antitoneOn_div_sq_of_deriv_le,MorganTianLib.monotoneOn_div_sq_of_le_deriv}
\leanok
Fix $r_0>0$ and let $s,h\colon(0,r_0)\to\mathbb R$ be differentiable
with $s>0$. If $h'\le 2\frac{s'}{s}\,h$ on $(0,r_0)$, then $h/s^2$
is non-increasing there; if instead $h'\ge2\frac{s'}{s}\,h$, then
$h/s^2$ is non-decreasing.
\end{lemma}

\begin{proof}
\leanok
By the quotient rule, on $(0,r_0)$,
$$\left(\frac{h}{s^2}\right)'
=\frac{h's^2-2ss'\,h}{s^4}
=\frac{h'-2\frac{s'}{s}h}{s^2},$$
which is $\le0$ (resp.\ $\ge0$) under the first (resp.\ second)
hypothesis; a differentiable function with nonpositive (resp.\
nonnegative) derivative on an interval is non-increasing (resp.\
non-decreasing).
\end{proof}

\begin{lemma}[Scalar Sturm comparison]
\label{lem:scalar-sturm-comparison}
\lean{MorganTianLib.scalar_sturm_comparison,MorganTianLib.scalar_sturm_pos}
\leanok
Fix $K\ge0$ and $t_1>0$ with $\sqrt K\,t_1<\pi$ (no condition when
$K=0$), and write $s_K$ for the solution of $s_K''+K\,s_K=0$,
$s_K(0)=0$, $s_K'(0)=1$ (the function $s_\lambda$ with $\lambda=K$
from the subsection on local isometries and constant curvature).
Let $f\colon[0,t_1]\to\mathbb{R}$ be continuous and twice
differentiable on $(0,t_1)$, and suppose that: $f''\ge -K f$ on
$(0,t_1)$; $f(0)=0$ and $f>0$ on $(0,t_1)$; $f'$ is bounded near
$0^+$; and $f(t)/t\to c$ as $t\to0^+$. Then
$$f(t)\ \ge\ c\,s_K(t)\qquad\text{for all }t\in(0,t_1].$$
In particular, if $c>0$ then $f(t_1)>0$.
\end{lemma}

\begin{proof}
\leanok
On $[0,t_1]$ we have $\sqrt K\,t\le\sqrt K\,t_1<\pi$, so $s_K\ge0$
on $[0,t_1]$ and $s_K>0$ on $(0,t_1]$. Consider the Wronskian
$W=f'\,s_K-f\,s_K'$ on $(0,t_1)$. Since $s_K''=-K\,s_K$,
$$W'=f''\,s_K-f\,s_K''=(f''+K f)\,s_K\ \ge\ 0\qquad
\text{on }(0,t_1),$$
so $W$ is non-decreasing. As $t\to0^+$: $f(t)=(f(t)/t)\cdot t\to
c\cdot0=0$; $s_K(t)\to0$ while $f'$ is bounded near $0^+$, so
$f'\,s_K\to0$; and $|s_K'|\le1$, so $f\,s_K'\to0$. Hence
$W(t)\to0$ as $t\to0^+$, and monotonicity gives $W\ge0$ on
$(0,t_1)$. Therefore
$$\left(\frac f{s_K}\right)'=\frac W{s_K^2}\ \ge\ 0
\qquad\text{on }(0,t_1),$$
so $f/s_K$ is non-decreasing on $(0,t_1)$. Since
$s_K(t)/t\to s_K'(0)=1$ as $t\to0^+$,
$$\frac{f(t)}{s_K(t)}=\frac{f(t)}{t}\cdot\frac{t}{s_K(t)}
\longrightarrow c\qquad(t\to0^+),$$
and monotonicity gives $f/s_K\ge c$ on $(0,t_1)$, i.e.
$f\ge c\,s_K$ there; the case $t=t_1$ follows by continuity of $f$
and $s_K$. Finally, if $c>0$ then
$f(t_1)\ge c\,s_K(t_1)>0$, since $0<t_1$ and $\sqrt K\,t_1<\pi$
give $s_K(t_1)>0$.
\end{proof}

The conjugate-point argument below applies the scalar comparison to
the norm $|Y|$ of a Jacobi field. We isolate that step in a
frame-independent, vector-valued form: expressed in a parallel
orthonormal frame along the geodesic, a Jacobi field becomes a curve
$V$ in a fixed inner product space, and the Jacobi equation together
with the curvature bound gives exactly the hypothesis
$\langle V'',V\rangle\ge-K\,|V|^2$ below.

\begin{lemma}[Vector-valued Sturm comparison]
\label{lem:vector-sturm-comparison}
\uses{lem:scalar-sturm-comparison}
\lean{MorganTianLib.vector_sturm_comparison,MorganTianLib.vector_sturm_ne_zero,MorganTianLib.tendsto_norm_div_self_of_hasDerivWithinAt}
\leanok
Fix $K\ge0$ and $T$ with $\sqrt K\,T<\pi$, and let $E$ be a real
inner product space. Let $V\colon[0,T]\to E$ be continuous on
$[0,T]$ and twice differentiable on $(0,T)$, and suppose that:
$\langle V''(t),V(t)\rangle\ge-K\,|V(t)|^2$ on $(0,T)$; $|V'|$ is
bounded near $0^+$; and $|V(t)|/t\to c$ as $t\to0^+$. Then
$$|V(t)|\ \ge\ c\,s_K(t)\qquad\text{for all }t\in(0,T].$$
In particular, if $c>0$ then $V$ has no zero on $(0,T]$; and if
$V(0)=0$ and $V$ is differentiable at $0$ with derivative $V'(0)$,
then the slope hypothesis holds automatically with $c=|V'(0)|$.
\end{lemma}

\begin{proof}
\uses{lem:scalar-sturm-comparison}
\leanok
Since $s_K\ge0$ on $[0,T]$ (as $\sqrt K\,T<\pi$), the bound is
trivial when $c\le0$; assume $c>0$. First suppose $t_1\in(0,T]$ is
such that $V\not=0$ on $(0,t_1)$. There the function
$f=|V|=\sqrt{\langle V,V\rangle}$ is positive and twice
differentiable, with
$$f'=\frac{\langle V',V\rangle}{|V|},\qquad
f''=\frac{\langle V'',V\rangle+|V'|^2}{|V|}
-\frac{\langle V',V\rangle^2}{|V|^3},$$
by the chain and quotient rules. The Cauchy--Schwarz inequality
$\langle V',V\rangle^2\le|V'|^2\,|V|^2$ gives
$$f''\ \ge\ \frac{\langle V'',V\rangle}{|V|}\ \ge\
\frac{-K|V|^2}{|V|}=-K f\qquad\text{on }(0,t_1),$$
using the hypothesis on $\langle V'',V\rangle$. Moreover
$|f'|\le|V'|$ is bounded near $0^+$ (Cauchy--Schwarz again), $f$ is
continuous on $[0,t_1]$, and $f(t)/t\to c$. By the scalar Sturm
comparison (Lemma~\ref{lem:scalar-sturm-comparison}, with
$\sqrt K\,t_1\le\sqrt K\,T<\pi$),
$$f(t)\ \ge\ c\,s_K(t)\qquad\text{on }(0,t_1].$$
Now suppose, for contradiction, that $V$ vanishes somewhere on
$(0,T]$. Since $|V(t)|/t\to c>0$, there is $\varepsilon>0$ with
$V\not=0$ on $(0,\varepsilon)$, so every zero lies in
$[\varepsilon,T]$; the set of zeros in $[\varepsilon,T]$ is closed
and non-empty, so it has a least element $t^\star$, and
$V\not=0$ on $(0,t^\star)$. The previous paragraph applied at
$t_1=t^\star$ gives $0=|V(t^\star)|\ge c\,s_K(t^\star)>0$ (as
$0<t^\star$ and $\sqrt K\,t^\star<\pi$), a contradiction. Hence $V$
has no zero on $(0,T]$, and the displayed bound holds on all of
$(0,T]$ by the first paragraph with $t_1=T$. Finally, if $V(0)=0$
and $V$ is differentiable at $0$ with derivative $V'(0)$, then
$|V(t)|/t=|V(t)/t|\to|V'(0)|$ as $t\to0^+$, which is the slope
hypothesis with $c=|V'(0)|$.
\end{proof}

The passage from a Jacobi field along a curve, measured in the
(time-dependent) inner products of the metric, to a Euclidean curve
eligible for the vector-valued Sturm comparison is a frame computation
that we record once and for all. It packages the parallel orthonormal
frame of Lemma~\ref{lem:parallel-frame} together with the Parseval
expansion and the endpoint slope analysis.

\begin{lemma}[Jacobi fields in a parallel orthonormal frame]
\label{lem:jacobi-frame-reduction}
\uses{lem:parallel-frame,lem:jacobi-field-coordinates,lem:vector-sturm-comparison,def:riemannian-metric}
\lean{MorganTianLib.jacobi_frame_sturm_comparison,MorganTianLib.jacobi_ne_zero_of_curvature_le,MorganTianLib.chartMetricInner_parseval,MorganTianLib.exists_chartMetricInner_orthonormal,MorganTianLib.chartMetricInner_pos}
\leanok
Fix $K\ge0$ and $T>0$ with $\sqrt K\,T<\pi$. Let $u\colon[0,T]\to U$ be a
$C^1$ curve in a coordinate chart $U$ (the Gram matrix of the metric being
differentiable along $u$), and let $(J,\nabla J)$ be a Jacobi field along
$u$ (Lemma~\ref{lem:jacobi-field-coordinates}) with $J(0)=0$. Suppose that
along $u$ the curvature term of the Jacobi equation satisfies
$$\langle{\mathcal R}(J,\dot u)\dot u,\,J\rangle\ \le\ K\,\langle J,J\rangle
\qquad\text{on }(0,T),$$
where $\langle\cdot,\cdot\rangle=\langle\cdot,\cdot\rangle_{u(t)}$ denotes
the inner product defined by the metric at $u(t)$. Then
$$\sqrt{\langle\nabla J(0),\nabla J(0)\rangle}\;s_K(t)\ \le\
\sqrt{\langle J(t),J(t)\rangle}\qquad\text{for all }t\in(0,T].$$
In particular, if $\nabla J(0)\not=0$ then $J$ has no zero on $(0,T]$.
\end{lemma}

\begin{proof}
\uses{lem:parallel-frame,lem:jacobi-field-coordinates,lem:vector-sturm-comparison}
\leanok
Write $n=\dim M$ and $B_t(\cdot,\cdot)=\langle\cdot,\cdot\rangle_{u(t)}$
for the Gram bilinear form of the metric at $u(t)$; each $B_t$ is
symmetric and positive definite.

{\em An orthonormal frame.} Diagonalizing the symmetric bilinear form
$B_0$ (over a field of characteristic $\not=2$ every symmetric bilinear
form admits an orthogonal basis) and rescaling each basis vector by the
square root of its positive $B_0$-square norm produces a basis
$e_1^0,\ldots,e_n^0$ with $B_0(e_i^0,e_j^0)=\delta_{ij}$. By
Lemma~\ref{lem:parallel-frame} there are fields $e_1,\ldots,e_n$ parallel
along $u$ on $[0,T]$ with $e_i(0)=e_i^0$, and their Gram matrix is
preserved: $B_t(e_i(t),e_j(t))=\delta_{ij}$ for all $t\in[0,T]$. In
particular the $e_i(t)$ are linearly independent (a vanishing linear
combination pairs to zero against each $e_j(t)$), hence a basis;
consequently, expanding both arguments in this basis, one obtains for any
vectors $x,z$ the Parseval identity
$$\sum_{i=1}^n B_t(x,e_i(t))\,B_t(z,e_i(t))=B_t(x,z).$$

{\em Frame coordinates.} Let $V(t)\in\mathbb{R}^n$ (with its standard
Euclidean inner product) have components $V_i(t)=B_t(J(t),e_i(t))$, and
set likewise $V'_i(t)=B_t(\nabla J(t),e_i(t))$ and
$V''_i(t)=-B_t({\mathcal R}(J,\dot u)\dot u(t),e_i(t))$. By the product
rule for the metric along the curve (metric compatibility, as in the
proof of Lemma~\ref{lem:parallel-frame}), at interior times
$$\frac{d}{dt}\,B_t(X(t),e_i(t))
=B_t(\nabla X(t),e_i(t))+B_t(X(t),\nabla e_i(t))
=B_t(\nabla X(t),e_i(t)),$$
the frame being parallel. Applied to $X=J$ and $X=\nabla J$ — whose
covariant derivatives along $u$ are $\nabla J$ and
$-{\mathcal R}(J,\dot u)\dot u$ by the Jacobi equation
(Lemma~\ref{lem:jacobi-field-coordinates}) — this shows that $V$ is twice
differentiable on $(0,T)$ with derivatives $V'$ and $V''$ there; all
three are continuous on $[0,T]$, being built from continuous functions.
By the Parseval identity,
$$|V(t)|^2=B_t(J(t),J(t)),\qquad
\langle V''(t),V(t)\rangle=-B_t({\mathcal R}(J,\dot u)\dot u,\,J)
\ \ge\ -K\,B_t(J,J)=-K\,|V(t)|^2$$
using the curvature hypothesis, and likewise
$|V'(0)|^2=B_0(\nabla J(0),\nabla J(0))$.

{\em Slope at $0$.} $V(0)=0$ since $J(0)=0$. We claim
$|V(t)|/t\to|V'(0)|$ as $t\to0^+$. Given $\varepsilon>0$, by continuity
of $V'$ at $0$ there is $\eta>0$ with $|V'(x)-V'(0)|\le\varepsilon$ for
$0\le x\le\eta$, $x\le T$. For $0<\delta\le t<\min(\eta,T)$ the mean
value inequality applied to $W(s)=V(s)-s\,V'(0)$ on $[\delta,t]$ (where
$W$ is differentiable with $|W'|=|V'-V'(0)|\le\varepsilon$) gives
$|W(t)-W(\delta)|\le\varepsilon(t-\delta)$; letting $\delta\to0^+$ and
using $W(\delta)\to W(0)=0$ by continuity yields
$|V(t)-t\,V'(0)|\le\varepsilon t$, whence
$\bigl||V(t)|/t-|V'(0)|\bigr|\le\varepsilon$. In particular $|V'|$ is
bounded near $0^+$, again by continuity on $[0,T]$.

{\em Conclusion.} The curve $V$ satisfies all the hypotheses of the
vector-valued Sturm comparison (Lemma~\ref{lem:vector-sturm-comparison})
on $[0,T]$ with $c=|V'(0)|$, so on $(0,T]$
$$\sqrt{\langle J(t),J(t)\rangle}=|V(t)|\ \ge\ c\,s_K(t)
=\sqrt{\langle\nabla J(0),\nabla J(0)\rangle}\,s_K(t).$$
If $\nabla J(0)\not=0$ then $c>0$ by positive definiteness of $B_0$, and
since $s_K>0$ on $(0,T]$ (as $\sqrt K\,T<\pi$), the field $J$ has no zero
on $(0,T]$.
\end{proof}

\begin{lemma}[No conjugate points under an upper curvature bound]
\label{lem:conjugate-sturm}
\uses{def:conjugate-point,def:sectional-curvature,def:geodesic}
\lean{MorganTianLib.not_isConjugatePointAt_of_sectionalCurvatureAt_le,MorganTianLib.IsJacobiFieldAlongOn.sqrt_metricInner_comparison,MorganTianLib.IsJacobiFieldAlongOn.ne_zero_of_sectionalCurvatureAt_le,MorganTianLib.IsJacobiFieldAlongOn.eqOn_zero,MorganTianLib.scalar_sturm_comparison_extend,MorganTianLib.scalar_sturm_comparison_of_start}
\leanok
Fix $K\ge 0$, with the convention $\pi/\sqrt K=+\infty$ for $K=0$.
Let $\gamma\colon[0,T]\to M$ be a unit-speed geodesic such that
$K(P)\le K$ for every $2$-plane $P$ tangent to $M$ at a point of
$\gamma$. Then there is no conjugate point along $\gamma$ at any
parameter $t_1$ with $0<t_1<\min(T,\pi/\sqrt K)$.
\end{lemma}

\begin{proof}
\leanok
\uses{def:conjugate-point,lem:jacobi-field-coordinates,lem:chart-curvature-coordinates,lem:jacobi-frame-reduction,lem:vector-sturm-comparison}
Suppose $Y$ is a non-zero Jacobi field along $\gamma$ with
$Y(0)=0=Y(t_1)$, $0<t_1<\min(T,\pi/\sqrt K)$. The geodesic may cross
several coordinate charts; the argument passes between charts using
only the chart-independent scalar pairings — no transformation law for
the Christoffel symbols is required.

{\em Step 1: $\nabla Y(0)\not=0$.} Otherwise $Y$ would vanish
identically: let $c\le t_1$ be the supremum of the times $s$ such that
$(Y,\nabla Y)$ vanishes on $[0,s]$. The pair vanishes on $[0,c)$, hence
at $c$ by continuity of the chart readings in a chart around
$\gamma(c)$; and if $c<t_1$, the forward Gr\"onwall uniqueness of the
chart Jacobi system (Lemma~\ref{lem:jacobi-field-coordinates}(2),
applied in that chart from the time $c$) extends the vanishing strictly
past $c$, a contradiction. So $c=t_1$, contradicting $Y\not=0$.

{\em Step 2: the chart-independent scalars.} Set
$$F=\langle Y,Y\rangle,\qquad G=\langle\nabla Y,Y\rangle,\qquad
H=\langle\nabla Y,\nabla Y\rangle,$$
the intrinsic metric pairings along $\gamma$. In any chart containing
the local piece of $\gamma$, these equal the chart Gram pairings of the
chart readings of $(Y,\nabla Y)$, so they are globally defined and
continuous on $[0,t_1]$, with $F(0)=0$, $F\ge0$, and $H$ bounded. At
every interior time, differentiating in a local chart with the metric
compatibility of the Levi-Civita connection and the Jacobi system gives
$$F'=2G,\qquad
G'=-\langle{\mathcal R}(Y,\dot\gamma)\dot\gamma,\,Y\rangle+H
\ \ge\ -K\,F+H,$$
the inequality by the sectional bound in chart coordinates
(Lemma~\ref{lem:chart-curvature-coordinates}(2)) together with
$\langle\dot\gamma,\dot\gamma\rangle=1$; and $G^2\le F\,H$ by
Cauchy--Schwarz for the metric pairing.

{\em Step 3: the first chart.} Choose $b_0\in(0,t_1]$ with
$\gamma([0,b_0])$ contained in a single chart around $\gamma(0)$. On
$(0,b_0]$ the frame reduction (Lemma~\ref{lem:jacobi-frame-reduction},
whose curvature hypothesis
$\langle{\mathcal R}(Y,\dot\gamma)\dot\gamma,Y\rangle\le K\langle
Y,Y\rangle$ is the bound of Step 2) gives
$$\sqrt{F(t)}\ \ge\ c\,s_K(t),\qquad
c=|\nabla Y(0)|>0 .$$

{\em Step 4: scalar Sturm continuation.} Suppose $F$ had a zero in
$(0,t_1]$, and let $t_0$ be the least one; by Step 3, $t_0>b_0$ and
$F>0$ on $(0,t_0)$. On $(0,t_0)$ set $v=\sqrt F$; then $v'=G/\sqrt F$
and, differentiating once more and using Step 2,
$$v''=\frac{G'}{\sqrt F}-\frac{G^2}{F^{3/2}}
\ \ge\ \frac{-KF+H}{\sqrt F}-\frac{H}{\sqrt F}=-K\,v .$$
The Wronskian $W=v's_K-vs_K'$ therefore satisfies
$W'=s_K(v''+Kv)\ge0$, so $W$ is nondecreasing on $(0,t_0)$; moreover
$W\to0$ as $t\to0^+$, since $|v'|=|G|/\sqrt F\le\sqrt H$ is bounded
(Cauchy--Schwarz again) while $s_K\to0$, and $v\to0$ while $s_K'$ is
bounded. Hence $W\ge0$, so $v/s_K$ is nondecreasing on $(0,t_0)$; by
Step 3 it is $\ge c$ at $b_0$, hence on all of $[b_0,t_0)$, and on
$(0,b_0]$ directly by Step 3. Letting $t\to t_0$,
$$0=\sqrt{F(t_0)}\ \ge\ c\,s_K(t_0)>0$$
(the last inequality since $0<t_0\le t_1<\pi/\sqrt K$), a
contradiction. So $F>0$ on $(0,t_1]$; in particular
$F(t_1)=|Y(t_1)|^2>0$, contradicting $Y(t_1)=0$. Hence no such $Y$
exists.

(The multi-chart bookkeeping differs from the classical argument via a
parallel orthonormal frame along all of $\gamma$ — which would need
parallel transport glued across charts: the present route carries only
the scalars $F,G,H$ across charts, and recovers the full vector-valued
Sturm comparison of Lemma~\ref{lem:vector-sturm-comparison} in the
first chart through Lemma~\ref{lem:jacobi-frame-reduction}. When the
whole piece of $\gamma$ lies in a single chart the two arguments
coincide.)
\end{proof}

The comparison theorems below bound the radial geometry from above
under a lower sectional curvature bound. The proof of the
volume--injectivity radius estimate at the end of the chapter also
requires the mirror-image bounds from below, under an upper
curvature bound; we record them here. Recall the function
$s_\lambda$ (with $s_\lambda''+\lambda s_\lambda=0$,
$s_\lambda(0)=0$, $s_\lambda'(0)=1$) and $D_\lambda$ from the
subsection on local isometries and constant curvature.

The analytic core of the mirror bound is the following
manifold-free comparison for operator-valued solutions of a Riccati
differential inequality. Unlike the upper bound of
Lemma~\ref{lem:scalar-riccati-comparison}, it cannot be reduced to
a scalar inequality along a fixed direction (the Cauchy--Schwarz
inequality points the wrong way); instead one works with the
smallest eigenvalue of the operator, which is merely locally
Lipschitz, through one-sided upper support functions and a barrier
argument for its left difference quotients.

\begin{lemma}[Operator Riccati comparison, lower form]
\label{lem:operator-riccati-lower}
\lean{MorganTianLib.operator_riccati_nonneg,MorganTianLib.minRayleigh,MorganTianLib.le_of_left_slope_eventually_lt,MorganTianLib.apply_eq_minRayleigh_smul}
\leanok
Let $E$ be a nontrivial finite-dimensional real inner product space
and fix $r_0>0$. Let $s\colon(0,r_0)\to\mathbb R$ be differentiable
and positive, bounded near $0^+$, and write $a=s'/s$. Let $U(r)$,
$r\in(0,r_0)$, be a differentiable family of symmetric
endomorphisms of $E$ with $U(r)\to0$ as $r\to0^+$, satisfying the
Riccati differential inequality
$$\langle U'(r)X,X\rangle\ \ge\
-\langle U(r)^2X,X\rangle-2a(r)\,\langle U(r)X,X\rangle
\qquad\text{for all }X\in E,\ r\in(0,r_0).$$
Then $U(r)\ge0$ for all $r\in(0,r_0)$.
\end{lemma}

\begin{proof}
\leanok
Let $m(r)$ be the smallest eigenvalue of $U(r)$, i.e.\ the minimum
of the quadratic form $X\mapsto\langle U(r)X,X\rangle$ over the
unit sphere of $E$, attained by compactness. For unit $X$ the
values $\langle U(r)X,X\rangle$ and $\langle U(r')X,X\rangle$
differ by at most $\|U(r)-U(r')\|$, so comparing minima through
their minimizers gives $|m(r)-m(r')|\le\|U(r)-U(r')\|$; in
particular $m$ is continuous on $(0,r_0)$, and
$|m(r)|\le\|U(r)\|\to0$ as $r\to0^+$.

\emph{Step 1: differential inequality for $m$ via upper support
functions.} Fix $x\in(0,r_0)$ and let $X_0$ be a unit vector
realizing $m(x)$. A minimizer on the unit sphere of the quadratic
form of a symmetric endomorphism is an eigenvector, with eigenvalue
the minimum value: $U(x)X_0=m(x)X_0$. The differentiable function
$\varphi(y)=\langle U(y)X_0,X_0\rangle$ satisfies $\varphi\ge m$
everywhere (definition of the minimum as an infimum over unit
vectors), $\varphi(x)=m(x)$, and, by the hypothesis at $X_0$
together with the eigenvector relation,
$$\varphi'(x)=\langle U'(x)X_0,X_0\rangle\ \ge\
-\langle U(x)^2X_0,X_0\rangle-2a(x)\langle U(x)X_0,X_0\rangle
=-m(x)^2-2a(x)\,m(x).$$
Hence for $y<x$,
$$\frac{m(x)-m(y)}{x-y}\ \ge\ \frac{\varphi(x)-\varphi(y)}{x-y}
\ \xrightarrow[\ y\to x^-\ ]{}\ \varphi'(x),$$
so for every $\delta>0$ the left difference quotients of $m$ at
$x$ eventually exceed $-m(x)^2-2a(x)m(x)-\delta$.

\emph{Step 2: barrier lemma.} If $h$ is continuous on $[c,d]$ and
at every $x\in(c,d]$ its left difference quotients
$\frac{h(x)-h(y)}{x-y}$ are, as $y\to x^-$, eventually smaller
than every positive $\delta$, then $h(d)\le h(c)$. Indeed, the
reflection $g(t)=-h(-t)$ on $[-d,-c]$ has, at every
$t\in[-d,-c)$, right difference quotients eventually below every
positive $\delta$ (the quotients coincide:
$\frac{g(z)-g(t)}{z-t}=\frac{h(-t)-h(-z)}{(-t)-(-z)}$), so the
standard fencing argument for continuous functions — comparing $g$
with the linear barriers $g(-d)+\varepsilon(t+d)$, $\varepsilon>0$,
whose slope strictly exceeds the limit inferior of the difference
quotients of $g$ at any contact point — gives
$g(-c)\le g(-d)+\varepsilon(d-c)$ for every $\varepsilon>0$, i.e.\
$h(d)\le h(c)$.

\emph{Step 3: conclusion.} Suppose for contradiction that
$m(r_2)<0$ for some $r_2\in(0,r_0)$. Since $\|U\|$ is continuous
on $(0,r_2]$ and tends to $0$ at $0^+$, it is bounded there: fix
$W\ge0$ with $\|U(r)\|\le W$, hence $|m(r)|\le W$, on $(0,r_2]$.
Let
$$r_3=\sup\bigl(\{0\}\cup\{r\in(0,r_2]:m(r)\ge0\}\bigr),$$
so $0\le r_3<r_2$ (by continuity $m<0$ on a neighborhood of
$r_2$), $m<0$ on $(r_3,r_2]$, and $m(r_3)\ge0$ if $r_3>0$
(continuity again, as $r_3$ is then a limit of points where
$m\ge0$). Put $w=-m$ and, on $(r_3,r_2]$,
$$h(r)=w(r)\,\rho(r),\qquad \rho(r)=s(r)^2e^{-Wr},$$
so $0<w\le W$ on $(r_3,r_2]$ and $\rho>0$ is differentiable with
$\rho'=\rho\,(2a-W)$. For $x\in(r_3,r_2]$ and $y<x$,
$$\frac{h(x)-h(y)}{x-y}=\rho(x)\,\frac{w(x)-w(y)}{x-y}
+w(y)\,\frac{\rho(x)-\rho(y)}{x-y};$$
as $y\to x^-$ the second summand tends to $w(x)\rho'(x)$ while, by
Step~1, the first is eventually less than
$\rho(x)\bigl(w(x)^2-2a(x)w(x)+\delta\bigr)$ for every
$\delta>0$. Since
$$\rho\,\bigl(w^2-2aw\bigr)+w\rho'
=\rho\,\bigl(w^2-2aw\bigr)+w\rho\,(2a-W)
=\rho\,w\,(w-W)\ \le\ 0\qquad\text{on }(r_3,r_2]$$
(using $0\le w\le W$ and $\rho>0$), the left difference quotients
of $h$ at each $x\in(r_3,r_2]$ are eventually smaller than every
positive $\delta$. By Step~2 applied on $[\bar r,r_2]$,
$$h(r_2)\ \le\ h(\bar r)\qquad\text{for every }\bar
r\in(r_3,r_2).$$
Finally, as $\bar r\to r_3^+$: if $r_3>0$ then
$h(\bar r)\to w(r_3)s(r_3)^2e^{-Wr_3}\le0$ (as $w(r_3)\le0$),
while if $r_3=0$ then $h(\bar r)\to0$ (as $w\to0$ while $s$ is
bounded near $0^+$ and $e^{-Wr}\le1$). Either way $h(r_2)\le0$,
contradicting $h(r_2)=w(r_2)s(r_2)^2e^{-Wr_2}>0$. Hence $m\ge0$ on
$(0,r_0)$, i.e.\ $U(r)\ge0$ for all $r\in(0,r_0)$.
\end{proof}

In the opposite direction — an upper bound on the operator under a
lower bound on the curvature term — the Cauchy--Schwarz inequality
points the right way, and the operator bound follows from the scalar
comparison along each fixed direction.

\begin{lemma}[Operator Riccati comparison, upper form]
\label{lem:operator-riccati-upper}
\uses{lem:scalar-riccati-comparison}
\lean{MorganTianLib.operator_riccati_le}
\leanok
Let $E$ be a real inner product space and fix $k\ge0$, $r_0>0$. Let
$A(r)$, $r\in(0,r_0)$, be a differentiable family of symmetric
endomorphisms of $E$ with
$A(r)-\tfrac1r\,\mathrm{Id}\to0$ in operator norm as $r\to0^+$,
satisfying
$$\langle A'(r)X,X\rangle\ \le\ k\,|X|^2-\langle A(r)^2X,X\rangle
\qquad\text{for all }X\in E,\ r\in(0,r_0).$$
Then for all $r\in(0,r_0)$ and $X\in E$,
$$\langle A(r)X,X\rangle\ \le\
\frac{\sn_k'(r)}{\sn_k(r)}\,|X|^2 .$$
\end{lemma}

\begin{proof}
\uses{lem:scalar-riccati-comparison}
\leanok
By homogeneity it suffices to consider a unit vector $X$. Set
$\phi(r)=\langle A(r)X,X\rangle$. By symmetry
$\langle A^2X,X\rangle=|AX|^2$, and by the Cauchy--Schwarz
inequality $\phi^2\le|AX|^2$, so the hypothesis gives
$$\phi'\ \le\ k-|AX|^2\ \le\ k-\phi^2 .$$
Moreover
$|\phi(r)-\tfrac1r|=|\langle(A(r)-\tfrac1r\,\mathrm{Id})X,X\rangle|
\le\|A(r)-\tfrac1r\,\mathrm{Id}\|\to0$ as $r\to0^+$. The scalar
Riccati comparison (Lemma~\ref{lem:scalar-riccati-comparison})
yields $\phi(r)\le\sn_k'(r)/\sn_k(r)$ on $(0,r_0)$.
\end{proof}

\begin{lemma}[Rauch lower bounds under an upper curvature bound]
\label{lem:rauch-lower}
\uses{lem:geodesic-polar-form,def:sectional-curvature}
Fix $\bar K\ge0$. In the setting of
Lemma~\ref{lem:geodesic-polar-form} (a cone $C$ of radius $r_0$ on
which $\exp_p$ is defined and non-singular), suppose in addition
that $r_0\le D_{\bar K}=\pi/\sqrt{\bar K}$ and that $K(P)\le\bar K$
for every $2$-plane $P$ tangent at points of $C$'s image. Then for
$0<r<r_0$, writing $s=s_{\bar K}$:
\begin{enumerate}
\item $A(r)\ge\frac{s'(r)}{s(r)}\,\mathrm{Id}$ as symmetric
operators on the tangent space of the coordinate spheres;
\item $g_{ij}(r,\theta)\ge s(r)^2\,\hat g_{ij}(\theta)$ as
bilinear forms;
\item the function $\frac12r^2$ satisfies
$$\Hess\left(\tfrac12r^2\right)\ \ge\
\min\left(1,\ r\,\frac{s'(r)}{s(r)}\right)\tilde g .$$
\end{enumerate}
\end{lemma}

\begin{proof}
\uses{lem:geodesic-polar-form,lem:hessian-symmetric,lem:operator-riccati-lower,lem:ratio-monotone}
(1) Fix a radial geodesic and let $a(r)=s'(r)/s(r)$, so
$a'+a^2=-\bar K$ and, by the Taylor expansion of $s$ at $0$,
$a(r)-1/r\to0$ as $r\to0^+$. Let $U(r)=A(r)-a(r)\,\mathrm{Id}$, a
symmetric operator field along the geodesic with $U(r)\to0$ as
$r\to0^+$ (Lemma~\ref{lem:geodesic-polar-form}(3)). By the Riccati
equation of Lemma~\ref{lem:geodesic-polar-form}(2) and the
curvature hypothesis $R_{\partial_r}\le\bar K\,\mathrm{Id}$ (its
quadratic form on unit vectors is a sectional curvature),
$$\nabla_{\partial_r}U=-A^2-R_{\partial_r}-a'\,\mathrm{Id}
=-\left(A^2-a^2\,\mathrm{Id}\right)-\left(R_{\partial_r}
-\bar K\,\mathrm{Id}\right)\cdot 1
\ \ge\ -\left(A^2-a^2\,\mathrm{Id}\right)
=-\left(U^2+2a\,U\right),$$
where we used that $A=U+a\,\mathrm{Id}$ commutes with $U$. Now
express $U$ in a parallel orthonormal frame of $\partial_r^\perp$
along the radial geodesic: parallel transport preserves the metric
and $\partial_r$, so this identifies $U$ with a differentiable
family of symmetric endomorphisms of a fixed euclidean space of
dimension $n-1$ (nontrivial, since the polar-coordinate framework
presupposes $n\ge2$; for $n=1$ the assertion is vacuous), with
$\nabla_{\partial_r}U$ becoming the ordinary derivative $U'$, and
the displayed inequality becomes, as quadratic forms,
$$\langle U'(r)X,X\rangle\ \ge\
-\langle U(r)^2X,X\rangle-2a(r)\langle U(r)X,X\rangle .$$
Since $r_0\le D_{\bar K}$, the function $s=s_{\bar K}$ is positive
and differentiable on $(0,r_0)$ and bounded near $0^+$, and
$U(r)\to0$ as $r\to0^+$; so
Lemma~\ref{lem:operator-riccati-lower} applies and gives $U\ge0$
on $(0,r_0)$, i.e.\ $A\ge\frac{s'}{s}\,\mathrm{Id}$, which is
(1).

(2) Fix a sphere-tangent coordinate vector $w$ with constant
components and set $h(r)=g_{ij}(r,\theta)w^iw^j$. By
Lemma~\ref{lem:geodesic-polar-form}(2) and part (1),
$h'=2S(w,w)\ge2\frac{s'}{s}h$, so $h/s^2$ is non-decreasing
(Lemma~\ref{lem:ratio-monotone}); by
Lemma~\ref{lem:geodesic-polar-form}(3) its limit as $r\to0^+$ is
$\hat g_{ij}w^iw^j$, whence $h(r)\ge s(r)^2\hat g_{ij}w^iw^j$.

(3) The gradient of $\frac12r^2$ is $r\partial_r$, so by the
Hessian formula (Lemma~\ref{lem:hessian-symmetric}) and the Leibniz
rule,
$$\Hess\left(\tfrac12r^2\right)(X,Y)
=\left\langle\nabla_X\left(r\partial_r\right),Y\right\rangle
=dr(X)\,dr(Y)+r\,S(X,Y).$$
Decomposing $X$ into radial and sphere-tangent parts and using
$S(\partial_r,\cdot)=0$ together with part (1) on the
sphere-tangent part gives the assertion, since on the radial part
the form is $dr\otimes dr$ and on the spherical part it is $r\,S\ge
r\frac{s'}{s}\,g$.
\end{proof}

Before passing to polar coordinates we isolate the comparison estimate
itself, as a statement about Jacobi fields along a single geodesic.  This
is the analytic content of the Sectional Curvature Comparison theorem; the
passage to the metric $g_{ij}(r,\theta)$ afterwards is the Gauss Lemma
(Lemma~\ref{lem:geodesic-polar-form}(1)) applied to the polar coordinate
fields, which are exactly Jacobi fields vanishing at $p$.

\begin{lemma}[Jacobi field comparison along a geodesic]
\label{lem:jacobi-metric-comparison}
\uses{def:geodesic,def:conjugate-point,def:sectional-curvature,
def:sinh-comparison-function,lem:parallel-frame,lem:jacobi-frame-reduction,
lem:jacobi-operator-symmetric,lem:second-order-linear-ode,lem:radial-comparison}
\lean{MorganTianLib.exists_isRadialJacobi,
MorganTianLib.IsJacobiSolOn.apply,
MorganTianLib.continuousOn_frameCurv,
MorganTianLib.frameCurvOp_symm,
MorganTianLib.isJacobiSolOn_frameVec,
MorganTianLib.exists_isRadialJacobi_of_geodesic,
MorganTianLib.curvatureFormAt_jacobi_le_of_sectionalCurvatureAt_ge,
MorganTianLib.metricInner_jacobi_le_snK,
MorganTianLib.metricInner_jacobi_le_snK_of_sectionalCurvature}
\leanok
Fix $k\ge 0$. Let $\gamma$ be a geodesic of $(M,g)$ parameterized at unit
speed, defined on an interval containing $[0,r_0]$ in its interior, with
$\gamma(0)=p$, and suppose
$$-k\le K(P)\qquad\text{for every $2$-plane }P\subset T_{\gamma(r)}M,\quad
0<r<r_0.$$
Suppose moreover that $\gamma$ has no conjugate point of $p$ on $(0,r_0)$,
i.e.\ that the matrix Jacobi field $\mathcal J$ of $\gamma$ (the solution of
$\mathcal J''+R\mathcal J=0$ with $\mathcal J(0)=0$, $\mathcal J'(0)=\Id$ in a
parallel orthonormal frame) is invertible on $(0,r_0)$. Then every Jacobi
field $J$ along $\gamma$ with $J(0)=0$ satisfies
$$|J(r)|_g\ \le\ \sn_k(r)\,|\nabla J(0)|_g,\qquad 0<r<r_0.$$
\end{lemma}

\begin{proof}
\uses{lem:parallel-frame,lem:jacobi-frame-reduction,lem:jacobi-operator-symmetric,
lem:second-order-linear-ode,lem:radial-comparison,def:sinh-comparison-function}
\leanok
Choose a parallel $g$-orthonormal frame $E_1,\ldots,E_n$ along $\gamma$
(Lemma~\ref{lem:parallel-frame}: parallel transport is a $g$-isometry, so
orthonormality at $\gamma(0)$ propagates). Write $c_i(r)=\langle
J(r),E_i(r)\rangle$ and $d_i(r)=\langle\nabla J(r),E_i(r)\rangle$. Because the
frame is parallel, $c_i'=d_i$; and expanding $J=\sum_j c_jE_j$ in the curvature
slot of the Jacobi equation turns it into the closed linear scalar system
$$c_i''=\sum_j\mathcal R_{ij}c_j,\qquad
\mathcal R_{ij}(r)={\mathcal R}\big(E_j,\gamma',\gamma',E_i\big)(r)$$
(Lemma~\ref{lem:jacobi-frame-reduction}), whose coefficient matrix is symmetric
by self-adjointness of the Jacobi operator
(Lemma~\ref{lem:jacobi-operator-symmetric}). Put $R=-\mathcal R$, so that the
system reads $c''+Rc=0$.

The matrix $R(r)$ depends continuously on $r$: continuity is a local property,
and each $r$ has a neighbourhood on which $\gamma$ stays inside a single chart;
there $R$ is the chart Gram pairing of the chart curvature (a smooth function of
the chart point) with the chart readings of $\gamma'$ and of the frame, and those
readings are continuous because they solve the geodesic, resp.\ the parallel
transport, ODE in that chart. Being continuous on the compact $[0,r_0]$, $R$ is
bounded there. Hence the matrix equation $\mathcal J''+R\mathcal J=0$,
$\mathcal J(0)=0$, $\mathcal J'(0)=\Id$, is a linear ODE with continuous bounded
coefficient, and so has a solution on all of $[0,r_0]$
(Lemma~\ref{lem:second-order-linear-ode}); its columns solve the vector equation
$c''+Rc=0$. Both $c$ and $\mathcal J\cdot c'(0)$ solve that linear system with
the same initial data $(0,c'(0))$ (here $c(0)=0$ because $J(0)=0$), so by
uniqueness for linear ODEs
$$c(r)=\mathcal J(r)\,c'(0),\qquad 0\le r\le r_0 .$$

Next, the curvature hypothesis. For $X\in T_{\gamma(r)}M$, antisymmetry of the
curvature tensor in its last pair gives
${\mathcal R}(X,\gamma',\gamma',X)=-{\mathcal R}(X,\gamma',X,\gamma')$, and the
definition of the sectional curvature gives
${\mathcal R}(X,\gamma',X,\gamma')=K(X\wedge\gamma')\,|X\wedge\gamma'|^2\ge
-k\,|X\wedge\gamma'|^2$ (both sides vanish when $X$ and $\gamma'$ are dependent).
Since $\gamma$ has unit speed, $|X\wedge\gamma'|^2=|X|^2-\langle
X,\gamma'\rangle^2\le|X|^2$, and $k\ge0$, so
$${\mathcal R}(X,\gamma',\gamma',X)\le k\,|X|^2,\qquad\text{i.e.}\qquad
\langle R(r)x,x\rangle\ge -k\,|x|^2$$
for the coefficient vector $x$ of $X$ in the frame (the frame being orthonormal,
the coefficient map is a linear isometry of $(T_{\gamma(r)}M,g)$ onto
$\mathbb R^n$).

The hypotheses of Lemma~\ref{lem:radial-comparison} are now met on $(0,r_0)$, and
its second conclusion gives $|\mathcal J(r)w|^2\le\sn_k^2(r)|w|^2$ for every $w$.
Taking $w=c'(0)$, the coefficient vector of $\nabla J(0)$, and using the isometry
of the coefficient map once more,
$$|J(r)|_g^2=|c(r)|^2=|\mathcal J(r)c'(0)|^2\le\sn_k^2(r)\,|c'(0)|^2
=\sn_k^2(r)\,|\nabla J(0)|_g^2 . \qedhere$$
\end{proof}

Now we can compare manifolds of varying sectional curvature with
those of constant curvature.


\begin{theorem}[Sectional curvature comparison theorem]
\label{thm:sectional-curvature-comparison}
\label{SCC}
\uses{def:sectional-curvature,def:sinh-comparison-function}
(Sectional Curvature Comparison) Fix $k\ge 0$. Let $(M, g)$ be a
Riemannian manifold with the property that $-k \leq K(P)$ for every
$2$-plane $P$ in $TM$. Fix a minimizing geodesic $\gamma\colon
[0,r_0)\to M$ parameterized at unit speed with $\gamma(0)=p$. Impose
Gaussian polar coordinates $(r,\theta^1,\ldots,\theta^{n-1})$ on a
neighborhood of $\gamma$ so that $g=dr^2+g_{ij}\theta^i\otimes
\theta^j$. Then for all $0<r<r_0$, at $\gamma(r)$ we have
$$ (g_{ij}(r,\theta))_{1 \leq
i,j \leq n-1} \leq \sn_{k}^{2}(r),$$ and the shape operator
associated to  the distance function from $p$, $f$, satisfies
$$(S_{ij}(r,\theta))_{1 \leq i,j \leq n-1} \leq \sqrt{k}\ct_{k}(r).$$
\end{theorem}

The two displayed inequalities are to be read as inequalities of
symmetric bilinear forms on the tangent space to the coordinate
sphere at the point $\gamma(r)$: the first says $g_{ij}(r,\theta)w^iw^j\le
\sn_k^2(r)\,\hat g_{ij}(\theta)w^iw^j$ for every $w$, where $\hat g$
is the round metric of the unit sphere of $T_pM$ in the coordinates
$\theta$, and the second says $S_{ij}(r,\theta)w^iw^j\le
\sqrt{k}\,\ct_k(r)\,g_{ij}(r,\theta)w^iw^j$. (Choosing the
coordinates $\theta$ so that at the point under consideration
$\hat g_{ij}=\delta_{ij}$, respectively $g_{ij}=\delta_{ij}$, one
recovers the literal matrix inequalities displayed.) Note that
$\sqrt k\,\ct_k(r)=\sn_k'(r)/\sn_k(r)$.

% Formalization status: the analytic content of this proof -- the shape-operator
% and metric estimates for the matrix Jacobi field in a parallel frame -- is
% formalized and checked as lem:radial-comparison(1),(2), and the MANIFOLD-TO-FRAME
% BRIDGE is now closed as well: lem:jacobi-metric-comparison (above, \leanok) proves
% the metric estimate |J(r)| <= sn_k(r)|nabla J(0)| for every Jacobi field vanishing
% at p along a unit-speed geodesic with -k <= K(P) and no conjugate point, directly
% on the manifold, from Morgan-Tian's own curvature hypothesis
% (MorganTianLib.metricInner_jacobi_le_snK_of_sectionalCurvature).
%
% What remains for \leanok on SCC itself:
%   (a) the POLAR-COORDINATE reading: identifying the coordinate fields
%       d/dtheta^i of Gaussian polar coordinates with Jacobi fields vanishing at p,
%       i.e. lem:exponential-differential-jacobi (d exp_p = Jacobi field) plus
%       lem:geodesic-polar-form(1).  Only this turns lem:jacobi-metric-comparison's
%       conclusion into the displayed g_{ij}(r,theta) <= sn_k^2(r).
%   (b) the SHAPE-OPERATOR half, S_ij <= sqrt(k) ct_k(r): the estimate itself is
%       lem:radial-comparison(1) (MorganTianLib.shapeOp_inner_le, \leanok), but it is
%       stated for the operator A of the matrix Jacobi field; identifying A with
%       S = Hess(r) is lem:geodesic-polar-form(2), still open.
%   (c) discharging the no-conjugate-point hypothesis from minimality, i.e.
%       prop:minimal-geodesic-no-conjugate.
\begin{proof}
\uses{prop:minimal-geodesic-no-conjugate,lem:geodesic-polar-form,lem:radial-comparison,lem:jacobi-metric-comparison,lem:operator-riccati-upper,lem:ratio-monotone,lem:hessian-symmetric,def:cut-locus,lem:cut-locus-properties}
Since $\gamma$ is minimizing on $[0,r_0)$, for every $t<r_0$ the
restriction $\gamma|_{[0,t]}$ extends to a minimizing geodesic on a
larger interval, so by
Proposition~\ref{prop:minimal-geodesic-no-conjugate} it is the
unique minimal geodesic between its endpoints and there are no
conjugate points along $\gamma|_{[0,r_0)}$. By the description of
the differential of $\exp_p$ in terms of Jacobi fields, $\exp_p$
is non-singular at $tv$ for $0\le t<r_0$, where $v=\gamma'(0)$;
since non-singularity is an open condition and $\exp_p$ is
defined on an open set, there are an open neighborhood $V$ of $v$
in the unit sphere of $T_pM$ and, for each $t_1<r_0$, a cone
$C=\{tw\,|\,w\in V,\ 0\le t<t_1\}$ on which $\exp_p$ is defined
and non-singular. This provides the Gaussian polar coordinates of
the statement, and Lemma~\ref{lem:geodesic-polar-form} applies on
$C$. Moreover, along $\gamma$ the coordinate $r$ agrees with the
distance function $f=d(p,\cdot)$ near each point $\gamma(t)$,
$0<t<r_0$: by Proposition~\ref{prop:minimal-geodesic-no-conjugate}
and the openness of $U_p$
(Lemma~\ref{lem:cut-locus-properties}(1)), the
points $tv$, $t<r_0$, lie in $U_p$ together with a neighborhood, and
on $\exp_p(U_p)$ the distance to $p$ of $\exp_p(y)$ equals
$|y|$, which is the polar coordinate $r$. Hence the shape operator
of the statement is the tensor $S$ of
Lemma~\ref{lem:geodesic-polar-form}(2).

\emph{The shape operator estimate.} Write $\theta_0$ for the
angular coordinate of $\gamma$, so that $\gamma$ is the ray
$\theta=\theta_0$. Express the $g$-symmetric operator $A$ of
Lemma~\ref{lem:geodesic-polar-form}(2) in a parallel orthonormal
frame of $\partial_r^\perp$ along $\gamma$, so that $A$ becomes a
differentiable family of symmetric endomorphisms of a fixed
euclidean space and $\nabla_{\partial_r}A$ becomes the ordinary
derivative $A'$. The Riccati equation of
Lemma~\ref{lem:geodesic-polar-form}(2) and the curvature
hypothesis $K\ge-k$ give, as quadratic forms on sphere-tangent
vectors,
$$\langle A'X,X\rangle=-\langle A^2X,X\rangle
-K(X\wedge\partial_r)\,|X|^2
\ \le\ k\,|X|^2-\langle A^2X,X\rangle,$$
while $A(r)-\tfrac1r\,\mathrm{Id}\to0$ as $r\to0^+$ by
Lemma~\ref{lem:geodesic-polar-form}(3). The operator Riccati
comparison (Lemma~\ref{lem:operator-riccati-upper}) yields
$$\langle A(r)X,X\rangle\ \le\
\frac{\sn_k'(r)}{\sn_k(r)}\,|X|^2\qquad(0<r<r_0),$$
which is exactly the asserted inequality $S\le\sqrt
k\,\ct_k(r)\,g$ on the tangent space of the coordinate sphere,
since $\sn_k'/\sn_k=\sqrt k\,\ct_k$.

\emph{The metric estimate.} Fix a tangent vector $w=w^i
\partial_{\theta^i}$ to the coordinate sphere, with constant
components $w^i$, and set $h(r)=g_{ij}(r,\theta_0)w^iw^j$ along
$\gamma$. By Lemma~\ref{lem:geodesic-polar-form}(2) and the shape
operator estimate just proved,
$$h'(r)=\left(\partial_rg_{ij}\right)w^iw^j=2S_{ij}w^iw^j
\le2\frac{\sn_k'(r)}{\sn_k(r)}\,h(r),$$
so $h/\sn_k^2$ is non-increasing on $(0,r_0)$ by
Lemma~\ref{lem:ratio-monotone}. By Lemma~\ref{lem:geodesic-polar-form}(3),
$h(r)/r^2\to\hat g_{ij}(\theta_0)w^iw^j$ and
$\sn_k^2(r)/r^2\to 1$ as $r\to0$, so
$h(r)/\sn_k^2(r)\to\hat g_{ij}w^iw^j$. Monotonicity then gives
$$g_{ij}(r,\theta_0)w^iw^j=h(r)\le\sn_k^2(r)\,\hat
g_{ij}(\theta_0)w^iw^j\qquad(0<r<r_0),$$
which is the first inequality.
\end{proof}

The analogous comparisons in the opposite direction, under a
positive upper bound for the sectional curvature, were recorded in
Lemma~\ref{lem:rauch-lower} above. From an upper curvature bound we
shall also need the following local diffeomorphism property of the
exponential mapping.



% Formalization status (run 0110, session 0014).  The MATHEMATICAL CONTENT of this lemma is
% formalized and kernel-checked (axiom-clean, sorry-free).  The node is still NOT \leanok, for
% one reason only -- a hypothesis-form mismatch, recorded as (b) below.
%
% PROVED.  For every v with sqrt(K)*|v|_g < pi -- the statement's |v| < pi/sqrt K, with the
% pi/sqrt 0 = +infty convention built into the product form -- and an upper sectional bound K
% on the closed ball of radius |v|_g about p, exp_p is a LOCAL DIFFEOMORPHISM AT v:
%   * expDifferential_isEquiv_of_sectionalCurvatureAt_le  (ExpLocalDiffeo.lean)
%       d(exp_p)_v is a continuous linear ISOMORPHISM.  This is what the proof below actually
%       establishes, and what the applications consume.
%   * expMapGlobal_locallyInjective_of_sectionalCurvatureAt_le  (ExpLocalDiffeo.lean)
%       exp_p is INJECTIVE on a neighbourhood of v.
%   * expMapGlobal_map_nhds_of_sectionalCurvatureAt_le  (ExpContinuity.lean)
%       exp_p maps neighbourhoods of v ONTO neighbourhoods of exp_p(v) -- so the image is open
%       in M, and exp_p is a local HOMEOMORPHISM at v (packaged as
%       expMapGlobal_isLocalHomeomorphAt_of_sectionalCurvatureAt_le).  The chart-level inverse
%       is smooth by the inverse function theorem, whence "diffeomorphism".
% Supporting infrastructure, new and reusable:
%   * not_isConjugatePointAt_one_of_sectionalCurvatureAt_le -- the no-conjugate-point step, via
%     lem:conjugate-sturm applied to the unit-speed reparametrization and rescaled back
%     (JacobiRescale.lean: Jacobi fields and conjugate points transport across t |-> c*t, the
%     chart pair system scaling by c and c^2 -- the Christoffel contraction being linear and
%     the curvature quadratic in the velocity).
%   * hasStrictFDerivAt_chartReading_expMapGlobal -- the STRICT form of
%     lem:exponential-differential-jacobi (the flow chain already produced a strict derivative;
%     the strict form is what the inverse function theorem consumes).
%   * continuousAt_expMapGlobal / continuous_expMapGlobal -- exp_p is CONTINUOUS, from
%     lem:geodesic-continuous-dependence.  Reusable far beyond this lemma (normal balls, the
%     cut locus, prop:exponential-diffeomorphism-cut-locus).
%
% (b) THE REMAINING GAP -- hypothesis form and quantification, not mathematics.  The Lean
%     statements take a SECTIONAL upper bound (sectionalCurvatureAt <= K on the closed ball of
%     radius |v|_g) rather than the curvature-operator bound |Rm| <= K, and they are stated per
%     tangent vector v rather than for the whole ball at once.  Two things are needed:
%       (b1) DONE this session, at the level of a single fiber:
%            abs_sectionalCurvature_le_of_hasCurvatureOperatorNormLe (CurvatureNormSectional.lean)
%            proves "|Rm| <= K ==> |K(P)| <= K for EVERY pair x, y" -- the final claim of
%            def:curvature-operator-norm, which was previously available only for ORTHONORMAL
%            pairs.  Gram-Schmidt (exists_orthonormal_changeBasis) plus the basis-invariance of
%            the sectional curvature; the degenerate (dependent) case is the junk value 0 <= K.
%       (b2) STILL OPEN, routine: lift (b1) from a fiber to a point of M -- i.e. state it for
%            sectionalCurvatureAt g nabla x, which carries a `Bundle.RiemannianBundle` local
%            instance (see isAlgCurvatureForm_curvatureFormAt) -- and then re-quantify the
%            results above over all v in B(0, pi/sqrt K), adding the image claim
%            exp_p(B(0,pi/sqrt K)) is contained in B(p, pi/sqrt K) (immediate: the geodesic from
%            p to exp_p(v) has length |v| < pi/sqrt K).  That earns \leanok.
\begin{lemma}[Exponential map local diffeomorphism for bounded curvature]
\label{lem:local-diffeomorphism-bounded-curvature}
\label{localdiffeo}
\uses{def:curvature-operator-norm,def:exponential-map}
\lean{MorganTianLib.hasStrictFDerivAt_chartReading_expMapGlobal,
MorganTianLib.speedSq_globalGeodesic,
MorganTianLib.globalGeodesic_smul,
MorganTianLib.not_isConjugatePointAt_one_of_sectionalCurvatureAt_le,
MorganTianLib.expDifferential_isEquiv_of_sectionalCurvatureAt_le,
MorganTianLib.expMapGlobal_locallyInjective_of_sectionalCurvatureAt_le,
MorganTianLib.continuousAt_expMapGlobal,
MorganTianLib.continuous_expMapGlobal,
MorganTianLib.expMapGlobal_map_nhds_of_sectionalCurvatureAt_le,
MorganTianLib.expMapGlobal_isLocalHomeomorphAt_of_sectionalCurvatureAt_le}
Fix $K\ge 0$, with the convention $\pi/\sqrt K=+\infty$ for $K=0$.
Assume that $|\Rm(x)|\le K$ (in the sense of
Definition~\ref{def:curvature-operator-norm}) for all $x\in
B(p,\pi/\sqrt{K})$, and that $\exp_p$ is defined on the ball
$B(0,\pi/\sqrt{K})$ in $T_pM$, i.e., that every geodesic from $p$
extends to length $\pi/\sqrt K$ (as holds, for instance, if $M$ is
complete, or if $B(p,\pi/\sqrt K)$ has compact closure). Then
$\exp_p$ is a local diffeomorphism from $B(0,\pi/\sqrt{K})$ to
the  ball $B(p,\pi/\sqrt{K})$ in
$M$.
\end{lemma}

\begin{proof}
\uses{lem:conjugate-sturm,def:curvature-operator-norm,def:exponential-map,def:sectional-curvature}
By Definition~\ref{def:curvature-operator-norm}, the hypothesis
implies $|K(P)|\le K$ for every $2$-plane $P$ tangent to $M$ at a
point of $B(p,\pi/\sqrt K)$.

Let $v\in T_pM$ with $0<|v|<\pi/\sqrt K$ and let
$\gamma(t)=\exp_p(tv/|v|)$ be the unit-speed radial geodesic.
Choose $T$ with $|v|<T<\pi/\sqrt K$; by the extension assumption,
$\gamma$ is defined on $[0,T]$. For $t\le T$ we have
$d(p,\gamma(t))\le t\le T<\pi/\sqrt{K}$, so
$\gamma([0,T])\subset B(p,\pi/\sqrt K)$ and the sectional
curvatures along $\gamma|_{[0,T]}$ are bounded above by $K$.
Applying Lemma~\ref{lem:conjugate-sturm} to $\gamma|_{[0,T]}$, and
noting that $\min(T,\pi/\sqrt K)=T>|v|$, there is no conjugate
point along
$\gamma$ at any parameter $t\le|v|$. By the description
of the differential of the exponential map in terms of Jacobi
fields, the kernel of $d(\exp_p)_v$ is identified with the
space of Jacobi fields along $\gamma$ vanishing at $0$ and at
$|v|$; since only the zero field does so, $d(\exp_p)_v$ is an
isomorphism. The same holds trivially at $v=0$. Hence $\exp_p$
is a local diffeomorphism on $B(0,\pi/\sqrt K)$, and its image lies
in $B(p,\pi/\sqrt K)$ because the geodesic from $p$ to
$\exp_p(v)$ has length $|v|<\pi/\sqrt K$.
\end{proof}

There is a crucial comparison result for volume which involves the
Ricci curvature. Its algebraic-analytic core is the following traced
form of the Riccati comparison, which converts the matrix Riccati
equation of the shape operator into the scalar comparison via the
Cauchy--Schwarz inequality for traces.

\begin{lemma}[Trace Riccati comparison]
\label{lem:trace-riccati-comparison}
\uses{def:sinh-comparison-function}
\lean{MorganTianLib.trace_riccati_comparison,MorganTianLib.sq_trace_le_finrank_mul_trace_comp_self}
\leanok
Fix $k\ge0$ and $r_0>0$. Let $E$ be a real inner product space of
finite dimension $m\ge1$ and let $A(r)$, $0<r<r_0$, be a
differentiable family of symmetric endomorphisms of $E$ such that
$$\Tr A'(r)+\Tr\!\left(A(r)^2\right)\ \le\ m\,k
\qquad(0<r<r_0),$$
and $\Tr A(r)-m/r\to0$ as $r\to0^+$. Then
$$\Tr A(r)\ \le\ m\,\frac{\sn_k'(r)}{\sn_k(r)}\qquad(0<r<r_0).$$
\end{lemma}

\begin{proof}
\uses{lem:scalar-riccati-comparison}
\leanok
First, for any symmetric endomorphism $A$ of $E$,
$$(\Tr A)^2\ \le\ m\,\Tr(A^2):$$
in an orthonormal basis $(e_i)_{i=1}^m$ of $E$, symmetry gives
$\Tr(A^2)=\sum_i\langle A^2e_i,e_i\rangle=\sum_i|Ae_i|^2$, the
Cauchy--Schwarz inequality in $E$ gives
$|Ae_i|^2\ge\langle Ae_i,e_i\rangle^2$ for each $i$ (as $|e_i|=1$),
and the Cauchy--Schwarz inequality in $\mathbb{R}^m$ gives
$\left(\sum_i\langle Ae_i,e_i\rangle\right)^2
\le m\sum_i\langle Ae_i,e_i\rangle^2$; combining,
$(\Tr A)^2\le m\,\Tr(A^2)$. Now set $\phi(r)=\Tr A(r)/m$. Then
$\phi$ is differentiable on $(0,r_0)$ with $\phi'=\Tr A'/m$ (the
trace is linear), and
$$\phi'+\phi^2=\frac{\Tr A'}{m}+\frac{(\Tr A)^2}{m^2}
\ \le\ \frac{\Tr A'+\Tr(A^2)}{m}\ \le\ k,$$
using the trace inequality and the hypothesis. Moreover
$\phi(r)-1/r=(\Tr A(r)-m/r)/m\to0$ as $r\to0^+$. By the scalar
Riccati comparison (Lemma~\ref{lem:scalar-riccati-comparison}),
$\phi(r)\le\sn_k'(r)/\sn_k(r)$ on $(0,r_0)$, which is the claim.
\end{proof}

We can now state the monotonicity of the volume element in geodesic
polar coordinates, which we isolate because it will be used again in
the proofs of the Bishop--Gromov inequality and of the
volume--injectivity radius estimate.

\begin{lemma}[Monotonicity of the volume element]
\label{lem:volume-element-comparison}
\uses{lem:geodesic-polar-form,def:ricci-curvature,def:sinh-comparison-function}
Fix $k\ge0$. Let $p\in M$ and let $\gamma(t)=\exp_p(tv)$ be a
unit-speed radial geodesic, defined and free of conjugate points on
$[0,r_0)$, along which $\Ric\ge-(n-1)k$. Let $\lambda(t,\theta)$ be
the volume element of Lemma~\ref{lem:geodesic-polar-form}(4) in
polar coordinates around $p$ (defined near the direction of
$\gamma$ for $t<r_0$). Then along $\gamma$,
$$\Tr(A)(t)\le(n-1)\frac{\sn_k'(t)}{\sn_k(t)},\qquad
\frac{\lambda(t)}{\sn_k^{n-1}(t)}\ \text{is non-increasing on
}(0,r_0),\qquad \lim_{t\to0^+}\frac{\lambda(t)}{\sn_k^{n-1}(t)}=1.$$
In particular $\lambda(t)\le\sn_k^{n-1}(t)$ for $t\in(0,r_0)$.
\end{lemma}

% Formalization status: the first assertion, Tr(A)(t) <= (n-1) sn_k'/sn_k in a
% parallel orthonormal frame, is formalized and checked as
% lem:radial-comparison(3), and the volume-element derivative
% lambda' = lambda Tr(A) as lem:radial-volume-element. What remains for \leanok
% here is (i) the positivity lambda > 0 on (0,r_0) and the normalization
% lambda(t)/t^{n-1} -> 1, and (ii) the identification of lambda with the
% polar-coordinate volume density, i.e. lem:geodesic-polar-form(4).
\begin{proof}
\uses{lem:geodesic-polar-form,lem:radial-comparison,lem:radial-volume-element,lem:trace-riccati-comparison,def:ricci-curvature}
As in the proof of Theorem~\ref{thm:sectional-curvature-comparison},
the absence of conjugate points means that $\exp_p$ is
non-singular on a cone around the direction $v$, so
Lemma~\ref{lem:geodesic-polar-form} applies. Let $A(t)$ be the shape
operator along $\gamma$, viewed, in a parallel orthonormal frame
$E_1,\ldots,E_{n-1}$ of $\gamma'^\perp$, as a differentiable family
of symmetric endomorphisms of a fixed $(n-1)$-dimensional inner
product space. Taking the trace of the Riccati equation of
Lemma~\ref{lem:geodesic-polar-form}(2) over this frame (parallelism
makes differentiating traces commute with $\nabla_{\partial_r}$),
$$\left(\Tr A\right)'+\Tr\left(A^2\right)
+\sum_{i=1}^{n-1}K(E_i\wedge\gamma')=0,$$
and $\sum_iK(E_i\wedge\gamma')=\Ric(\gamma',\gamma')\ge-(n-1)k$ by
the expression of the Ricci tensor in an orthonormal basis, so
$$\left(\Tr A\right)'+\Tr\left(A^2\right)\ \le\ (n-1)k.$$
By Lemma~\ref{lem:geodesic-polar-form}(3), $A(t)=\frac1t\,
\mathrm{Id}+O(t)$, so $\Tr A(t)-(n-1)/t\to0$ as $t\to0^+$. By the
trace Riccati comparison
(Lemma~\ref{lem:trace-riccati-comparison}, with $m=n-1$),
$$\Tr A(t)\le(n-1)\frac{\sn_k'(t)}{\sn_k(t)}\qquad(0<t<r_0),$$
which is the first assertion. By
Lemma~\ref{lem:geodesic-polar-form}(4),
$$\partial_t\log\frac{\lambda(t)}{\sn_k^{n-1}(t)}
=\Tr(A)(t)-(n-1)\frac{\sn_k'(t)}{\sn_k(t)}\le0,$$
so $\lambda/\sn_k^{n-1}$ is non-increasing; its limit as $t\to0^+$
is $1$ because $\lambda(t)/t^{n-1}\to1$
(Lemma~\ref{lem:geodesic-polar-form}(4)) and
$\sn_k(t)/t\to1$. The final inequality follows from monotonicity
and the value of the limit.
\end{proof}

\begin{theorem}[Ricci curvature comparison theorem]
\label{thm:ricci-curvature-comparison}
\label{riccurvcomp}
\uses{def:ricci-curvature,def:sinh-comparison-function}
{\bf (Ricci curvature comparison)} Fix $k\ge 0$. Assume that $(M,
g)$ satisfies $\Ric \geq -(n-1)k$. Let $\gamma\colon [0,r_0)\to
M$ be a minimal geodesic of unit speed. Then for any $r<r_0$ at
$\gamma(r)$ we have
$$\sqrt{\det\, g(r,\theta)} \leq \sn^{n-1}_{k}(r)$$ and
$$\Tr(S)(r,\theta) \leq (n-1)\frac{\sn_{k}^{'}(r)}{\sn_{k}(r)}.$$
\end{theorem}

As in Theorem~\ref{thm:sectional-curvature-comparison}, the first
inequality is normalized against the round metric: it reads
$\sqrt{\det g(r,\theta)/\det\hat g(\theta)}\le\sn_k^{n-1}(r)$, and
takes the displayed form when the coordinates $\theta$ are chosen
so that $\hat g_{ij}=\delta_{ij}$ at the point in question.

\begin{proof}
\uses{prop:minimal-geodesic-no-conjugate,lem:volume-element-comparison,lem:geodesic-polar-form}
Since $\gamma$ is minimal on $[0,r_0)$, there are no conjugate
points along $\gamma|_{[0,r_0)}$
(Proposition~\ref{prop:minimal-geodesic-no-conjugate}, applied as in
the proof of Theorem~\ref{thm:sectional-curvature-comparison}), so
Lemma~\ref{lem:volume-element-comparison} applies to $\gamma$. Its
first assertion is precisely
$$\Tr(S)(r,\theta)=\Tr(A)(r)\le(n-1)\frac{\sn_k'(r)}{\sn_k(r)},$$
and its final assertion is
$\lambda(r)=\sqrt{\det g(r,\theta)/\det\hat g(\theta)}
\le\sn_k^{n-1}(r)$, which is the first displayed inequality in the
normalization explained above.
\end{proof}

Note that the pointwise inequality $\triangle f\le(n-1)/f$ away
from the cut locus, which underlies
Example~\ref{ex:distance-function-laplacian}, follows from this
theorem with $k=0$.

The comparison result in Theorem~\ref{thm:ricci-curvature-comparison} holds out to
every radius, a fact that will be used repeatedly in our arguments.
This result evolved over the period 1964-1980 and now is referred to
as the Bishop-Gromov inequality; see
Proposition 4.1 of \cite{CheegerGromovTaylor}.

Since the Bishop--Gromov inequality is stated for a ball with
compact closure in a manifold that need not be complete, we first
record that the cut-locus apparatus localizes to such a ball.

\begin{lemma}[Localized cut-locus apparatus]
\label{lem:localized-cut-locus}
\uses{def:exponential-map,def:geodesic,def:conjugate-point}
Let $(M,g)$ be a Riemannian manifold, not assumed complete, let
$p\in M$ and $R>0$, and suppose that $B(p,R)$ has compact closure
in $M$. Then:
\begin{enumerate}
\item every unit-speed geodesic from $p$ extends to every length
less than $R$; in particular
$B(0,R)\subset T_pM$ lies in the maximal domain of $\exp_p$;
\item every $q\in B(p,R)$ is joined to $p$ by a minimizing geodesic
of length $d(p,q)$;
\item the set $U_p^{<R}$, consisting of $0$ together with all $v\in
T_pM$ with $0<|v|<R$ such that $\gamma_v$ is the unique minimizing
geodesic from $p$ to $\exp_p(v)$ and $\exp_p$ is a local
diffeomorphism at $v$, is open and star-shaped about $0$: for each
unit vector $v$ there is $c^R(v)\in(0,R]$ with $\{0<t<R\,|\,tv\in
U_p^{<R}\}=(0,c^R(v))$; moreover $d(p,\exp_p(w))=|w|$ for $w\in
U_p^{<R}$, and along each ray the geodesic $\gamma_v$ is minimizing
and free of conjugate points on $[0,c^R(v))$;
\item for every $r<R$ the set $B(p,r)\setminus
\exp_p\left(U_p^{<R}\cap B(0,r)\right)$ has measure zero, and
$\exp_p$ restricted to $U_p^{<R}\cap B(0,r)$ is a
diffeomorphism onto its image in $B(p,r)$.
\end{enumerate}
\end{lemma}

\begin{proof}
\uses{prop:minimal-geodesic-no-conjugate,def:conjugate-point,lem:geodesic-polar-form,lem:normal-ball,lem:exponential-differential-jacobi,lem:minimizing-geodesic-exists}
(1) Let $\gamma$ be a unit-speed geodesic from $p$, defined on a
maximal interval $[0,T)$ with $T<R$. Then
$\gamma([0,T))\subset \overline{B(p,T)}$, a closed subset of the
compact set $\overline{B(p,R')}$ for $T<R'<R$. Since $|\gamma'|
\equiv1$, $\gamma$ is uniformly continuous and converges as $t\to
T$ to a point of this compact set; by the local existence and
uniqueness theorem for the geodesic equation at that point,
$\gamma$ extends beyond $T$, contradicting maximality. Hence
$T\ge R$, so $\gamma$ is defined on $[0,t]$ for every $t<R$.

(2) Let $q\in B(p,R)$, $d=d(p,q)<R$, and fix $d<R'<R$. By part
(1), every unit-speed geodesic from $p$ is defined on $[0,R']$;
hence Lemma~\ref{lem:minimizing-geodesic-exists}, applied with
$D=R'$, provides a geodesic from $p$ to $q$ of length $d(p,q)$.

(3) First, for every $w\in U_p^{<R}$ the geodesic $\gamma_w$ is
minimizing by condition (i), so
$$d\left(p,\exp_p(w)\right)=L(\gamma_w)=|w|.$$

\emph{$U_p^{<R}$ is a neighborhood of $0$.} Let $\delta>0$ be the
radius provided by Lemma~\ref{lem:normal-ball} for $K=\{p\}$, and
let $0<|w|<\min(\delta,R)$. Then $\exp_p$ is non-singular at $w$
(being a diffeomorphism on $B(0,\delta)$), which is condition
(ii); and any minimizing geodesic from $p$ to $\exp_p(w)$,
parameterized by $[0,1]$, is a piecewise smooth curve of length
$d(p,\exp_p(w))=|w|$, hence by
Lemma~\ref{lem:normal-ball}(3) a monotone reparameterization of
the radial segment; having constant speed, it equals $\gamma_w$.
This is condition (i). Hence
$B(0,\min(\delta,R))\subseteq U_p^{<R}$.

\emph{$U_p^{<R}$ is star-shaped.} Let $tv\in U_p^{<R}$ with
$|v|=1$, $0<t<R$, and let $0<s<t$. By condition (i),
$\gamma_{tv}\colon[0,1]\to M$ is the unique minimizing geodesic
from $p$ to its endpoint. By
Proposition~\ref{prop:minimal-geodesic-no-conjugate}, its
restriction to $[0,s/t]$ --- which is $\gamma_{sv}$ after
reparameterization --- is the unique minimizing geodesic between
its endpoints, so condition (i) holds at $sv$; and there are no
conjugate points along $\gamma_{tv}$ at any parameter in $[0,1)$,
so $\exp_p$ is a local diffeomorphism at $sv$ by
Lemma~\ref{lem:exponential-differential-jacobi}, which is
condition (ii). Hence $sv\in U_p^{<R}$, and, setting
$c^R(v)=\sup\{t\in(0,R)\,|\,tv\in U_p^{<R}\}\in(0,R]$, we get
$\{0<t<R\,|\,tv\in U_p^{<R}\}=(0,c^R(v))$ provided $c^R(v)v\notin
U_p^{<R}$ when $c^R(v)<R$; and indeed the set $\{t\,|\,tv\in
U_p^{<R}\}$ is open in $(0,R)$ once we prove $U_p^{<R}$ open,
which we do now.

\emph{$U_p^{<R}$ is open.} Let $v\in U_p^{<R}$; the case $v=0$ is
covered by the neighborhood of $0$ above, so assume $v\not=0$ and
fix $|v|<R'<R$. Suppose, for a contradiction, that $v_j\to v$ with
$0<|v_j|<R'$ and $v_j\notin U_p^{<R}$ for all $j$. The set where
$d(\exp_p)$ is singular is closed and misses $v$ (condition
(ii)), so for large $j$ condition (ii) holds at $v_j$; hence
condition (i) fails at $v_j$: writing $q_j=\exp_p(v_j)$, the
geodesic $\gamma_{v_j}$ is not the unique minimizing geodesic from
$p$ to $q_j$. Note $d(p,q_j)\le|v_j|<R'$, so by part (2) there is
a minimizing geodesic from $p$ to $q_j$; parameterized by $[0,1]$
it is $\gamma_{w_j}$ with $|w_j|=d(p,q_j)\le|v_j|$, and we may
choose $w_j\not=v_j$: if $\gamma_{v_j}$ is not minimizing then
$|w_j|<|v_j|$ forces $w_j\not=v_j$, while if $\gamma_{v_j}$ is
minimizing but not unique we take for $\gamma_{w_j}$ a minimizing
geodesic different from $\gamma_{v_j}$ (two distinct geodesics
parameterized by $[0,1]$ have distinct initial velocities, by
uniqueness for the geodesic equation). The vectors $w_j$ lie in
the compact ball $\overline{B(0,R')}\subset T_pM$; passing to a
subsequence, $w_j\to w$. By continuity
$\exp_p(w)=\lim q_j=\exp_p(v)$ and
$|w|=\lim d(p,q_j)=d\left(p,\exp_p(v)\right)=|v|$, so
$\gamma_w$ is a minimizing geodesic from $p$ to $\exp_p(v)$ of
length $|v|$; by the uniqueness in condition (i) at $v$,
$w=v$. Thus $v_j\to v$ and $w_j\to v$ with
$v_j\not=w_j$ and $\exp_p(v_j)=\exp_p(w_j)$ for all $j$,
contradicting the injectivity of $\exp_p$ on a neighborhood of
$v$ furnished by condition (ii). Hence $U_p^{<R}$ is open.

\emph{No conjugate points along the rays.} Let $|v|=1$ and
$t<c^R(v)$; pick $t<t'<c^R(v)$. Then $t'v\in U_p^{<R}$, so by
Proposition~\ref{prop:minimal-geodesic-no-conjugate} the geodesic
$\gamma_v$ is minimizing on $[0,t']\supset[0,t]$ and has no
conjugate point at any parameter in $[0,t')\supseteq[0,t]$.

(4) Injectivity of $\exp_p$ on $U_p^{<R}$: if $v,w\in U_p^{<R}$
satisfy $\exp_p(v)=\exp_p(w)=q$, then by condition (i) each
of $\gamma_v,\gamma_w$ is \emph{the} unique minimizing geodesic
from $p$ to $q$, so $\gamma_v=\gamma_w$, and differentiating at
$t=0$ gives $v=w$. With the local diffeomorphism condition (ii)
this makes $\exp_p|_{U_p^{<R}\cap B(0,r)}$ an injective local
diffeomorphism, hence a diffeomorphism onto its (open)
image. For the measure-zero claim, let $q\in B(p,r)$ be outside the
image. By (2) there is a minimizing geodesic $\gamma_w$,
$|w|=d(p,q)<r$, from $p$ to $q$; since $w\notin U_p^{<R}$, either
$\exp_p$ is singular at $w$, so that $q$ is a critical value of
$\exp_p|_{B(0,r)}$, or there are two distinct minimizing
geodesics from $p$ to $q$. The set of critical values has measure
zero by Sard's theorem.

And at a point $q$ with two distinct minimizing geodesics, the
distance function $d(p,\cdot)$ is not differentiable. To see this,
apply Lemma~\ref{lem:normal-ball} to the compact set
$K=\overline{B(p,r)}$ (a closed subset of the compact
$\overline{B(p,R)}$), obtaining a radius $\delta>0$ valid at every
point of $K$. Let $\gamma\colon[0,1]\to M$ be a minimizing
geodesic from $p$ to $q$, of length $\ell=d(p,q)$, and for
$0<\epsilon<\min(\delta/\ell,1)/2$ set
$y_\epsilon=\gamma(1-\epsilon)\in K$. Then
$d(y_\epsilon,q)=\epsilon\ell<\delta$ (subsegments of minimizing
geodesics are minimizing), so $q$ lies in the metric ball
$B(y_\epsilon,\delta)$, on which, by
Lemma~\ref{lem:normal-ball}(1),(2), the function
$d(y_\epsilon,\cdot)=|\exp_{y_\epsilon}^{-1}(\cdot)|$ is smooth
away from the center $y_\epsilon\not=q$. The function
$$u_\epsilon(x):=d\left(p,y_\epsilon\right)
+d\left(y_\epsilon,x\right)$$
is therefore smooth near $q$; it is an upper support function
for $d(p,\cdot)$ at $q$: $u_\epsilon\ge d(p,\cdot)$ by the
triangle inequality, with equality at $q$ since $\gamma$ realizes
both distances. Its differential at $q$ is the differential of
$d(y_\epsilon,\cdot)$, which by
Lemma~\ref{lem:geodesic-polar-form}(2) (applied to the polar
coordinates of $\exp_{y_\epsilon}$ on $B(0,\delta)$, where
$\exp_{y_\epsilon}$ is non-singular) is the unit radial covector
at $q$; and the radial segment of $B(y_\epsilon,\delta)$ ending at
$q$ is the reparameterized $\gamma|_{[1-\epsilon,1]}$ (a
minimizing curve of length $d(y_\epsilon,q)$, hence radial by
Lemma~\ref{lem:normal-ball}(3)), so this covector is $\langle
T,\cdot\rangle$ with $T=\gamma'(1)/|\gamma'(1)|$ the unit terminal
velocity of $\gamma$. If $d(p,\cdot)$ were differentiable at $q$,
its differential would agree with that of each smooth upper
support function touching it at $q$ (the difference has a local
minimum at $q$), giving $\langle
T_1,\cdot\rangle=\langle T_2,\cdot\rangle$ for the terminal
velocities of the two minimizing geodesics; but $T_1\not=T_2$,
since two geodesics of the same length, the same terminal point,
and the same terminal velocity coincide by (backwards) uniqueness
for the geodesic equation --- a contradiction. Since $d(p,\cdot)$
is Lipschitz, by Rademacher's theorem it is differentiable almost
everywhere, so the set of such $q$ has measure zero as well.
\end{proof}

\begin{theorem}[Bishop-Gromov relative volume comparison]
\label{thm:bishop-gromov}
\label{BishopGromov}
\uses{def:ricci-curvature,thm:ricci-curvature-comparison}
(Relative Volume Comparison, Bishop-Gromov 1964-1980) Suppose $(M,
g)$ is a Riemannian manifold. Fix a point $p\in M$, and suppose that
$B(p,R)$ has compact closure in $M$. Suppose that for some $k\ge 0$
we have $\Ric \geq -(n-1)k$ on $B(p,R)$. Recall that $H^n_k$ is
the simply connected, complete manifold of constant curvature $-k$
and $q_k\in H^n_k$ is a point. Then
\[ \frac{\Vol\,B(p,r)}{\Vol\, B_{H^n_k}(q_k,r)} \]
is a non-increasing function of $r$ for $r<R$, whose limit as $r\rightarrow 0$
is 1. In particular, if the Ricci curvature of $(M,g)$ is $\ge 0$ on $B(p,R)$,
then $\Vol\,B(p,r)/r^n$ is a non-increasing function of $r$ for $r<R$.
\end{theorem}

\begin{proof}
\uses{lem:volume-element-comparison,lem:geodesic-polar-form,lem:localized-cut-locus,lem:model-polar-isometry,lem:normal-ball}
\emph{The volume of $B(p,r)$ in polar coordinates.} Let
$U:=U_p^{<R}\subset T_pM$ be the star-shaped open set of
Lemma~\ref{lem:localized-cut-locus} and, for $v$ in the unit sphere
$S\subset T_pM$, let $c(v):=c^R(v)\in(0,R]$ be its cut time, so
that $\{0<t<R\,|\,tv\in U\}=(0,c(v))$. By parts (3) and (4) of that
lemma, $d(p,\exp_p(w))=|w|$ for $w\in U$, so that
$$\exp_p\left(U\cap B(0,r)\right)
=\exp_p(U)\cap B(p,r)\qquad(r<R),$$
the complement $B(p,r)\setminus\exp_p(U\cap B(0,r))$ has
measure zero, and $\exp_p$ is a diffeomorphism of $U\cap
B(0,r)$ onto its image. Since $U$ is open and star-shaped, the cut
time $c$ is lower semicontinuous, so $U$ is a union of cones
$\{tw\,|\,w\in V,\ 0\le t<r_0\}$ over open sets $V$ of unit
vectors, on each of which $\exp_p$ is non-singular; the polar
volume element $\lambda(t,v)$ of
Lemma~\ref{lem:geodesic-polar-form}(4) is therefore defined at
every point of $U$. The change of variables formula and
integration in polar coordinates therefore give,
for $r<R$,
$$\Vol B(p,r)=\int_{S}\int_0^r
\mathbf{1}_{\{t<c(v)\}}\,\lambda(t,v)\,dt\,d\mathrm{vol}_{S}(v).$$

\emph{The model volume.} By
Lemma~\ref{lem:model-polar-isometry} applied to $H^n_k$ (complete,
simply connected, of constant curvature $-k$), the exponential map
of $H^n_k$ at $q_k$ is an isometry from $T_{q_k}H^n_k$ equipped
with the polar metric $dr^2+\sn_k^2(r)g_{S^{n-1}}$ onto $H^n_k$
(note that $s_{\lambda}=\sn_k$ for $\lambda=-k$). Hence
$$\Vol B_{H^n_k}(q_k,r)=\omega_{n-1}\int_0^r\sn_k^{n-1}(t)\,dt
=:V_k(r),$$
where $\omega_{n-1}$ is the volume of the unit $(n-1)$-sphere.

\emph{Monotonicity.} For $v\in S$ and $0<t<R$ define
$$h(t,v)=\mathbf{1}_{\{t<c(v)\}}\,
\frac{\lambda(t,v)}{\sn_k^{n-1}(t)},$$
so that $\Vol B(p,r)=\int_S\int_0^rh(t,v)\sn_k^{n-1}(t)\,dt\,dv$.
We claim each $h(\cdot,v)$ is non-increasing on $(0,R)$. For
$t<c(v)$ the geodesic $\gamma_v|_{[0,t]}$ lies in $B(p,R)$
(as $d(p,\gamma_v(s))=s$ by
Lemma~\ref{lem:localized-cut-locus}(3)), is free of conjugate
points (ibid.), and $\Ric\ge-(n-1)k$ along it; so by
Lemma~\ref{lem:volume-element-comparison} the ratio
$\lambda(\cdot,v)/\sn_k^{n-1}$ is non-increasing on
$(0,c(v))$ with limit $1$ at $0^+$; in particular $0\le
h\le1$. For $t\ge c(v)$ we have $h(t,v)=0$. Hence $h(\cdot,v)$ is
non-increasing, and therefore so is
$H(t):=\int_Sh(t,v)\,d\mathrm{vol}_S(v)$ (all functions here are
measurable and bounded: $U$ is open and $\lambda$ is continuous
on it). Now write, for $0<r_1<r_2<R$, with
$w(t)=\sn_k^{n-1}(t)$,
$$\Vol B(p,r_2)\,V_k(r_1)-\Vol B(p,r_1)\,V_k(r_2)
=\omega_{n-1}\iint_{[0,r_2]\times[0,r_1]}
\left(H(t)-H(s)\right)w(t)w(s)\,dt\,ds.$$
The contribution of the square $[0,r_1]\times[0,r_1]$ vanishes,
the integrand being antisymmetric in $(t,s)$; on the remaining
region $t\in(r_1,r_2]$, $s\in[0,r_1]$ we have $t>s$, hence
$H(t)\le H(s)$ and the integrand is non-positive. Therefore
$$\frac{\Vol B(p,r_2)}{V_k(r_2)}\le\frac{\Vol B(p,r_1)}{V_k(r_1)},$$
which is the asserted monotonicity.

\emph{The limit at $r=0$.} From $h\le1$ we get $\Vol B(p,r)\le
V_k(r)$ for all $r<R$. Conversely, by
Lemma~\ref{lem:normal-ball} (applied with $K=\{p\}$) there is
$r_0>0$ such that $\exp_p$ is a diffeomorphism of $B(0,r_0)$ onto
the metric ball $B(p,r_0)$, with $\exp_p(B(0,r))=B(p,r)$ for
$r\le r_0$. For $r<r_0$ the change of
variables formula gives directly
$\Vol B(p,r)=\int_S\int_0^r\lambda(t,v)\,dt\,dv$, and by
Lemma~\ref{lem:geodesic-polar-form}(4) (with compactness of $S$
providing uniformity) $\lambda(t,v)/t^{n-1}\to1$ uniformly in $v$
as $t\to0$; since also $\sn_k(t)/t\to1$, for every $\epsilon>0$ we
have $\lambda(t,v)\ge(1-\epsilon)\sn_k^{n-1}(t)$ for all
sufficiently small $t$ and all $v$, whence $\Vol B(p,r)\ge
(1-\epsilon)V_k(r)$ for small $r$. Thus $\Vol B(p,r)/V_k(r)\to1$
as $r\to0$.

Finally, if $\Ric\ge0$ on $B(p,R)$ we may take $k=0$, for which
$V_0(r)=\omega_{n-1}r^n/n$; the monotonicity of $\Vol
B(p,r)/V_0(r)$ is then that of $\Vol B(p,r)/r^n$ up to the constant
factor $\omega_{n-1}/n$.
\end{proof}

\section{Local volume and the injectivity radius}

As the following results show, in the presence of bounded curvature the volume
of a ball $B(p,r)$ in $M$ is bounded away from zero if and only if the
injectivity radius of $M$ at $p$ is bounded away from zero.
The arguments in this section repeatedly rescale the metric by a
constant; the effect of such a rescaling on the geometric
quantities in play is recorded in the following lemma.

\begin{lemma}[Rescaling the metric]
\label{lem:metric-rescaling}
\uses{def:riemannian-metric,def:sectional-curvature,def:curvature-operator-norm,def:ricci-curvature,def:injectivity-radius,def:exponential-map}
Let $(M,g)$ be a Riemannian manifold and let $c>0$ be a constant;
set $g'=c\,g$. Then:
\begin{enumerate}
\item $g$ and $g'$ have the same Christoffel symbols, the same
geodesics (as parameterized curves), the same exponential maps, and
the same $(1,3)$ curvature tensor; the $(0,4)$ tensor satisfies
$R'_{ijkl}=c\,R_{ijkl}$;
\item $K'(P)=K(P)/c$ for every $2$-plane $P$, and
$|\Rm'(x)|=|\Rm(x)|/c$ for the curvature operator norm of
Definition~\ref{def:curvature-operator-norm};
\item $\Ric'=\Ric$ as $(0,2)$-tensors; in particular
$\Ric\ge(n-1)k\,g$ holds if and only if
$\Ric'\ge(n-1)(k/c)\,g'$;
\item $d'=\sqrt{c}\,d$; hence $B_{g'}(p,\sqrt{c}\,r)=B_g(p,r)$,
and $(M,g')$ is complete if and only if $(M,g)$ is;
\item $\Vol_{g'}=c^{n/2}\Vol_g$;
\item $\inj_{(M,g')}(p)=\sqrt{c}\,\inj_{(M,g)}(p)$.
\end{enumerate}
\end{lemma}

\begin{proof}
\uses{thm:levi-civita-connection,def:riemann-curvature-tensor,def:geodesic,def:sectional-curvature,def:curvature-operator,def:curvature-operator-norm,def:ricci-curvature,def:injectivity-radius,def:exponential-map}
Throughout, primed quantities are those computed with respect to
$g'=c\,g$. In local coordinates $g'_{ij}=c\,g_{ij}$, hence
$g'^{kl}=c^{-1}g^{kl}$, and, since $c$ is constant,
$\partial_a g'_{ij}=c\,\partial_a g_{ij}$.

(1) In Equation~(\ref{Gamma}) the factor $c$ coming from the
derivatives of the metric cancels against the factor $c^{-1}$
coming from the inverse metric:
$$\Gamma'^{k}_{ij}
=\frac{1}{2}g'^{kl}(\partial_i g'_{lj}+\partial_j g'_{il}
-\partial_l g'_{ij})
=\frac{1}{2}c^{-1}g^{kl}\cdot c\,(\partial_i g_{lj}
+\partial_j g_{il}-\partial_l g_{ij})
=\Gamma^{k}_{ij}.$$
Since the Levi-Civita connection is determined on each chart by its
Christoffel symbols, $\nabla'=\nabla$. The geodesic equation
$\nabla_{\dot\gamma}\dot\gamma=0$ involves only the connection, so
a curve is a geodesic for $g'$ if and only if it is a geodesic for
$g$, with the same parameterization. For $v\in T_pM$ the initial
condition $(\gamma(0),\dot\gamma(0))=(p,v)$ does not involve the
metric, so the geodesic $\gamma_v$ is the same curve for both
metrics, with the same maximal interval of definition; hence the
maximal domains $O_p\subset T_pM$ of the two exponential maps
coincide and $\exp'_p=\exp_p$ as maps. The $(1,3)$ curvature
tensor ${\mathcal R}(X,Y)Z=\nabla_X\nabla_YZ-\nabla_Y\nabla_XZ
-\nabla_{[X,Y]}Z$ is built from the connection alone, so
${\mathcal R}'={\mathcal R}$ as $(1,3)$-tensors; in local
coordinates ${{R'_{ij}}^l}_k={{R_{ij}}^l}_k$. Lowering the index
with $g'$ gives
$$R'_{ijkl}=g'_{ks}{{R_{ij}}^s}_l=c\,g_{ks}{{R_{ij}}^s}_l
=c\,R_{ijkl},$$
or invariantly ${\mathcal R}'(X,Y,Z,W)=g'({\mathcal R}(X,Y)W,Z)
=c\,{\mathcal R}(X,Y,Z,W)$ for the $(0,4)$-tensors.

(2) Let $P\subset T_xM$ be a $2$-plane and let $\{X,Y\}$ be a
$g$-orthonormal basis of $P$. Since $g'(X/\sqrt{c},X/\sqrt{c})
=c^{-1}g'(X,X)=g(X,X)=1$, and similarly for the other products,
$\{X/\sqrt{c},Y/\sqrt{c}\}$ is a $g'$-orthonormal basis of $P$.
By multilinearity of the $(0,4)$-tensor and (1),
$$K'(P)={\mathcal R}'\!\left(\tfrac{X}{\sqrt{c}},
\tfrac{Y}{\sqrt{c}},\tfrac{X}{\sqrt{c}},\tfrac{Y}{\sqrt{c}}\right)
=\frac{1}{c^2}\,{\mathcal R}'(X,Y,X,Y)
=\frac{c}{c^2}\,{\mathcal R}(X,Y,X,Y)=\frac{K(P)}{c}.$$
For the operator norm, the inner product of
Definition~\ref{def:curvature-operator-norm} on $\wedge^2T_xM$
computed with $g'$ satisfies
$$\langle X\wedge Y,Z\wedge W\rangle'
=g'(X,Z)g'(Y,W)-g'(X,W)g'(Y,Z)
=c^2\langle X\wedge Y,Z\wedge W\rangle,$$
so by bilinearity $\langle\cdot,\cdot\rangle'
=c^2\langle\cdot,\cdot\rangle$ on all of $\wedge^2T_xM$. The
bilinear form of Definition~\ref{def:curvature-operator} satisfies
$$\Rm'(\varphi,\psi)=R'_{ijkl}\varphi^{ij}\psi^{kl}
=c\,R_{ijkl}\varphi^{ij}\psi^{kl}=c\,\Rm(\varphi,\psi),$$
the components $\varphi^{ij},\psi^{kl}$ of sections of
$\wedge^2TM$ being independent of the metric. Let $T$ (resp.\ $T'$)
be the symmetric operator on $\wedge^2T_xM$ associated to $\Rm$ by
$\langle\cdot,\cdot\rangle$ (resp.\ to $\Rm'$ by
$\langle\cdot,\cdot\rangle'$). For all
$\varphi,\psi\in\wedge^2T_xM$,
$$\langle T'\varphi,\psi\rangle'=\Rm'(\varphi,\psi)
=c\,\Rm(\varphi,\psi)=c\,\langle T\varphi,\psi\rangle
=c^{-1}\langle T\varphi,\psi\rangle',$$
and since $\langle\cdot,\cdot\rangle'$ is non-degenerate this
forces $T'=c^{-1}T$ (which is indeed symmetric for
$\langle\cdot,\cdot\rangle'$, the two inner products being
proportional). Thus $T'$ has the same eigenvectors as $T$, with
each eigenvalue divided by $c$; consequently, for $K\ge0$, the
eigenvalues of $T'$ lie in $[-K/c,K/c]$ if and only if those of
$T$ lie in $[-K,K]$. Taking the least such bounds gives
$|\Rm'(x)|=|\Rm(x)|/c$.

(3) By the local-coordinate definition of the Ricci tensor and (1),
$$\Ric'(X,Y)=g'^{kl}R'(X,\partial_k,Y,\partial_l)
=c^{-1}g^{kl}\cdot c\,R(X,\partial_k,Y,\partial_l)=\Ric(X,Y),$$
so $\Ric'=\Ric$. Since $\Ric'=\Ric$ and
$(n-1)k\,g=(n-1)(k/c)\,g'$, for every tangent vector $X$ the
inequality $\Ric(X,X)\ge(n-1)k\,g(X,X)$ is literally the same
inequality as $\Ric'(X,X)\ge(n-1)(k/c)\,g'(X,X)$; hence one holds
for all $X$ if and only if the other does.

(4) For every smooth path $\gamma\colon[a,b]\to M$ the lengths with
respect to the two metrics satisfy
$$L'(\gamma)=\int_a^b\sqrt{g'(\dot\gamma(t),\dot\gamma(t))}\,dt
=\sqrt{c}\int_a^b\sqrt{g(\dot\gamma(t),\dot\gamma(t))}\,dt
=\sqrt{c}\,L(\gamma).$$
Taking the infimum over the family of smooth paths from $p$ to $q$,
which does not depend on the metric, gives
$d'(p,q)=\sqrt{c}\,d(p,q)$. Consequently $d'(p,q)<\sqrt{c}\,r$ if
and only if $d(p,q)<r$, i.e.,
$B_{g'}(p,\sqrt{c}\,r)=B_g(p,r)$. Moreover, since the two distance
functions differ by the constant positive factor $\sqrt{c}$, a
sequence in $M$ is Cauchy for $d'$ if and only if it is Cauchy for
$d$, and it converges for $d'$ if and only if it converges for
$d$ (to the same limit); hence $(M,d')$ is complete if and only if
$(M,d)$ is.

(5) In any local coordinate chart the Riemannian measure of $g$ is
$\sqrt{\det(g_{ij})}\,dx^1\cdots dx^n$. Since
$\det(g'_{ij})=\det(c\,g_{ij})=c^n\det(g_{ij})$, we have
$\sqrt{\det(g'_{ij})}=c^{n/2}\sqrt{\det(g_{ij})}$ on every chart,
with the same constant factor $c^{n/2}$ in all charts. Decomposing
a measurable set by a partition of unity subordinate to a covering
by charts and integrating, we conclude
$\Vol_{g'}(A)=c^{n/2}\Vol_g(A)$ for every measurable $A\subset M$,
i.e., $\Vol_{g'}=c^{n/2}\Vol_g$.

(6) For $v\in T_pM$ we have $\sqrt{g'(v,v)}=\sqrt{c}\sqrt{g(v,v)}$,
so the ball of radius $r$ about the origin in $(T_pM,g'_p)$ is the
ball of radius $r/\sqrt{c}$ in $(T_pM,g_p)$:
$$B_{g'|T_pM}(0,r)=B_{g|T_pM}(0,r/\sqrt{c}).$$
By (1) the two exponential maps at $p$ coincide, with the same
maximal domain $O_p$; hence the restriction of $\exp'_p$ to
$B_{g'|T_pM}(0,r)$ is defined and is a diffeomorphism into $M$ if
and only if the restriction of $\exp_p$ to
$B_{g|T_pM}(0,r/\sqrt{c})$ is. Therefore $r$ lies in the set of
radii whose supremum defines $\inj_{(M,g')}(p)$ if and only if
$r/\sqrt{c}$ lies in the corresponding set for $(M,g)$; the first
set is the image of the second under multiplication by the
constant $\sqrt{c}$, so the suprema satisfy
$\inj_{(M,g')}(p)=\sqrt{c}\,\inj_{(M,g)}(p)$.
\end{proof}


\begin{proposition}[Bounded curvature implies volume-injectivity relation]
\label{prop:injectivity-radius-volume}
\label{injvol}
\uses{def:injectivity-radius,def:curvature-operator-norm}
Fix an integer $n>0$. For every $\epsilon>0$ there is $\delta>0$
depending on $n$ and $\epsilon$ such that the following holds.
Suppose that $(M^n,g)$ is a complete connected Riemannian
manifold of dimension $n$ and that $p\in M$. Suppose that $|\Rm(x)|\le
r^{-2}$ for all $x\in B(p,r)$. If the injectivity radius of $M$ at
$p$ is at least $\epsilon r$, then $\Vol(B(p,r))\ge \delta
r^n$.
\end{proposition}

\begin{proof}
\uses{lem:rauch-lower,def:curvature-operator-norm,lem:geodesic-polar-form}
Suppose that $|\Rm(x)|\le r^{-2}$ for all $x\in B(p,r)$.
Replacing $g$ by $r^{-2} g$ (which multiplies all sectional
curvatures by $r^{2}$, distances by $r^{-1}$ and volumes by
$r^{-n}$, as in Step 1 of the proof of
Theorem~\ref{thm:uniformization}) allows us to assume that $r=1$.
Without loss of generality we can assume that $\epsilon\le 1$. The
map $\exp_p$ is a diffeomorphism on the ball $B(0,\epsilon)$ in
the tangent space, and by
Definition~\ref{def:curvature-operator-norm} the sectional
curvatures on $B(p,1)$ satisfy $K(P)\le 1$; note
$\epsilon\le1<\pi=D_1$. Hence Lemma~\ref{lem:rauch-lower} applies
on the cone over the whole unit sphere of $T_pM$ of radius
$\epsilon$, and gives $g_{ij}(t,\theta)\ge s_1(t)^2\hat
g_{ij}(\theta)$ with $s_1(t)=\sin t$. Since the determinant is
monotone on positive definite symmetric matrices, the volume
element of Lemma~\ref{lem:geodesic-polar-form}(4) satisfies
$\lambda(t,\theta)\ge\sin^{n-1}(t)$, and integrating in polar
coordinates over the diffeomorphic image,
$$\Vol B(p,\epsilon)=\int_{S^{n-1}}\int_0^\epsilon
\lambda(t,\theta)\,dt\,d\mathrm{vol}_{S^{n-1}}
\ge\omega_{n-1}\int_0^\epsilon\sin^{n-1}(t)\,dt,$$
which is the volume of the ball of radius $\epsilon$ in
the $n$-sphere of radius $1$. This gives a lower bound to the volume
of $B(p,\epsilon)$, and a fortiori to $B(p,1)$, in terms of $n$ and
$\epsilon$.
\end{proof}

 We shall normally work with volume, which behaves nicely under Ricci
flow, but in order to take limits we need to bound the injectivity
radius away from zero. Thus, the more important, indeed crucial,
result for our purposes is the converse to the previous proposition;
see Theorem 4.3, especially Inequality (4.22), on page 46 of
\cite{CheegerGromovTaylor}, or see Theorem 5.8 on page 96 of
\cite{CheegerEbin}.

Our proof follows Section 4 of \cite{CheegerGromovTaylor}
(Theorem 4.3 there, via their Lemmas 4.4--4.6), with all steps
carried out in full; the one point which is asserted without proof
there --- the strict convexity of squared-distance functions in the
ball with the pulled-back metric --- is established below
(Lemma~\ref{lem:pullback-unique-minimizers}) by a homotopy argument
special to the pulled-back ball. The mechanism is a count of
preimages of points under the exponential map: a short geodesic
loop at $p$ forces $p$ (and then every nearby point) to have many
preimages in a ball of the tangent space on which
$\exp_p$ is non-singular, and each preimage carries its own copy
of the local volume; an upper bound for the total volume upstairs
then forces the volume downstairs to be small.

\begin{lemma}[The pulled-back ball and lifts of short curves]
\label{lem:pullback-ball}
\uses{def:exponential-map,def:geodesic,lem:geodesic-polar-form,def:local-isometry}
\source{cheeger-gromov-taylor:page-0032}
Let $p\in M$, let $0<\rho_*\le\infty$, and suppose that
$\exp_p$ is defined and non-singular on
$B=B(0,\rho_*)\subset T_pM$. Let $\tilde g=\exp_p^*g$ on $B$.
Then:
\begin{enumerate}
\item $\exp_p\colon (B,\tilde g)\to M$ is a local isometry; the
rays $t\mapsto tv$ are unit-speed $\tilde g$-geodesics (for unit
$v$); every piecewise smooth curve $c$ in $B$ satisfies
$L_{\tilde g}(c)\ge\bigl||c(1)|-|c(0)|\bigr|$; and the
$\tilde g$-geodesics starting at $0$ are exactly the rays.
\item (Path lifting) For every continuous curve $c\colon[0,1]\to M$
of finite length $L(c)$ with $c(0)=\exp_p(y_0)$ for some
$y_0\in B$ with $|y_0|+L(c)<\rho_*$, there is a unique continuous
curve $\tilde c$ (the \emph{lift of $c$ at $y_0$}) with $\tilde
c(0)=y_0$ and $\exp_p\circ\tilde c=c$; it satisfies $|\tilde
c(t)|\le|y_0|+L(c|_{[0,t]})$. The lift of the reversed curve at
$\tilde c(1)$ is the reversal of $\tilde c$, and lifts of
concatenations are the concatenations of the successive lifts.
\item (Homotopy invariance) Let $H\colon[0,1]\times[0,1]\to M$ be
continuous, write $c_u=H(\cdot,u)$, and suppose that all $c_u$ have
the same endpoints $q_0=\exp_p(y_0)$ and $q_1$, and finite
lengths $L(c_u)\le\Lambda$ with $|y_0|+\Lambda<\rho_*$. Then the
lift endpoints $\tilde c_u(1)$ coincide for all $u$.
\item (Endpoint map and cancellation) For a continuous curve $c$
from $p$ of finite length $L(c)<\rho_*$ set
$E(c):=\tilde c(1)\in\exp_p^{-1}(c(1))$, the endpoint of the
lift of $c$ at $0$; write $\mathrm{rad}_y$, $y\in B$, for the curve
$t\mapsto\exp_p(ty)$, so that $E(\mathrm{rad}_y)=y$. Then, with
all length sums below assumed $<\rho_*$ and $\cup$ denoting
concatenation (first curve first):
\begin{enumerate}
\item[(a)] $c$ is homotopic to $\mathrm{rad}_{E(c)}$ with fixed
endpoints through curves of length $\le L(c)$;
\item[(b)] if $E(\theta_1\cup c)=E(\theta_2\cup c)$ for curves
$\theta_1,\theta_2$ from $p$ to $q'$ and $c$ from $q'$, then
$E(\theta_1)=E(\theta_2)$;
\item[(c)] $E(\theta\cup c)$ is the endpoint of the lift of $c$ at
$E(\theta)$; in particular it depends on $\theta$ only through
$E(\theta)$;
\item[(d)] if $\sigma$ is a curve from $p$ with
$E(\sigma)=y$, then $E(\theta\cup\sigma)
=E(\theta\cup\mathrm{rad}_y)$ for every loop $\theta$ at $p$.
\end{enumerate}
\item (Deck motions) For a loop $\theta$ at $p$ of length
$\ell_\theta$ and $a>0$ with $a+\ell_\theta<\rho_*$, define
$$\pi_\theta\colon B(0,a)\to B(0,a+\ell_\theta),\qquad
\pi_\theta(y):=E(\theta\cup\mathrm{rad}_y).$$
Then $\exp_p\circ\pi_\theta=\exp_p$; the map
$\pi_\theta$ is smooth with everywhere isometric differential, so
$L_{\tilde g}(\pi_\theta\circ\omega)=L_{\tilde g}(\omega)$ for
every curve $\omega$ in $B(0,a)$; moreover
$\pi_{\theta'}\circ\pi_\theta=\pi_{\theta'\cup\theta}$ and
$\pi_{\bar\theta}\circ\pi_\theta=\mathrm{id}$ on suitable balls
(with $\bar\theta$ the reversed loop), and
$\pi_\theta(E(\sigma))=E(\theta\cup\sigma)$ for curves $\sigma$
from $p$.
\end{enumerate}
\end{lemma}

\begin{proof}
\uses{lem:geodesic-polar-form,def:local-isometry,def:exponential-map,lem:local-isometry-exponential,lem:local-isometry-normal-balls,lem:normal-ball}
(1) $\exp_p^*g=\tilde g$ says precisely that $d\exp_p$ is a
linear isometry at each point of $B$, so $\exp_p$ is a local
isometry. A local isometry carries geodesics to geodesics
(Lemma~\ref{lem:local-isometry-exponential}, first assertion),
and near each point it restricts to an isometry onto its image
whose inverse is again a local isometry
(Definition~\ref{def:local-isometry}); since being a geodesic is a
local property, a curve in $B$ is a $\tilde g$-geodesic if and
only if its image is a geodesic of $M$. In particular the rays
$t\mapsto tv$, whose images are the radial geodesics
$t\mapsto\exp_p(tv)$, are unit-speed $\tilde g$-geodesics. The
length estimate is
Lemma~\ref{lem:geodesic-polar-form}(1) (the polar form of $\tilde
g$): away from the origin $|c'|_{\tilde g}\ge|\frac{d}{dt}|c||$,
and if $c$ passes through the origin one applies this on the
complementary subintervals. A $\tilde g$-geodesic starting at $0$
has the same initial data as a ray, hence equals it as long as both
are defined.

(2) Uniqueness: if $\tilde c_1,\tilde c_2$ are lifts with $\tilde
c_1(0)=\tilde c_2(0)$, the set $\{t\,|\,\tilde c_1(t)=\tilde
c_2(t)\}$ is closed, non-empty, and open (near an agreement point,
both lifts lie in a neighborhood on which $\exp_p$ is
injective, and they project to the same curve), hence everything.
Existence: fix $\rho$ with $|y_0|+L(c)<\rho<\rho_*$ and let
$K_\rho=\{|x|\le\rho\}$, a compact subset of $B$; by
Lemma~\ref{lem:local-isometry-normal-balls}(3), applied to the
local isometry $\exp_p\colon(B,\tilde g)\to M$ and the compact
set $K_\rho$, every continuous partial lift of $c$ with image in
$K_\rho$ has the same length as the corresponding piece of $c$.
Let $T\in(0,1]$ be maximal such that a lift with the
stated radial bound exists on $[0,t]$ for every $t<T$ (a local lift
exists near $0$ through a local inverse of $\exp_p$; and the
radial bound propagates: while the lift stays in $K_\rho$ its
$\tilde g$-length equals that of $c$, so by (1)
$|\tilde c(t)|\le
|y_0|+L_{\tilde g}(\tilde c|_{[0,t]})=|y_0|+L(c|_{[0,t]})
\le|y_0|+L(c)<\rho$, which keeps it strictly inside $K_\rho$ and
extends the bound by continuity). By uniqueness these lifts agree
on common domains and
define a lift on $[0,T)$; as $t\to T$, $\tilde c(t)$ has a limit
point $x^*$ in the compact set $\{|x|\le|y_0|+L(c)\}\subset B$
(by the radial bound), and in fact converges to $x^*$: on a
neighborhood $V\ni x^*$ carried homeomorphically by $\exp_p$
onto a neighborhood of $c(T)$, the lift must eventually agree with
(local inverse)$\circ\,c$, which extends continuously to $T$. The
same local inverse extends the lift beyond $T$ if $T<1$,
contradicting maximality; so $T=1$ and the lift extends to $[0,1]$.
The statements about reversals and concatenations follow from
uniqueness.

(3) Fix $\rho$ with $|y_0|+\Lambda<\rho<\rho_*$. All lifts
$\tilde c_u$ (which exist by (2)) have image in the compact set
$K^*=\{|x|\le|y_0|+\Lambda\}$, contained in the interior of the
compact set $K_\rho=\{|x|\le\rho\}\subset B$. Apply
Lemma~\ref{lem:local-isometry-normal-balls}(1),(2) to
$\exp_p\colon(B,\tilde g)\to M$ and $K_\rho$: there is
$\delta>0$ such that for every $x\in K_\rho$, $\exp_p$ maps the
metric ball $B_{\tilde g}(x,\delta)$ diffeomorphically onto
$B(\exp_p(x),\delta)$, and
$d_{\tilde g}(a,b)=d\left(\exp_p(a),\exp_p(b)\right)$ for
$a,b\in B_{\tilde g}(x,\delta/4)$. For $x\in K_\rho$ write
$\Phi_x=\left(\exp_p|_{B_{\tilde g}(x,\delta)}\right)^{-1}$, the
\emph{sheet at $x$}. By the naturality identity of
Lemma~\ref{lem:local-isometry-normal-balls}(1),
$\exp_p|_{B_{\tilde g}(x,\delta)}
=\exp_{\exp_p(x)}\circ\,d(\exp_p)_x\circ\widetilde\exp_x^{-1}$,
where $\widetilde\exp_x$ is the exponential map of $(B,\tilde g)$
at $x$; since by Lemma~\ref{lem:normal-ball}(1) (which underlies
the cited lemma) the two exponential maps here carry $s$-balls of
the tangent spaces onto metric $s$-balls for every $0<s\le\delta$,
and $d(\exp_p)_x$ is a linear isometry, the sheet $\Phi_x$ maps
$B(\exp_p(x),s)$ onto $B_{\tilde g}(x,s)$ for every
$0<s\le\delta$.

By uniform continuity of $H$ on the compact square, choose a
partition $0=t_0<\cdots<t_m=1$ and $\eta>0$ such that
$$d\left(H(t,w),H(t_j,u)\right)<\delta/8\qquad\text{for }
t\in[t_j,t_{j+1}],\ |w-u|\le\eta,\ w,u\in[0,1].$$
Fix $u$ and $u'$ with $|u'-u|\le\eta$. We show by induction on
$j$ that
$$\tilde c_{u'}(t_j)\in B_{\tilde g}\left(\tilde
c_u(t_j),\delta/4\right)\qquad(0\le j\le m).$$
For $j=0$ this holds since both lifts start at $y_0$. Assume it
at $j$, and write $x=\tilde c_u(t_j)\in K^*\subset K_\rho$. For
$t\in[t_j,t_{j+1}]$ both $c_u(t)$ and $c_{u'}(t)$ lie in
$B(c_u(t_j),\delta/8)\subset B(\exp_p(x),\delta)$. The curves
$t\mapsto\Phi_x\left(c_u(t)\right)$ and
$t\mapsto\Phi_x\left(c_{u'}(t)\right)$ on $[t_j,t_{j+1}]$ are
continuous lifts of $c_u$ and $c_{u'}$; they agree with $\tilde
c_u$ and $\tilde c_{u'}$ at $t_j$ ($\Phi_x$ inverts $\exp_p$ on
$B_{\tilde g}(x,\delta)$, which contains $\tilde c_u(t_j)=x$ and,
by the inductive hypothesis, $\tilde c_{u'}(t_j)$), so by the
uniqueness of lifts (proved in (2), and valid for any continuous
lift) they agree with $\tilde c_u$, $\tilde c_{u'}$ on all of
$[t_j,t_{j+1}]$. In particular
$$\tilde c_u(t_{j+1})=\Phi_x\left(c_u(t_{j+1})\right),\qquad
\tilde c_{u'}(t_{j+1})=\Phi_x\left(c_{u'}(t_{j+1})\right),$$
and both lie in $B_{\tilde g}(x,\delta/4)$, since $c_u(t_{j+1})$
and $c_{u'}(t_{j+1})$ lie in $B(\exp_p(x),\delta/8)$ and $\Phi_x$
maps $s$-balls to $s$-balls. By
Lemma~\ref{lem:local-isometry-normal-balls}(2),
$$d_{\tilde g}\left(\tilde c_u(t_{j+1}),\tilde
c_{u'}(t_{j+1})\right)
=d\left(c_u(t_{j+1}),c_{u'}(t_{j+1})\right)<\delta/4,$$
which is the inductive claim at $j+1$. At $j=m$: $\tilde
c_{u'}(1)$ lies in $B_{\tilde g}(\tilde c_u(1),\delta/4)$ and
both lift endpoints project to $q_1$; since $\exp_p$ is
injective on $B_{\tilde g}(\tilde c_u(1),\delta)$, we get $\tilde
c_{u'}(1)=\tilde c_u(1)$. Thus
$u\mapsto\tilde c_u(1)$ is locally constant, hence constant. (The
same induction shows, for a continuous family of curves with
uniformly bounded lengths but \emph{varying} endpoints, that the
lift endpoint $\tilde c_u(1)=\Phi_{\tilde c_{u_0}(1)}
\left(c_u(1)\right)$ for $u$ near $u_0$ depends continuously on
$u$; this is used in the proof of (5) below.)

(4)(a) Let $\tilde c$ be the lift of $c$ at $0$ and define, for
$u\in[0,1]$, the curve $\tilde c_u$ which follows the radial
segment from $0$ to $\tilde c(u)$ and then $\tilde c|_{[u,1]}$;
each $\tilde c_u$ runs from $0$ to $\tilde c(1)$, depends
continuously on $u$, and
$$L_{\tilde g}(\tilde c_u)=|\tilde c(u)|
+L_{\tilde g}(\tilde c|_{[u,1]})
\le L_{\tilde g}(\tilde c|_{[0,u]})+L_{\tilde g}(\tilde
c|_{[u,1]})=L(c)$$
by (1). The curves $c_u:=\exp_p\circ\tilde c_u$ give the
required homotopy from $c$ ($u=0$) to $\mathrm{rad}_{E(c)}$
($u=1$).

(b) By (2), the lift of $\theta_i\cup c\cup\bar c$ at $0$ is the
lift of $\theta_i\cup c$ followed by the reversal of the lift of
$c$ at $E(\theta_i)$; hence
$E(\theta_i\cup c\cup\bar c)=E(\theta_i)$. On the other hand
$\theta_1\cup c\cup\bar c$ and $\theta_2\cup c\cup\bar c$ have
lifts whose final thirds start at the common point
$E(\theta_1\cup c)=E(\theta_2\cup c)$ and lift the common curve
$\bar c$; so they have the same endpoint, i.e.,
$E(\theta_1)=E(\theta_2)$.

(c) is immediate from the concatenation property in (2). For (d),
by (a) the curve $\sigma$ is fixed-endpoint homotopic to
$\mathrm{rad}_y$ through curves of length $\le L(\sigma)$;
concatenating with the constant homotopy of $\theta$ gives a
fixed-endpoint homotopy from $\theta\cup\sigma$ to
$\theta\cup\mathrm{rad}_y$ through curves of bounded length, and
(3) applies.

(5) $\exp_p(\pi_\theta(y))=\exp_p(y)$ holds since
$\pi_\theta(y)$ is a preimage of the endpoint
$\exp_p(y)$ of $\theta\cup\mathrm{rad}_y$. Continuity of
$\pi_\theta$: the curves $\theta\cup\mathrm{rad}_y$ depend
continuously on $y$ with uniformly bounded lengths
$\le\ell_\theta+a$, and the continuity remark at the end of the
proof of (3) shows the lift endpoint depends continuously on
$y$. Smoothness and the
isometric differential: near a fixed $y_1$, let $V$ be a
neighborhood of $\pi_\theta(y_1)$ on which $\exp_p$ is
invertible onto its image; by continuity
$\pi_\theta(y)\in V$ for $y$ near $y_1$, and then
$\pi_\theta=(\exp_p|_V)^{-1}\circ\exp_p$ there, a
composition of a local isometry and the inverse of one. The
composition rule: by (4c) and (4d),
$$\pi_{\theta'}(\pi_\theta(y))
=E\left(\theta'\cup\mathrm{rad}_{E(\theta\cup\mathrm{rad}_y)}\right)
=E\left(\theta'\cup\theta\cup\mathrm{rad}_y\right)
=\pi_{\theta'\cup\theta}(y),$$
and $\pi_{\bar\theta\cup\theta}=\mathrm{id}$ because
$\bar\theta\cup\theta\cup\mathrm{rad}_y$ retracts to
$\mathrm{rad}_y$ (reverse--retrace the initial loop, a
fixed-endpoint homotopy through curves of length $\le
2\ell_\theta+|y|$, then apply (3)). Finally
$\pi_\theta(E(\sigma))=E(\theta\cup\sigma)$ is (4d) combined with
(4c).
\end{proof}

\begin{lemma}[Unique minimizers and convexity in the pulled-back ball]
\label{lem:pullback-unique-minimizers}
\uses{lem:pullback-ball,def:sectional-curvature,lem:rauch-lower,lem:conjugate-sturm}
In the setting of Lemma~\ref{lem:pullback-ball} with $\rho_*=1$,
suppose in addition that $|K(P)|\le1$ for every $2$-plane tangent
to $(B,\tilde g)$. Write $\widetilde\exp_z$ for the exponential
map of $(B,\tilde g)$ at $z$. Then for every $z\in
\overline{B(0,1/32)}$:
\begin{enumerate}
\item $\widetilde\exp_z$ is defined and non-singular on
$B(0,1/2)\subset T_z B$, and injective on $B(0,1/8)$;
\item every $y\in B$ with $d_{\tilde g}(z,y)\le1/10$ lies in
$\Omega_z:=\widetilde\exp_z(B(0,1/8))$; in particular any two
points $z,y\in\overline{B(0,1/32)}$ are joined by a
$\tilde g$-minimizing geodesic, of length $\le1/8$, unique among
geodesics from $z$ to $y$ of length $\le 1/8$; and on
$\Omega_z$ the distance
$d_z:=d_{\tilde g}(z,\cdot)$ satisfies
$d_z=|\widetilde\exp_z^{-1}|$ and is smooth away from $z$;
\item $\Hess_{\tilde g}(d_z^2)\ge\tilde g$ on $\Omega_z$.
\end{enumerate}
(Here $d_{\tilde g}$ denotes the infimum of $\tilde g$-lengths of
curves in $B$.)
\end{lemma}

\begin{proof}
\uses{lem:pullback-ball,lem:rauch-lower,lem:conjugate-sturm,lem:geodesic-polar-form,lem:scalar-riccati-comparison,lem:localized-cut-locus,lem:exponential-differential-jacobi}
\emph{Two-sided comparison with the Euclidean metric.} On the cone
over the full unit sphere with $r_0=1/2$,
Lemma~\ref{lem:geodesic-polar-form} applies to $\tilde g$ around
$0$ (non-singularity is our hypothesis). The hypothesis $K\ge-1$
gives, exactly as in the shape-operator estimate in the proof of
Theorem~\ref{thm:sectional-curvature-comparison} (the parallel
eigenvector trick together with
Lemma~\ref{lem:scalar-riccati-comparison}), the bound $A(r)\le
\frac{\sn_1'(r)}{\sn_1(r)}\mathrm{Id}$ and hence, integrating
$\partial_rg_{ij}=2S_{ij}$ as there,
$g_{ij}(r,\theta)\le\sn_1^2(r)\hat g_{ij}(\theta)$; the hypothesis
$K\le1$ gives $g_{ij}(r,\theta)\ge s_1^2(r)\hat g_{ij}(\theta)$ by
Lemma~\ref{lem:rauch-lower}(2). Since
$\sn_1^2(r)\le r^2(1+r^2)$ and $s_1^2(r)\ge r^2(1-r^2/3)$ for
$0\le r\le 1/2$, comparing with
$g_{\mathrm{euc}}=dr^2+r^2g_{S^{n-1}}$ yields
\begin{equation}\label{twosided}
0.9\,g_{\mathrm{euc}}\ \le\ \tilde g\ \le\ 1.3\,g_{\mathrm{euc}}
\qquad\text{on }B(0,1/2)
\end{equation}
(as quadratic forms; the radial parts agree and
$1-r^2/3\ge0.9$, $1+r^2\le1.3$ there). In particular, for
$z,y\in\overline{B(0,1/32)}$, the straight segment gives
$d_{\tilde g}(z,y)\le1.2\,|z-y|\le1.2\cdot\tfrac{2}{32}\le0.08$.

\emph{Geodesics from $z$.} We apply
Lemma~\ref{lem:localized-cut-locus}(1) to the Riemannian manifold
$(B,\tilde g)$ at the point $z$, with $R=0.6$. Its hypothesis
holds: every point of the metric ball $B_{\tilde g}(z,0.6)$ is
joined to $z$ by a curve of $\tilde g$-length $<0.6$, hence lies
in $\{|x|\le|z|+0.6\}$ by the radial estimate of
Lemma~\ref{lem:pullback-ball}(1); as $|z|+0.6\le1/32+0.6<1$, the
closure of $B_{\tilde g}(z,0.6)$ is a closed subset of this
compact subset of $B$, hence compact. Therefore every unit-speed
$\tilde g$-geodesic from $z$ extends to every length less than
$0.6$; in particular $\widetilde\exp_z$ is defined on $B(0,1/2)$.
It is non-singular there: by Lemma~\ref{lem:conjugate-sturm}
($K\le1$ along all these geodesics and $1/2<\pi$) there is no
conjugate point along any of these geodesics at parameters
below $1/2$, so $d(\widetilde\exp_z)_v$ is non-singular for
$|v|<1/2$ by Lemma~\ref{lem:exponential-differential-jacobi}.

\emph{Injectivity of $\widetilde\exp_z$ on $B(0,1/8)$.} Suppose
$\sigma_1,\sigma_2$ are distinct $\tilde g$-geodesics from $z$ to a
common point $y'$, of lengths $L_1,L_2\le1/8$. Both lie in
$\{|x|\le5/32\}$. Consider the loop $\ell=\sigma_1\cup\bar\sigma_2$
at $z$, of length $\le1/4$. Using the linear structure of $B\subset
T_pM$, the map $(u,t)\mapsto u\,\ell(t)$ is a free homotopy from
$\ell$ to the constant loop at $0$, through loops based at $uz$
contained in $\{|x|\le5/32\}$ and of $\tilde g$-length at most
$1.2\cdot L_{\mathrm{euc}}(u\ell)=1.2\,u\,L_{\mathrm{euc}}(\ell)
\le1.2\cdot\left(1.1\cdot\tfrac14\right)\le0.33$
by (\ref{twosided}). Converting to a based homotopy at $z$ by
conjugating with the track $u\mapsto uz$ (of $\tilde g$-length
$\le1.2\,|z|\le0.04$) exhibits $\ell$ as fixed-endpoint homotopic,
through loops at $z$ of length $\le0.33+2\cdot0.04\le0.45$, to the
loop $\bar\tau\cup\tau$ ($\tau$ the track, run from $0$ to $z$),
which retracts to the constant loop through loops of length
$\le0.08$. Now apply
Lemma~\ref{lem:pullback-ball} \emph{to the manifold $(B,\tilde g)$
at the point $z$} with $\rho_*=1/2$: its hypotheses hold by the
previous paragraph, all the loops of the homotopy have length
$\le0.45<1/2$, so by parts (3) and (2) of that lemma the lift of
$\ell$ at $0\in T_zB$ is closed:
$$E_z(\sigma_1\cup\bar\sigma_2)=E_z(\text{constant})=0.$$
But $\sigma_1,\sigma_2$ are geodesics \emph{starting at $z$}, so
their lifts at $0$ are the rays to $w_i:=L_i\,\sigma_i'(0)$
(Lemma~\ref{lem:pullback-ball}(1) applied to $(B,\tilde g)$ at
$z$); and the lift of $\ell$ is the lift of $\sigma_1$ (ending at
$w_1$) followed by the lift of $\bar\sigma_2$ at $w_1$. The lift of
$\ell$ being closed means that the lift of $\bar\sigma_2$ at $w_1$
ends at $0$; reversing (Lemma~\ref{lem:pullback-ball}(2)), the lift
of $\sigma_2$ at $0$, which is the ray to $w_2$, ends at $w_1$;
hence $w_2=w_1$, so $L_1=L_2$ and
$\sigma_1'(0)=\sigma_2'(0)$, and $\sigma_1=\sigma_2$ ---
a contradiction. This proves both the injectivity in (1) and the
uniqueness in (2).

\emph{Existence of minimizers, and $d_z$.} We apply
Lemma~\ref{lem:localized-cut-locus}(2) to $(B,\tilde g)$ at $z$,
now with $R=1/9$: as above, the closure of the metric ball
$B_{\tilde g}(z,1/9)$ is compact (it lies in
$\{|x|\le1/32+1/9\}$, by the radial estimate). Hence every $y\in
B$ with $d_{\tilde g}(z,y)\le1/10<1/9$ is joined to $z$ by a
minimizing $\tilde g$-geodesic $\sigma$, of length $d_{\tilde
g}(z,y)\le1/10<1/8$; writing $w=L(\sigma)\,\sigma'(0)$ (for
unit-speed $\sigma$) we get $y=\widetilde\exp_z(w)$ with
$|w|<1/8$, i.e., $y\in\Omega_z$. For
$z,y\in\overline{B(0,1/32)}$ we have $d_{\tilde
g}(z,y)\le0.08\le1/10$, so such $z,y$ are joined by a minimizing
geodesic of length $\le0.08<1/8$; its uniqueness among geodesics
of length $\le1/8$ was proved above.

Next, $d_z=|\widetilde\exp_z^{-1}|$ on all of $\Omega_z$: let
$y=\widetilde\exp_z(w)$ with $|w|<1/8$. The radial geodesic gives
$d_z(y)\le|w|$. Conversely, no curve $c$ from $z$ to $y$ has
$L(c)<|w|$: for such a $c$, apply
Lemma~\ref{lem:pullback-ball}(2) to $(B,\tilde g)$ at $z$ with
$\rho_*=1/2$ (its hypotheses --- $\widetilde\exp_z$ defined and
non-singular on $B(0,1/2)$ --- hold by the second paragraph, and
$0+L(c)<1/8<1/2$): the lift of $c$ at $0\in T_zB$ ends at a
preimage of $y$ in $\overline{B(0,L(c))}\subset B(0,1/8)$, which
by the injectivity of $\widetilde\exp_z$ on $B(0,1/8)$ must be
$w$; but then the radial estimate
(Lemma~\ref{lem:pullback-ball}(1) for $(B,\tilde g)$ at $z$)
gives $L(c)\ge|w|$, a contradiction. Therefore $d_z(y)=|w|$,
i.e., $d_z=|\widetilde\exp_z^{-1}|$ on $\Omega_z$, which is
smooth away from $z$ since $\widetilde\exp_z$ is a diffeomorphism
of $B(0,1/8)$ onto $\Omega_z$ (non-singular and injective there).

(3) In geodesic polar coordinates around $z$ in $(B,\tilde g)$
(furnished on $\Omega_z\setminus\{z\}$ by the diffeomorphism just
established), $d_z$ is the radial coordinate, so
Lemma~\ref{lem:rauch-lower}(3) with $\bar K=1$ gives
$$\Hess_{\tilde g}\left(\tfrac12d_z^2\right)
\ge\min\left(1,\ d_z\,\frac{s_1'(d_z)}{s_1(d_z)}\right)\tilde g
=\min\left(1,\,d_z\cot d_z\right)\tilde g
\ge\left(1-\tfrac{d_z^2}{2}\right)\tilde g\ge\tfrac12\,\tilde g$$
on $\Omega_z\setminus\{z\}$ (where $d_z<1/8$). Since
$d_z^2=|\widetilde\exp_z^{-1}|^2$ is smooth on all of $\Omega_z$
(the squared norm is smooth at the origin) and both sides of the
inequality $\Hess_{\tilde g}(d_z^2)\ge\tilde g$ are continuous,
the bound extends from the dense subset
$\Omega_z\setminus\{z\}$ to all of $\Omega_z$, which is assertion
(3).
\end{proof}

\begin{lemma}[Distinctness of the lifted multiples of a geodesic loop]
\label{lem:loop-multiples-distinct}
\uses{lem:pullback-ball,lem:pullback-unique-minimizers}
\source{cheeger-gromov-taylor:page-0033}
In the setting of Lemma~\ref{lem:pullback-unique-minimizers}
($\rho_*=1$, $|K(P)|\le1$ on $(B,\tilde g)$), let $\gamma$ be a
geodesic loop at $p$ of length $2\ell>0$, smooth on the open
parameter interval, and let $N\ge1$ be an integer with
$2\ell N\le1/256$. Then the points
$$E(\gamma^0)=0,\ E(\gamma^1),\ \ldots,\ E(\gamma^N)$$
(where $\gamma^m$ denotes the $m$-fold concatenation) are pairwise
distinct points of $\overline{B(0,1/256)}$, all mapping to $p$
under $\exp_p$.
\end{lemma}

\begin{proof}
\uses{lem:pullback-ball,lem:pullback-unique-minimizers}
Each $\gamma^m$, $m\le N$, has length $2\ell m\le1/256<1$, so
$x_m:=E(\gamma^m)$ is defined, $|x_m|\le2\ell m\le1/256$
(Lemma~\ref{lem:pullback-ball}(2)) and
$\exp_p(x_m)=\gamma^m(\text{end})=p$.

First note that $E(\gamma)\not=0$: since $\gamma$ is smooth on the
open interval, its lift at $0$ is an unbroken $\tilde g$-geodesic
starting at $0$, hence a ray
(Lemma~\ref{lem:pullback-ball}(1)), so
$E(\gamma)=2\ell\,\gamma'(0)\not=0$.

Suppose now that $x_m=x_{m'}$ with $0\le m<m'\le N$. Writing
$\gamma^{m'}=\gamma^{i}\cup\gamma^{m}$ with $i:=m'-m\ge 1$ and
$\gamma^m$ as the common tail, cancellation
(Lemma~\ref{lem:pullback-ball}(4b)) gives $E(\gamma^i)=E(\text{constant
loop})=0$.

Let $\pi:=\pi_\gamma$ be the deck motion of
Lemma~\ref{lem:pullback-ball}(5); we shall use it, together with
its powers and inverse, on the ball $\overline{B(0,0.45)}$ (all
defined, since $0.45+2\ell N\le0.45+1/256<1$). We derive a
contradiction from $E(\gamma^i)=0$ in three steps. Throughout,
$d=d_{\tilde g}$, and we use the radial estimate of
Lemma~\ref{lem:pullback-ball}(1): every curve $c$ in $B$ satisfies
$L_{\tilde g}(c)\ge\bigl||c(1)|-|c(0)|\bigr|$, so in particular
$d(y_1,y_2)\ge\bigl||y_1|-|y_2|\bigr|$.

\emph{$\pi$ has no fixed point.} If $\pi(y)=y$ then
$E(\gamma\cup\mathrm{rad}_y)=y=E(\mathrm{rad}_y)$, so cancellation
with the common tail $\mathrm{rad}_y$ gives $E(\gamma)=0$,
contradicting the first paragraph.

\emph{$\pi$ permutes a finite orbit.} By
Lemma~\ref{lem:pullback-ball}(5) and (4d),
$$\pi(x_j)=E\left(\gamma\cup\mathrm{rad}_{E(\gamma^j)}\right)
=E\left(\gamma\cup\gamma^j\right)=x_{j+1},$$
and $x_i=E(\gamma^i)=0=x_0$; so $\pi$ permutes the finite set
$O:=\{x_0,\ldots,x_{i-1}\}\subset\overline{B(0,1/256)}$
cyclically. Moreover, by the composition rule,
$\pi^{j}=\pi_{\gamma^{j}}$ for $1\le j\le i$, and for
$y\in\overline{B(0,0.4)}$,
$$\pi_{\gamma^i}(y)=E\left(\gamma^i\cup\mathrm{rad}_y\right)$$
is the endpoint of the lift of $\mathrm{rad}_y$'s image at
$E(\gamma^i)=0$ (Lemma~\ref{lem:pullback-ball}(4c)), i.e., equals
$y$: thus $\pi^i=\mathrm{id}$ on $\overline{B(0,0.4)}$. The same
applies to the reversed loop: $\pi^{-1}=\pi_{\bar\gamma}$ and
$\pi_{\bar\gamma^{\,j}}=\pi^{-j}$, with
$\pi^{-i}=\mathrm{id}$ there as well.

In particular the powers of $\pi$ do not increase distances at the
scales we use: for $y_1,y_2\in\overline{B(0,0.2)}$ with
$d(y_1,y_2)\le0.2$ and $|j|\le i$,
$$d\left(\pi^{j}(y_1),\pi^{j}(y_2)\right)\le d(y_1,y_2).$$
Indeed, for $\epsilon\in(0,0.01)$ take a curve
$\omega$ from $y_1$ to $y_2$ of $\tilde g$-length $\le
d(y_1,y_2)+\epsilon\le0.21$; by the radial estimate $\omega$
stays in $B(0,0.2+0.22)\subset B(0,0.45)$, where $\pi^{j}$ is
defined with everywhere isometric differential
(Lemma~\ref{lem:pullback-ball}(5), applied to the loop
$\gamma^{|j|}$ or $\bar\gamma^{\,|j|}$, of length
$2\ell|j|\le1/256$). Hence $\pi^{j}\circ\omega$ joins
$\pi^{j}(y_1)$ to $\pi^{j}(y_2)$ with the same length, so
$d(\pi^{j}y_1,\pi^{j}y_2)\le d(y_1,y_2)+\epsilon$, and letting
$\epsilon\to0$ the claim follows. Combining the two directions
gives the equality
\begin{equation}\label{deckisometry}
d\left(\pi^{j}(y_1),\pi^{j}(y_2)\right)
=d(y_1,y_2),\qquad
y_1,y_2\in\overline{B(0,1/7)},\ d(y_1,y_2)\le0.2,\
|j|\le i:
\end{equation}
here $\le$ is the contraction inequality; and since
$|\pi^{j}(y_k)|\le2\ell|j|+|y_k|\le1/256+1/7\le0.2$
(Lemma~\ref{lem:pullback-ball}(2)) and
$d(\pi^{j}y_1,\pi^{j}y_2)\le0.2$, the contraction inequality
applied to the points $\pi^{j}(y_k)$ with the power $-j$, together
with $\pi^{-j}\circ\pi^{j}=\pi_{\bar\gamma^{\,j}\cup\gamma^{j}}
=\mathrm{id}$ on $\overline{B(0,1/7)}$
(the composition rule and the reverse--retrace homotopy of
Lemma~\ref{lem:pullback-ball}(5)), gives the reverse inequality
$d(y_1,y_2)\le d(\pi^{j}y_1,\pi^{j}y_2)$.

\emph{The center of mass.} Set
$$F(y):=\sum_{j=0}^{i-1}d\left(x_j,y\right)^2,
\qquad y\in\overline{B(0,1/16)},$$
a continuous function on a compact set; let $y^*$ be a minimum
point, and write $m=F(y^*)$. Since all
$x_j\in\overline{B(0,1/256)}$, and since for any $y\in B$ the ray
$t\mapsto ty$ is a curve of $\tilde g$-length $|y|$
(Lemma~\ref{lem:pullback-ball}(1)), so that $d(0,y)\le|y|$, we
have $d(0,x_j)\le|x_j|\le1/256\le0.005$, so $F(0)\le
i\,(0.005)^2$; while for $|y|=1/16$ the radial estimate gives
$d(x_j,y)\ge|y|-|x_j|\ge1/16-1/256\ge 0.058$, so
$F(y)\ge i\,(0.058)^2>F(0)$: the minimum is attained in the open
ball $B(0,1/16)$. The minimum point lies close to the orbit: since
$d(x_j,y^*)\ge d(0,y^*)-d(0,x_j)\ge d(0,y^*)-0.005$ for every
$j$, the inequality
$i\left(d(0,y^*)-0.005\right)^2\le F(y^*)\le F(0)\le
i\,(0.005)^2$ (valid whenever $d(0,y^*)\ge0.005$) forces
$d(0,y^*)\le0.01$; in particular $|y^*|\le0.01$ by the radial
estimate. By (\ref{deckisometry}) (applicable: the points
$x_{j}$ and any $y\in\overline{B(0,1/16)}$ lie in
$\overline{B(0,1/7)}$ at mutual distances
$d(x_j,y)\le d(x_j,0)+d(0,y)\le|x_j|+|y|
\le1/256+1/16\le0.2$, using $d(0,\cdot)\le|\cdot|$ as above) and
the permutation property,
$$F(\pi(y))=\sum_jd\left(x_j,\pi(y)\right)^2
=\sum_jd\left(\pi\left(\pi^{-1}(x_j)\right),\pi(y)\right)^2
=\sum_jd\left(\pi^{-1}(x_j),y\right)^2=F(y)$$
for $y\in\overline{B(0,1/16)}$ with $\pi(y)\in
\overline{B(0,1/16)}$ (the inverse $\pi^{-1}=\pi_{\bar\gamma}$
permutes $O$ as well). Now
$d(0,\pi(y^*))\le d(0,x_1)+d(x_1,\pi(y^*))
=d(0,x_1)+d(x_0,y^*)\le0.005+0.01\le0.015$ (using
(\ref{deckisometry}) and $\pi(x_0)=x_1$), so
$|\pi(y^*)|\le0.015$ and $\pi(y^*)$ lies in the open ball
$B(0,1/16)$ as well, with $F(\pi(y^*))=F(y^*)=m$: it is also a
minimum point.

If $y^*\not=\pi(y^*)$, note first that both points lie in
$\overline{B(0,1/32)}$ (their norms are $\le0.015<1/32$), so by
Lemma~\ref{lem:pullback-unique-minimizers}(2) there is a
minimizing geodesic $\sigma\colon[0,1]\to B$ from $y^*$ to
$\pi(y^*)$, of length $d(y^*,\pi(y^*))\le
d(y^*,0)+d(0,\pi(y^*))\le0.01+0.015=0.025$. Every point of
$\sigma$ satisfies $d(y^*,\sigma(t))\le0.025$, hence
$d(0,\sigma(t))\le0.035$ and, by the radial estimate,
$|\sigma(t)|\le0.035<1/16$: the whole of $\sigma$ lies in the open
ball $B(0,1/16)$, the interior of the domain over which $F$ was
minimized. Moreover, for each $j$,
$d\left(x_j,\sigma(t)\right)\le d(x_j,0)+d(0,\sigma(t))
\le0.005+0.035\le1/10$, so by the first clause of
Lemma~\ref{lem:pullback-unique-minimizers}(2) (applicable, as
$x_j\in\overline{B(0,1/256)}\subset\overline{B(0,1/32)}$) we get
$\sigma(t)\in\Omega_{x_j}$, on which $d_{x_j}^2$ is smooth with
$\Hess(d_{x_j}^2)\ge\tilde g$
(Lemma~\ref{lem:pullback-unique-minimizers}(2),(3)). Hence $t\mapsto
F(\sigma(t))$ is smooth with second derivative
$\sum_j\Hess(d_{x_j}^2)(\sigma',\sigma')\ge
i\,|\sigma'|^2>0$ ($\sigma$ is non-constant), so it is strictly
convex on $[0,1]$; consequently, for $0<t<1$,
$F(\sigma(t))<\max\left(F(\sigma(0)),F(\sigma(1))\right)=m$. But
$\sigma(t)\in B(0,1/16)$, where $F\ge m$ --- a contradiction.
Therefore $\pi(y^*)=y^*$, contradicting fixed-point freeness.
Hence the $x_m$, $0\le m\le N$, are pairwise distinct.
\end{proof}

\begin{theorem}[Volume lower bound implies injectivity radius lower bound]
\label{thm:volume-injectivity-radius}
\label{volinj}
\uses{def:injectivity-radius,def:curvature-operator-norm}
Fix an integer $n>0$. For every $\epsilon>0$ there is $\delta>0$
depending on $n$ and $\epsilon$ such that the following holds.
Suppose that $(M^n,g)$ is a complete connected Riemannian
manifold of dimension $n$ and that $p\in M$. Suppose that $|\Rm(x)|\le
r^{-2}$ for all $x\in B(p,r)$. If $Vol(B(p,r))\ge \epsilon r^n$ then
the injectivity radius of $M$ at $p$ is at least $\delta r$.
\end{theorem}

\begin{proof}
\uses{def:curvature-operator-norm,thm:hopf-rinow,lem:conjugate-sturm,lem:klingenberg-loop,lem:pullback-ball,lem:pullback-unique-minimizers,lem:loop-multiples-distinct,lem:volume-element-comparison,lem:geodesic-polar-form,thm:bishop-gromov,def:injectivity-radius,def:ricci-curvature,lem:metric-rescaling,lem:exponential-differential-jacobi,lem:model-polar-isometry}
\source{cheeger-gromov-taylor:page-0032,cheeger-gromov-taylor:page-0033}
Replacing $g$ by $r^{-2}g$, which by
Lemma~\ref{lem:metric-rescaling} (with $c=r^{-2}$) multiplies
$|\Rm|$ by $r^{2}$,
distances by $r^{-1}$, volumes by $r^{-n}$ and the injectivity
radius by $r^{-1}$, while preserving completeness, we may assume
$r=1$: thus $M$ is
complete, $|\Rm|\le1$ on $B(p,1)$, and $\Vol B(p,1)\ge\epsilon$; we
must bound $\inj_M(p)$ below by a constant $\delta(n,\epsilon)>0$.

Throughout, by Definition~\ref{def:curvature-operator-norm} the
hypothesis gives $|K(P)|\le1$ for all $2$-planes tangent at points
of $B(p,1)$, hence also $\Ric\ge-(n-1)$ there (expressing
$\Ric(X,X)$ as a sum of $n-1$ sectional curvatures in an
orthonormal basis). Every unit-speed geodesic from $p$
restricted to $[0,T]$, $T<1$, lies in
$\overline{B(p,T)}\subset B(p,1)$, where $K(P)\le1$; applying
Lemma~\ref{lem:conjugate-sturm} to it (with $K=1$, and
$T<1<\pi=\pi/\sqrt{1}$)
shows: \emph{no geodesic from $p$ has a conjugate point at
distance less than $1$}. By
completeness and the Hopf--Rinow theorem
(Theorem~\ref{thm:hopf-rinow}(1)), $\exp_p$ is defined
on all of $T_pM$, and by
Lemma~\ref{lem:exponential-differential-jacobi} it is
non-singular on $B(0,1)$ (a singular point $v$, $|v|<1$, would
produce a conjugate point at distance $|v|<1$ along
$\gamma_v$). Thus
Lemma~\ref{lem:pullback-ball} applies with $\rho_*=1$; the
pulled-back metric $\tilde g=\exp_p^*g$ on $B=B(0,1)$
satisfies $|K_{\tilde g}(P)|\le1$ (the curvature tensor of a
pullback under a local isometry is the pullback of the curvature
tensor), so Lemmas~\ref{lem:pullback-unique-minimizers} and
\ref{lem:loop-multiples-distinct} apply as well.

If $\inj_M(p)\ge\tfrac{1}{1024}$, we are done with
$\delta=\tfrac1{1024}$. So suppose
$\ell:=\inj_M(p)<\tfrac1{1024}$. Since there are no conjugate
points at distance $\le\ell<1$ along geodesics from $p$,
Lemma~\ref{lem:klingenberg-loop} provides a geodesic loop $\gamma$
at $p$ of length exactly $2\ell>0$, smooth on the open parameter
interval.

\emph{Counting preimages.} Let
$$N:=\left\lfloor\frac{1}{512\,\ell}\right\rfloor\ \ge\ 2,
\qquad\text{so that}\qquad 2\ell N\le\frac{1}{256}.$$
By Lemma~\ref{lem:loop-multiples-distinct} the points
$x_m=E(\gamma^m)$, $0\le m\le N$, are $N+1$ distinct points of
$\overline{B(0,1/256)}$ with $\exp_p(x_m)=p$. Now fix any
$q\in B(p,s)$ with $s:=\tfrac1{256}$, and let $\sigma$ be a
minimizing geodesic from $p$ to $q$, which exists by completeness
and the Hopf--Rinow theorem (Theorem~\ref{thm:hopf-rinow}(2)),
of length $<s$. The points
$$y_m:=E\left(\gamma^m\cup\sigma\right),\qquad 0\le m\le N,$$
satisfy $\exp_p(y_m)=q$ and $|y_m|\le2\ell m+s\le\tfrac1{128}$
(Lemma~\ref{lem:pullback-ball}(2)), and they are pairwise
distinct: if $y_m=y_{m'}$ then cancellation of the common tail
$\sigma$ (Lemma~\ref{lem:pullback-ball}(4b)) gives
$x_m=x_{m'}$, so $m=m'$. Hence \emph{every point of $B(p,s)$ has
at least $N+1$ preimages under $\exp_p$ in
$\overline{B(0,1/128)}$}.

\emph{Volume transfer.} Since $\exp_p$ is a local
diffeomorphism on $B(0,1)$, the open set $B(0,1/64)$ can be written
as a countable disjoint union of Borel sets on each of which
$\exp_p$ is injective; summing the change of variables formula
over the pieces gives
$$\int_M\#\left(\exp_p^{-1}(q)\cap
B(0,1/64)\right)d\vol(q)
=\Vol_{\tilde g}\left(B(0,1/64)\right),$$
where $\Vol_{\tilde g}$ denotes the volume for the pulled-back
metric. The left side is at least $(N+1)\Vol B(p,s)$ by the
previous paragraph (note $\overline{B(0,1/128)}\subset B(0,1/64)$).
For the right side, along every ray the hypotheses of
Lemma~\ref{lem:volume-element-comparison} hold with $k=1$ and
$r_0=1$ (no conjugate points, $\Ric\ge-(n-1)$), so the volume
element satisfies $\lambda(t,v)\le\sn_1^{n-1}(t)$ and, integrating
in polar coordinates (Lemma~\ref{lem:geodesic-polar-form}(4)),
$$\Vol_{\tilde g}\left(B(0,1/64)\right)
\le\omega_{n-1}\int_0^{1/64}\sn_1^{n-1}(t)\,dt=:V_1(1/64),$$
a constant depending only on $n$. Hence
\begin{equation}\label{volinjcount}
(N+1)\,\Vol B(p,s)\ \le\ V_1(1/64).
\end{equation}

\emph{The volume lower bound.} Closed balls in the complete
manifold $M$ are compact, by the Hopf--Rinow theorem
(Theorem~\ref{thm:hopf-rinow}(3)). Thus
Theorem~\ref{thm:bishop-gromov} applies on $B(p,R)$ for every
$R\le1$, with $k=1$: for $s<r'<1$,
$$\frac{\Vol B(p,s)}{V_1(s)}\ \ge\ \frac{\Vol B(p,r')}{V_1(r')},$$
where $V_1(t)=\omega_{n-1}\int_0^t\sn_1^{n-1}$ is the volume of
the $t$-ball in the model $H^n_1$ (computed via
Lemma~\ref{lem:model-polar-isometry}, as in the proof of
Theorem~\ref{thm:bishop-gromov}). Letting $r'\to1$ and using
$\Vol B(p,r')\to\Vol B(p,1)\ge\epsilon$,
$$\Vol B(p,s)\ \ge\ \frac{V_1(s)}{V_1(1)}\,\epsilon
\ =:\ c_1(n)\,\epsilon.$$

\emph{Conclusion.} Combining this with (\ref{volinjcount}),
$$\frac{1}{512\,\ell}\ \le\ N+1\ \le\
\frac{V_1(1/64)}{c_1(n)\,\epsilon},$$
so $\ell\ge\delta_0(n,\epsilon):=
\dfrac{c_1(n)\,\epsilon}{512\,V_1(1/64)}$. Taking
$\delta=\min\left(\tfrac1{1024},\ \delta_0(n,\epsilon)\right)$
proves the theorem (after undoing the initial rescaling, which
multiplies both $\inj_M(p)$ and the required bound by $r$).
\end{proof}

